\documentclass[11pt,a4paper,twoside,openany]{book}

% ---- encoding and fonts (all from texlive-fonts-recommended, present on arXiv)
\usepackage[T1]{fontenc}
\usepackage[utf8]{inputenc}
\usepackage{textcomp}
\usepackage{mathpazo}                 % text: URW Palladio (Palatino)
\usepackage[scaled=0.90]{helvet}      % headings: Nimbus Sans
\usepackage{amssymb}
\usepackage{pmboxdraw}                % box-drawing glyphs for pdflatex
\usepackage{microtype}
\linespread{1.06}

% ---- non-ASCII characters that appear in the manuscript
\DeclareUnicodeCharacter{2192}{\ensuremath{\rightarrow}}
\DeclareUnicodeCharacter{2193}{\ensuremath{\downarrow}}
\DeclareUnicodeCharacter{25B6}{\ensuremath{\blacktriangleright}}
\DeclareUnicodeCharacter{25C0}{\ensuremath{\blacktriangleleft}}
\DeclareUnicodeCharacter{25BC}{\ensuremath{\blacktriangledown}}
\DeclareUnicodeCharacter{2713}{\ensuremath{\checkmark}}
\DeclareUnicodeCharacter{2717}{\ensuremath{\times}}
\DeclareUnicodeCharacter{23F8}{\textbf{||}}
\DeclareUnicodeCharacter{00B7}{\textperiodcentered}
\DeclareUnicodeCharacter{00A7}{\S}
\DeclareUnicodeCharacter{2020}{\textdagger}   % effort-table footnote marker

% ---- page geometry
\usepackage[a4paper,inner=32mm,outer=28mm,top=25mm,bottom=30mm,
            headsep=7mm,footskip=12mm]{geometry}

% ---- colour identity
%
% Indigo and brass, unchanged as an identity. What changed is its application.
% Brass measures 2.42:1 on white, so it can carry meaning only on a dark ground;
% every light-ground use of it was previously substituted by an olive
% #8A6D1E, which is the muddiest point on the gold ramp. brassdark is now a
% warmer bronze at 5.85:1, and the light-ground accent role moved to indigo.
% Every value here is gated by src/check_palette.py.
\usepackage{xcolor}
\definecolor{deepindigo}{HTML}{0F1B33}   % dark ground
\definecolor{indigo}{HTML}{1E2E56}       % section + subsection headings, note rule
\definecolor{brass}{HTML}{C9A227}        % accent on DARK grounds only (2.42:1 on white)
\definecolor{brassdark}{HTML}{8C5A16}    % bronze: kickers and the caution rule (5.85:1)
\definecolor{cream}{HTML}{F4F1E8}        % type on dark grounds
\definecolor{bodyink}{HTML}{16202B}      % running text; was unset, i.e. pure black
\definecolor{paleindigo}{HTML}{DFE5F0}   % practitioner-note fill (1.26:1)
\definecolor{palebrass}{HTML}{F6E8CC}    % caution fill (1.21:1)
\definecolor{palegrey}{HTML}{ECEFF5}     % generic callout fill
\definecolor{quoterule}{HTML}{6E7A8B}    % generic callout rule
\definecolor{inkgrey}{HTML}{5A6472}      % running heads, TOC sections, code rule
\definecolor{codebg}{HTML}{EAEDF3}       % code ground (1.17:1; was 1.07:1, invisible)

% ---- structure
\usepackage{titlesec}
\usepackage{fancyhdr}
\usepackage{enumitem}
\usepackage{longtable}
\usepackage{array}
\usepackage{booktabs}
\usepackage{ragged2e}
\usepackage{needspace}
\usepackage{tikz}
\usepackage[most]{tcolorbox}
\usepackage{eso-pic}
\usepackage{placeins}
\usepackage{multicol}

% ---- definition lists
%
% The glossary and the index are definition lists, and were being set as single
% run-on paragraphs: 32 glossary entries and 40-odd index entries flowed
% together, so a locator list ran straight into the next term. A hanging indent
% keeps the term at the margin and its text indented under it, which is what
% makes such a list scannable at all.
\newenvironment{deflist}{%
  \begin{list}{}{%
    \setlength{\leftmargin}{6mm}\setlength{\itemindent}{-6mm}%
    \setlength{\listparindent}{0pt}\setlength{\labelwidth}{0pt}%
    \setlength{\labelsep}{0pt}\setlength{\topsep}{2mm}%
    \setlength{\itemsep}{0.9mm}\setlength{\parsep}{0pt}}}%
  {\end{list}}

% ---- float placement
%
% Figures are floats so they can never be split across a page boundary. The
% measured tallest figure occupies 0.58 of the text block, so a top fraction of
% 0.85 guarantees every one of them can take a top slot on the page it belongs
% to. build_tex.py emits an explicit \FloatBarrier before every section and
% every chapter, part and appendix opener, so a figure can never drift out of
% its own section or land on a full-bleed part page. The [section] option of
% placeins is deliberately not used: it patches \section at load time and
% titlesec redefines \section afterwards, which makes the result depend on
% package order.
\setcounter{topnumber}{2}
\setcounter{bottomnumber}{1}
\setcounter{totalnumber}{3}
\renewcommand{\topfraction}{0.85}
\renewcommand{\bottomfraction}{0.60}
\renewcommand{\textfraction}{0.08}
\renewcommand{\floatpagefraction}{0.70}
\setlength{\floatsep}{6mm plus 2mm minus 1mm}
\setlength{\textfloatsep}{7mm plus 2mm minus 2mm}
\setlength{\intextsep}{6mm plus 2mm minus 1mm}
\usepackage[hidelinks,bookmarksnumbered,pdfusetitle]{hyperref}
\hypersetup{
  pdfsubject={A verified practitioner's reference for agentic work. Version 1.0. Verification baseline: Claude Code 2.1.241, 23 Aug. 2026; re-verified against 2.1.246, 26 Aug. 2026.},
  pdfkeywords={agentic artificial intelligence, AI-assisted software engineering,
    Claude Code, large language model agents, Model Context Protocol, prompt injection,
    zero trust, software supply-chain security, human oversight, AI governance,
    verification and validation}}
\usepackage{bookmark}

% ---- running heads
\pagestyle{fancy}
\fancyhf{}
\renewcommand{\headrulewidth}{0pt}
\renewcommand{\footrulewidth}{0pt}
\fancyhead[LE]{\footnotesize\sffamily\color{inkgrey}\nouppercase{\leftmark}}
\fancyhead[RO]{\footnotesize\sffamily\color{inkgrey}\nouppercase{\rightmark}}
\fancyfoot[LE,RO]{\footnotesize\sffamily\color{deepindigo}\thepage}
\fancypagestyle{plain}{\fancyhf{}\fancyfoot[LE,RO]{%
  \footnotesize\sffamily\color{deepindigo}\thepage}}
\renewcommand{\chaptermark}[1]{\markboth{#1}{#1}}
\renewcommand{\sectionmark}[1]{\markright{\thesection\ \ #1}}

% ---- the brass rule-stack, echoing the cover
\newcommand{\rulestack}[1][brass]{%
  \noindent
  \begin{tikzpicture}[baseline=0pt]
    \fill[#1]        (0,0)     rectangle (74mm,1.5mm);
    \fill[#1!78!white] (0,-3.6mm) rectangle (52mm,-2.1mm);
    \fill[#1!56!white] (0,-7.2mm) rectangle (34mm,-5.7mm);
    \fill[#1!36!white] (0,-10.8mm) rectangle (20mm,-9.3mm);
  \end{tikzpicture}}

% ---- part pages: full-bleed indigo
\newcommand{\partbackground}{%
  \AddToShipoutPictureBG*{%
    \AtPageLowerLeft{\tikz[remember picture,overlay]{%
      \fill[left color=deepindigo,right color=indigo]
        (0,0) rectangle (\paperwidth,\paperheight);}}}}

% ---- title page: the same field, held clear of the left paper edge
%
% arXiv stamps the identifier down the left margin of page 1. Its own FAQ names
% "a figure that goes to the edge of the page" as a reason the stamp is lost or
% unreadable, so the title-page field starts one gutter in. Part pages, which
% arXiv never stamps, stay full-bleed. The body text block begins at inner=32mm,
% well clear of the 20mm gutter, so nothing else moves.
\newlength{\arxivstampgutter}
\setlength{\arxivstampgutter}{20mm}
\newcommand{\titlebackground}{%
  \AddToShipoutPictureBG*{%
    \AtPageLowerLeft{\tikz[remember picture,overlay]{%
      \fill[left color=deepindigo,right color=indigo]
        (\arxivstampgutter,0) rectangle (\paperwidth,\paperheight);}}}}

\newcommand{\parttitle}[2]{%
  \FloatBarrier\clearpage
  \thispagestyle{empty}%
  \partbackground
  \begingroup
  \sffamily\color{cream}
  \vspace*{0.62\textheight}
  \noindent\rulestack\par
  \vspace{7mm}
  {\fontsize{10}{12}\selectfont\bfseries\color{brass}\MakeUppercase{#1}\par}
  \vspace{5mm}
  {\fontsize{27}{31}\selectfont\bfseries\raggedright #2\par}
  \endgroup
  \clearpage
  \markboth{#1\ \ \textperiodcentered\ \ #2}{#1\ \ \textperiodcentered\ \ #2}}

% ---- chapter openings: a title page of their own
\titleformat{\chapter}[display]
  {\sffamily\color{deepindigo}}
  {\vspace*{0.30\textheight}%
   \parbox{\textwidth}{\rulestack\\[6mm]%
     \fontsize{9.6}{12}\selectfont\bfseries\color{brassdark}%
     \MakeUppercase{\chaptertitlelabel}}}
  {4mm}
  {\fontsize{25}{29}\selectfont\bfseries\raggedright}
  [\vspace*{\fill}\clearpage]
\titlespacing*{\chapter}{0pt}{0pt}{0pt}
\newcommand{\chaptertitlelabel}{}

% front- and back-matter openers: a plain macro, not a titlesec shape
\newcommand{\frontchapter}[1]{%
  \FloatBarrier\clearpage
  \thispagestyle{plain}%
  \noindent\rulestack\par
  \vspace{7mm}
  {\sffamily\fontsize{21}{25}\selectfont\bfseries\color{deepindigo}\raggedright #1\par}
  \vspace{7mm}
  \addcontentsline{toc}{chapter}{#1}%
  \markboth{#1}{#1}}

% The contents list sets its own opener. \tableofcontents emits the heading as
% \chapter*{\contentsname}, and titlesec applies the chapter shape's after-code
% -- \vspace*{\fill}\clearpage -- to the starred form as well as the numbered
% one. That flung "Contents" onto a page of its own and started the entries on
% the next: a heading separated from the thing it introduces, the same defect
% class as a table stranding its heading. Setting the opener here in the
% \frontchapter idiom and calling \@starttoc directly keeps the two together,
% and matches the other front-matter openers. No \addcontentsline: a contents
% list does not list itself.
\makeatletter
\newcommand{\handbooktoc}{%
  \FloatBarrier\clearpage
  \thispagestyle{plain}%
  \noindent\rulestack\par
  \vspace{7mm}
  {\sffamily\fontsize{21}{25}\selectfont\bfseries\color{deepindigo}\raggedright
    \contentsname\par}
  \vspace{7mm}
  \markboth{\contentsname}{\contentsname}%
  \@starttoc{toc}}
\makeatother

% appendices: a rule-topped opener, not a page of their own
\newcommand{\appendixchapter}[2]{%
  \FloatBarrier\clearpage
  \thispagestyle{plain}%
  \noindent\rulestack\par
  \vspace{6mm}
  {\sffamily\fontsize{9.6}{12}\selectfont\bfseries\color{brassdark}%
    \MakeUppercase{Appendix #1}\par}
  \vspace{3.4mm}
  {\sffamily\fontsize{21}{25}\selectfont\bfseries\color{deepindigo}\raggedright #2\par}
  \vspace{7mm}
  \addcontentsline{toc}{chapter}{\protect\numberline{#1}#2}%
  \markboth{Appendix #1\ \ \textperiodcentered\ \ #2}{Appendix #1\ \ \textperiodcentered\ \ #2}}

\titleformat{\section}{\sffamily\large\bfseries\color{indigo}}{\thesection}{0.6em}{}
\titlespacing*{\section}{0pt}{5.2mm plus 1mm minus .6mm}{1.9mm}
% 57 subsections set in olive gold on white read as cheap and measured only
% 4.90:1. The light-ground accent role belongs to indigo; brass stays on dark.
\titleformat{\subsection}{\sffamily\normalsize\bfseries\color{indigo}}{}{0em}{}
\titlespacing*{\subsection}{0pt}{4.2mm plus .8mm minus .5mm}{1.4mm}

% ---- table of contents
\setcounter{tocdepth}{1}
\setcounter{secnumdepth}{2}
\usepackage{titletoc}
\titlecontents{part}[0pt]{\addvspace{5mm}\sffamily\bfseries\color{deepindigo}}
  {}{}{}
\titlecontents{chapter}[3.4em]{\addvspace{1.4mm}\sffamily}
  {\contentslabel{3.4em}}{\hspace*{-3.4em}}
  {\titlerule*[0.7pc]{.}\contentspage}
\titlecontents{section}[6.6em]{\small\color{inkgrey}}
  {\contentslabel{3.2em}}{}{\titlerule*[0.7pc]{.}\contentspage}

% ---- lists
\setlist[itemize]{leftmargin=5.5mm,itemsep=0.5mm,topsep=1.4mm,parsep=0pt}
\setlist[enumerate]{leftmargin=6.5mm,itemsep=0.5mm,topsep=1.4mm,parsep=0pt}

% ---- code blocks: verbatim rendered from pre-escaped lines
%
% The four block types are separated by BOTH accent hue and rule weight. Weight
% is the cue that survives greyscale printing, which hue alone does not: the
% code rule and the caution rule sit at almost the same luminance (0.125 and
% 0.129) and are told apart by 0.7mm against 2.0mm.
\newtcolorbox{codeblock}[2]{breakable, enhanced, sharp corners, boxrule=0pt,
  colback=codebg, colframe=codebg,
  borderline west={0.7mm}{0pt}{inkgrey},
  left=3.4mm, right=1.5mm, top=2.2mm, bottom=2.2mm,
  before skip=2.8mm, after skip=2.8mm, parbox=false,
  fontupper=\fontsize{#1}{#2}\selectfont\ttfamily}

% The figure variant is explicitly unbreakable. A figure that fits inside the
% text block must never be split across a page boundary, which is what the
% breakable code box did to all seven of them.
\newtcolorbox{codefig}[2]{breakable=false, enhanced, sharp corners, boxrule=0pt,
  colback=codebg, colframe=codebg,
  borderline west={0.7mm}{0pt}{inkgrey},
  left=3.4mm, right=1.5mm, top=2.2mm, bottom=2.2mm,
  before skip=0mm, after skip=0mm, parbox=false,
  fontupper=\fontsize{#1}{#2}\selectfont\ttfamily}
\newcommand{\cl}[1]{\noindent\mbox{#1}\par}

% ---- callouts: fill, rule colour, rule weight
\newtcolorbox{calloutbox}[3]{breakable, enhanced, sharp corners, boxrule=0pt,
  colback=#1, colframe=#1, borderline west={#3}{0pt}{#2},
  left=3.4mm, right=3mm, top=2.4mm, bottom=2.4mm,
  before skip=3mm, after skip=3mm, fontupper=\small}

% ---- captions for tables and figures
% \tabcaption  free-standing, above a longtable that is allowed to break
% \boxcaption  above a table already inside an atomic box: no leading skip and
%              no \needspace, both of which are meaningless inside a minipage
% \figcaption  below a figure, which is where IEEE puts a figure caption
\newcommand{\tabcaption}[1]{\par\addvspace{2.6mm}\needspace{4\baselineskip}%
  {\sffamily\footnotesize\bfseries\color{deepindigo}#1\par}\addvspace{1.2mm}}
\newcommand{\boxcaption}[1]{%
  {\sffamily\footnotesize\bfseries\color{deepindigo}#1\par}\addvspace{1.4mm}}
\newcommand{\figcaption}[1]{\par\addvspace{2.6mm}%
  {\sffamily\footnotesize\color{deepindigo}\RaggedRight #1\par}}

% ---- ornamental break
\newcommand{\breakrule}{\par\addvspace{3mm}\noindent
  \hspace*{\fill}{\color{brass}\rule{18mm}{0.5mm}}\hspace*{\fill}\par\addvspace{3mm}}

% ---- table column helper
% \tblbox holds a whole table while its set height is tested against the page.
% build_tex.py emits the test; LaTeX takes the decision, because a table's real
% height depends on how ragged-right columns wrap and cannot be predicted from
% the source text.
\newsavebox{\tblbox}
% \kpbox does the same for a Key points block, which is a heading bound to the
% list it introduces and must be placed as one unit.
\newsavebox{\kpbox}
% Two ragged-right column types, differing only in what a hyphen costs.
%
% L is the default. \hyphenpenalty was 150 for every column, which makes a
% hyphen almost free, so TeX split words rather than breaking at a space:
% "Official documenta-/tion", "Third-party compo-/nents". At 2500 a space-break
% wins wherever one exists.
%
% H is for a column too narrow to hold its own longest word -- which happens in
% the wide crosswalk tables, where six columns of standards vocabulary need
% more measure than A4 has. There a hyphen is not gratuitous, it is the only
% way to avoid overflowing the measure, so the old cheap penalty applies.
% build_tex.py decides per column, by comparing the allocated width against the
% width the longest token actually needs.
\newcolumntype{L}[1]{>{\RaggedRight\arraybackslash\hyphenpenalty=2500\relax}p{#1}}
\newcolumntype{H}[1]{>{\RaggedRight\arraybackslash\hyphenpenalty=150\relax}p{#1}}
\renewcommand{\arraystretch}{1.18}
\setlength{\tabcolsep}{1.6mm}

% Retained defensively. Under openany nothing calls it: chapter, part, appendix
% and front-matter openers all use \clearpage, so no blank verso is emitted at
% all. If a future change reinstates openright, this keeps those versos free of
% a running head and a folio.
\makeatletter
\def\cleardoublepage{\clearpage\if@twoside\ifodd\c@page\else
  \null\thispagestyle{empty}\newpage\if@twocolumn\null\newpage\fi\fi\fi}
\makeatother

% Running text was previously left at the default pure black while the HTML and
% DOCX editions used #16202B. The typeset edition now matches the identity.
\AtBeginDocument{\color{bodyink}}

\hyphenpenalty=1200
\tolerance=1800
\emergencystretch=3em
\widowpenalty=10000
\clubpenalty=10000
\raggedbottom

\newenvironment{refslist}{\begin{list}{}{\leftmargin=8mm\labelwidth=6mm\labelsep=2mm\itemindent=-8mm\listparindent=0pt\itemsep=1.1mm\parsep=0pt\topsep=2mm}}{\end{list}}
\title{Claude Code Complete User Handbook}
\author{David Soldani}
\date{26 August 2026}
\begin{document}
\frontmatter
\begin{titlepage}
\thispagestyle{empty}
\titlebackground
\begingroup\sffamily\color{cream}
\vspace*{0.30\textheight}
\rulestack
\vspace{9mm}\par
{\fontsize{9}{11}\selectfont\bfseries\color{brass}\MakeUppercase{A verified practitioner\textquotesingle s reference}\par}
\vspace{12mm}
{\fontsize{34}{38}\selectfont\bfseries\raggedright Claude Code Complete User Handbook\par}
\vspace{7mm}
{\fontsize{15}{19}\selectfont\color{paleindigo}\raggedright A verified practitioner's reference for agentic work\par}
\vspace*{\fill}
\noindent{\color{brass}\rule{\textwidth}{0.25mm}}\par\vspace{5mm}
{\fontsize{12.5}{16}\selectfont\bfseries David Soldani\par}
\vspace{5mm}
{\fontsize{9.4}{13}\selectfont\color{paleindigo}Version 1.0\ \ \textperiodcentered\ \ 26 August 2026\par}
\endgroup
\end{titlepage}
\clearpage
\pagestyle{fancy}
\handbooktoc
\clearpage

\textbf{Verification baseline:} Claude Code 2.1.241. Every statement of product behaviour in this book was adjudicated against Anthropic's official documentation on 23 August 2026, enumerated from the published documentation index [1].


\textbf{Re-verified 26 August 2026 against Claude Code 2.1.246.} All 81 cited documentation pages were re-fetched and hashed on that date and the published index re-enumerated: no cited page had been removed or renamed, and the 57 pages added since the baseline are all sub-pages of sections this book already covers. All 96 unique URLs in the text resolved, one of them behind an authentication challenge. Ten of the most heavily cited pages were re-adjudicated claim by claim; the effort vocabulary was re-established against a running installation, and the permission-mode vocabulary against that installation and the documentation that governs it, which disagreed. The corrections this produced are listed in Appendix J. The re-verification was a reconciliation and a targeted re-adjudication, not a second full pass: every claim not named in the evidence ledger rests on the 23 August 2026 adjudication.


Interfaces, limits, plan entitlements and preview features change on a weekly cadence; Appendix K states the re-verification procedure and should be run before this book is relied upon for a consequential decision.


\breakrule

\FloatBarrier
\section*{Abstract}
\markboth{Abstract}{Abstract}

Claude Code is an agentic work environment: a language model operating in a loop with filesystem access, shell execution, browser control, scheduled and cloud execution, external tool connections through the Model Context Protocol, and multi-agent orchestration [2], [3]. Its capability envelope now exceeds what a single practitioner can supervise by attention alone, and its failure modes are systemic rather than local: an unreviewed hook, an over-scoped connector, a stale completion condition, an autonomous routine inheriting every credential on an account.


This book is a task-oriented reference for operating that system safely and productively, written for practitioners who are accountable for the result. It advances four propositions. First, that capability without a defined and observable completion condition is not productivity. Second, that instruction, permission enforcement, sandboxing and operating-system isolation are four distinct layers of a control stack, only two of which are enforced, and that conflating them is the most common cause of loss of control. Third, that third-party skills, plugins, marketplaces, channels and MCP servers are software supply-chain dependencies and must be governed as such. Fourth, that the correct unit of trust in agentic work is observed evidence, not an agent's closing statement.


Thirty-four chapters take the reader from installation to a fully verified capstone, with a governance part addressing managed policy, data residency and retention, observability and accessibility. Every product claim carries a citation to a primary source, and the evidence ledger of Appendix J records the adjudication wherever a claim in circulation was found wrong, what the re-verification changed, and what remains unverified; controls are mapped to seventeen external frameworks in the crosswalk of Appendix H; and an organisational adoption maturity model is proposed in Appendix L. Claims that could not be confirmed from primary sources are labelled UNVERIFIED rather than softened.


\textbf{Keywords:} agentic artificial intelligence, AI-assisted software engineering, Claude Code, large language model agents, Model Context Protocol, prompt injection, zero trust, software supply-chain security, human oversight, AI governance, verification and validation.


\FloatBarrier
\section*{Contributions of this work}
\markboth{Contributions of this work}{Contributions of this work}

The product documentation this book cites is authoritative and extensive. This book exists because authoritative documentation answers \emph{what a feature does} and leaves four questions unanswered for anyone accountable for the outcome: what may be trusted, what must be enforced, what evidence would demonstrate correctness, and what an organisation must be able to show. Its original contributions are:


\begin{enumerate}
\item \textbf{The four-element brief} — Outcome, Context, Boundaries, Evidence — proposed in Chapter 1 as the invariant structure of an agentic instruction, and threaded through the plan contract (Chapter 6), the delegation brief (Chapter 16), the completion condition (Chapter 21) and the release gate (Chapter 34).
\item \textbf{The agentic control stack} (Figure 1) — a four-layer model separating advisory guidance from enforced control, which resolves the most persistent category error in current practice: the belief that an instruction in a project file constrains an agent.
\item \textbf{Evidence-Gated Delivery} (Chapter 34) — an eight-stage method with named gates for producing an agentic deliverable that can be released on evidence rather than assertion.
\item \textbf{A verification methodology for fast-decaying technical subjects} (“Scope and method”) — raw-source adjudication, a machine-checked citation and cross-reference apparatus, an explicit UNVERIFIED class, and a published re-verification procedure (Appendix K).
\item \textbf{A threat model and control map} (Appendix F) for agentic development environments, mapped to OWASP LLM Top 10 and MITRE ATLAS [4], [5].
\item \textbf{A standards and regulatory crosswalk} (Appendix H) mapping the book's controls to NIST AI RMF 1.0 and its generative profile, NIST SP 800-207, SP 800-218 and CSF 2.0, ISO/IEC 42001, 27001, 23894 and 25010, IEEE Std 1012-2016 and 7000-2021, the EU AI Act, and the ASD Essential Eight [6]–[18].
\item \textbf{An adoption maturity model} (Appendix L) — five levels of organisational capability, each with its enforced controls, required evidence and accountable owner.
\item \textbf{A consolidated correction of superseded guidance} (Part 0) — twenty-one widely repeated claims about Claude Code that are wrong, incomplete or inverted at the verification date, each with the primary source that settles it.
\end{enumerate}

\FloatBarrier
\section*{Executive summary}
\markboth{Executive summary}{Executive summary}

\emph{For readers who will decide whether and how an organisation adopts this class of tool, rather than operate it.}


\textbf{What the system is.} Claude Code is not a code-completion assistant. It is an agent that reads and writes files, executes shell commands, drives a browser using your authenticated sessions, connects to external services under your identity, publishes web pages, schedules its own future work, runs unattended in cloud infrastructure, and delegates to as many as a thousand copies of itself in a single run [2], [19]–[21].


\textbf{Where the risk concentrates.} Not in the model's competence, which is high and improving, but in four structural places. \emph{Delegation}, where a subagent created casually inherits the conversation it was meant to audit independently (Chapter 16). \emph{Untrusted input}, where a web page, a tool result, a pull-request comment or a peer message carries instructions to the agent (Chapter 20;Chapter 13) — the failure mode documented in the peer-reviewed literature as indirect prompt injection [22], [23]. \emph{Autonomy}, where a scheduled routine runs with every connected credential and no approval prompt (Chapter 24). And \emph{supply chain}, where a plugin may place executables on the shell path, start background processes and replace the system prompt of the main session on enable (Chapter 18).


\textbf{What must be enforced rather than requested.} An instruction in a project file is a preference; a permission rule is a control; a sandbox is an operating-system boundary; a managed policy is the only one of the four a user cannot unset (Figure 1). Organisations that write their standards into \texttt{CLAUDE.\allowbreak{}md} have documented an intention. Organisations that write them into managed settings have deployed a control (Chapter 31).


\textbf{What can be evidenced.} Permission decisions, mode changes, connector connections, plugin loads, authentication events and model refusals are all exportable through OpenTelemetry, and every event carries a prompt identifier that reconstructs the full activity of a single instruction (Chapter 33). An organisation that enables this can answer an auditor. One that does not cannot.


\textbf{What the reader should decide.} Six questions, in this order: which permission mode is the default for each cohort and who may change it; whether sandboxing is enforced or optional; which connectors, plugins and marketplaces are permitted; whether autonomous execution — routines, goals, unattended loops — is allowed at all and under whose approval; what leaves the boundary and under which retention terms (Chapter 32); and who is accountable when an agent acts. Appendix L sequences these into five maturity levels; Appendix E is the operational checklist that follows from them.


\textbf{The single-sentence version.} Adopt the capability at the pace at which you can evidence its output, and not faster.


\FloatBarrier
\section*{How to cite this work}
\markboth{How to cite this work}{How to cite this work}

\begin{calloutbox}{palegrey}{quoterule}{1.2mm}
D. Soldani, \emph{Claude Code Complete User Handbook: A Verified Practitioner's Reference for Agentic Work}, Version 1.0. Sydney, Australia, Aug. 2026. Verification baseline: Claude Code 2.1.241, 23 Aug. 2026; re-verified against 2.1.246, 26 Aug. 2026.
\end{calloutbox}

\begin{codeblock}{6.0}{7.1}
\cl{@book\{soldani2026claudecode,}
\cl{~~author~~~~=~\{Soldani,~David\},}
\cl{~~title~~~~~=~\{Claude~Code~Complete~User~Handbook:}
\cl{~~~~~~~~~~~~~~~A~Verified~Practitioner\textquotesingle{}s~Reference~for~Agentic~Work\},}
\cl{~~version~~~=~\{1.0\},}
\cl{~~year~~~~~~=~\{2026\},}
\cl{~~month~~~~~=~\{8\},}
\cl{~~address~~~=~\{Sydney,~Australia\},}
\cl{~~note~~~~~~=~\{Verification~baseline:~Claude~Code~2.1.241,~23~Aug.~2026;~re-verified~against~2.1.246,~26~Aug.~2026\}}
\cl{\}}
\end{codeblock}

\FloatBarrier
\section*{Scope and method}
\markboth{Scope and method}{Scope and method}

\textbf{Scope.} The book covers Claude Code as delivered through the terminal CLI, the VS Code and JetBrains extensions, the desktop application, the web and mobile surfaces, and Remote Control, at the verification baseline stated above. It covers the extension mechanisms — MCP, skills, subagents, hooks, plugins, output styles — the autonomy mechanisms — goals, workflows, three distinct schedulers, agent teams — and the governance mechanisms — managed policy, data retention, telemetry and accessibility. It does not cover the Agent SDK as a development platform beyond the points at which it changes the behaviour of an interactive session, nor cloud-provider deployment topologies beyond their effect on feature availability.


\textbf{Method.} Four rules were applied, and are stated here so that a reader can judge the evidence rather than trust the author.


\emph{Raw-source adjudication.} Every atomic product claim was adjudicated against the raw text of the primary documentation page that governs it, never against a summary of that page. The rule earned its keep twice. An automated summarisation pass reported three real product capabilities as non-existent; and, during the reference audit, a second summarisation pass asserted that an international secure-AI guideline had been issued by a single agency, when the issuing bodies' own announcements record it as a joint publication of two national agencies with seventeen further international partners. Both were false negatives, and either would have introduced an error had it been accepted. Where a summary and the source disagree, the source is authoritative and the summary is discarded.


\emph{Machine-checked apparatus.} The manuscript is authored with symbolic citation keys and symbolic cross-references, resolved to IEEE numbering and section numbers by a build step. The build fails on an unresolved key, a dangling cross-reference, a reference cited but not listed, a reference listed but not cited, or a cited documentation page absent from the official index. Two hundred and one automated checks gate the released artefacts. This is a modest application of the principle behind IEEE Std 1012-2016: verification is an activity with defined inputs, defined outputs and an independent result, not an assertion of diligence [15].


\emph{An explicit unverified class.} Claims that could be settled only against a running installation — and therefore could not be settled against documentation — are labelled UNVERIFIED and listed in Appendix J with the local check that closes them. They were not softened into plausible prose. The temptation to do so is precisely what produces confident, wrong technical writing.


\emph{Version stamping.} Behaviour that was introduced or changed at a specific release carries that release inline. Thirty-five release-specific behaviour notes appear in the pages this book cites, across a span of roughly ninety releases. A reference work on this subject has a half-life measured in weeks; the honest response is a stated procedure rather than a disclaimer, and that procedure is Appendix K.


\textbf{Independence.} This book is not published by, endorsed by, or affiliated with Anthropic PBC. "Claude", "Claude Code" and related marks are the property of their respective owners and are used descriptively. Quotations from official documentation are short, attributed, and used for technical criticism and review.


\textbf{Tools used in preparation.} This manuscript was drafted, critically reviewed, revised and typeset with the assistance of AI coding agents (Claude Code, Anthropic; Codex, OpenAI). Their use spans the writing of chapter prose from the author's outlines and source material; the critical review of that prose against the documentation it cites; the toolchain that resolves citations and cross-references and renders the four editions; the automated audit that gates every release; and the typographic revision of the typeset edition. Every claim of product behaviour is cited to the primary documentation that governs it, and the author has reviewed and accepted each one; the evidence ledger of Appendix J records the outcome. Generative AI tools are not authors of this work and are not listed as such. Responsibility for every statement, correction and error in it rests solely with the author, however the text was produced.


\FloatBarrier
\section*{Limitations}
\markboth{Limitations}{Limitations}

The method in the preceding section is stated so that a reader can judge the evidence. This section states what that evidence does not establish. It is placed in the front matter rather than in an appendix because a reader deciding how much weight to give the book should not have to reach page two hundred to find out.


\textbf{Documentation is not a running system.} The verification method adjudicates claims against primary documentation. Documentation describes intended behaviour; a running installation exhibits actual behaviour, and the two diverge — in defaults that differ by platform, in features gated by plan or organisation, and in the gap between a documented flag and its effect. Claims that could be settled only against a running installation are labelled UNVERIFIED and listed in Appendix J with the local check that closes each one. Three remain open at this release, all of them narrowing what the \texttt{/\allowbreak{}effort} selector offers on a surface or a plan that was not captured. A separate class could not be settled from publisher sources rather than from a running installation: the ISO/IEC 27001:2022 Annex A control identifiers cited in the crosswalk, the page range of one conference paper whose publisher returns 403 to automated clients, and the official expansion of the MITRE ATLAS acronym, for which two forms circulate in the maintainers' own materials. Both classes are recorded in Appendix J.


\textbf{Paywalled standards were verified structurally, not textually.} ISO and IEC standards are sold, not published. The ISO/IEC 42001 clause structure used in the crosswalk was corrected against the issuing body's own public catalogue record and clause listing, and the correction was material — seven crosswalk rows and one worked example carried wrong clause numbers before it. But the \textasciitilde{}25 Annex A control identifiers cited from ISO/IEC 27001:2022 were reachable only through secondary summaries, and a caveat to that effect is recorded in the crosswalk rather than implying they were checked. Any verification of ISO-family standards performed without purchased access has this limitation; stating it is more useful than concealing it.


\textbf{Link liveness is hygiene, not attestation.} The reference list holds 92 distinct URLs. A sweep from an unproxied network on 26 August 2026 found every one of the 96 unique URLs in the text reachable, one of them behind an authentication challenge. That sweep is not part of the default check suite, which still treats reachability as UNCHECKED, and reachability is in any case not a claim about content: a URL can resolve and no longer say what was adjudicated. A reader who finds a dead link should treat the citation's bibliographic data as the authoritative record.


\textbf{The false-negative finding cuts both ways.} Two automated summarisation passes were caught asserting things that the raw sources contradict, and both were false negatives — a capability or a fact reported as absent when it was present. False negatives are detectable by this method, because the raw source is consulted and the discrepancy surfaces. False \emph{positives} — a summariser asserting a capability that does not exist, in a case where the raw source was also consulted and happened to be ambiguous — would not surface the same way. No such case was found, and the method provides no strong assurance that none occurred.


\textbf{One author, one reviewer.} Verification was performed by the author. There was no independent adjudicator, no second reader reproducing the checks from the sources, and no editorial review of the kind a publisher supplies. Drafting and review were AI-assisted, as recorded in “Scope and method”; that is not independent adjudication, because the same class of tool both produced and checked the prose, and the author remains the only reviewer. The apparatus compensates in one direction only: the automated checks catch mechanical failures — a dangling cross-reference, an uncited reference, a stale count — and cannot catch a claim that is well-formed and wrong. Readers with access to a running installation are invited to close the UNVERIFIED items and to report discrepancies; the re-verification procedure in Appendix K is written to make that reproducible.


\textbf{The maturity model and the exercises are proposals, not findings.} The five-level adoption model in Appendix L, the four-element brief, the eight-point plan contract and the thirteen exercises with their worked solutions in Appendix M are the author's own constructions, offered as sequencing and teaching aids. They are informed by practice but have not been validated empirically against a population of organisations or learners, and no claim of external validity is made for them. Where the book cites a standard, the standard is the authority; where the book proposes a framework, the proposal is the author's and should be treated as such.


\textbf{Half-life.} The verification baseline is a single release. Thirty-five release-specific behaviour notes appear in the cited pages across roughly ninety releases, which is the measured rate of change of the subject rather than an impression of it. Statements of product behaviour in this book should be assumed to decay from the date on the title page, and Appendix K states the procedure and the cadence rather than offering a disclaimer in place of one.


\FloatBarrier
\section*{Disclaimer}
\markboth{Disclaimer}{Disclaimer}

This book describes the operation of a system that can modify files, execute commands, transact with external services and act autonomously under a user's own credentials. Procedures are provided for instruction. The author accepts no liability for data loss, disclosure of confidential information, financial loss, regulatory breach or service disruption arising from their use. Nothing here constitutes legal, financial, security-assurance or compliance advice; where a control is described in relation to a regulatory or standards framework, the description is informational and the reader remains responsible for their own obligations. Test every procedure in a disposable environment before applying it to systems that matter.


\FloatBarrier
\section*{Conventions used in this book}
\markboth{Conventions used in this book}{Conventions used in this book}

\subsection*{Evidence labels}

\begingroup
\def\tblrows{%
\textbf{VERIFIED} & Confirmed in Anthropic's official documentation, cited at the point of use, as of 23 August 2026. \\
\textbf{PRACTITIONER NOTE} & A technique from the author's practice or the wider community. Not a product guarantee. \\
\textbf{UNVERIFIED} & Plausible and in circulation, but not confirmable from primary sources at the verification date. Check locally before relying on it. \\
\textbf{CAUTION} & An action with material security, privacy, financial, regulatory or data-loss consequence. \\
\textbf{CHECKPOINT} & Evidence the reader should observe before continuing. \\
}%
\def\tblbody{\begin{minipage}{\textwidth}
{\footnotesize\begin{tabular}{L{24.6mm}L{118.4mm}}
\toprule
\textbf{Label} & \textbf{Meaning} \\
\midrule
\tblrows
\bottomrule\end{tabular}}\end{minipage}}%
\begingroup
\def\sloppy{\tolerance 9999\emergencystretch 3em\hfuzz 200pt\vfuzz 200pt}%
\hbadness=10000\vbadness=10000\hfuzz=200pt\vfuzz=200pt
\global\setbox\tblbox=\hbox{\tblbody}%
\endgroup
\par\addvspace{2.6mm}
\ifdim\dimexpr\ht\tblbox+\dp\tblbox\relax>0.55\textheight
  \typeout{HANDBOOK-TABLE broken \the\dimexpr\ht\tblbox+\dp\tblbox\relax}%
  
  {\footnotesize\begin{longtable}{L{24.6mm}L{118.4mm}}
  \toprule
\textbf{Label} & \textbf{Meaning} \\
\midrule\endfirsthead
  \multicolumn{2}{@{}l@{}}{%
  \sffamily\footnotesize\itshape\color{inkgrey}Continued}\\[1.2mm]
  \toprule
\textbf{Label} & \textbf{Meaning} \\
\midrule\endhead
  \bottomrule\endfoot
  \bottomrule\endlastfoot
  \tblrows
  \end{longtable}}%
\else
  \typeout{HANDBOOK-TABLE atomic \the\dimexpr\ht\tblbox+\dp\tblbox\relax}%
  \noindent\tblbody
\fi
\par\addvspace{2.6mm}
\endgroup

\subsection*{Citation, cross-reference and version conventions}

Citations are IEEE numeric: a bracketed number resolves to the reference list, numbered in order of first appearance. Cross-references name a chapter or a numbered section, for example Chapter 7 or Section 7.2. Figures and tables are numbered sequentially and listed after the contents.


Where a behaviour was introduced or changed at a specific release, the release is stated inline in the form \emph{(requires 2.1.219 or later)}. Where a behaviour is plan-dependent, surface-dependent or provider-dependent, that dependency is stated rather than assumed away.


Command syntax is shown as typed. A leading \texttt{/\allowbreak{}} marks an in-session command; a leading \texttt{claude} marks a shell invocation. The two are different interfaces and are indexed separately in Appendix A and Appendix B. Each chapter closes with \textbf{Key points} — the three to five statements a reader should retain.


\FloatBarrier
\section*{Release and re-verification}
\markboth{Release and re-verification}{Release and re-verification}

\begingroup
\def\tblrows{%
1.0 & 23 Aug. 2026 & Claude Code 2.1.241 & First release. 34 chapters, 14 appendices, 7 figures. \\
1.0, revised & 26 Aug. 2026 & Re-verified against 2.1.246 & Cited corpus re-fetched, hashed and reconciled against the published index; ten pages re-adjudicated; the effort and permission-mode vocabularies generated from one source. 201 automated build checks, all passing. \\
}%
\def\tblbody{\begin{minipage}{\textwidth}
{\footnotesize\begin{tabular}{L{16.9mm}L{10.2mm}L{22.6mm}L{86.9mm}}
\toprule
\textbf{Version} & \textbf{Date} & \textbf{Baseline} & \textbf{Status} \\
\midrule
\tblrows
\bottomrule\end{tabular}}\end{minipage}}%
\begingroup
\def\sloppy{\tolerance 9999\emergencystretch 3em\hfuzz 200pt\vfuzz 200pt}%
\hbadness=10000\vbadness=10000\hfuzz=200pt\vfuzz=200pt
\global\setbox\tblbox=\hbox{\tblbody}%
\endgroup
\par\addvspace{2.6mm}
\ifdim\dimexpr\ht\tblbox+\dp\tblbox\relax>0.55\textheight
  \typeout{HANDBOOK-TABLE broken \the\dimexpr\ht\tblbox+\dp\tblbox\relax}%
  
  {\footnotesize\begin{longtable}{L{16.9mm}L{10.2mm}L{22.6mm}L{86.9mm}}
  \toprule
\textbf{Version} & \textbf{Date} & \textbf{Baseline} & \textbf{Status} \\
\midrule\endfirsthead
  \multicolumn{4}{@{}l@{}}{%
  \sffamily\footnotesize\itshape\color{inkgrey}Continued}\\[1.2mm]
  \toprule
\textbf{Version} & \textbf{Date} & \textbf{Baseline} & \textbf{Status} \\
\midrule\endhead
  \bottomrule\endfoot
  \bottomrule\endlastfoot
  \tblrows
  \end{longtable}}%
\else
  \typeout{HANDBOOK-TABLE atomic \the\dimexpr\ht\tblbox+\dp\tblbox\relax}%
  \noindent\tblbody
\fi
\par\addvspace{2.6mm}
\endgroup

Re-verify quarterly, or before relying on the book for a consequential decision, following Appendix K. Record what changed in the ledger of Appendix J.


\FloatBarrier
\section*{Acknowledgements}
\markboth{Acknowledgements}{Acknowledgements}

Anthropic's documentation team maintains an unusually candid public record of limitations, deprecations and version-specific behaviour [1]; a book of this kind would not be possible without it. Responsibility for every judgement, correction and error rests with the author.


\FloatBarrier
\frontchapter{List of figures and tables}

\FloatBarrier
\section*{Figures}
\markboth{Figures}{Figures}

\begingroup
\def\tblrows{%
1 & The agentic control stack. Layers 1 and 2 are the product's; layers 3 and 4 are the operating system's… & Ch. 1 \\
2 & How a tool call is decided. Deny is evaluated first and cannot be widened from below… & Ch. 7 \\
3 & What survives each context operation. The column that matters is the last one: only files and version control survive everything. & Ch. 11 \\
4 & Where untrusted content enters an agentic session. Every arrow is a documented injection channel… & Ch. 12 \\
5 & Five delegation topologies and what each one inherits. Independence decreases from left to right along the top row. & Ch. 16 \\
6 & A cross-checked research workflow with a human release gate & Ch. 22 \\
7 & Choosing a scheduling mechanism. Start from what must be true when the work runs, not from how often it runs. & Ch. 24 \\
}%
\def\tblbody{\begin{minipage}{\textwidth}
{\footnotesize\begin{tabular}{L{9.7mm}L{110.8mm}L{19.4mm}}
\toprule
\textbf{Fig.} & \textbf{Caption} & \textbf{Location} \\
\midrule
\tblrows
\bottomrule\end{tabular}}\end{minipage}}%
\begingroup
\def\sloppy{\tolerance 9999\emergencystretch 3em\hfuzz 200pt\vfuzz 200pt}%
\hbadness=10000\vbadness=10000\hfuzz=200pt\vfuzz=200pt
\global\setbox\tblbox=\hbox{\tblbody}%
\endgroup
\par\addvspace{2.6mm}
\ifdim\dimexpr\ht\tblbox+\dp\tblbox\relax>0.55\textheight
  \typeout{HANDBOOK-TABLE broken \the\dimexpr\ht\tblbox+\dp\tblbox\relax}%
  
  {\footnotesize\begin{longtable}{L{9.7mm}L{110.8mm}L{19.4mm}}
  \toprule
\textbf{Fig.} & \textbf{Caption} & \textbf{Location} \\
\midrule\endfirsthead
  \multicolumn{3}{@{}l@{}}{%
  \sffamily\footnotesize\itshape\color{inkgrey}Continued}\\[1.2mm]
  \toprule
\textbf{Fig.} & \textbf{Caption} & \textbf{Location} \\
\midrule\endhead
  \bottomrule\endfoot
  \bottomrule\endlastfoot
  \tblrows
  \end{longtable}}%
\else
  \typeout{HANDBOOK-TABLE atomic \the\dimexpr\ht\tblbox+\dp\tblbox\relax}%
  \noindent\tblbody
\fi
\par\addvspace{2.6mm}
\endgroup

\FloatBarrier
\section*{Tables}
\markboth{Tables}{Tables}

\begingroup
\def\tblrows{%
1 & Superseded guidance and common misconceptions, with the verified position & front matter \\
2 & Claude Code surfaces and their trade-offs & Ch. 1 \\
3 & Effort levels accepted by surface, observed on Claude Code 2.1.246, 26 August 2026 & Ch. 4 \\
4 & Suggested starting posture by work type (practitioner guidance) & Ch. 4 \\
5 & Permission modes, the cycle position that reaches each, and appropriate use & Ch. 7 \\
6 & Instruction file locations, in load order from broadest to most specific & Ch. 9 \\
7 & Where a given piece of knowledge belongs & Ch. 10 \\
8 & Mechanisms that change the shape of a conversation & Ch. 11 \\
9 & Controls over transcript location, retention and writing & Ch. 12 \\
10 & MCP configuration scopes & Ch. 13 \\
11 & Skill frontmatter fields with a security consequence & Ch. 15 \\
12 & Subagent definition precedence & Ch. 16 \\
13 & Choosing a delegation mechanism & Ch. 16 \\
14 & Plugin components, all at the plugin root & Ch. 18 \\
15 & Built-in output styles & Ch. 19 \\
16 & Dynamic workflow runtime constraints & Ch. 22 \\
17 & What \texttt{/\allowbreak{}loop} does with each combination of arguments & Ch. 23 \\
18 & Comparing the three scheduling mechanisms & Ch. 24 \\
19 & Where channels sit among the ways to reach a session & Ch. 27 \\
20 & Artifact page constraints & Ch. 28 \\
21 & Ways to run more than one thing at a time & Ch. 29 \\
22 & How a managed policy reaches a machine & Ch. 31 \\
23 & Non-inference data flows and their opt-outs & Ch. 32 \\
24 & Retention by account type & Ch. 32 \\
25 & Exported metrics & Ch. 33 \\
26 & Book controls mapped to external frameworks & App. H \\
27 & Misconception adjudication, with verdict and primary source & App. J \\
28 & Corrections applied at the 26 August 2026 re-verification & App. J \\
29 & Items closed by direct observation on Claude Code 2.1.246, 26 August 2026 & App. J \\
30 & Adoption maturity levels for agentic development tooling & App. L \\
}%
\def\tblbody{\begin{minipage}{\textwidth}
{\footnotesize\begin{tabular}{L{12.1mm}L{108.4mm}L{19.4mm}}
\toprule
\textbf{Table} & \textbf{Caption} & \textbf{Location} \\
\midrule
\tblrows
\bottomrule\end{tabular}}\end{minipage}}%
\begingroup
\def\sloppy{\tolerance 9999\emergencystretch 3em\hfuzz 200pt\vfuzz 200pt}%
\hbadness=10000\vbadness=10000\hfuzz=200pt\vfuzz=200pt
\global\setbox\tblbox=\hbox{\tblbody}%
\endgroup
\par\addvspace{2.6mm}
\ifdim\dimexpr\ht\tblbox+\dp\tblbox\relax>0.55\textheight
  \typeout{HANDBOOK-TABLE broken \the\dimexpr\ht\tblbox+\dp\tblbox\relax}%
  
  {\footnotesize\begin{longtable}{L{12.1mm}L{108.4mm}L{19.4mm}}
  \toprule
\textbf{Table} & \textbf{Caption} & \textbf{Location} \\
\midrule\endfirsthead
  \multicolumn{3}{@{}l@{}}{%
  \sffamily\footnotesize\itshape\color{inkgrey}Continued}\\[1.2mm]
  \toprule
\textbf{Table} & \textbf{Caption} & \textbf{Location} \\
\midrule\endhead
  \bottomrule\endfoot
  \bottomrule\endlastfoot
  \tblrows
  \end{longtable}}%
\else
  \typeout{HANDBOOK-TABLE atomic \the\dimexpr\ht\tblbox+\dp\tblbox\relax}%
  \noindent\tblbody
\fi
\par\addvspace{2.6mm}
\endgroup

\breakrule

\FloatBarrier
\mainmatter
\parttitle{Part 0}{Superseded guidance and common misconceptions}
\addcontentsline{toc}{part}{Part 0 \textemdash\ Superseded guidance and common misconceptions}

Read this part before following any procedure in the book, and before following any procedure found elsewhere.


Claude Code changes weekly. Guidance written against a release three months old is not merely dated: in several documented cases it is inverted, and a reader who follows it does the opposite of what they intend. Twenty-one such claims are in wide circulation at the verification date.


This part is a register, not a treatment. Each row states the claim as it is commonly made, the position established from primary documentation, and the chapter that owns the correction and develops it in full. \textbf{The chapter is authoritative}; this table exists so that a reader who already holds one of these beliefs recognises the collision when they meet the correct statement. The four printed in bold are the ones whose consequences are operational rather than cosmetic.


Every entry appears again in the evidence ledger of Appendix J, with the verdict recording \emph{how} the claim failed and the primary source that settles it. Both tables are generated from one record, so the claim is written once.


\begingroup
\def\tblrows{%
\textbf{M01} & Run \texttt{/\allowbreak{}loop} with no arguments to see your scheduled tasks & A bare \texttt{/\allowbreak{}loop} starts an autonomous maintenance loop and lists nothing. To list tasks, ask Claude in natural language or use \texttt{Cron\allowbreak{}List}. [24] & Chapter 23 \\
\textbf{M02} & Closing the session stops a loop & Backgrounding carries loop tasks into a background session with no terminal attached, and resuming restores unexpired ones. Only a fresh conversation clears them. [24] & Chapter 23 \\
\textbf{M03} & A subagent gives you an independent second opinion & Only if you make it one. Fork mode is on by default in interactive sessions and a fork inherits the entire parent conversation; a parent in \texttt{accept\allowbreak{}Edits} or \texttt{bypass\allowbreak{}Permissions} also overrides the subagent's own mode. [25] & Chapter 16 \\
\textbf{M04} & \texttt{.\allowbreak{}env} keeps credentials out of Claude's reach & False as a guarantee. Permission deny rules, sandbox credential controls, network restriction, output discipline and least-privilege credentials are the controls that hold. [26]–[28] & Chapter 14 \\
M05 & Scheduling jitter is unpredictable & It is specified and deterministic, and the offset derives from the task ID. Schedule at a minute that is neither \texttt{:\allowbreak{}00} nor \texttt{:\allowbreak{}30}. [24] & Chapter 23 \\
M06 & \texttt{CLAUDE\_\allowbreak{}CODE\_\allowbreak{}DISABLE\_\allowbreak{}AUTO\_\allowbreak{}MEMORY} is how you turn auto memory off & The \texttt{auto\allowbreak{}Memory\allowbreak{}Enabled} setting is the primary control, settable per project; the environment variable is the fallback. [26] & Chapter 10 \\
M07 & The 30-day cleanup period governs local Claude Code state & Session transcripts are swept, but auto-memory files are excluded from that sweep and persist until edited or deleted. [26], [29] & Chapter 10 \\
M08 & Session history is capped at one month & Thirty days is a configurable default, set by \texttt{cleanup\allowbreak{}Period\allowbreak{}Days}. [30] & Chapter 12 \\
M09 & \texttt{bypass\allowbreak{}Permissions} is a reasonable productivity setting & Isolated containers or virtual machines only. It propagates to workflow agents and agent-team teammates, who cannot decline it. [19], [29], [31] & Chapter 7 \\
M10 & Every MCP tool definition always consumes context & Tool Search defers many definitions until they are needed. Measure occupancy with \texttt{/\allowbreak{}context} rather than assuming it. [32], [33] & Chapter 13 \\
M11 & Copying \texttt{.\allowbreak{}mcp.\allowbreak{}json} transfers an authenticated server & Shared project configuration must be credential-free, and project-scoped servers additionally require per-user approval and workspace trust. [32] & Chapter 13 \\
M12 & A popular third-party skill is effectively safe & Skills execute inline shell, grant tools for the invoking turn and can register hooks that persist for the session. Stars are not a security review. [34] & Chapter 15 \\
M13 & \texttt{/\allowbreak{}agents} opens an agent editor & Changed at release 2.1.198: it now directs you to ask Claude to manage agents, or to edit the agent files directly. [25], [35] & Chapter 16 \\
M14 & Workflows are beta and cannot be monitored in an IDE & Dynamic workflows run in the CLI, the desktop app, the IDE extensions, \texttt{claude -\allowbreak{}p} and the Agent SDK, with \texttt{/\allowbreak{}workflows} as the progress view. [19] & Chapter 22 \\
M15 & Workflow agents inherit the session's permission mode & They always run in \texttt{accept\allowbreak{}Edits} and inherit your tool allowlist, whatever the session's mode. Your mode governs only the launch prompt. [19] & Chapter 22 \\
M16 & Claude in Chrome is a beta feature & Generally available. It is unavailable under WSL and to API-key sessions, and its access to authenticated browser state raises the stakes on permission hygiene. [21], [36] & Chapter 20 \\
M17 & \texttt{/\allowbreak{}loop} is equivalent to operating-system cron & It is session-scoped, idle-dependent and jittered, capped at 50 tasks per session, expiring after seven days, with no catch-up for missed fires. [24] & Chapter 23 \\
M18 & Plugins are convenient feature packages & A plugin may ship executables placed on the Bash tool's \texttt{PATH}, background monitors, hooks, MCP and LSP servers, and a settings key that activates one of its agents as the main thread. [37], [38] & Chapter 18 \\
M19 & Routines are simply sessions on a schedule & A research preview that also triggers on API calls and GitHub events, includes all connected connectors by default, acts under your identity, and has no approval prompts. [20] & Chapter 24 \\
M20 & A green status on a routine run means the task succeeded & It means the session started and exited without an infrastructure error. Blocked requests, missing connector tools and task failures appear only in the transcript. [20] & Chapter 24 \\
M21 & Remote Control keeps everything local & Execution is local, but while connected the transcript is stored on Anthropic servers to synchronise across devices — which is why it is disabled under Zero Data Retention. [39]–[41] & Chapter 25 \\
}%
\def\tblbody{\begin{minipage}{\textwidth}\boxcaption{Table 1 — Superseded guidance and common misconceptions, with the verified position}
{\footnotesize\begin{tabular}{L{6.1mm}L{63.5mm}L{50.0mm}L{16.9mm}}
\toprule
\textbf{} & \textbf{Claim as commonly made} & \textbf{The position at 23 August 2026} & \textbf{Treated in} \\
\midrule
\tblrows
\bottomrule\end{tabular}}\end{minipage}}%
\begingroup
\def\sloppy{\tolerance 9999\emergencystretch 3em\hfuzz 200pt\vfuzz 200pt}%
\hbadness=10000\vbadness=10000\hfuzz=200pt\vfuzz=200pt
\global\setbox\tblbox=\hbox{\tblbody}%
\endgroup
\par\addvspace{2.6mm}
\ifdim\dimexpr\ht\tblbox+\dp\tblbox\relax>0.55\textheight
  \typeout{HANDBOOK-TABLE broken \the\dimexpr\ht\tblbox+\dp\tblbox\relax}%
  \tabcaption{Table 1 — Superseded guidance and common misconceptions, with the verified position}
  {\footnotesize\begin{longtable}{L{6.1mm}L{63.5mm}L{50.0mm}L{16.9mm}}
  \toprule
\textbf{} & \textbf{Claim as commonly made} & \textbf{The position at 23 August 2026} & \textbf{Treated in} \\
\midrule\endfirsthead
  \multicolumn{4}{@{}l@{}}{%
  \sffamily\footnotesize\itshape\color{inkgrey}Table 1 — Superseded guidance and common misconceptions, with the verified position \textemdash\ continued}\\[1.2mm]
  \toprule
\textbf{} & \textbf{Claim as commonly made} & \textbf{The position at 23 August 2026} & \textbf{Treated in} \\
\midrule\endhead
  \bottomrule\endfoot
  \bottomrule\endlastfoot
  \tblrows
  \end{longtable}}%
\else
  \typeout{HANDBOOK-TABLE atomic \the\dimexpr\ht\tblbox+\dp\tblbox\relax}%
  \noindent\tblbody
\fi
\par\addvspace{2.6mm}
\endgroup

Several of these claims are traceable to a widely viewed introductory video course published on 10 July 2026 [42], which remains a reasonable orientation to the product's surface but was recorded against an earlier release. Where a reader has learned one of the positions above from that source or from material derived from it, the entry above supersedes it.


\begin{calloutbox}{palebrass}{brassdark}{2.0mm}
\textbf{CAUTION — non-negotiable human gates} Require explicit human confirmation before publication, payment, credential or account change, production deployment, destructive deletion, merge, or direct push to a protected branch. An agent reporting "done" is a claim, not evidence. This principle is operationalised as Evidence-Gated Delivery in Chapter 34 and as a checklist in Appendix E.
\end{calloutbox}

\breakrule

\FloatBarrier
\frontchapter{How to use this book}

The book has four reading paths.


The \textbf{essential path} is Chapters 1–12. It takes a reader from installation to a controlled working method: a project with bounded context, a plan reviewed before execution, permissions that hold, version control underneath, and durable project knowledge. Complete it before adding autonomous or third-party components. Nothing in Chapters 13 onward is safe without it.


The \textbf{extension path} is Chapters 13–19. It covers the mechanisms that change what Claude Code can reach: connectors, secrets, reusable procedures, delegated agents, event-driven automation, packaged plugins, and the system-prompt layer. Every chapter in this part is also a supply-chain chapter, and each carries a threat model.


The \textbf{autonomy path} is Chapters 20–30. Browser control, persistent goals, multi-agent workflows, three distinct scheduling mechanisms, remote and cloud execution, event push, shared artifacts, and team orchestration. The organising question in this part is not "can it do this" but "how will I know what it did".


The \textbf{governance path} is Chapters 31–33 and Appendices D, F and K. It is written for readers who must answer to somebody else: managed policy, data residency and retention, observability, audit, accessibility, and the re-verification procedure that keeps a document like this one honest.


Chapter 34 combines the whole system into a single verified deliverable. Appendix E is the release checklist that goes with it.


Use a disposable practice folder. Do not learn permissions, hooks, browser actions or scheduling in a production repository, and do not open a home directory as a project.


\FloatBarrier
\section*{The running project}
\markboth{The running project}{The running project}

Worked examples build a \textbf{Research Brief Workspace}. It is deliberately not a software project: the argument of this book is that agentic work is a general operating discipline, and a non-programming example makes the discipline visible without the distraction of a build system.


\begin{codeblock}{9.0}{10.6}
\cl{research-brief-workspace/}
\cl{\pmboxdrawuni{251C}\pmboxdrawuni{2500}\pmboxdrawuni{2500}~CLAUDE.md}
\cl{\pmboxdrawuni{251C}\pmboxdrawuni{2500}\pmboxdrawuni{2500}~README.md}
\cl{\pmboxdrawuni{251C}\pmboxdrawuni{2500}\pmboxdrawuni{2500}~sources/~~~~~~~~~~\#~immutable~source~material}
\cl{\pmboxdrawuni{251C}\pmboxdrawuni{2500}\pmboxdrawuni{2500}~notes/~~~~~~~~~~~~\#~working~extraction}
\cl{\pmboxdrawuni{251C}\pmboxdrawuni{2500}\pmboxdrawuni{2500}~drafts/~~~~~~~~~~~\#~unapproved~manuscripts}
\cl{\pmboxdrawuni{251C}\pmboxdrawuni{2500}\pmboxdrawuni{2500}~deliverables/~~~~~\#~human-approved~outputs~only}
\cl{\pmboxdrawuni{251C}\pmboxdrawuni{2500}\pmboxdrawuni{2500}~references/~~~~~~~\#~policies~the~agent~retrieves~on~demand}
\cl{\pmboxdrawuni{251C}\pmboxdrawuni{2500}\pmboxdrawuni{2500}~templates/~~~~~~~~\#~reusable~structures}
\cl{\pmboxdrawuni{2514}\pmboxdrawuni{2500}\pmboxdrawuni{2500}~.claude/}
\cl{~~~~\pmboxdrawuni{251C}\pmboxdrawuni{2500}\pmboxdrawuni{2500}~settings.json}
\cl{~~~~\pmboxdrawuni{251C}\pmboxdrawuni{2500}\pmboxdrawuni{2500}~rules/}
\cl{~~~~\pmboxdrawuni{251C}\pmboxdrawuni{2500}\pmboxdrawuni{2500}~skills/}
\cl{~~~~\pmboxdrawuni{2514}\pmboxdrawuni{2500}\pmboxdrawuni{2500}~agents/}
\end{codeblock}

By Chapter 34 the workspace will research a bounded question, keep an evidence ledger, draft a brief, review it with an agent that did not write it, render a deliverable, and stop at a human release gate that it cannot pass on its own.


\breakrule

\FloatBarrier
\parttitle{Part I}{Foundations}
\addcontentsline{toc}{part}{Part I \textemdash\ Foundations}

\FloatBarrier
\renewcommand{\chaptertitlelabel}{Part I \textperiodcentered\ Chapter 1}
\setcounter{chapter}{0}
\chapter{What Claude Code is, and what it is not}

Claude Code is an agentic work environment. A chat interface receives messages and returns messages. Claude Code additionally inspects a working directory, edits files, executes tools, requests permissions, reaches configured services, publishes pages, schedules its own future turns, and delegates to other instances of itself [2], [3].


That capability does not make it an employee, and it does not make it an expert. A productive mental model is this:


\begin{calloutbox}{palegrey}{quoterule}{1.2mm}
Claude Code is a fast operator with variable judgement, broad tool access, no inherent knowledge of your unstated intent, and a strong disposition to report success.
\end{calloutbox}

The last clause is the operationally important one. Every control described in this book exists because the agent's own account of what it did is the least reliable artefact it produces.


\FloatBarrier
\setcounter{section}{0}
\section{Four things that must be explicit}

Work goes well when four elements are stated rather than inferred:


\begin{enumerate}
\item \textbf{Outcome} — the finished state you need, in terms someone else could check.
\item \textbf{Context} — files, constraints, prior decisions, and what is out of scope.
\item \textbf{Boundaries} — what may change without asking, and what requires approval.
\item \textbf{Evidence} — how both parties will know the outcome is correct.
\end{enumerate}

Chapters 5 and 6 turn these into a brief and a plan; Chapter 21 turns the fourth into a machine-checked completion condition; Appendix E turns all four into a release checklist.


\FloatBarrier
\setcounter{section}{1}
\section{The control stack}

Most loss of control in agentic work is a category error: treating a layer that \emph{shapes intent} as though it \emph{constrains action}. The four layers of Figure 1 are not interchangeable, and only two of them are enforced.


\begin{figure}[tbp]
\begin{codefig}{9.0}{10.6}
\cl{~~MANAGED~POLICY~\pmboxdrawuni{2500}~deployed~by~an~organisation;~the~only~tier~a~user~cannot~unset}
\cl{~~\pmboxdrawuni{2550}\pmboxdrawuni{2550}\pmboxdrawuni{2550}\pmboxdrawuni{2550}\pmboxdrawuni{2550}\pmboxdrawuni{2550}\pmboxdrawuni{2550}\pmboxdrawuni{2550}\pmboxdrawuni{2550}\pmboxdrawuni{2550}\pmboxdrawuni{2550}\pmboxdrawuni{2550}\pmboxdrawuni{2550}\pmboxdrawuni{2550}\pmboxdrawuni{2550}\pmboxdrawuni{2550}\pmboxdrawuni{2550}\pmboxdrawuni{2550}\pmboxdrawuni{2550}\pmboxdrawuni{2550}\pmboxdrawuni{2550}\pmboxdrawuni{2550}\pmboxdrawuni{2550}\pmboxdrawuni{2550}\pmboxdrawuni{2550}\pmboxdrawuni{2550}\pmboxdrawuni{2550}\pmboxdrawuni{2550}\pmboxdrawuni{2550}\pmboxdrawuni{2550}\pmboxdrawuni{2550}\pmboxdrawuni{2550}\pmboxdrawuni{2550}\pmboxdrawuni{2550}\pmboxdrawuni{2550}\pmboxdrawuni{2550}\pmboxdrawuni{2550}\pmboxdrawuni{2550}\pmboxdrawuni{2550}\pmboxdrawuni{2550}\pmboxdrawuni{2550}\pmboxdrawuni{2550}\pmboxdrawuni{2550}\pmboxdrawuni{2550}\pmboxdrawuni{2550}\pmboxdrawuni{2550}\pmboxdrawuni{2550}\pmboxdrawuni{2550}\pmboxdrawuni{2550}\pmboxdrawuni{2550}\pmboxdrawuni{2550}\pmboxdrawuni{2550}\pmboxdrawuni{2550}\pmboxdrawuni{2550}\pmboxdrawuni{2550}\pmboxdrawuni{2550}\pmboxdrawuni{2550}\pmboxdrawuni{2550}\pmboxdrawuni{2550}\pmboxdrawuni{2550}\pmboxdrawuni{2550}\pmboxdrawuni{2550}\pmboxdrawuni{2550}\pmboxdrawuni{2550}\pmboxdrawuni{2550}\pmboxdrawuni{2550}\pmboxdrawuni{2550}\pmboxdrawuni{2550}\pmboxdrawuni{2550}\pmboxdrawuni{2550}\pmboxdrawuni{2550}\pmboxdrawuni{2550}\pmboxdrawuni{2550}\pmboxdrawuni{2550}\pmboxdrawuni{2550}\pmboxdrawuni{2550}\pmboxdrawuni{2550}\pmboxdrawuni{2550}\pmboxdrawuni{2550}}
\cl{}
\cl{~~~ADVISORY~~~\pmboxdrawuni{250C}\pmboxdrawuni{2500}\pmboxdrawuni{2500}\pmboxdrawuni{2500}\pmboxdrawuni{2500}\pmboxdrawuni{2500}\pmboxdrawuni{2500}\pmboxdrawuni{2500}\pmboxdrawuni{2500}\pmboxdrawuni{2500}\pmboxdrawuni{2500}\pmboxdrawuni{2500}\pmboxdrawuni{2500}\pmboxdrawuni{2500}\pmboxdrawuni{2500}\pmboxdrawuni{2500}\pmboxdrawuni{2500}\pmboxdrawuni{2500}\pmboxdrawuni{2500}\pmboxdrawuni{2500}\pmboxdrawuni{2500}\pmboxdrawuni{2500}\pmboxdrawuni{2500}\pmboxdrawuni{2500}\pmboxdrawuni{2500}\pmboxdrawuni{2500}\pmboxdrawuni{2500}\pmboxdrawuni{2500}\pmboxdrawuni{2500}\pmboxdrawuni{2500}\pmboxdrawuni{2500}\pmboxdrawuni{2500}\pmboxdrawuni{2500}\pmboxdrawuni{2500}\pmboxdrawuni{2500}\pmboxdrawuni{2500}\pmboxdrawuni{2500}\pmboxdrawuni{2500}\pmboxdrawuni{2500}\pmboxdrawuni{2500}\pmboxdrawuni{2500}\pmboxdrawuni{2500}\pmboxdrawuni{2500}\pmboxdrawuni{2500}\pmboxdrawuni{2500}\pmboxdrawuni{2500}\pmboxdrawuni{2500}\pmboxdrawuni{2500}\pmboxdrawuni{2500}\pmboxdrawuni{2500}\pmboxdrawuni{2500}\pmboxdrawuni{2500}\pmboxdrawuni{2500}\pmboxdrawuni{2500}\pmboxdrawuni{2500}\pmboxdrawuni{2500}\pmboxdrawuni{2500}\pmboxdrawuni{2500}\pmboxdrawuni{2500}\pmboxdrawuni{2500}\pmboxdrawuni{2500}\pmboxdrawuni{2500}\pmboxdrawuni{2500}\pmboxdrawuni{2510}}
\cl{~~~shapes~~~~~\pmboxdrawuni{2502}~1~~INSTRUCTION~~~CLAUDE.md~·~rules~·~skills~·~the~prompt~~~~~\pmboxdrawuni{2502}}
\cl{~~~intent~~~~~\pmboxdrawuni{2502}~~~~Read~as~context~and~generally~followed.~No~guarantee,~~~~~\pmboxdrawuni{2502}}
\cl{~~~~~~~~~~~~~~\pmboxdrawuni{2502}~~~~especially~where~two~instructions~conflict.~~~~~~~~~~~~~~~\pmboxdrawuni{2502}}
\cl{~~~~~~~~~~~~~~\pmboxdrawuni{2514}\pmboxdrawuni{2500}\pmboxdrawuni{2500}\pmboxdrawuni{2500}\pmboxdrawuni{2500}\pmboxdrawuni{2500}\pmboxdrawuni{2500}\pmboxdrawuni{2500}\pmboxdrawuni{2500}\pmboxdrawuni{2500}\pmboxdrawuni{2500}\pmboxdrawuni{2500}\pmboxdrawuni{2500}\pmboxdrawuni{2500}\pmboxdrawuni{2500}\pmboxdrawuni{2500}\pmboxdrawuni{2500}\pmboxdrawuni{2500}\pmboxdrawuni{2500}\pmboxdrawuni{2500}\pmboxdrawuni{2500}\pmboxdrawuni{2500}\pmboxdrawuni{2500}\pmboxdrawuni{2500}\pmboxdrawuni{2500}\pmboxdrawuni{2500}\pmboxdrawuni{2500}\pmboxdrawuni{2500}\pmboxdrawuni{2500}\pmboxdrawuni{2500}\pmboxdrawuni{2500}\pmboxdrawuni{2500}\pmboxdrawuni{2500}\pmboxdrawuni{2500}\pmboxdrawuni{2500}\pmboxdrawuni{2500}\pmboxdrawuni{2500}\pmboxdrawuni{2500}\pmboxdrawuni{2500}\pmboxdrawuni{2500}\pmboxdrawuni{2500}\pmboxdrawuni{2500}\pmboxdrawuni{2500}\pmboxdrawuni{2500}\pmboxdrawuni{2500}\pmboxdrawuni{2500}\pmboxdrawuni{2500}\pmboxdrawuni{2500}\pmboxdrawuni{2500}\pmboxdrawuni{2500}\pmboxdrawuni{2500}\pmboxdrawuni{2500}\pmboxdrawuni{2500}\pmboxdrawuni{2500}\pmboxdrawuni{2500}\pmboxdrawuni{2500}\pmboxdrawuni{2500}\pmboxdrawuni{2500}\pmboxdrawuni{2500}\pmboxdrawuni{2500}\pmboxdrawuni{2500}\pmboxdrawuni{2500}\pmboxdrawuni{2500}\pmboxdrawuni{2518}}
\cl{~~~~~~~~~~~~~~~~~~~~~~~~~~~~~~~~~~~~~~~\pmboxdrawuni{2502}}
\cl{~~~ENFORCED~~~\pmboxdrawuni{250C}\pmboxdrawuni{2500}\pmboxdrawuni{2500}\pmboxdrawuni{2500}\pmboxdrawuni{2500}\pmboxdrawuni{2500}\pmboxdrawuni{2500}\pmboxdrawuni{2500}\pmboxdrawuni{2500}\pmboxdrawuni{2500}\pmboxdrawuni{2500}\pmboxdrawuni{2500}\pmboxdrawuni{2500}\pmboxdrawuni{2500}\pmboxdrawuni{2500}\pmboxdrawuni{2500}\pmboxdrawuni{2500}\pmboxdrawuni{2500}\pmboxdrawuni{2500}\pmboxdrawuni{2500}\pmboxdrawuni{2500}\pmboxdrawuni{2500}\pmboxdrawuni{2500}\pmboxdrawuni{2500}\pmboxdrawuni{2500}\pmboxdrawuni{2500}\pmboxdrawuni{2500}\pmboxdrawuni{2500}\pmboxdrawuni{2500}\pmboxdrawuni{2500}\pmboxdrawuni{2500}\pmboxdrawuni{2500}\pmboxdrawuni{2500}\pmboxdrawuni{2500}\pmboxdrawuni{2500}\pmboxdrawuni{2500}\pmboxdrawuni{2500}\pmboxdrawuni{2500}\pmboxdrawuni{2500}\pmboxdrawuni{2500}\pmboxdrawuni{2500}\pmboxdrawuni{2500}\pmboxdrawuni{2500}\pmboxdrawuni{2500}\pmboxdrawuni{2500}\pmboxdrawuni{2500}\pmboxdrawuni{2500}\pmboxdrawuni{2500}\pmboxdrawuni{2500}\pmboxdrawuni{2500}\pmboxdrawuni{2500}\pmboxdrawuni{2500}\pmboxdrawuni{2500}\pmboxdrawuni{2500}\pmboxdrawuni{2500}\pmboxdrawuni{2500}\pmboxdrawuni{2500}\pmboxdrawuni{2500}\pmboxdrawuni{2500}\pmboxdrawuni{2500}\pmboxdrawuni{2500}\pmboxdrawuni{2500}\pmboxdrawuni{2500}\pmboxdrawuni{2510}}
\cl{~~~blocks~~~~~\pmboxdrawuni{2502}~2~~PERMISSION~~~~deny~→~ask~→~allow~·~permission~modes~~~~~~~\pmboxdrawuni{2502}}
\cl{~~~action~~~~~\pmboxdrawuni{2502}~~~~Product-enforced.~A~deny~cannot~be~widened~from~below.~~~~\pmboxdrawuni{2502}}
\cl{~~~~~~~~~~~~~~\pmboxdrawuni{251C}\pmboxdrawuni{2500}\pmboxdrawuni{2500}\pmboxdrawuni{2500}\pmboxdrawuni{2500}\pmboxdrawuni{2500}\pmboxdrawuni{2500}\pmboxdrawuni{2500}\pmboxdrawuni{2500}\pmboxdrawuni{2500}\pmboxdrawuni{2500}\pmboxdrawuni{2500}\pmboxdrawuni{2500}\pmboxdrawuni{2500}\pmboxdrawuni{2500}\pmboxdrawuni{2500}\pmboxdrawuni{2500}\pmboxdrawuni{2500}\pmboxdrawuni{2500}\pmboxdrawuni{2500}\pmboxdrawuni{2500}\pmboxdrawuni{2500}\pmboxdrawuni{2500}\pmboxdrawuni{2500}\pmboxdrawuni{2500}\pmboxdrawuni{2500}\pmboxdrawuni{2500}\pmboxdrawuni{2500}\pmboxdrawuni{2500}\pmboxdrawuni{2500}\pmboxdrawuni{2500}\pmboxdrawuni{2500}\pmboxdrawuni{2500}\pmboxdrawuni{2500}\pmboxdrawuni{2500}\pmboxdrawuni{2500}\pmboxdrawuni{2500}\pmboxdrawuni{2500}\pmboxdrawuni{2500}\pmboxdrawuni{2500}\pmboxdrawuni{2500}\pmboxdrawuni{2500}\pmboxdrawuni{2500}\pmboxdrawuni{2500}\pmboxdrawuni{2500}\pmboxdrawuni{2500}\pmboxdrawuni{2500}\pmboxdrawuni{2500}\pmboxdrawuni{2500}\pmboxdrawuni{2500}\pmboxdrawuni{2500}\pmboxdrawuni{2500}\pmboxdrawuni{2500}\pmboxdrawuni{2500}\pmboxdrawuni{2500}\pmboxdrawuni{2500}\pmboxdrawuni{2500}\pmboxdrawuni{2500}\pmboxdrawuni{2500}\pmboxdrawuni{2500}\pmboxdrawuni{2500}\pmboxdrawuni{2500}\pmboxdrawuni{2500}\pmboxdrawuni{2524}}
\cl{~~~~~~~~~~~~~~\pmboxdrawuni{2502}~3~~SANDBOX~~~~~~~filesystem~and~network,~per~command~~~~~~~~~\pmboxdrawuni{2502}}
\cl{~~~~~~~~~~~~~~\pmboxdrawuni{2502}~~~~OS-enforced:~Seatbelt~on~macOS,~Linux~and~WSL2.~~~~~~~~~~~\pmboxdrawuni{2502}}
\cl{~~~~~~~~~~~~~~\pmboxdrawuni{2502}~~~~Not~available~on~native~Windows.~~~~~~~~~~~~~~~~~~~~~~~~~~\pmboxdrawuni{2502}}
\cl{~~~~~~~~~~~~~~\pmboxdrawuni{251C}\pmboxdrawuni{2500}\pmboxdrawuni{2500}\pmboxdrawuni{2500}\pmboxdrawuni{2500}\pmboxdrawuni{2500}\pmboxdrawuni{2500}\pmboxdrawuni{2500}\pmboxdrawuni{2500}\pmboxdrawuni{2500}\pmboxdrawuni{2500}\pmboxdrawuni{2500}\pmboxdrawuni{2500}\pmboxdrawuni{2500}\pmboxdrawuni{2500}\pmboxdrawuni{2500}\pmboxdrawuni{2500}\pmboxdrawuni{2500}\pmboxdrawuni{2500}\pmboxdrawuni{2500}\pmboxdrawuni{2500}\pmboxdrawuni{2500}\pmboxdrawuni{2500}\pmboxdrawuni{2500}\pmboxdrawuni{2500}\pmboxdrawuni{2500}\pmboxdrawuni{2500}\pmboxdrawuni{2500}\pmboxdrawuni{2500}\pmboxdrawuni{2500}\pmboxdrawuni{2500}\pmboxdrawuni{2500}\pmboxdrawuni{2500}\pmboxdrawuni{2500}\pmboxdrawuni{2500}\pmboxdrawuni{2500}\pmboxdrawuni{2500}\pmboxdrawuni{2500}\pmboxdrawuni{2500}\pmboxdrawuni{2500}\pmboxdrawuni{2500}\pmboxdrawuni{2500}\pmboxdrawuni{2500}\pmboxdrawuni{2500}\pmboxdrawuni{2500}\pmboxdrawuni{2500}\pmboxdrawuni{2500}\pmboxdrawuni{2500}\pmboxdrawuni{2500}\pmboxdrawuni{2500}\pmboxdrawuni{2500}\pmboxdrawuni{2500}\pmboxdrawuni{2500}\pmboxdrawuni{2500}\pmboxdrawuni{2500}\pmboxdrawuni{2500}\pmboxdrawuni{2500}\pmboxdrawuni{2500}\pmboxdrawuni{2500}\pmboxdrawuni{2500}\pmboxdrawuni{2500}\pmboxdrawuni{2500}\pmboxdrawuni{2500}\pmboxdrawuni{2524}}
\cl{~~~~~~~~~~~~~~\pmboxdrawuni{2502}~4~~ISOLATION~~~~~container~·~virtual~machine~·~dev~container~\pmboxdrawuni{2502}}
\cl{~~~~~~~~~~~~~~\pmboxdrawuni{2502}~~~~The~only~layer~that~contains~a~failure~of~layers~2~or~3.~~\pmboxdrawuni{2502}}
\cl{~~~~~~~~~~~~~~\pmboxdrawuni{2514}\pmboxdrawuni{2500}\pmboxdrawuni{2500}\pmboxdrawuni{2500}\pmboxdrawuni{2500}\pmboxdrawuni{2500}\pmboxdrawuni{2500}\pmboxdrawuni{2500}\pmboxdrawuni{2500}\pmboxdrawuni{2500}\pmboxdrawuni{2500}\pmboxdrawuni{2500}\pmboxdrawuni{2500}\pmboxdrawuni{2500}\pmboxdrawuni{2500}\pmboxdrawuni{2500}\pmboxdrawuni{2500}\pmboxdrawuni{2500}\pmboxdrawuni{2500}\pmboxdrawuni{2500}\pmboxdrawuni{2500}\pmboxdrawuni{2500}\pmboxdrawuni{2500}\pmboxdrawuni{2500}\pmboxdrawuni{2500}\pmboxdrawuni{2500}\pmboxdrawuni{2500}\pmboxdrawuni{2500}\pmboxdrawuni{2500}\pmboxdrawuni{2500}\pmboxdrawuni{2500}\pmboxdrawuni{2500}\pmboxdrawuni{2500}\pmboxdrawuni{2500}\pmboxdrawuni{2500}\pmboxdrawuni{2500}\pmboxdrawuni{2500}\pmboxdrawuni{2500}\pmboxdrawuni{2500}\pmboxdrawuni{2500}\pmboxdrawuni{2500}\pmboxdrawuni{2500}\pmboxdrawuni{2500}\pmboxdrawuni{2500}\pmboxdrawuni{2500}\pmboxdrawuni{2500}\pmboxdrawuni{2500}\pmboxdrawuni{2500}\pmboxdrawuni{2500}\pmboxdrawuni{2500}\pmboxdrawuni{2500}\pmboxdrawuni{2500}\pmboxdrawuni{2500}\pmboxdrawuni{2500}\pmboxdrawuni{2500}\pmboxdrawuni{2500}\pmboxdrawuni{2500}\pmboxdrawuni{2500}\pmboxdrawuni{2500}\pmboxdrawuni{2500}\pmboxdrawuni{2500}\pmboxdrawuni{2500}\pmboxdrawuni{2500}\pmboxdrawuni{2518}}
\cl{}
\cl{~~~HOOKS~cut~across~layers~1~and~2~in~both~directions:~a~PreToolUse~hook~exiting~2}
\cl{~~~blocks~a~call~before~permission~rules~are~evaluated,~and~no~hook~decision~can}
\cl{~~~bypass~a~matching~deny~rule.}
\end{codefig}
\figcaption{Figure 1 — The agentic control stack. Layers 1 and 2 are the product's; layers 3 and 4 are the operating system's. Only layers 2 to 4 constrain what can happen.}
\end{figure}

An instruction that must hold belongs at layer 2 or below. An instruction that merely should hold belongs at layer 1, where it costs context and buys a tendency. Chapters 7, 17 and 31 develop the enforced layers; Chapters 9 and 10 develop the advisory one.



\FloatBarrier
\setcounter{section}{2}
\section{Where the work runs}

The same engine is reachable through several surfaces, and your repository's \texttt{CLAUDE.\allowbreak{}md}, settings and MCP servers apply across all of them [2], [43].


\begingroup
\def\tblrows{%
Terminal CLI & Earliest feature access, scripting, remote hosts, full command surface & Requires comfort with a terminal \\
VS Code extension & Inline diffs, \texttt{@}-mentions, plan review, conversation history beside files [44] & Some controls are terminal-only \\
JetBrains plugin & IntelliJ, PyCharm, WebStorm and related IDEs [45] & Requires the CLI installed separately \\
Desktop app & Parallel sessions with git isolation, visual diff review, scheduled local tasks, computer use [46] & A different interaction model from the terminal \\
Web (claude.ai/code) & Cloud sandboxes, work on repositories you do not have locally, long tasks you check later [47] & No access to local files; research preview \\
Mobile & Starting, monitoring and steering tasks from a phone [48] & Small control surface; poor for reviewing diffs \\
Remote Control & Supervising a session that is executing on your own machine [39] & Transcript is relayed and stored server-side while connected \\
}%
\def\tblbody{\begin{minipage}{\textwidth}\boxcaption{Table 2 — Claude Code surfaces and their trade-offs}
{\footnotesize\begin{tabular}{L{28.7mm}L{67.8mm}L{43.3mm}}
\toprule
\textbf{Surface} & \textbf{Best for} & \textbf{Trade-off} \\
\midrule
\tblrows
\bottomrule\end{tabular}}\end{minipage}}%
\begingroup
\def\sloppy{\tolerance 9999\emergencystretch 3em\hfuzz 200pt\vfuzz 200pt}%
\hbadness=10000\vbadness=10000\hfuzz=200pt\vfuzz=200pt
\global\setbox\tblbox=\hbox{\tblbody}%
\endgroup
\par\addvspace{2.6mm}
\ifdim\dimexpr\ht\tblbox+\dp\tblbox\relax>0.55\textheight
  \typeout{HANDBOOK-TABLE broken \the\dimexpr\ht\tblbox+\dp\tblbox\relax}%
  \tabcaption{Table 2 — Claude Code surfaces and their trade-offs}
  {\footnotesize\begin{longtable}{L{28.7mm}L{67.8mm}L{43.3mm}}
  \toprule
\textbf{Surface} & \textbf{Best for} & \textbf{Trade-off} \\
\midrule\endfirsthead
  \multicolumn{3}{@{}l@{}}{%
  \sffamily\footnotesize\itshape\color{inkgrey}Table 2 — Claude Code surfaces and their trade-offs \textemdash\ continued}\\[1.2mm]
  \toprule
\textbf{Surface} & \textbf{Best for} & \textbf{Trade-off} \\
\midrule\endhead
  \bottomrule\endfoot
  \bottomrule\endlastfoot
  \tblrows
  \end{longtable}}%
\else
  \typeout{HANDBOOK-TABLE atomic \the\dimexpr\ht\tblbox+\dp\tblbox\relax}%
  \noindent\tblbody
\fi
\par\addvspace{2.6mm}
\endgroup

\textbf{PRACTITIONER NOTE:} the source course prefers VS Code because it puts an explorer, an editor, a terminal and several Claude tabs on one screen [42]. That is a reasonable default for a beginner. It is not a requirement, and several controls in Part V are easier to observe in the CLI.


\FloatBarrier
\setcounter{section}{3}
\section{Appropriate first tasks}

Choose work that is reversible and cheap to inspect:


\begin{itemize}
\item organise copies of documents;
\item produce a report outline from supplied notes;
\item compare two files and return a difference table;
\item research a bounded question and save a source ledger;
\item explain an unfamiliar repository without changing it.
\end{itemize}

Avoid, as a first task, anything that touches an account you cannot afford to have altered. Capability should expand only after verification habits do.


\subsection*{Exercise}

Write a four-line brief and keep it. You will run it in Chapter 5.


\begin{codeblock}{8.0}{9.4}
\cl{Outcome:~~~~A~one-page~brief~built~from~the~files~in~notes/.}
\cl{Context:~~~~Use~only~local~files.~Identify~evidence~gaps~rather~than~filling~them.}
\cl{Boundaries:~Do~not~edit~or~delete~anything~in~notes/~or~sources/.~No~network~access.}
\cl{Evidence:~~~Save~to~drafts/brief-v1.md~with~a~claim-to-source~table.}
\end{codeblock}

\emph{A worked solution is given in Appendix M.}


\begingroup
\def\kprows{%
\item Claude Code is an agent, not an assistant: it edits files, runs commands, reaches accounts and schedules its own work.
\item Four elements must be explicit in every instruction — Outcome, Context, Boundaries, Evidence. Omitting the fourth is what turns capability into rework.
\item The control stack has four layers and only the lower three are enforced. An instruction is a preference; a permission rule is a control.
\item Choose first tasks that are reversible and cheap to inspect. Capability should expand only after verification habits do.
}%
\def\kpbody{\begin{minipage}{\textwidth}\subsection*{Key points}\begin{itemize}\kprows\end{itemize}\end{minipage}}%
\begingroup
\def\sloppy{\tolerance 9999\emergencystretch 3em\hfuzz 200pt\vfuzz 200pt}%
\hbadness=10000\vbadness=10000\hfuzz=200pt\vfuzz=200pt
\global\setbox\kpbox=\hbox{\kpbody}%
\endgroup
\par\addvspace{4.2mm}
\ifdim\dimexpr\ht\kpbox+\dp\kpbox\relax>0.30\textheight
  \typeout{HANDBOOK-KEYPOINTS broken \the\dimexpr\ht\kpbox+\dp\kpbox\relax}%
  \subsection*{Key points}
  \begin{itemize}\kprows\end{itemize}
\else
  \typeout{HANDBOOK-KEYPOINTS atomic \the\dimexpr\ht\kpbox+\dp\kpbox\relax}%
  \noindent\kpbody
\fi
\par\addvspace{1.4mm}
\endgroup

\FloatBarrier
\renewcommand{\chaptertitlelabel}{Part I \textperiodcentered\ Chapter 2}
\setcounter{chapter}{1}
\chapter{Install, verify and authenticate}

\FloatBarrier
\setcounter{section}{0}
\section{Requirements}

As verified on 23 August 2026, Claude Code supports macOS 13 and later, Windows 10 1809 or Windows Server 2019 and later, Ubuntu 20.04+, Debian 10+ and Alpine 3.19+, with at least 4 GB of RAM, network access, and Bash, Zsh, PowerShell or CMD. A Pro, Max, Team, Enterprise or Console account is required, or a configured cloud provider [2], [49].


Not every feature follows the platform matrix. Sandboxing does not run on native Windows [27]; Chrome integration does not run under WSL [21]; artifacts, cloud sessions, routines and channels are unavailable on third-party model providers [20], [47], [50], [51]. Appendix D lists the dependencies that most often surprise people.


\FloatBarrier
\setcounter{section}{1}
\section{Installation}

The native installer is the official default and updates itself in the background [2], [49].


\begin{codeblock}{9.0}{10.6}
\cl{\#~macOS,~Linux,~WSL}
\cl{curl~-fsSL~https://claude.ai/install.sh~|~bash}
\end{codeblock}

\begin{codeblock}{9.0}{10.6}
\cl{\#~Windows~PowerShell}
\cl{irm~https://claude.ai/install.ps1~|~iex}
\end{codeblock}

\begin{codeblock}{8.0}{9.4}
\cl{REM~Windows~CMD}
\cl{curl~-fsSL~https://claude.ai/install.cmd~-o~install.cmd~\&\&~install.cmd~\&\&~del~install.cmd}
\end{codeblock}

Package-manager installations do \textbf{not} auto-update and must be upgraded explicitly [2]:


\begin{codeblock}{8.0}{9.4}
\cl{brew~install~--cask~claude-code~~~~~~~~\#~stable~channel;~claude-code@latest~tracks~latest}
\cl{winget~install~Anthropic.ClaudeCode~~~~\#~upgrade~with:~winget~upgrade~Anthropic.ClaudeCode}
\end{codeblock}

Debian, Fedora, RHEL and Alpine packages are also available through \texttt{apt}, \texttt{dnf} and \texttt{apk} [2].


\begin{calloutbox}{palebrass}{brassdark}{2.0mm}
\textbf{CAUTION:} a shell installer executes code on your machine with your privileges. Never substitute a command copied from a video description, an advertisement, a forum post or an unverified repository for the documented one. Verify the domain character by character.
\end{calloutbox}

\FloatBarrier
\setcounter{section}{2}
\section{Verify before you log in}

\begin{codeblock}{9.0}{10.6}
\cl{claude~--version}
\cl{claude~doctor}
\end{codeblock}

\texttt{claude doct\allowbreak{}or} runs a read-only installation and configuration checkup; from 2.1.206 it also proposes trims for a checked-in \texttt{CLAUDE.\allowbreak{}md} [26]. Then start the product and complete browser authentication:


\begin{codeblock}{9.0}{10.6}
\cl{claude}
\end{codeblock}

Inside a session, \texttt{/\allowbreak{}login} reauthenticates or switches account and \texttt{/\allowbreak{}logout} signs out [35], [52]. Credentials are stored in the macOS Keychain where available and protected by file permissions on Windows and Linux [53].


\textbf{CAUTION:} several features require a claude.ai subscription login specifically and are silently unavailable when \texttt{ANTHROPIC\_\allowbreak{}API\_\allowbreak{}KEY}, \texttt{ANTHROPIC\_\allowbreak{}AUTH\_\allowbreak{}TOKEN} or an \texttt{api\allowbreak{}Key\allowbreak{}Helper} takes precedence — Chrome integration, artifacts, routines, cloud sessions and teleport among them [20], [21], [47], [50]. If a documented command appears to be missing, check the authentication path before you check the version.


\FloatBarrier
\setcounter{section}{3}
\section{Windows: native or WSL 2}

\begin{itemize}
\item \textbf{Native Windows} suits Windows-native projects and tooling. Sandboxing is unavailable [27]. Git for Windows is recommended but no longer mandatory; without it, Claude Code uses PowerShell as its shell tool [2].
\item \textbf{WSL 2} suits Linux toolchains and is required for the sandboxed Bash tool [27]. Chrome integration is not supported under WSL [21].
\end{itemize}

\begin{calloutbox}{palebrass}{brassdark}{2.0mm}
\textbf{CAUTION — Windows WebDAV.} Anthropic recommends against enabling WebDAV or allowing access to paths such as \texttt{\textbackslash{}\textbackslash{}*} that may contain WebDAV subdirectories, because doing so can let Claude Code trigger network requests to remote hosts outside the permission system [53].
\end{calloutbox}

\subsection*{Checkpoint}

\begin{itemize}
\item \texttt{claude -\allowbreak{}-\allowbreak{}version} prints a version.
\item \texttt{claude doct\allowbreak{}or} reports no unresolved critical issue.
\item \texttt{claude} starts and shows the account you intended.
\item You know whether your installation auto-updates.
\item You know whether you are authenticated with a subscription login or an API key.
\end{itemize}

\begingroup
\def\kprows{%
\item The native installer auto-updates; Homebrew and WinGet installations do not.
\item Verify with \texttt{claude -\allowbreak{}-\allowbreak{}version} and \texttt{claude doct\allowbreak{}or} before you log in, and know which authentication path you are on.
\item Several features — Chrome, artifacts, routines, cloud sessions, teleport — require a claude.ai subscription login and are silently unavailable behind an API key.
\item Sandboxing does not run on native Windows; Chrome integration does not run under WSL. Choose the Windows path accordingly.
}%
\def\kpbody{\begin{minipage}{\textwidth}\subsection*{Key points}\begin{itemize}\kprows\end{itemize}\end{minipage}}%
\begingroup
\def\sloppy{\tolerance 9999\emergencystretch 3em\hfuzz 200pt\vfuzz 200pt}%
\hbadness=10000\vbadness=10000\hfuzz=200pt\vfuzz=200pt
\global\setbox\kpbox=\hbox{\kpbody}%
\endgroup
\par\addvspace{4.2mm}
\ifdim\dimexpr\ht\kpbox+\dp\kpbox\relax>0.30\textheight
  \typeout{HANDBOOK-KEYPOINTS broken \the\dimexpr\ht\kpbox+\dp\kpbox\relax}%
  \subsection*{Key points}
  \begin{itemize}\kprows\end{itemize}
\else
  \typeout{HANDBOOK-KEYPOINTS atomic \the\dimexpr\ht\kpbox+\dp\kpbox\relax}%
  \noindent\kpbody
\fi
\par\addvspace{1.4mm}
\endgroup


\FloatBarrier
\renewcommand{\chaptertitlelabel}{Part I \textperiodcentered\ Chapter 3}
\setcounter{chapter}{2}
\chapter{Surfaces, projects and boundaries}

\FloatBarrier
\setcounter{section}{0}
\section{Set up the IDE}

Install Visual Studio Code, open \textbf{Extensions}, search for \textbf{Claude Code}, and confirm the publisher is \textbf{Anthropic} before installing [44]. From the Command Palette, \textbf{Claude Code: Open in New Tab} gives the conversation a full editor tab and \textbf{Open in Terminal} launches the CLI inside the IDE. In VS Code, Chrome integration is available whenever the browser extension is installed, with no additional flag [21].


\FloatBarrier
\setcounter{section}{1}
\section{Create the practice workspace}

Create \texttt{research-\allowbreak{}brief-\allowbreak{}workspace}, open it with \textbf{File → Open Folder}, and create the directories of “The running project”. Add two harmless text files to \texttt{notes/\allowbreak{}}.


The opened folder is the default working boundary. In Manual mode Claude Code can write only inside that folder and its subfolders, and asks before reading paths outside it with the Read, Grep and Glob tools; in auto mode it reads outside the boundary without asking [53]. Extend the boundary deliberately with \texttt{-\allowbreak{}-\allowbreak{}add-\allowbreak{}dir} or \texttt{/\allowbreak{}add-\allowbreak{}dir}, not reflexively [28].


\begin{calloutbox}{palebrass}{brassdark}{2.0mm}
\textbf{CAUTION:} an opened folder is not a security sandbox. Claude Code can request additional directories, run commands, reach the network and use connected services subject to permissions. Do not open your home directory as a project: trust acceptance there is held for the session only and is never written to disk, so the prompt reappears at every launch, and the blast radius of a permission mistake is your entire personal filesystem [53].
\end{calloutbox}

\FloatBarrier
\setcounter{section}{2}
\section{Monorepos and large trees}

In a monorepo, nested \texttt{CLAUDE.\allowbreak{}md} files from other teams are picked up from every ancestor directory. Use \texttt{claude\allowbreak{}Md\allowbreak{}Excludes} to skip them, path-scoped rules to load instructions only for matching files, and per-package skills to keep guidance local [26], [54]. Sparse worktrees and code intelligence are covered in [54].


\subsection*{Exercise}

In Plan mode, ask:


\begin{codeblock}{7.0}{8.3}
\cl{Inspect~this~workspace~without~modifying~it.~Describe~its~structure,~identify~what~is~missing}
\cl{for~a~research-brief~workflow,~and~propose~the~smallest~safe~next~step.~Do~not~access~anything}
\cl{outside~this~folder.}
\end{codeblock}

\emph{A worked solution is given in Appendix M.}


\subsection*{Checkpoint}

The response names only files inside the practice folder and proposes — rather than performs — changes.


\begingroup
\def\kprows{%
\item The opened folder is the default working boundary, not a security sandbox.
\item In Manual mode Claude writes only inside that folder; in auto mode it reads outside it without asking.
\item Never open a home directory as a project: trust acceptance there is never written to disk and the blast radius is your whole filesystem.
\item In a monorepo, use \texttt{claude\allowbreak{}Md\allowbreak{}Excludes} and path-scoped rules before you use a bigger context window.
}%
\def\kpbody{\begin{minipage}{\textwidth}\subsection*{Key points}\begin{itemize}\kprows\end{itemize}\end{minipage}}%
\begingroup
\def\sloppy{\tolerance 9999\emergencystretch 3em\hfuzz 200pt\vfuzz 200pt}%
\hbadness=10000\vbadness=10000\hfuzz=200pt\vfuzz=200pt
\global\setbox\kpbox=\hbox{\kpbody}%
\endgroup
\par\addvspace{4.2mm}
\ifdim\dimexpr\ht\kpbox+\dp\kpbox\relax>0.30\textheight
  \typeout{HANDBOOK-KEYPOINTS broken \the\dimexpr\ht\kpbox+\dp\kpbox\relax}%
  \subsection*{Key points}
  \begin{itemize}\kprows\end{itemize}
\else
  \typeout{HANDBOOK-KEYPOINTS atomic \the\dimexpr\ht\kpbox+\dp\kpbox\relax}%
  \noindent\kpbody
\fi
\par\addvspace{1.4mm}
\endgroup

\FloatBarrier
\renewcommand{\chaptertitlelabel}{Part I \textperiodcentered\ Chapter 4}
\setcounter{chapter}{3}
\chapter{Modes, models, effort and usage}

\FloatBarrier
\setcounter{section}{0}
\section{The controls that matter first}

\begin{itemize}
\item \textbf{Permission mode} — switched with \texttt{Shift+Tab} in the CLI, the mode indicator in VS Code, or the mode selector in Desktop [31]. The cycle runs \texttt{default} → \texttt{accept\allowbreak{}Edits} → \texttt{plan} in every session, with two further modes slotting in after \texttt{plan} when they are enabled; \texttt{dont\allowbreak{}Ask} never appears in it at all. Table 5 gives the full picture. Covered in Chapter 7.
\item \textbf{Model} — \texttt{/\allowbreak{}model}, or the picker. Availability changes; do not memorise a model name from a recording [55].
\item \textbf{Effort} — \texttt{/\allowbreak{}effort}, with levels \texttt{low}, \texttt{medium}, \texttt{high}, \texttt{xhigh} and \texttt{max} on every surface, plus \texttt{ultracode} and \texttt{auto}, whose availability differs by surface [35], [55]. Section 4.2 records what each surface accepted. \texttt{ultracode} is not merely an effort level; see Chapter 22.
\item \textbf{Fast mode} — \texttt{/\allowbreak{}fast}, trading depth for latency on models that support it [56].
\item \textbf{Advisor} — \texttt{/\allowbreak{}advisor} pairs the working model with a stronger model consulted at key moments [57].
\item \textbf{Context attachment} — \texttt{@} mentions and the attachment control.
\item \textbf{Interrupt} — \texttt{Esc} stops the current turn; \texttt{Ctrl+C} is stronger [58].
\item \textbf{Settings from the prompt} — \texttt{/\allowbreak{}config} opens settings or sets an individual key, including keys with no dedicated command [35].
\item \textbf{Command menu} — type \texttt{/\allowbreak{}} at the start of a message.
\end{itemize}

\FloatBarrier
\setcounter{section}{1}
\section{The effort vocabulary, by surface}

\texttt{/\allowbreak{}effort} and \texttt{claude -\allowbreak{}-\allowbreak{}effort} are two interfaces onto one control, and they do not accept the same words. The table records what each surface did when the value was offered to it, on Claude Code 2.1.246 on 26 August 2026. It is a configuration-stamped observation, not a product guarantee: the menu a given plan and model offer may differ, and the entry in Appendix J states the residual [35], [55].


\begingroup
\def\tblrows{%
\texttt{low} & yes & yes & yes \\
\texttt{medium} & yes & yes & yes \\
\texttt{high} & yes & yes & yes \\
\texttt{xhigh} & yes & yes & yes \\
\texttt{max} & yes & yes & yes \\
\texttt{auto} & — & yes & no \\
\texttt{ultracode} & yes & — † & yes † \\
}%
\def\tblbody{\begin{minipage}{\textwidth}\boxcaption{Table 3 — Effort levels accepted by surface, observed on Claude Code 2.1.246, 26 August 2026}
{\footnotesize\begin{tabular}{L{36.2mm}L{32.5mm}L{33.5mm}L{34.4mm}}
\toprule
\textbf{Level} & \textbf{\texttt{/\allowbreak{}effort} (terminal)} & \textbf{\texttt{/\allowbreak{}effort} (VS Code)} & \textbf{\texttt{claude -\allowbreak{}-\allowbreak{}effort}} \\
\midrule
\tblrows
\bottomrule\end{tabular}}\end{minipage}}%
\begingroup
\def\sloppy{\tolerance 9999\emergencystretch 3em\hfuzz 200pt\vfuzz 200pt}%
\hbadness=10000\vbadness=10000\hfuzz=200pt\vfuzz=200pt
\global\setbox\tblbox=\hbox{\tblbody}%
\endgroup
\par\addvspace{2.6mm}
\ifdim\dimexpr\ht\tblbox+\dp\tblbox\relax>0.55\textheight
  \typeout{HANDBOOK-TABLE broken \the\dimexpr\ht\tblbox+\dp\tblbox\relax}%
  \tabcaption{Table 3 — Effort levels accepted by surface, observed on Claude Code 2.1.246, 26 August 2026}
  {\footnotesize\begin{longtable}{L{36.2mm}L{32.5mm}L{33.5mm}L{34.4mm}}
  \toprule
\textbf{Level} & \textbf{\texttt{/\allowbreak{}effort} (terminal)} & \textbf{\texttt{/\allowbreak{}effort} (VS Code)} & \textbf{\texttt{claude -\allowbreak{}-\allowbreak{}effort}} \\
\midrule\endfirsthead
  \multicolumn{4}{@{}l@{}}{%
  \sffamily\footnotesize\itshape\color{inkgrey}Table 3 — Effort levels accepted by surface, observed on Claude Code 2.1.246, 26 August 2026 \textemdash\ continued}\\[1.2mm]
  \toprule
\textbf{Level} & \textbf{\texttt{/\allowbreak{}effort} (terminal)} & \textbf{\texttt{/\allowbreak{}effort} (VS Code)} & \textbf{\texttt{claude -\allowbreak{}-\allowbreak{}effort}} \\
\midrule\endhead
  \bottomrule\endfoot
  \bottomrule\endlastfoot
  \tblrows
  \end{longtable}}%
\else
  \typeout{HANDBOOK-TABLE atomic \the\dimexpr\ht\tblbox+\dp\tblbox\relax}%
  \noindent\tblbody
\fi
\par\addvspace{2.6mm}
\endgroup

\textbf{yes} accepted · \textbf{no} refused · \textbf{—} not established on that surface · \textbf{†} absent from that surface's own help or usage text.


Invocation differs too. In the terminal, a bare \texttt{/\allowbreak{}effort} shows the current level and a selector; in the VS Code extension it requires an explicit level and answers a bare call with its usage string.


Two consequences matter in practice. \textbf{\texttt{claude -\allowbreak{}-\allowbreak{}effort auto} is refused}: the flag warns that \texttt{auto} is not a valid value and the session then runs at the default effort, so a script that sets it neither gets what it asked for nor fails. And \textbf{\texttt{ultracode} is accepted by the flag although \texttt{claude -\allowbreak{}-\allowbreak{}help} does not list it} — which is why this book treats a help string as evidence of what is supported and never as evidence of what is refused (Appendix J).


\FloatBarrier
\setcounter{section}{2}
\section{A decision rule for model and effort}

\begingroup
\def\tblrows{%
Locate a file; reformat known text & Lower-cost model, normal effort \\
Summarise supplied material & Normal effort \\
Design a project; resolve genuine ambiguity & Stronger model, high effort \\
Execute an already approved, explicit plan & Normal to high effort; verify completeness rather than creativity \\
Multi-source research; architecture trade-off & \texttt{xhigh}, or a dynamic workflow (Chapter 22) \\
}%
\def\tblbody{\begin{minipage}{\textwidth}\boxcaption{Table 4 — Suggested starting posture by work type (practitioner guidance)}
{\footnotesize\begin{tabular}{L{70.4mm}L{72.6mm}}
\toprule
\textbf{Work} & \textbf{Suggested starting posture} \\
\midrule
\tblrows
\bottomrule\end{tabular}}\end{minipage}}%
\begingroup
\def\sloppy{\tolerance 9999\emergencystretch 3em\hfuzz 200pt\vfuzz 200pt}%
\hbadness=10000\vbadness=10000\hfuzz=200pt\vfuzz=200pt
\global\setbox\tblbox=\hbox{\tblbody}%
\endgroup
\par\addvspace{2.6mm}
\ifdim\dimexpr\ht\tblbox+\dp\tblbox\relax>0.55\textheight
  \typeout{HANDBOOK-TABLE broken \the\dimexpr\ht\tblbox+\dp\tblbox\relax}%
  \tabcaption{Table 4 — Suggested starting posture by work type (practitioner guidance)}
  {\footnotesize\begin{longtable}{L{70.4mm}L{72.6mm}}
  \toprule
\textbf{Work} & \textbf{Suggested starting posture} \\
\midrule\endfirsthead
  \multicolumn{2}{@{}l@{}}{%
  \sffamily\footnotesize\itshape\color{inkgrey}Table 4 — Suggested starting posture by work type (practitioner guidance) \textemdash\ continued}\\[1.2mm]
  \toprule
\textbf{Work} & \textbf{Suggested starting posture} \\
\midrule\endhead
  \bottomrule\endfoot
  \bottomrule\endlastfoot
  \tblrows
  \end{longtable}}%
\else
  \typeout{HANDBOOK-TABLE atomic \the\dimexpr\ht\tblbox+\dp\tblbox\relax}%
  \noindent\tblbody
\fi
\par\addvspace{2.6mm}
\endgroup

\textbf{PRACTITIONER NOTE.} This table is judgement, not a product guarantee. Raising effort does not make an answer correct, and lowering it does not make an operation safe. Effort changes how much the model deliberates; it changes nothing about what it is permitted to do.


\FloatBarrier
\setcounter{section}{3}
\section{Measure usage and context separately}

\texttt{/\allowbreak{}usage} reports session cost, plan limits and activity, and can break usage down by skill, subagent, plugin and MCP server [35], [59]. \texttt{/\allowbreak{}cost} and \texttt{/\allowbreak{}stats} are both documented aliases of it, and \texttt{/\allowbreak{}stats} opens on the Stats tab [35]. \texttt{/\allowbreak{}context} answers a different question: what currently occupies the conversation window [33]. Allocation and occupancy are related to tokens but they are not the same quantity, and optimising one at the expense of the other is a common and expensive mistake.


Two further facts belong here because they explain most surprising bills. First, usage climbs superlinearly in a long session because each request carries the accumulated history [59]. Second, prompt caching means a model switch, a \texttt{/\allowbreak{}compact}, or an output-style change forces an uncached turn, and a \texttt{CLAUDE.\allowbreak{}md} edit does not take effect mid-session [60]. If a change to project instructions appears to have been ignored, that is usually why.


\FloatBarrier
\setcounter{section}{4}
\section{What to optimise first}

\begin{enumerate}
\item Stop passing irrelevant material.
\item Persist durable outputs to files instead of regenerating them.
\item Delegate only bounded, independent work (Chapter 16).
\item Remove unused plugins, connectors and agents — after measuring their contribution, not before.
\item Start a clean session when the task changes.
\end{enumerate}

Do not optimise by skipping verification. A cheaper wrong answer is still waste, and in this domain it is waste that arrives looking like success.



\begingroup
\def\kprows{%
\item Permission mode, model, effort and output style are four independent controls; changing one does not change the others.
\item \texttt{/\allowbreak{}usage} measures your allocation; \texttt{/\allowbreak{}context} measures what occupies the window. They are different questions.
\item Usage climbs superlinearly in a long session because every request carries the accumulated history.
\item A \texttt{CLAUDE.\allowbreak{}md} edit does not take effect mid-session — the system context is read once at session start.
}%
\def\kpbody{\begin{minipage}{\textwidth}\subsection*{Key points}\begin{itemize}\kprows\end{itemize}\end{minipage}}%
\begingroup
\def\sloppy{\tolerance 9999\emergencystretch 3em\hfuzz 200pt\vfuzz 200pt}%
\hbadness=10000\vbadness=10000\hfuzz=200pt\vfuzz=200pt
\global\setbox\kpbox=\hbox{\kpbody}%
\endgroup
\par\addvspace{4.2mm}
\ifdim\dimexpr\ht\kpbox+\dp\kpbox\relax>0.30\textheight
  \typeout{HANDBOOK-KEYPOINTS broken \the\dimexpr\ht\kpbox+\dp\kpbox\relax}%
  \subsection*{Key points}
  \begin{itemize}\kprows\end{itemize}
\else
  \typeout{HANDBOOK-KEYPOINTS atomic \the\dimexpr\ht\kpbox+\dp\kpbox\relax}%
  \noindent\kpbody
\fi
\par\addvspace{1.4mm}
\endgroup

\breakrule

\FloatBarrier
\parttitle{Part II}{A safe working method}
\addcontentsline{toc}{part}{Part II \textemdash\ A safe working method}

\FloatBarrier
\renewcommand{\chaptertitlelabel}{Part II \textperiodcentered\ Chapter 5}
\setcounter{chapter}{4}
\chapter{Start and steer a session}

\FloatBarrier
\setcounter{section}{0}
\section{Begin with an inspectable brief}

A first message needs no special syntax. It needs enough structure that omissions are visible:


\begin{codeblock}{7.0}{8.3}
\cl{Objective}
\cl{Create~a~one-page~research~brief~from~the~files~in~notes/.}
\cl{}
\cl{Scope}
\cl{Use~only~the~supplied~notes.~Identify~evidence~gaps~instead~of~filling~them~from~memory~or}
\cl{the~web.}
\cl{}
\cl{Output}
\cl{Save~drafts/brief-v1.md~with~an~executive~summary,~three~findings,~and~a~claim-to-source~table.}
\cl{}
\cl{Boundaries}
\cl{Do~not~edit~source~files.~Do~not~access~other~folders~or~the~internet.}
\cl{}
\cl{Done~when}
\cl{Every~factual~statement~maps~to~a~named~source~note,~and~the~file~passes~a~completeness~review.}
\end{codeblock}

The headings serve the reader as much as the model. They make a missing boundary conspicuous before it becomes an incident.


\FloatBarrier
\setcounter{section}{1}
\section{Watch the beginning of unfamiliar work}

Claude Code renders progress and tool calls. Early in a task, read them and ask:


\begin{itemize}
\item Is it reading the files you intended?
\item Is it attempting an external action you did not request?
\item Is its interpretation of "done" the same as yours?
\item Did it delegate, and if so with what brief? (See Chapter 16.)
\end{itemize}

Once the direction is stable you need not watch every line. Unattended work, by contrast, needs tighter boundaries and a stronger completion test, not looser ones.


\FloatBarrier
\setcounter{section}{2}
\section{Steer without restarting}

Type a message and press Enter to queue it. If tools are running, Claude receives it when those calls finish and can adjust within the same turn; anything further becomes a subsequent turn [58]. Current Desktop accepts a correction after the active action rather than requiring a restart.


When the correction cannot wait, press \texttt{Esc} to stop the current response or tool call. Completed work remains; queued content is sent; the Up arrow retrieves a queued entry for editing. Say explicitly whether to resume:


\begin{codeblock}{8.0}{9.4}
\cl{Stop.~Do~not~use~the~internet;~I~omitted~that~boundary.~Keep~the~local~analysis~already}
\cl{completed,~discard~any~external~material,~and~resume~from~the~approved~notes~only.}
\end{codeblock}

\begin{calloutbox}{palebrass}{brassdark}{2.0mm}
\textbf{CAUTION:} interruption does not undo completed file changes and cannot recall an external action. Inspect the diff or the checkpoint state (Chapter 7).
\end{calloutbox}

\FloatBarrier
\setcounter{section}{3}
\section{Side questions, focus and recall}

Three small controls repay learning early. \texttt{/\allowbreak{}btw} asks a side question without adding it to the conversation history. \texttt{/\allowbreak{}focus} collapses the view to the last prompt and response. \texttt{/\allowbreak{}recap} generates a one-line summary of the session, and \texttt{/\allowbreak{}export} writes a rendered plain-text transcript for the record [30], [35]. The last of these matters for audit: it is the supported way to retain evidence of what happened, and it is preferable to parsing the raw JSONL transcript, whose format is internal and changes between releases [30].


\FloatBarrier
\setcounter{section}{4}
\section{Voice input}

Voice dictation makes long briefs easier to compose [61]. Read the transcription before sending. Filenames, command flags, negations and numbers are where dictation fails, and a dropped "not" in a boundary statement is the most expensive single-word error in this book.


\subsection*{Exercise}

Run the brief from Chapter 1. Mid-task, queue one non-urgent steering message requiring the findings to be ordered by decision relevance. Interrupt only if a boundary is crossed.


\emph{A worked solution is given in Appendix M.}


\begingroup
\def\kprows{%
\item A brief with headings makes omissions visible before they become incidents.
\item Queue a correction and it applies after the active tool call; press \texttt{Esc} only when it cannot wait.
\item Interruption does not undo completed file changes and cannot recall an external action.
\item \texttt{/\allowbreak{}export} is the supported way to retain evidence of a session; the raw JSONL transcript format is internal and changes between releases.
}%
\def\kpbody{\begin{minipage}{\textwidth}\subsection*{Key points}\begin{itemize}\kprows\end{itemize}\end{minipage}}%
\begingroup
\def\sloppy{\tolerance 9999\emergencystretch 3em\hfuzz 200pt\vfuzz 200pt}%
\hbadness=10000\vbadness=10000\hfuzz=200pt\vfuzz=200pt
\global\setbox\kpbox=\hbox{\kpbody}%
\endgroup
\par\addvspace{4.2mm}
\ifdim\dimexpr\ht\kpbox+\dp\kpbox\relax>0.30\textheight
  \typeout{HANDBOOK-KEYPOINTS broken \the\dimexpr\ht\kpbox+\dp\kpbox\relax}%
  \subsection*{Key points}
  \begin{itemize}\kprows\end{itemize}
\else
  \typeout{HANDBOOK-KEYPOINTS atomic \the\dimexpr\ht\kpbox+\dp\kpbox\relax}%
  \noindent\kpbody
\fi
\par\addvspace{1.4mm}
\endgroup

\FloatBarrier
\renewcommand{\chaptertitlelabel}{Part II \textperiodcentered\ Chapter 6}
\setcounter{chapter}{5}
\chapter{Plan, comment, approve, execute}

\FloatBarrier
\setcounter{section}{0}
\section{When to plan}

Use Plan mode when the work has several stages, unclear requirements, expensive external actions, or broad file impact. Enter it with \texttt{/\allowbreak{}plan}, the mode selector, or \texttt{claude -\allowbreak{}-\allowbreak{}permission-\allowbreak{}mode plan} [31], [62]. In Plan mode Claude reads and explores but does not edit source; browser calls that only read run without prompting, while state-changing browser calls still ask [21].


A plan should reduce uncertainty before edits, not restate your request in longer form.


\FloatBarrier
\setcounter{section}{1}
\section{The minimum complete plan}

A plan is ready for approval when it identifies:


\begin{enumerate}
\item the intended outcome and the explicit exclusions;
\item inputs and the source-of-truth files;
\item every file or system that will change;
\item ordered implementation steps;
\item human decisions still outstanding;
\item tests or inspection evidence for each material result;
\item external side effects;
\item a rollback or recovery route.
\end{enumerate}

For a report, item 6 means source coverage, link checks, spelling and rendering. For software it means automated tests plus direct behavioural observation.


\FloatBarrier
\setcounter{section}{2}
\section{Review the plan like a contract}

Do not approve a plan because it is long or confident. Comment on specific weaknesses:


\begin{codeblock}{8.0}{9.4}
\cl{Step~3~is~too~broad.~Separate~source~extraction~from~interpretation,~and~require~one}
\cl{evidence-ledger~row~per~material~claim~before~any~drafting~begins.}
\end{codeblock}

\begin{codeblock}{9.0}{10.6}
\cl{Add~a~human~approval~gate~before~anything~is~published~or~sent~externally.}
\end{codeblock}

\begin{codeblock}{9.0}{10.6}
\cl{Define~what~"complete"~means~for~the~final~file,~and~how~you~will~verify~it.}
\end{codeblock}

In surfaces that render an editable plan, select text and attach comments; otherwise quote the step. Continue until the high-impact uncertainties are resolved.


Accepting a plan also gives the session a generated title based on the plan, unless you have already named it [30] — a small but useful audit artefact.


\FloatBarrier
\setcounter{section}{3}
\section{Approve the right scope}

Approving a plan is not blanket permission for any method. Keep the permission mode you meant to be in. A good plan combined with auto mode can support long unattended execution, but auto mode is a classifier-assisted permission mode: it is not a sandbox and not a guarantee [31]. See Chapter 7.


\FloatBarrier
\setcounter{section}{4}
\section{Require a completion report, then check it}

\begin{codeblock}{8.0}{9.4}
\cl{Report~the~files~changed,~the~tests~or~checks~actually~executed~with~their~output,~unresolved}
\cl{limitations,~and~every~external~side~effect.~Do~not~describe~a~check~that~did~not~run.}
\end{codeblock}

Then inspect the evidence yourself: \texttt{/\allowbreak{}diff} opens the interactive diff viewer for uncommitted changes [35], and \texttt{git diff} is the version you can archive.


\subsection*{Exercise}

Plan a second brief that adds web research. Require a query strategy, source-quality criteria, a claim ledger, a source register, a contradiction pass, and explicit approval before any external publication. Do not execute until the plan satisfies all eight points of Section 6.2.


\emph{A worked solution is given in Appendix M.}


\begingroup
\def\kprows{%
\item Plan mode reduces uncertainty before edits; it is not a longer restatement of your request.
\item A plan is ready for approval when it names outcomes, inputs, changes, steps, open decisions, tests, side effects and a rollback route.
\item Approving a plan is not blanket permission for any method — keep the permission mode you intended.
\item Require a completion report that distinguishes checks that ran from checks that were described, then inspect the evidence yourself.
}%
\def\kpbody{\begin{minipage}{\textwidth}\subsection*{Key points}\begin{itemize}\kprows\end{itemize}\end{minipage}}%
\begingroup
\def\sloppy{\tolerance 9999\emergencystretch 3em\hfuzz 200pt\vfuzz 200pt}%
\hbadness=10000\vbadness=10000\hfuzz=200pt\vfuzz=200pt
\global\setbox\kpbox=\hbox{\kpbody}%
\endgroup
\par\addvspace{4.2mm}
\ifdim\dimexpr\ht\kpbox+\dp\kpbox\relax>0.30\textheight
  \typeout{HANDBOOK-KEYPOINTS broken \the\dimexpr\ht\kpbox+\dp\kpbox\relax}%
  \subsection*{Key points}
  \begin{itemize}\kprows\end{itemize}
\else
  \typeout{HANDBOOK-KEYPOINTS atomic \the\dimexpr\ht\kpbox+\dp\kpbox\relax}%
  \noindent\kpbody
\fi
\par\addvspace{1.4mm}
\endgroup

\FloatBarrier
\renewcommand{\chaptertitlelabel}{Part II \textperiodcentered\ Chapter 7}
\setcounter{chapter}{6}
\chapter{Permissions, sandboxing and checkpoints}

Permissions are product-enforced rules governing tool use. Instructions in a prompt or in \texttt{CLAUDE.\allowbreak{}md} influence what Claude attempts; they do not grant or revoke access [26], [28]. This distinction is the foundation of the whole book: an instruction is a preference, a permission rule is a control, and a hook is executable enforcement (Chapter 17).


\FloatBarrier
\setcounter{section}{0}
\section{Permission modes}

On Pro, Max and Team plans the built-in starting mode is \textbf{auto mode} [31]. That is a categorical statement in current documentation, not a tendency. Organisation settings, the surface, and a \texttt{default\allowbreak{}Mode} in settings can change it, so inspect the selector rather than assuming.


\begingroup
\def\tblrows{%
Manual & \texttt{default} (alias \texttt{manual}, 2.1.200+) & 1 & Asks before most actions that edit files, run commands or reach the network & Learning, sensitive repositories, unfamiliar tasks \\
Accept Edits & \texttt{accept\allowbreak{}Edits} & 2 & Auto-approves file edits and a fixed set of filesystem Bash commands (\texttt{mkdir}, \texttt{touch}, \texttt{rm}, \texttt{mv}, \texttt{cp}, \texttt{sed}) inside the working directory & Well-bounded local editing with review \\
Plan & \texttt{plan} & 3 & Reads and explores; no source edits & Analysis and design before execution \\
Bypass Permissions & \texttt{bypass\allowbreak{}Permissions} & 4 & Skips ordinary prompts & Isolated disposable containers or VMs only \\
Auto & \texttt{auto} & 5 & A classifier model reviews actions in place of you and blocks those it judges unsafe & Bounded work with sandboxing and verification \\
Don't Ask & \texttt{dont\allowbreak{}Ask} & — & Denies anything not pre-approved & Locked-down repeatable workflows \\
}%
\def\tblbody{\begin{minipage}{\textwidth}\boxcaption{Table 5 — Permission modes, the cycle position that reaches each, and appropriate use}
{\footnotesize\begin{tabular}{H{22.0mm}H{34.1mm}L{21.3mm}L{28.0mm}L{28.0mm}}
\toprule
\textbf{Mode} & \textbf{Config value} & \textbf{\texttt{Shift+Tab}} & \textbf{Behaviour} & \textbf{Sensible use} \\
\midrule
\tblrows
\bottomrule\end{tabular}}\end{minipage}}%
\begingroup
\def\sloppy{\tolerance 9999\emergencystretch 3em\hfuzz 200pt\vfuzz 200pt}%
\hbadness=10000\vbadness=10000\hfuzz=200pt\vfuzz=200pt
\global\setbox\tblbox=\hbox{\tblbody}%
\endgroup
\par\addvspace{2.6mm}
\ifdim\dimexpr\ht\tblbox+\dp\tblbox\relax>0.55\textheight
  \typeout{HANDBOOK-TABLE broken \the\dimexpr\ht\tblbox+\dp\tblbox\relax}%
  \tabcaption{Table 5 — Permission modes, the cycle position that reaches each, and appropriate use}
  {\footnotesize\begin{longtable}{H{22.0mm}H{34.1mm}L{21.3mm}L{28.0mm}L{28.0mm}}
  \toprule
\textbf{Mode} & \textbf{Config value} & \textbf{\texttt{Shift+Tab}} & \textbf{Behaviour} & \textbf{Sensible use} \\
\midrule\endfirsthead
  \multicolumn{5}{@{}l@{}}{%
  \sffamily\footnotesize\itshape\color{inkgrey}Table 5 — Permission modes, the cycle position that reaches each, and appropriate use \textemdash\ continued}\\[1.2mm]
  \toprule
\textbf{Mode} & \textbf{Config value} & \textbf{\texttt{Shift+Tab}} & \textbf{Behaviour} & \textbf{Sensible use} \\
\midrule\endhead
  \bottomrule\endfoot
  \bottomrule\endlastfoot
  \tblrows
  \end{longtable}}%
\else
  \typeout{HANDBOOK-TABLE atomic \the\dimexpr\ht\tblbox+\dp\tblbox\relax}%
  \noindent\tblbody
\fi
\par\addvspace{2.6mm}
\endgroup

\textbf{The keystroke does not reach every mode, and what it reaches depends on your session.} In the CLI the cycle runs \texttt{default} → \texttt{accept\allowbreak{}Edits} → \texttt{plan} and back, in every session. Two further modes slot in after \texttt{plan}, in this order: \textbf{\texttt{bypass\allowbreak{}Permissions}}, when the session was started with \texttt{-\allowbreak{}-\allowbreak{}permission-\allowbreak{}mode bypass\allowbreak{}Permissions}, \texttt{-\allowbreak{}-\allowbreak{}dangerously-\allowbreak{}skip-\allowbreak{}permissions}, \texttt{-\allowbreak{}-\allowbreak{}allow-\allowbreak{}dangerously-\allowbreak{}skip-\allowbreak{}permissions} or a \texttt{bypass\allowbreak{}Permissions} default in settings; then \textbf{\texttt{auto}}, when the account meets the auto-mode requirements. So the keystroke reaches three modes by default, and at most five in a session that has enabled both. \textbf{\texttt{dont\allowbreak{}Ask} never appears in the cycle at all}: it is set with \texttt{claude -\allowbreak{}-\allowbreak{}permission-\allowbreak{}mode dont\allowbreak{}Ask} [31], [62].


The numbering in the table is the cycle order, not a ranking. Two consequences follow for a reader who takes the table and reaches for the keystroke. If \texttt{bypass\allowbreak{}Permissions} and \texttt{auto} are not enabled, the keystroke moves between the first three rows only, and the fourth and fifth are reached by flag or by settings. And \texttt{dont\allowbreak{}Ask} is never reached by cycling in any configuration, which is the one case where the table and the keystroke disagree about what exists.


\texttt{claude -\allowbreak{}-\allowbreak{}permission-\allowbreak{}mode} names six values when it refuses one — \texttt{accept\allowbreak{}Edits}, \texttt{auto}, \texttt{bypass\allowbreak{}Permissions}, \texttt{dont\allowbreak{}Ask}, \texttt{manual} and \texttt{plan} — and additionally accepts \texttt{default}, which that enumeration omits. That is not two vocabularies but one: \texttt{manual} is a documented alias for \texttt{default}, from 2.1.200 onward, so the validator lists the alias and not the canonical value the settings files and the SDK use [31]. It is the plainest case of the rule in Appendix J — an enumeration is evidence of what is accepted, never of what is not.


Two mode facts have operational consequences. \texttt{plan} and \texttt{bypass\allowbreak{}Permissions} are \textbf{never} restored when a session resumes; a resumed session returns to the mode a new session would start in [30]. And \texttt{bypass\allowbreak{}Permissions} propagates: agent-team teammates inherit a lead running without checks [29], and a child agent inherits its parent's mode, which a \texttt{permission\allowbreak{}Mode} of its own cannot override when the parent is in \texttt{bypass\allowbreak{}Permissions}, \texttt{accept\allowbreak{}Edits} or \texttt{auto} [25]. Workflow subagents are the exception, and they run the other way. The subagents a workflow spawns always run in \texttt{accept\allowbreak{}Edits} and inherit your tool allowlist, whatever the session's mode [19]. Your mode governs only the launch prompt. The session's mode is therefore not a reliable guide to what an agent it spawns may do, in either direction; Chapter 22 treats the workflow case in full.


\FloatBarrier
\setcounter{section}{1}
\section{Rules, and the order they are evaluated in}

Run \texttt{/\allowbreak{}permissions} to inspect and edit Allow, Ask and Deny rules. \textbf{Rules are evaluated in the order deny, then ask, then allow; the first match in that order decides, and specificity does not change the order} [28]. A broad deny such as \texttt{Bash(aws *)} therefore blocks every matching call, including one that also matches a narrower allow like \texttt{Bash(aws s3\allowbreak{} ls)}. A deny rule cannot carry allowlist exceptions.


Two subtleties are worth committing to memory. A \textbf{bare tool name} in \texttt{deny} removes the tool from Claude's context entirely, so Claude never sees it; a \textbf{scoped rule} such as \texttt{Bash(rm *)} leaves the tool available and blocks matching calls [28]. And a deny or ask rule matches past a leading environment assignment, while an allow rule does not: \texttt{Bash(rm *)} in \texttt{deny} still matches \texttt{FOO=\allowbreak{}bar rm -\allowbreak{}rf tmp/\allowbreak{}} [28].


\begin{codeblock}{9.0}{10.6}
\cl{\{}
\cl{~~"permissions":~\{}
\cl{~~~~"deny":~[}
\cl{~~~~~~"Read(./.env)",}
\cl{~~~~~~"Read(./.env.*)",}
\cl{~~~~~~"Read(./secrets/**)",}
\cl{~~~~~~"Read(\textasciitilde{}/.ssh/**)",}
\cl{~~~~~~"Bash(git~push~*)",}
\cl{~~~~~~"Bash(curl~*)",}
\cl{~~~~~~"Bash(wget~*)"}
\cl{~~~~],}
\cl{~~~~"ask":~[}
\cl{~~~~~~"WebFetch",}
\cl{~~~~~~"WebSearch"}
\cl{~~~~],}
\cl{~~~~"allow":~[}
\cl{~~~~~~"WebFetch(domain:code.claude.com)"}
\cl{~~~~]}
\cl{~~\}}
\cl{\}}
\end{codeblock}

Note what that example does \emph{not} do: because deny precedes allow, the \texttt{Web\allowbreak{}Fetch} allow rule does not defeat any deny rule, and the \texttt{Bash(curl *\allowbreak{})} deny is the reason to route web access through the \texttt{Web\allowbreak{}Fetch} tool with a domain specifier rather than through the shell [28], [53]. Wildcards in a domain specifier match only between dots unless they lead — \texttt{Web\allowbreak{}Fetch(domai\allowbreak{}n:\allowbreak{}*.\allowbreak{}example.\allowbreak{}com)} matches subdomains but not \texttt{example.\allowbreak{}com}, and \texttt{Web\allowbreak{}Fetch(domai\allowbreak{}n:\allowbreak{}example.\allowbreak{}*)} matches \texttt{example.\allowbreak{}org} but not \texttt{example.\allowbreak{}evil.\allowbreak{}com} [28].


Symlinks are resolved against both allow and deny: a link inside an allowed directory pointing at \texttt{\textasciitilde{}/\allowbreak{}.\allowbreak{}ssh/\allowbreak{}id\_\allowbreak{}rsa} is blocked because the target fails the allow rule and matches the deny rule [28].


Share team rules in \texttt{.\allowbreak{}claude/\allowbreak{}settings.\allowbreak{}json}; keep personal overrides and saved approvals in \texttt{.\allowbreak{}claude/\allowbreak{}settings.\allowbreak{}local.\allowbreak{}json}, which is not committed; user-wide settings live in \texttt{\textasciitilde{}/\allowbreak{}.\allowbreak{}claude/\allowbreak{}settings.\allowbreak{}json} [63]. Managed settings sit above all of them and cannot be overridden by any other level, including command-line arguments [28], [64]. See Chapter 31. Figure 2 traces a single tool call through the whole chain, in the order the decision is actually taken.


\begin{figure}[tbp]
\begin{codefig}{9.0}{10.6}
\cl{~~~Claude~proposes~a~tool~call}
\cl{~~~~~~~~~~~~~\pmboxdrawuni{2502}}
\cl{~~~~~~~~~~~~~▼}
\cl{~~~\pmboxdrawuni{250C}\pmboxdrawuni{2500}\pmboxdrawuni{2500}\pmboxdrawuni{2500}\pmboxdrawuni{2500}\pmboxdrawuni{2500}\pmboxdrawuni{2500}\pmboxdrawuni{2500}\pmboxdrawuni{2500}\pmboxdrawuni{2500}\pmboxdrawuni{2500}\pmboxdrawuni{2500}\pmboxdrawuni{2500}\pmboxdrawuni{2500}\pmboxdrawuni{2500}\pmboxdrawuni{2500}\pmboxdrawuni{2500}\pmboxdrawuni{2500}\pmboxdrawuni{2500}\pmboxdrawuni{2500}\pmboxdrawuni{2500}\pmboxdrawuni{2500}\pmboxdrawuni{2500}\pmboxdrawuni{2500}\pmboxdrawuni{2510}~~~exit~2~~~~~~\pmboxdrawuni{250C}\pmboxdrawuni{2500}\pmboxdrawuni{2500}\pmboxdrawuni{2500}\pmboxdrawuni{2500}\pmboxdrawuni{2500}\pmboxdrawuni{2500}\pmboxdrawuni{2500}\pmboxdrawuni{2500}\pmboxdrawuni{2500}\pmboxdrawuni{2500}\pmboxdrawuni{2500}\pmboxdrawuni{2500}\pmboxdrawuni{2500}\pmboxdrawuni{2500}\pmboxdrawuni{2500}\pmboxdrawuni{2500}\pmboxdrawuni{2500}\pmboxdrawuni{2500}\pmboxdrawuni{2500}\pmboxdrawuni{2500}\pmboxdrawuni{2500}\pmboxdrawuni{2500}\pmboxdrawuni{2500}\pmboxdrawuni{2500}\pmboxdrawuni{2500}\pmboxdrawuni{2500}\pmboxdrawuni{2500}\pmboxdrawuni{2500}\pmboxdrawuni{2500}\pmboxdrawuni{2500}\pmboxdrawuni{2510}}
\cl{~~~\pmboxdrawuni{2502}~PreToolUse~hook~~~~~~~\pmboxdrawuni{2502}\pmboxdrawuni{2500}\pmboxdrawuni{2500}\pmboxdrawuni{2500}\pmboxdrawuni{2500}\pmboxdrawuni{2500}\pmboxdrawuni{2500}\pmboxdrawuni{2500}\pmboxdrawuni{2500}\pmboxdrawuni{2500}\pmboxdrawuni{2500}\pmboxdrawuni{2500}\pmboxdrawuni{2500}\pmboxdrawuni{2500}\pmboxdrawuni{2500}▶\pmboxdrawuni{2502}~BLOCKED~before~rules~are~~~~~\pmboxdrawuni{2502}}
\cl{~~~\pmboxdrawuni{2502}~(if~one~matches)~~~~~~\pmboxdrawuni{2502}~~~~~~~~~~~~~~~\pmboxdrawuni{2502}~evaluated~—~overrides~allow~~\pmboxdrawuni{2502}}
\cl{~~~\pmboxdrawuni{2514}\pmboxdrawuni{2500}\pmboxdrawuni{2500}\pmboxdrawuni{2500}\pmboxdrawuni{2500}\pmboxdrawuni{2500}\pmboxdrawuni{2500}\pmboxdrawuni{2500}\pmboxdrawuni{2500}\pmboxdrawuni{2500}\pmboxdrawuni{2500}\pmboxdrawuni{252C}\pmboxdrawuni{2500}\pmboxdrawuni{2500}\pmboxdrawuni{2500}\pmboxdrawuni{2500}\pmboxdrawuni{2500}\pmboxdrawuni{2500}\pmboxdrawuni{2500}\pmboxdrawuni{2500}\pmboxdrawuni{2500}\pmboxdrawuni{2500}\pmboxdrawuni{2500}\pmboxdrawuni{2500}\pmboxdrawuni{2518}~~~~~~~~~~~~~~~\pmboxdrawuni{2514}\pmboxdrawuni{2500}\pmboxdrawuni{2500}\pmboxdrawuni{2500}\pmboxdrawuni{2500}\pmboxdrawuni{2500}\pmboxdrawuni{2500}\pmboxdrawuni{2500}\pmboxdrawuni{2500}\pmboxdrawuni{2500}\pmboxdrawuni{2500}\pmboxdrawuni{2500}\pmboxdrawuni{2500}\pmboxdrawuni{2500}\pmboxdrawuni{2500}\pmboxdrawuni{2500}\pmboxdrawuni{2500}\pmboxdrawuni{2500}\pmboxdrawuni{2500}\pmboxdrawuni{2500}\pmboxdrawuni{2500}\pmboxdrawuni{2500}\pmboxdrawuni{2500}\pmboxdrawuni{2500}\pmboxdrawuni{2500}\pmboxdrawuni{2500}\pmboxdrawuni{2500}\pmboxdrawuni{2500}\pmboxdrawuni{2500}\pmboxdrawuni{2500}\pmboxdrawuni{2500}\pmboxdrawuni{2518}}
\cl{~~~~~~~~~~~~~~\pmboxdrawuni{2502}~exit~0,~or~a~JSON~decision}
\cl{~~~~~~~~~~~~~~▼}
\cl{~~~\pmboxdrawuni{250C}\pmboxdrawuni{2500}\pmboxdrawuni{2500}\pmboxdrawuni{2500}\pmboxdrawuni{2500}\pmboxdrawuni{2500}\pmboxdrawuni{2500}\pmboxdrawuni{2500}\pmboxdrawuni{2500}\pmboxdrawuni{2500}\pmboxdrawuni{2500}\pmboxdrawuni{2500}\pmboxdrawuni{2500}\pmboxdrawuni{2500}\pmboxdrawuni{2500}\pmboxdrawuni{2500}\pmboxdrawuni{2500}\pmboxdrawuni{2500}\pmboxdrawuni{2500}\pmboxdrawuni{2500}\pmboxdrawuni{2500}\pmboxdrawuni{2500}\pmboxdrawuni{2500}\pmboxdrawuni{2500}\pmboxdrawuni{2510}~~~match~~~~~~~\pmboxdrawuni{250C}\pmboxdrawuni{2500}\pmboxdrawuni{2500}\pmboxdrawuni{2500}\pmboxdrawuni{2500}\pmboxdrawuni{2500}\pmboxdrawuni{2500}\pmboxdrawuni{2500}\pmboxdrawuni{2500}\pmboxdrawuni{2500}\pmboxdrawuni{2500}\pmboxdrawuni{2500}\pmboxdrawuni{2500}\pmboxdrawuni{2500}\pmboxdrawuni{2500}\pmboxdrawuni{2500}\pmboxdrawuni{2500}\pmboxdrawuni{2500}\pmboxdrawuni{2500}\pmboxdrawuni{2500}\pmboxdrawuni{2500}\pmboxdrawuni{2500}\pmboxdrawuni{2500}\pmboxdrawuni{2500}\pmboxdrawuni{2500}\pmboxdrawuni{2500}\pmboxdrawuni{2500}\pmboxdrawuni{2500}\pmboxdrawuni{2500}\pmboxdrawuni{2500}\pmboxdrawuni{2500}\pmboxdrawuni{2510}}
\cl{~~~\pmboxdrawuni{2502}~DENY~rules~~~~~~~~~~~~\pmboxdrawuni{2502}\pmboxdrawuni{2500}\pmboxdrawuni{2500}\pmboxdrawuni{2500}\pmboxdrawuni{2500}\pmboxdrawuni{2500}\pmboxdrawuni{2500}\pmboxdrawuni{2500}\pmboxdrawuni{2500}\pmboxdrawuni{2500}\pmboxdrawuni{2500}\pmboxdrawuni{2500}\pmboxdrawuni{2500}\pmboxdrawuni{2500}\pmboxdrawuni{2500}▶\pmboxdrawuni{2502}~BLOCKED.~A~hook~returning~~~~\pmboxdrawuni{2502}}
\cl{~~~\pmboxdrawuni{2502}~(managed~first)~~~~~~~\pmboxdrawuni{2502}~~~~~~~~~~~~~~~\pmboxdrawuni{2502}~"allow"~does~not~override~it~\pmboxdrawuni{2502}}
\cl{~~~\pmboxdrawuni{2514}\pmboxdrawuni{2500}\pmboxdrawuni{2500}\pmboxdrawuni{2500}\pmboxdrawuni{2500}\pmboxdrawuni{2500}\pmboxdrawuni{2500}\pmboxdrawuni{2500}\pmboxdrawuni{2500}\pmboxdrawuni{2500}\pmboxdrawuni{2500}\pmboxdrawuni{252C}\pmboxdrawuni{2500}\pmboxdrawuni{2500}\pmboxdrawuni{2500}\pmboxdrawuni{2500}\pmboxdrawuni{2500}\pmboxdrawuni{2500}\pmboxdrawuni{2500}\pmboxdrawuni{2500}\pmboxdrawuni{2500}\pmboxdrawuni{2500}\pmboxdrawuni{2500}\pmboxdrawuni{2500}\pmboxdrawuni{2518}~~~~~~~~~~~~~~~\pmboxdrawuni{2514}\pmboxdrawuni{2500}\pmboxdrawuni{2500}\pmboxdrawuni{2500}\pmboxdrawuni{2500}\pmboxdrawuni{2500}\pmboxdrawuni{2500}\pmboxdrawuni{2500}\pmboxdrawuni{2500}\pmboxdrawuni{2500}\pmboxdrawuni{2500}\pmboxdrawuni{2500}\pmboxdrawuni{2500}\pmboxdrawuni{2500}\pmboxdrawuni{2500}\pmboxdrawuni{2500}\pmboxdrawuni{2500}\pmboxdrawuni{2500}\pmboxdrawuni{2500}\pmboxdrawuni{2500}\pmboxdrawuni{2500}\pmboxdrawuni{2500}\pmboxdrawuni{2500}\pmboxdrawuni{2500}\pmboxdrawuni{2500}\pmboxdrawuni{2500}\pmboxdrawuni{2500}\pmboxdrawuni{2500}\pmboxdrawuni{2500}\pmboxdrawuni{2500}\pmboxdrawuni{2500}\pmboxdrawuni{2518}}
\cl{~~~~~~~~~~~~~~\pmboxdrawuni{2502}~no~match}
\cl{~~~~~~~~~~~~~~▼}
\cl{~~~\pmboxdrawuni{250C}\pmboxdrawuni{2500}\pmboxdrawuni{2500}\pmboxdrawuni{2500}\pmboxdrawuni{2500}\pmboxdrawuni{2500}\pmboxdrawuni{2500}\pmboxdrawuni{2500}\pmboxdrawuni{2500}\pmboxdrawuni{2500}\pmboxdrawuni{2500}\pmboxdrawuni{2500}\pmboxdrawuni{2500}\pmboxdrawuni{2500}\pmboxdrawuni{2500}\pmboxdrawuni{2500}\pmboxdrawuni{2500}\pmboxdrawuni{2500}\pmboxdrawuni{2500}\pmboxdrawuni{2500}\pmboxdrawuni{2500}\pmboxdrawuni{2500}\pmboxdrawuni{2500}\pmboxdrawuni{2500}\pmboxdrawuni{2510}~~~match~~~~~~~\pmboxdrawuni{250C}\pmboxdrawuni{2500}\pmboxdrawuni{2500}\pmboxdrawuni{2500}\pmboxdrawuni{2500}\pmboxdrawuni{2500}\pmboxdrawuni{2500}\pmboxdrawuni{2500}\pmboxdrawuni{2500}\pmboxdrawuni{2500}\pmboxdrawuni{2500}\pmboxdrawuni{2500}\pmboxdrawuni{2500}\pmboxdrawuni{2500}\pmboxdrawuni{2500}\pmboxdrawuni{2500}\pmboxdrawuni{2500}\pmboxdrawuni{2500}\pmboxdrawuni{2500}\pmboxdrawuni{2500}\pmboxdrawuni{2500}\pmboxdrawuni{2500}\pmboxdrawuni{2500}\pmboxdrawuni{2500}\pmboxdrawuni{2500}\pmboxdrawuni{2500}\pmboxdrawuni{2500}\pmboxdrawuni{2500}\pmboxdrawuni{2500}\pmboxdrawuni{2500}\pmboxdrawuni{2500}\pmboxdrawuni{2510}}
\cl{~~~\pmboxdrawuni{2502}~ASK~rules~~~~~~~~~~~~~\pmboxdrawuni{2502}\pmboxdrawuni{2500}\pmboxdrawuni{2500}\pmboxdrawuni{2500}\pmboxdrawuni{2500}\pmboxdrawuni{2500}\pmboxdrawuni{2500}\pmboxdrawuni{2500}\pmboxdrawuni{2500}\pmboxdrawuni{2500}\pmboxdrawuni{2500}\pmboxdrawuni{2500}\pmboxdrawuni{2500}\pmboxdrawuni{2500}\pmboxdrawuni{2500}▶\pmboxdrawuni{2502}~PROMPT,~even~when~a~narrower~\pmboxdrawuni{2502}}
\cl{~~~\pmboxdrawuni{2514}\pmboxdrawuni{2500}\pmboxdrawuni{2500}\pmboxdrawuni{2500}\pmboxdrawuni{2500}\pmboxdrawuni{2500}\pmboxdrawuni{2500}\pmboxdrawuni{2500}\pmboxdrawuni{2500}\pmboxdrawuni{2500}\pmboxdrawuni{2500}\pmboxdrawuni{252C}\pmboxdrawuni{2500}\pmboxdrawuni{2500}\pmboxdrawuni{2500}\pmboxdrawuni{2500}\pmboxdrawuni{2500}\pmboxdrawuni{2500}\pmboxdrawuni{2500}\pmboxdrawuni{2500}\pmboxdrawuni{2500}\pmboxdrawuni{2500}\pmboxdrawuni{2500}\pmboxdrawuni{2500}\pmboxdrawuni{2518}~~~~~~~~~~~~~~~\pmboxdrawuni{2502}~allow~rule~also~matches~~~~~~\pmboxdrawuni{2502}}
\cl{~~~~~~~~~~~~~~\pmboxdrawuni{2502}~no~match~~~~~~~~~~~~~~~~~~~\pmboxdrawuni{2514}\pmboxdrawuni{2500}\pmboxdrawuni{2500}\pmboxdrawuni{2500}\pmboxdrawuni{2500}\pmboxdrawuni{2500}\pmboxdrawuni{2500}\pmboxdrawuni{2500}\pmboxdrawuni{2500}\pmboxdrawuni{2500}\pmboxdrawuni{2500}\pmboxdrawuni{2500}\pmboxdrawuni{2500}\pmboxdrawuni{2500}\pmboxdrawuni{2500}\pmboxdrawuni{2500}\pmboxdrawuni{2500}\pmboxdrawuni{2500}\pmboxdrawuni{2500}\pmboxdrawuni{2500}\pmboxdrawuni{2500}\pmboxdrawuni{2500}\pmboxdrawuni{2500}\pmboxdrawuni{2500}\pmboxdrawuni{2500}\pmboxdrawuni{2500}\pmboxdrawuni{2500}\pmboxdrawuni{2500}\pmboxdrawuni{2500}\pmboxdrawuni{2500}\pmboxdrawuni{2500}\pmboxdrawuni{2518}}
\cl{~~~~~~~~~~~~~~▼}
\cl{~~~\pmboxdrawuni{250C}\pmboxdrawuni{2500}\pmboxdrawuni{2500}\pmboxdrawuni{2500}\pmboxdrawuni{2500}\pmboxdrawuni{2500}\pmboxdrawuni{2500}\pmboxdrawuni{2500}\pmboxdrawuni{2500}\pmboxdrawuni{2500}\pmboxdrawuni{2500}\pmboxdrawuni{2500}\pmboxdrawuni{2500}\pmboxdrawuni{2500}\pmboxdrawuni{2500}\pmboxdrawuni{2500}\pmboxdrawuni{2500}\pmboxdrawuni{2500}\pmboxdrawuni{2500}\pmboxdrawuni{2500}\pmboxdrawuni{2500}\pmboxdrawuni{2500}\pmboxdrawuni{2500}\pmboxdrawuni{2500}\pmboxdrawuni{2510}~~~match~~~~~~~\pmboxdrawuni{250C}\pmboxdrawuni{2500}\pmboxdrawuni{2500}\pmboxdrawuni{2500}\pmboxdrawuni{2500}\pmboxdrawuni{2500}\pmboxdrawuni{2500}\pmboxdrawuni{2500}\pmboxdrawuni{2500}\pmboxdrawuni{2500}\pmboxdrawuni{2500}\pmboxdrawuni{2500}\pmboxdrawuni{2500}\pmboxdrawuni{2500}\pmboxdrawuni{2500}\pmboxdrawuni{2500}\pmboxdrawuni{2500}\pmboxdrawuni{2500}\pmboxdrawuni{2500}\pmboxdrawuni{2500}\pmboxdrawuni{2500}\pmboxdrawuni{2500}\pmboxdrawuni{2500}\pmboxdrawuni{2500}\pmboxdrawuni{2500}\pmboxdrawuni{2500}\pmboxdrawuni{2500}\pmboxdrawuni{2500}\pmboxdrawuni{2500}\pmboxdrawuni{2500}\pmboxdrawuni{2500}\pmboxdrawuni{2510}}
\cl{~~~\pmboxdrawuni{2502}~ALLOW~rules~~~~~~~~~~~\pmboxdrawuni{2502}\pmboxdrawuni{2500}\pmboxdrawuni{2500}\pmboxdrawuni{2500}\pmboxdrawuni{2500}\pmboxdrawuni{2500}\pmboxdrawuni{2500}\pmboxdrawuni{2500}\pmboxdrawuni{2500}\pmboxdrawuni{2500}\pmboxdrawuni{2500}\pmboxdrawuni{2500}\pmboxdrawuni{2500}\pmboxdrawuni{2500}\pmboxdrawuni{2500}▶\pmboxdrawuni{2502}~RUNS~without~a~prompt~~~~~~~~\pmboxdrawuni{2502}}
\cl{~~~\pmboxdrawuni{2514}\pmboxdrawuni{2500}\pmboxdrawuni{2500}\pmboxdrawuni{2500}\pmboxdrawuni{2500}\pmboxdrawuni{2500}\pmboxdrawuni{2500}\pmboxdrawuni{2500}\pmboxdrawuni{2500}\pmboxdrawuni{2500}\pmboxdrawuni{2500}\pmboxdrawuni{252C}\pmboxdrawuni{2500}\pmboxdrawuni{2500}\pmboxdrawuni{2500}\pmboxdrawuni{2500}\pmboxdrawuni{2500}\pmboxdrawuni{2500}\pmboxdrawuni{2500}\pmboxdrawuni{2500}\pmboxdrawuni{2500}\pmboxdrawuni{2500}\pmboxdrawuni{2500}\pmboxdrawuni{2500}\pmboxdrawuni{2518}~~~~~~~~~~~~~~~\pmboxdrawuni{2514}\pmboxdrawuni{2500}\pmboxdrawuni{2500}\pmboxdrawuni{2500}\pmboxdrawuni{2500}\pmboxdrawuni{2500}\pmboxdrawuni{2500}\pmboxdrawuni{2500}\pmboxdrawuni{2500}\pmboxdrawuni{2500}\pmboxdrawuni{2500}\pmboxdrawuni{2500}\pmboxdrawuni{2500}\pmboxdrawuni{2500}\pmboxdrawuni{2500}\pmboxdrawuni{2500}\pmboxdrawuni{2500}\pmboxdrawuni{2500}\pmboxdrawuni{2500}\pmboxdrawuni{2500}\pmboxdrawuni{2500}\pmboxdrawuni{2500}\pmboxdrawuni{2500}\pmboxdrawuni{2500}\pmboxdrawuni{2500}\pmboxdrawuni{2500}\pmboxdrawuni{2500}\pmboxdrawuni{2500}\pmboxdrawuni{2500}\pmboxdrawuni{2500}\pmboxdrawuni{2500}\pmboxdrawuni{2518}}
\cl{~~~~~~~~~~~~~~\pmboxdrawuni{2502}~no~match}
\cl{~~~~~~~~~~~~~~▼}
\cl{~~~\pmboxdrawuni{250C}\pmboxdrawuni{2500}\pmboxdrawuni{2500}\pmboxdrawuni{2500}\pmboxdrawuni{2500}\pmboxdrawuni{2500}\pmboxdrawuni{2500}\pmboxdrawuni{2500}\pmboxdrawuni{2500}\pmboxdrawuni{2500}\pmboxdrawuni{2500}\pmboxdrawuni{2500}\pmboxdrawuni{2500}\pmboxdrawuni{2500}\pmboxdrawuni{2500}\pmboxdrawuni{2500}\pmboxdrawuni{2500}\pmboxdrawuni{2500}\pmboxdrawuni{2500}\pmboxdrawuni{2500}\pmboxdrawuni{2500}\pmboxdrawuni{2500}\pmboxdrawuni{2500}\pmboxdrawuni{2500}\pmboxdrawuni{2500}\pmboxdrawuni{2500}\pmboxdrawuni{2500}\pmboxdrawuni{2500}\pmboxdrawuni{2500}\pmboxdrawuni{2500}\pmboxdrawuni{2500}\pmboxdrawuni{2500}\pmboxdrawuni{2500}\pmboxdrawuni{2500}\pmboxdrawuni{2500}\pmboxdrawuni{2500}\pmboxdrawuni{2500}\pmboxdrawuni{2500}\pmboxdrawuni{2500}\pmboxdrawuni{2500}\pmboxdrawuni{2500}\pmboxdrawuni{2500}\pmboxdrawuni{2500}\pmboxdrawuni{2500}\pmboxdrawuni{2500}\pmboxdrawuni{2500}\pmboxdrawuni{2500}\pmboxdrawuni{2500}\pmboxdrawuni{2500}\pmboxdrawuni{2500}\pmboxdrawuni{2500}\pmboxdrawuni{2500}\pmboxdrawuni{2500}\pmboxdrawuni{2500}\pmboxdrawuni{2500}\pmboxdrawuni{2500}\pmboxdrawuni{2500}\pmboxdrawuni{2500}\pmboxdrawuni{2500}\pmboxdrawuni{2500}\pmboxdrawuni{2500}\pmboxdrawuni{2500}\pmboxdrawuni{2500}\pmboxdrawuni{2500}\pmboxdrawuni{2500}\pmboxdrawuni{2500}\pmboxdrawuni{2500}\pmboxdrawuni{2500}\pmboxdrawuni{2500}\pmboxdrawuni{2510}}
\cl{~~~\pmboxdrawuni{2502}~PERMISSION~MODE~decides:~Manual~asks~·~auto~sends~it~to~the~~~~~~~~\pmboxdrawuni{2502}}
\cl{~~~\pmboxdrawuni{2502}~classifier~·~acceptEdits~auto-approves~edits~·~dontAsk~denies~·~~~~\pmboxdrawuni{2502}}
\cl{~~~\pmboxdrawuni{2502}~bypassPermissions~skips~the~prompt~entirely~~~~~~~~~~~~~~~~~~~~~~~~\pmboxdrawuni{2502}}
\cl{~~~\pmboxdrawuni{2514}\pmboxdrawuni{2500}\pmboxdrawuni{2500}\pmboxdrawuni{2500}\pmboxdrawuni{2500}\pmboxdrawuni{2500}\pmboxdrawuni{2500}\pmboxdrawuni{2500}\pmboxdrawuni{2500}\pmboxdrawuni{2500}\pmboxdrawuni{2500}\pmboxdrawuni{2500}\pmboxdrawuni{2500}\pmboxdrawuni{2500}\pmboxdrawuni{2500}\pmboxdrawuni{2500}\pmboxdrawuni{2500}\pmboxdrawuni{2500}\pmboxdrawuni{2500}\pmboxdrawuni{2500}\pmboxdrawuni{2500}\pmboxdrawuni{2500}\pmboxdrawuni{2500}\pmboxdrawuni{2500}\pmboxdrawuni{2500}\pmboxdrawuni{2500}\pmboxdrawuni{2500}\pmboxdrawuni{2500}\pmboxdrawuni{2500}\pmboxdrawuni{2500}\pmboxdrawuni{2500}\pmboxdrawuni{2500}\pmboxdrawuni{2500}\pmboxdrawuni{2500}\pmboxdrawuni{2500}\pmboxdrawuni{2500}\pmboxdrawuni{2500}\pmboxdrawuni{2500}\pmboxdrawuni{2500}\pmboxdrawuni{2500}\pmboxdrawuni{2500}\pmboxdrawuni{2500}\pmboxdrawuni{2500}\pmboxdrawuni{2500}\pmboxdrawuni{2500}\pmboxdrawuni{2500}\pmboxdrawuni{2500}\pmboxdrawuni{2500}\pmboxdrawuni{2500}\pmboxdrawuni{2500}\pmboxdrawuni{2500}\pmboxdrawuni{2500}\pmboxdrawuni{2500}\pmboxdrawuni{2500}\pmboxdrawuni{2500}\pmboxdrawuni{2500}\pmboxdrawuni{2500}\pmboxdrawuni{2500}\pmboxdrawuni{2500}\pmboxdrawuni{2500}\pmboxdrawuni{2500}\pmboxdrawuni{2500}\pmboxdrawuni{2500}\pmboxdrawuni{2500}\pmboxdrawuni{2500}\pmboxdrawuni{2500}\pmboxdrawuni{2500}\pmboxdrawuni{2500}\pmboxdrawuni{2500}\pmboxdrawuni{2518}}
\cl{~~~~~~~~~~~~~~\pmboxdrawuni{2502}~approved}
\cl{~~~~~~~~~~~~~~▼}
\cl{~~~\pmboxdrawuni{250C}\pmboxdrawuni{2500}\pmboxdrawuni{2500}\pmboxdrawuni{2500}\pmboxdrawuni{2500}\pmboxdrawuni{2500}\pmboxdrawuni{2500}\pmboxdrawuni{2500}\pmboxdrawuni{2500}\pmboxdrawuni{2500}\pmboxdrawuni{2500}\pmboxdrawuni{2500}\pmboxdrawuni{2500}\pmboxdrawuni{2500}\pmboxdrawuni{2500}\pmboxdrawuni{2500}\pmboxdrawuni{2500}\pmboxdrawuni{2500}\pmboxdrawuni{2500}\pmboxdrawuni{2500}\pmboxdrawuni{2500}\pmboxdrawuni{2500}\pmboxdrawuni{2500}\pmboxdrawuni{2500}\pmboxdrawuni{2500}\pmboxdrawuni{2500}\pmboxdrawuni{2500}\pmboxdrawuni{2500}\pmboxdrawuni{2500}\pmboxdrawuni{2500}\pmboxdrawuni{2500}\pmboxdrawuni{2500}\pmboxdrawuni{2500}\pmboxdrawuni{2500}\pmboxdrawuni{2500}\pmboxdrawuni{2500}\pmboxdrawuni{2500}\pmboxdrawuni{2500}\pmboxdrawuni{2500}\pmboxdrawuni{2500}\pmboxdrawuni{2500}\pmboxdrawuni{2500}\pmboxdrawuni{2500}\pmboxdrawuni{2500}\pmboxdrawuni{2500}\pmboxdrawuni{2500}\pmboxdrawuni{2500}\pmboxdrawuni{2500}\pmboxdrawuni{2500}\pmboxdrawuni{2500}\pmboxdrawuni{2500}\pmboxdrawuni{2500}\pmboxdrawuni{2500}\pmboxdrawuni{2500}\pmboxdrawuni{2500}\pmboxdrawuni{2500}\pmboxdrawuni{2500}\pmboxdrawuni{2500}\pmboxdrawuni{2500}\pmboxdrawuni{2500}\pmboxdrawuni{2500}\pmboxdrawuni{2500}\pmboxdrawuni{2500}\pmboxdrawuni{2500}\pmboxdrawuni{2500}\pmboxdrawuni{2500}\pmboxdrawuni{2500}\pmboxdrawuni{2500}\pmboxdrawuni{2500}\pmboxdrawuni{2510}}
\cl{~~~\pmboxdrawuni{2502}~SANDBOX~(if~enabled)~enforces~filesystem~and~network~at~the~OS~~~~~\pmboxdrawuni{2502}}
\cl{~~~\pmboxdrawuni{2514}\pmboxdrawuni{2500}\pmboxdrawuni{2500}\pmboxdrawuni{2500}\pmboxdrawuni{2500}\pmboxdrawuni{2500}\pmboxdrawuni{2500}\pmboxdrawuni{2500}\pmboxdrawuni{2500}\pmboxdrawuni{2500}\pmboxdrawuni{2500}\pmboxdrawuni{2500}\pmboxdrawuni{2500}\pmboxdrawuni{2500}\pmboxdrawuni{2500}\pmboxdrawuni{2500}\pmboxdrawuni{2500}\pmboxdrawuni{2500}\pmboxdrawuni{2500}\pmboxdrawuni{2500}\pmboxdrawuni{2500}\pmboxdrawuni{2500}\pmboxdrawuni{2500}\pmboxdrawuni{2500}\pmboxdrawuni{2500}\pmboxdrawuni{2500}\pmboxdrawuni{2500}\pmboxdrawuni{2500}\pmboxdrawuni{2500}\pmboxdrawuni{2500}\pmboxdrawuni{2500}\pmboxdrawuni{2500}\pmboxdrawuni{2500}\pmboxdrawuni{2500}\pmboxdrawuni{2500}\pmboxdrawuni{2500}\pmboxdrawuni{2500}\pmboxdrawuni{2500}\pmboxdrawuni{2500}\pmboxdrawuni{2500}\pmboxdrawuni{2500}\pmboxdrawuni{2500}\pmboxdrawuni{2500}\pmboxdrawuni{2500}\pmboxdrawuni{2500}\pmboxdrawuni{2500}\pmboxdrawuni{2500}\pmboxdrawuni{2500}\pmboxdrawuni{2500}\pmboxdrawuni{2500}\pmboxdrawuni{2500}\pmboxdrawuni{2500}\pmboxdrawuni{2500}\pmboxdrawuni{2500}\pmboxdrawuni{2500}\pmboxdrawuni{2500}\pmboxdrawuni{2500}\pmboxdrawuni{2500}\pmboxdrawuni{2500}\pmboxdrawuni{2500}\pmboxdrawuni{2500}\pmboxdrawuni{2500}\pmboxdrawuni{2500}\pmboxdrawuni{2500}\pmboxdrawuni{2500}\pmboxdrawuni{2500}\pmboxdrawuni{2500}\pmboxdrawuni{2500}\pmboxdrawuni{2500}\pmboxdrawuni{2518}}
\end{codefig}
\figcaption{Figure 2 — How a tool call is decided. Deny is evaluated first and cannot be widened from below; a blocking hook short-circuits the whole chain.}
\end{figure}


\FloatBarrier
\setcounter{section}{2}
\section{Understand a prompt before approving}

On supported Bash or PowerShell permission prompts, \texttt{Ctrl+E} asks for an explanation of the command and its risk without executing it. Deny is also an answer, and denying with a comment proposing a safer alternative is often the fastest route forward.


Before approving anything you did not write, ask:


\begin{itemize}
\item What exact files, account or host does it touch?
\item Is the target resolved, or does it rely on a broad wildcard?
\item Is the operation reversible?
\item Can it place a credential into output?
\item Can the same information be obtained read-only first?
\end{itemize}

Claude Code applies its own protections here — command-injection detection, fail-closed matching for unmatched commands in Manual mode, natural-language descriptions of complex commands, and a separate context window for web fetches so that fetched content cannot be injected directly into the conversation [53]. None of them removes the need for the five questions.


\FloatBarrier
\setcounter{section}{3}
\section{Sandboxing}

Run \texttt{/\allowbreak{}sandbox} to open the sandbox panel [27]. The sandbox constrains filesystem and network access \textbf{below} the model layer, enforced by the operating system for every Bash command and its child processes. It runs on macOS (via Seatbelt), Linux and WSL2; native Windows is not supported [27].


The panel writes to \texttt{.\allowbreak{}claude/\allowbreak{}settings.\allowbreak{}local.\allowbreak{}json}; set \texttt{sandbox.\allowbreak{}enabled} in user settings to enable it everywhere, or deploy it through managed settings for an organisation [27]. The controls worth knowing:


\begin{itemize}
\item \texttt{sandbox.\allowbreak{}filesystem.\allowbreak{}deny\allowbreak{}Write}, \texttt{deny\allowbreak{}Read} and \texttt{allow\allowbreak{}Read}. Where read rules overlap, the more specific path wins, and \textbf{a deny holds inside a wider allow} — so \texttt{allow\allowbreak{}Read:\allowbreak{} ["\textasciitilde{}/\allowbreak{}"]} with \texttt{deny\allowbreak{}Read:\allowbreak{} ["\textasciitilde{}/\allowbreak{}**/\allowbreak{}.\allowbreak{}env"]} leaves every \texttt{.\allowbreak{}env} blocked while the rest of the home directory is readable [27].
\item \texttt{sandbox.\allowbreak{}credentials.\allowbreak{}files} and \texttt{sandbox.\allowbreak{}credentials.\allowbreak{}env\allowbreak{}Vars}, each entry \texttt{deny} or \texttt{mask}. Denied files are unreadable inside the sandbox; denied environment variables are unset before each sandboxed command runs. Masking lets a command receive a value while redacting it from tool output [27].
\item \texttt{auto\allowbreak{}Allow\allowbreak{}Bash\allowbreak{}If\allowbreak{}Sandboxed} defaults to \texttt{true}, so sandboxed commands run without prompts; set it to \texttt{false} if you want prompts anyway [27].
\item Deny entries from every settings scope are merged, and a deny only ever narrows access: any scope can add one and no scope can remove one another scope added [27].
\end{itemize}

Sandboxing complements permissions. Neither is sufficient alone, and for a stronger boundary — untrusted code, unfamiliar dependencies, anything you would not run on your laptop — compare the options in [65] and consider a dev container [66].


\FloatBarrier
\setcounter{section}{4}
\section{Checkpoints and rewind}

Claude Code records checkpoints before file edits, and \texttt{/\allowbreak{}rewind} restores conversation, code or both [67]. It is useful. It is not a backup, and the gaps are systematic:


\begin{itemize}
\item external actions are not reversible by a local rewind;
\item changes made through Bash commands or external programs are not tracked as file-edit checkpoints;
\item most edits made by subagents are not in the parent's checkpoint history;
\item writes through symlinks or hard links may fall outside coverage;
\item untracked or ignored files need separate care;
\item browser actions, messages, published artifacts and pushed commits cannot be un-sent.
\end{itemize}

Use Git for durable history before you enable anything in Part V. Chapter 8 is not optional preparation.


\subsection*{Checkpoint}

\begin{itemize}
\item Sensitive paths carry deny rules, and you have tested that they actually block.
\item External writes require ask, or a human gate.
\item You know whether sandboxing is active and what it covers.
\item You can display the current diff and the checkpoint state.
\item \texttt{bypass\allowbreak{}Permissions} is not in use outside an isolated environment.
\end{itemize}

\begingroup
\def\kprows{%
\item Rules are evaluated deny, then ask, then allow. Specificity does not change that order, and a deny cannot carry exceptions.
\item A bare tool name in \texttt{deny} removes the tool from Claude's context; a scoped rule leaves it available and blocks matching calls.
\item Sandboxing constrains filesystem and network below the model layer, and a deny holds inside a wider allow.
\item Checkpoints are not a backup: they do not cover Bash-driven changes, most subagent edits, or anything already sent externally.
}%
\def\kpbody{\begin{minipage}{\textwidth}\subsection*{Key points}\begin{itemize}\kprows\end{itemize}\end{minipage}}%
\begingroup
\def\sloppy{\tolerance 9999\emergencystretch 3em\hfuzz 200pt\vfuzz 200pt}%
\hbadness=10000\vbadness=10000\hfuzz=200pt\vfuzz=200pt
\global\setbox\kpbox=\hbox{\kpbody}%
\endgroup
\par\addvspace{4.2mm}
\ifdim\dimexpr\ht\kpbox+\dp\kpbox\relax>0.30\textheight
  \typeout{HANDBOOK-KEYPOINTS broken \the\dimexpr\ht\kpbox+\dp\kpbox\relax}%
  \subsection*{Key points}
  \begin{itemize}\kprows\end{itemize}
\else
  \typeout{HANDBOOK-KEYPOINTS atomic \the\dimexpr\ht\kpbox+\dp\kpbox\relax}%
  \noindent\kpbody
\fi
\par\addvspace{1.4mm}
\endgroup

\FloatBarrier
\renewcommand{\chaptertitlelabel}{Part II \textperiodcentered\ Chapter 8}
\setcounter{chapter}{7}
\chapter{Protect work with Git and GitHub}

Git records local history. GitHub hosts a repository remotely and adds collaboration, review and access control. They are related but not interchangeable: a commit can exist entirely on your machine, and that is often exactly what you want while an agent is working.


\FloatBarrier
\setcounter{section}{0}
\section{A safe beginner sequence}

\begin{codeblock}{9.0}{10.6}
\cl{git~init}
\cl{git~status}
\end{codeblock}

Create \texttt{.\allowbreak{}gitignore} \textbf{before} the first commit:


\begin{codeblock}{9.0}{10.6}
\cl{.env}
\cl{.env.*}
\cl{!.env.example}
\cl{.claude/settings.local.json}
\cl{CLAUDE.local.md}
\cl{secrets/}
\cl{*.key}
\cl{*.pem}
\end{codeblock}

That pattern is a starting point, not a secret detector. Inspect before committing:


\begin{codeblock}{9.0}{10.6}
\cl{git~status}
\cl{git~diff~--staged}
\cl{git~add~.}
\cl{git~commit~-m~"Create~research~brief~workspace"}
\end{codeblock}

\FloatBarrier
\setcounter{section}{1}
\section{Add a remote deliberately}

\begin{codeblock}{9.0}{10.6}
\cl{gh~auth~login}
\cl{gh~repo~create~research-brief-workspace~--private~--source=.~--remote=origin}
\cl{git~remote~-v}
\end{codeblock}

Confirm owner, repository name, \textbf{visibility} and remote URL, then push only after a second secret and diff review:


\begin{codeblock}{9.0}{10.6}
\cl{git~push~-u~origin~main}
\end{codeblock}

If the branch is not \texttt{main}, use its real name. For collaborative work prefer a feature branch and a pull request over direct pushes to a default branch. Claude can propose all of these commands; account creation, authentication, visibility and release approval belong to you.


\FloatBarrier
\setcounter{section}{2}
\section{An agent-assisted Git workflow}

\begin{enumerate}
\item Claude summarises the intended change.
\item You approve the local edit plan.
\item Claude edits and runs verification.
\item You inspect \texttt{git diff}.
\item Claude proposes a commit message grounded in the diff.
\item You approve the commit.
\item Push or pull request is a \textbf{separate}, human-gated action.
\end{enumerate}

\begin{calloutbox}{palebrass}{brassdark}{2.0mm}
\textbf{CAUTION:} an automatic "commit and push at session end" hook can publish secrets, broken work, generated noise, or changes from the wrong branch. If you must automate, use a dedicated branch, a protected default branch, scoped credentials, pre-push checks and failure reporting (Chapter 17).
\end{calloutbox}

\FloatBarrier
\setcounter{section}{3}
\section{Automated review, and its limits}

\texttt{/\allowbreak{}security-\allowbreak{}review} runs a security pass over the changes on the current branch, and \texttt{/\allowbreak{}code-\allowbreak{}review} (alias \texttt{/\allowbreak{}review}) reviews a diff or pull request for bugs and cleanup [35], [68]. The \texttt{security-\allowbreak{}guidance} plugin has Claude review its own changes for vulnerabilities in the same session, and the Claude Security plugin scans a codebase and turns findings into patches you review [69], [70]. For a deeper pass, \texttt{/\allowbreak{}code-\allowbreak{}review ultr\allowbreak{}a} runs a multi-agent review in the cloud [71].


None of these replaces a human reviewer or a specialist assessment, and none should be a release gate on its own [53].


\FloatBarrier
\setcounter{section}{4}
\section{Recovery is broader than the repository}

Git does not cover untracked files, external databases and services, secret-manager values, browser actions, published artifacts, or local configuration excluded by \texttt{.\allowbreak{}gitignore}. Plan recovery for the workflow, not only for the source tree.


\subsection*{Exercise}

Initialise Git in the practice workspace. Ask Claude to draft a conservative \texttt{.\allowbreak{}gitignore} and to explain every pattern. Show the staged diff. Commit only after confirming that no sensitive or unrelated file is included.


\emph{A worked solution is given in Appendix M.}


\begingroup
\def\kprows{%
\item Create \texttt{.\allowbreak{}gitignore} before the first commit, and inspect \texttt{git diff -\allowbreak{}-\allowbreak{}staged} before every one.
\item Push is a separate, human-gated action — never a step inside an agent's plan.
\item An automatic commit-and-push hook can publish secrets, broken work or the wrong branch.
\item Git does not cover untracked files, external services, secret managers, browser actions or published artifacts. Plan recovery for the workflow.
}%
\def\kpbody{\begin{minipage}{\textwidth}\subsection*{Key points}\begin{itemize}\kprows\end{itemize}\end{minipage}}%
\begingroup
\def\sloppy{\tolerance 9999\emergencystretch 3em\hfuzz 200pt\vfuzz 200pt}%
\hbadness=10000\vbadness=10000\hfuzz=200pt\vfuzz=200pt
\global\setbox\kpbox=\hbox{\kpbody}%
\endgroup
\par\addvspace{4.2mm}
\ifdim\dimexpr\ht\kpbox+\dp\kpbox\relax>0.30\textheight
  \typeout{HANDBOOK-KEYPOINTS broken \the\dimexpr\ht\kpbox+\dp\kpbox\relax}%
  \subsection*{Key points}
  \begin{itemize}\kprows\end{itemize}
\else
  \typeout{HANDBOOK-KEYPOINTS atomic \the\dimexpr\ht\kpbox+\dp\kpbox\relax}%
  \noindent\kpbody
\fi
\par\addvspace{1.4mm}
\endgroup

\breakrule

\FloatBarrier
\parttitle{Part III}{Project knowledge and context}
\addcontentsline{toc}{part}{Part III \textemdash\ Project knowledge and context}

\FloatBarrier
\renewcommand{\chaptertitlelabel}{Part III \textperiodcentered\ Chapter 9}
\setcounter{chapter}{8}
\chapter{Build an effective \texttt{CLAUDE.\allowbreak{}md}}

\texttt{CLAUDE.\allowbreak{}md} is durable project guidance loaded at the start of every session. It is delivered as a user message after the system prompt, not as part of the system prompt, so Claude reads it and tries to follow it but there is no guarantee of compliance — especially for vague or conflicting instructions [26]. An instruction that \emph{must} hold belongs in a permission rule (Chapter 7) or a hook (Chapter 17).


\FloatBarrier
\setcounter{section}{0}
\section{Where instructions live}

\begingroup
\def\tblrows{%
Managed policy & macOS \texttt{/\allowbreak{}Library/\allowbreak{}Application\allowbreak{} Support/\allowbreak{}Claude\allowbreak{}Code/\allowbreak{}CLAUDE.\allowbreak{}md}; Linux and WSL \texttt{/\allowbreak{}etc/\allowbreak{}claude-\allowbreak{}code/\allowbreak{}CLAUDE.\allowbreak{}md}; Windows \texttt{C:\allowbreak{}\textbackslash{}Program Fil\allowbreak{}es\textbackslash{}Claude\allowbreak{}Code\textbackslash{}CLAUDE.\allowbreak{}md}; or the \texttt{claude\allowbreak{}Md} key in managed settings & Organisation-wide standards, security policy, compliance reminders. Cannot be excluded by any user setting \\
User & \texttt{\textasciitilde{}/\allowbreak{}.\allowbreak{}claude/\allowbreak{}CLAUDE.\allowbreak{}md} & Personal defaults that genuinely apply everywhere \\
Project & \texttt{.\allowbreak{}/\allowbreak{}CLAUDE.\allowbreak{}md} or \texttt{.\allowbreak{}/\allowbreak{}.\allowbreak{}claude/\allowbreak{}CLAUDE.\allowbreak{}md} & Purpose, structure, commands, rules, verification. Shared through version control \\
Local project & \texttt{.\allowbreak{}/\allowbreak{}CLAUDE.\allowbreak{}local.\allowbreak{}md} & Private project preferences. Add to \texttt{.\allowbreak{}gitignore} \\
Subdirectory & \texttt{path/\allowbreak{}to/\allowbreak{}CLAUDE.\allowbreak{}md} & Instructions relevant only within that subtree; loaded on demand when Claude reads files there \\
}%
\def\tblbody{\begin{minipage}{\textwidth}\boxcaption{Table 6 — Instruction file locations, in load order from broadest to most specific}
{\footnotesize\begin{tabular}{H{24.6mm}H{59.4mm}L{55.8mm}}
\toprule
\textbf{Scope} & \textbf{Location} & \textbf{Appropriate content} \\
\midrule
\tblrows
\bottomrule\end{tabular}}\end{minipage}}%
\begingroup
\def\sloppy{\tolerance 9999\emergencystretch 3em\hfuzz 200pt\vfuzz 200pt}%
\hbadness=10000\vbadness=10000\hfuzz=200pt\vfuzz=200pt
\global\setbox\tblbox=\hbox{\tblbody}%
\endgroup
\par\addvspace{2.6mm}
\ifdim\dimexpr\ht\tblbox+\dp\tblbox\relax>0.55\textheight
  \typeout{HANDBOOK-TABLE broken \the\dimexpr\ht\tblbox+\dp\tblbox\relax}%
  \tabcaption{Table 6 — Instruction file locations, in load order from broadest to most specific}
  {\footnotesize\begin{longtable}{H{24.6mm}H{59.4mm}L{55.8mm}}
  \toprule
\textbf{Scope} & \textbf{Location} & \textbf{Appropriate content} \\
\midrule\endfirsthead
  \multicolumn{3}{@{}l@{}}{%
  \sffamily\footnotesize\itshape\color{inkgrey}Table 6 — Instruction file locations, in load order from broadest to most specific \textemdash\ continued}\\[1.2mm]
  \toprule
\textbf{Scope} & \textbf{Location} & \textbf{Appropriate content} \\
\midrule\endhead
  \bottomrule\endfoot
  \bottomrule\endlastfoot
  \tblrows
  \end{longtable}}%
\else
  \typeout{HANDBOOK-TABLE atomic \the\dimexpr\ht\tblbox+\dp\tblbox\relax}%
  \noindent\tblbody
\fi
\par\addvspace{2.6mm}
\endgroup

Files are \textbf{concatenated}, not overridden. Across the tree, content is ordered from the filesystem root down to the working directory, so instructions closest to where you launched Claude are read last; within a directory, \texttt{CLAUDE.\allowbreak{}local.\allowbreak{}md} is appended after \texttt{CLAUDE.\allowbreak{}md} [26].


Claude Code reads \texttt{CLAUDE.\allowbreak{}md}, not \texttt{AGENTS.\allowbreak{}md}. If a repository already uses \texttt{AGENTS.\allowbreak{}md}, import it — \texttt{@AGENTS.\allowbreak{}md} on the first line of \texttt{CLAUDE.\allowbreak{}md} — rather than duplicating it [26].


\FloatBarrier
\setcounter{section}{1}
\section{A useful project file}

\begin{codeblock}{9.0}{10.6}
\cl{\#~Research~Brief~Workspace}
\cl{}
\cl{\#\#~Purpose}
\cl{Produce~concise,~evidence-led~briefs~from~approved~sources.}
\cl{}
\cl{\#\#~Directory~map}
\cl{-~\textasciigrave{}sources/\textasciigrave{}:~immutable~source~material}
\cl{-~\textasciigrave{}notes/\textasciigrave{}:~working~notes~and~extraction}
\cl{-~\textasciigrave{}drafts/\textasciigrave{}:~unapproved~manuscripts}
\cl{-~\textasciigrave{}deliverables/\textasciigrave{}:~human-approved~outputs~only}
\cl{-~\textasciigrave{}references/\textasciigrave{}:~retrieval-on-demand~policies}
\cl{-~\textasciigrave{}templates/\textasciigrave{}:~reusable~structures}
\cl{}
\cl{\#\#~Operating~rules}
\cl{-~Never~modify~files~in~\textasciigrave{}sources/\textasciigrave{}.}
\cl{-~Separate~source~claims~from~analyst~inference.}
\cl{-~Record~one~atomic~claim~per~evidence-ledger~row.}
\cl{-~Do~not~publish,~send,~push,~or~overwrite~a~deliverable~without~approval.}
\cl{-~Never~read~or~print~credentials.}
\cl{}
\cl{\#\#~Verification~commands}
\cl{-~Structure~check:~\textasciigrave{}python3~tools/check\_structure.py~drafts/<file>.md\textasciigrave{}}
\cl{-~Link~check:~~~~~~\textasciigrave{}python3~tools/check\_links.py~drafts/<file>.md\textasciigrave{}}
\cl{}
\cl{\#\#~Definition~of~done}
\cl{-~All~material~claims~carry~direct~source~references.}
\cl{-~Contradictions~and~gaps~are~visible~rather~than~resolved~silently.}
\cl{-~Links~and~rendering~have~been~checked~and~the~output~inspected.}
\cl{-~A~human~has~approved~promotion~from~\textasciigrave{}drafts/\textasciigrave{}~to~\textasciigrave{}deliverables/\textasciigrave{}.}
\cl{}
\cl{\#\#~References}
\cl{-~Source-quality~rules:~read~\textasciigrave{}references/source-policy.md\textasciigrave{}.}
\cl{-~House~style:~read~\textasciigrave{}references/editorial-style.md\textasciigrave{}~only~when~drafting.}
\end{codeblock}

\FloatBarrier
\setcounter{section}{2}
\section{Initialise, then edit}

\texttt{/\allowbreak{}init} generates a starter \texttt{CLAUDE.\allowbreak{}md} by analysing the project; if one exists it suggests improvements rather than overwriting [26]. It also reads Cursor and Copilot rule files and incorporates the relevant parts. Setting \texttt{CLAUDE\_\allowbreak{}CODE\_\allowbreak{}NEW\_\allowbreak{}INIT=\allowbreak{}1} enables an interactive multi-phase flow that explores with a subagent, asks follow-up questions and presents a reviewable proposal before writing anything [26].


The course warns against \texttt{/\allowbreak{}init} because generated files could become verbose [42]. Both positions reconcile: generate a draft, then edit it hard.


\begin{enumerate}
\item Run \texttt{/\allowbreak{}init}.
\item Read every line.
\item Delete generic description, changelog, transient state and duplicated documentation.
\item Keep operational facts, conventions that differ from tool defaults, pitfalls, and verification commands.
\item Run \texttt{/\allowbreak{}memory} to inspect what is loaded, and \texttt{/\allowbreak{}context} to confirm it.
\end{enumerate}

\textbf{Target under 200 lines.} Longer files consume context and reduce adherence. A file over 4 MiB is skipped entirely [26]. \texttt{/\allowbreak{}doctor} proposes trims for a checked-in file, cutting content Claude can derive from the codebase and keeping rationale and conventions (2.1.206 or later) [26].


\FloatBarrier
\setcounter{section}{3}
\section{An index, not a warehouse}

Long reference material belongs in separate files with a retrieval rule, as in the \texttt{References} section above. Imports with \texttt{@path/\allowbreak{}to/\allowbreak{}file} are supported to a maximum depth of four hops, but \textbf{imported files load at launch and do not reduce context occupancy} — they help organisation, not cost [26]. Import parsing skips code spans and fences, so `\texttt{ }@README\texttt{ }` in backticks stays literal.


An import in a \emph{project} memory file whose path resolves outside the working directory is external, and the first time Claude Code encounters one it shows an approval dialog listing the files; declining disables them permanently [26]. That dialog is a supply-chain control: it exists because someone else can commit an import into a shared repository.


Block-level HTML comments are stripped before injection, so maintainer notes cost no context [26].


\FloatBarrier
\setcounter{section}{4}
\section{Modular rules}

Split large instruction sets into \texttt{.\allowbreak{}claude/\allowbreak{}rules/\allowbreak{}*.\allowbreak{}md}, discovered recursively [26]. A rule with no \texttt{paths} frontmatter loads at launch with the same priority as \texttt{.\allowbreak{}claude/\allowbreak{}CLAUDE.\allowbreak{}md}. A rule \emph{with} \texttt{paths} loads only when Claude reads a matching file:


\begin{codeblock}{9.0}{10.6}
\cl{---}
\cl{paths:}
\cl{~~-~"sources/**/*.md"}
\cl{~~-~"notes/**/*.\{md,txt\}"}
\cl{---}
\cl{}
\cl{\#~Source~handling}
\cl{}
\cl{-~Never~edit~a~file~under~\textasciigrave{}sources/\textasciigrave{}.}
\cl{-~Record~provenance~as~\textasciigrave{}source-id~·~retrieved-date~·~URL~or~physical~location\textasciigrave{}.}
\end{codeblock}

User rules in \texttt{\textasciitilde{}/\allowbreak{}.\allowbreak{}claude/\allowbreak{}rules/\allowbreak{}} load before project rules, giving project rules higher priority. Symlinks are supported and circular links are handled. A rule's whole \texttt{paths} list shares a budget of 1,000 expanded patterns and 4 MiB; a pattern that would exceed it is used unexpanded and matches nothing [26].


\FloatBarrier
\setcounter{section}{5}
\section{Maintenance trigger}

Review \texttt{CLAUDE.\allowbreak{}md} when Claude repeatedly searches for a documented command, when you correct the same rule twice, when structure changes, when two instructions conflict, or when \texttt{/\allowbreak{}context} shows disproportionate instruction overhead. Ask for a proposed diff with a reason for every addition and removal — never a blind "optimise this".


\textbf{Note on compaction.} Project-root \texttt{CLAUDE.\allowbreak{}md} survives \texttt{/\allowbreak{}compact}: Claude Code re-reads it from disk and re-injects it [26]. Nested files and path-scoped rules reload as Claude reads the files they apply to. An instruction that disappears after compaction was given only in conversation — which is the argument for writing it down.


\subsection*{Exercise}

Create the project \texttt{CLAUDE.\allowbreak{}md}, start a new session, and ask Claude to state the directory rules without searching every file. Then ask which instructions are stable and which belong in a path-scoped rule or a skill.


\emph{A worked solution is given in Appendix M.}


\begingroup
\def\kprows{%
\item \texttt{CLAUDE.\allowbreak{}md} is delivered as context, not as enforced configuration; something that must hold belongs in a rule or a hook.
\item Files are concatenated from the filesystem root down, with \texttt{CLAUDE.\allowbreak{}local.\allowbreak{}md} appended last at each level.
\item Target under 200 lines. Imports organise content but do not reduce context, because they load at launch.
\item Project-root \texttt{CLAUDE.\allowbreak{}md} survives \texttt{/\allowbreak{}compact}; an instruction that disappears was given only in conversation.
}%
\def\kpbody{\begin{minipage}{\textwidth}\subsection*{Key points}\begin{itemize}\kprows\end{itemize}\end{minipage}}%
\begingroup
\def\sloppy{\tolerance 9999\emergencystretch 3em\hfuzz 200pt\vfuzz 200pt}%
\hbadness=10000\vbadness=10000\hfuzz=200pt\vfuzz=200pt
\global\setbox\kpbox=\hbox{\kpbody}%
\endgroup
\par\addvspace{4.2mm}
\ifdim\dimexpr\ht\kpbox+\dp\kpbox\relax>0.30\textheight
  \typeout{HANDBOOK-KEYPOINTS broken \the\dimexpr\ht\kpbox+\dp\kpbox\relax}%
  \subsection*{Key points}
  \begin{itemize}\kprows\end{itemize}
\else
  \typeout{HANDBOOK-KEYPOINTS atomic \the\dimexpr\ht\kpbox+\dp\kpbox\relax}%
  \noindent\kpbody
\fi
\par\addvspace{1.4mm}
\endgroup


\FloatBarrier
\renewcommand{\chaptertitlelabel}{Part III \textperiodcentered\ Chapter 10}
\setcounter{chapter}{9}
\chapter{Global instructions and auto memory}

Persistent knowledge serves three different purposes, and conflating them produces bloated context and stale behaviour:


\begin{enumerate}
\item \textbf{Instructions} tell Claude how to behave — \texttt{CLAUDE.\allowbreak{}md}, rules.
\item \textbf{Reference files} store facts and procedures Claude retrieves on demand.
\item \textbf{Auto memory} stores concise notes Claude writes for itself.
\end{enumerate}

\FloatBarrier
\setcounter{section}{0}
\section{Global instructions}

Use \texttt{\textasciitilde{}/\allowbreak{}.\allowbreak{}claude/\allowbreak{}CLAUDE.\allowbreak{}md} sparingly. Suitable: your preferred working language, a universal safety gate, a stable output convention, a pointer to a user rules directory. Unsuitable: project-specific commands, one client's house style, temporary priorities, or a catalogue of every tool you might one day use.


\FloatBarrier
\setcounter{section}{1}
\section{Auto memory}

Auto memory is on by default. As it works, Claude saves four kinds of note, recorded as a \texttt{type} field in the memory file's frontmatter [26]:


\begin{itemize}
\item \texttt{user} — your role, expertise and working preferences;
\item \texttt{feedback} — corrections you gave and approaches you confirmed;
\item \texttt{project} — ongoing work, deadlines and decisions not derivable from the code or history;
\item \texttt{reference} — where to find information outside the project.
\end{itemize}

Claude skips anything derivable from the codebase and anything your \texttt{CLAUDE.\allowbreak{}md} already says.


\textbf{Storage.} Each project gets \texttt{\textasciitilde{}/\allowbreak{}.\allowbreak{}claude/\allowbreak{}projects/\allowbreak{}<project>/\allowbreak{}memory/\allowbreak{}}, where \texttt{<project>} is derived \textbf{from the git repository}, so every worktree and subdirectory of the same repository shares one memory directory [26]. \texttt{MEMORY.\allowbreak{}md} is the index; topic files hold detail and are read on demand. Relocate storage with \texttt{auto\allowbreak{}Memory\allowbreak{}Directory}, an absolute or \texttt{\textasciitilde{}/\allowbreak{}}-prefixed path, honoured from any settings scope — though a value set in a project's \texttt{.\allowbreak{}claude/\allowbreak{}settings.\allowbreak{}json} or \texttt{.\allowbreak{}claude/\allowbreak{}settings.\allowbreak{}local.\allowbreak{}json} is honoured only under the same workspace-trust rule that governs hooks in settings files [26]. Auto memory is machine-local and is not shared across machines or cloud environments.


\textbf{Loading limit.} The first 200 lines of \texttt{MEMORY.\allowbreak{}md}, or the first 25 KB, whichever comes first, are loaded at the start of every conversation; content beyond that is not loaded [26]. When the index approaches a limit Claude Code prompts Claude to shorten it; when it exceeds one, the write succeeds but Claude Code returns an error instructing Claude to rewrite the index, because everything past the limit is dropped on the next load.


\textbf{Disabling.} Toggle auto memory in \texttt{/\allowbreak{}memory}, which writes \texttt{auto\allowbreak{}Memory\allowbreak{}Enabled} to \texttt{\textasciitilde{}/\allowbreak{}.\allowbreak{}claude/\allowbreak{}settings.\allowbreak{}json}; set \texttt{\{"auto\allowbreak{}Memory\allowbreak{}Enabled":\allowbreak{} false\}} in a project's settings to disable it for that project only; or set \texttt{CLAUDE\_\allowbreak{}CODE\_\allowbreak{}DISABLE\_\allowbreak{}AUTO\_\allowbreak{}MEMORY=\allowbreak{}1} [26]. This corrects M06 in Table 1.


\begin{calloutbox}{palebrass}{brassdark}{2.0mm}
\textbf{CAUTION — retention.} Claude Code deletes old session transcripts after the \texttt{cleanup\allowbreak{}Period\allowbreak{}Days} retention period but \textbf{excludes the memory directory from that sweep}. \texttt{MEMORY.\allowbreak{}md} and its topic files persist until you or Claude edits or deletes them [26]. If your retention obligations assume a 30-day local sweep, auto memory is outside it. This corrects E5.
\end{calloutbox}

Treat auto memory as a convenience, not a specification. Claude may remember something situational, or generalise a preference too far. Stable project truth belongs in reviewed files under version control.


\FloatBarrier
\setcounter{section}{2}
\section{A decision table}

\begingroup
\def\tblrows{%
"Never publish without approval" & Permission policy, plus a \texttt{CLAUDE.\allowbreak{}md} reminder \\
Project directory map & Project \texttt{CLAUDE.\allowbreak{}md} \\
Editorial style guide & \texttt{references/\allowbreak{}editorial-\allowbreak{}style.\allowbreak{}md}, retrieved on demand \\
Rules that apply only to one file type & \texttt{.\allowbreak{}claude/\allowbreak{}rules/\allowbreak{}*.\allowbreak{}md} with \texttt{paths} frontmatter \\
A multi-step procedure & A skill (Chapter 15) \\
Current open tasks & \texttt{notes/\allowbreak{}task-\allowbreak{}state.\allowbreak{}md} or an issue tracker \\
A correction useful only this session & The conversation \\
A recurring preference Claude inferred & Auto memory, after your review \\
Full change history & Git \\
Something that must hold regardless of what Claude decides & A \texttt{Pre\allowbreak{}Tool\allowbreak{}Use} hook (Chapter 17) \\
}%
\def\tblbody{\begin{minipage}{\textwidth}\boxcaption{Table 7 — Where a given piece of knowledge belongs}
{\footnotesize\begin{tabular}{L{74.0mm}L{69.0mm}}
\toprule
\textbf{Information} & \textbf{Best home} \\
\midrule
\tblrows
\bottomrule\end{tabular}}\end{minipage}}%
\begingroup
\def\sloppy{\tolerance 9999\emergencystretch 3em\hfuzz 200pt\vfuzz 200pt}%
\hbadness=10000\vbadness=10000\hfuzz=200pt\vfuzz=200pt
\global\setbox\tblbox=\hbox{\tblbody}%
\endgroup
\par\addvspace{2.6mm}
\ifdim\dimexpr\ht\tblbox+\dp\tblbox\relax>0.55\textheight
  \typeout{HANDBOOK-TABLE broken \the\dimexpr\ht\tblbox+\dp\tblbox\relax}%
  \tabcaption{Table 7 — Where a given piece of knowledge belongs}
  {\footnotesize\begin{longtable}{L{74.0mm}L{69.0mm}}
  \toprule
\textbf{Information} & \textbf{Best home} \\
\midrule\endfirsthead
  \multicolumn{2}{@{}l@{}}{%
  \sffamily\footnotesize\itshape\color{inkgrey}Table 7 — Where a given piece of knowledge belongs \textemdash\ continued}\\[1.2mm]
  \toprule
\textbf{Information} & \textbf{Best home} \\
\midrule\endhead
  \bottomrule\endfoot
  \bottomrule\endlastfoot
  \tblrows
  \end{longtable}}%
\else
  \typeout{HANDBOOK-TABLE atomic \the\dimexpr\ht\tblbox+\dp\tblbox\relax}%
  \noindent\tblbody
\fi
\par\addvspace{2.6mm}
\endgroup

\FloatBarrier
\setcounter{section}{3}
\section{Memory audit}

Monthly for active workspaces, or after any surprising behaviour:


\begin{codeblock}{8.0}{9.4}
\cl{List~every~instruction~and~memory~file~affecting~this~project.~For~each,~state~its~scope,}
\cl{approximate~context~cost,~last~relevant~use,~and~any~conflict~or~stale~entry.~Propose~changes}
\cl{as~a~diff.~Do~not~apply~them.}
\end{codeblock}

Then verify the file list yourself with \texttt{/\allowbreak{}memory} and \texttt{/\allowbreak{}context}, and use \texttt{/\allowbreak{}doctor} or \texttt{/\allowbreak{}hooks} when something configured is not taking effect [72]. The \texttt{Instruction\allowbreak{}s\allowbreak{}Loaded} hook logs exactly which instruction files loaded, when and why — the definitive answer when path-scoped rules misbehave [73].


\begingroup
\def\kprows{%
\item Instructions, reference files and auto memory serve three different purposes; conflating them bloats context and staleness.
\item Auto memory is keyed to the git repository and shared across all its worktrees.
\item Only the first 200 lines or 25 KB of \texttt{MEMORY.\allowbreak{}md} load at session start.
\item Auto-memory files are excluded from the \texttt{cleanup\allowbreak{}Period\allowbreak{}Days} retention sweep and persist until deleted.
}%
\def\kpbody{\begin{minipage}{\textwidth}\subsection*{Key points}\begin{itemize}\kprows\end{itemize}\end{minipage}}%
\begingroup
\def\sloppy{\tolerance 9999\emergencystretch 3em\hfuzz 200pt\vfuzz 200pt}%
\hbadness=10000\vbadness=10000\hfuzz=200pt\vfuzz=200pt
\global\setbox\kpbox=\hbox{\kpbody}%
\endgroup
\par\addvspace{4.2mm}
\ifdim\dimexpr\ht\kpbox+\dp\kpbox\relax>0.30\textheight
  \typeout{HANDBOOK-KEYPOINTS broken \the\dimexpr\ht\kpbox+\dp\kpbox\relax}%
  \subsection*{Key points}
  \begin{itemize}\kprows\end{itemize}
\else
  \typeout{HANDBOOK-KEYPOINTS atomic \the\dimexpr\ht\kpbox+\dp\kpbox\relax}%
  \noindent\kpbody
\fi
\par\addvspace{1.4mm}
\endgroup

\FloatBarrier
\renewcommand{\chaptertitlelabel}{Part III \textperiodcentered\ Chapter 11}
\setcounter{chapter}{10}
\chapter{Engineer context deliberately}

Context is what the model has for the current turn: system prompt, tool descriptions, project instructions, attached files, tool results and conversation history. A file on disk is not "known" merely because Claude could read it.


\FloatBarrier
\setcounter{section}{0}
\section{Inspect the window}

\begin{codeblock}{9.0}{10.6}
\cl{/context}
\end{codeblock}

Use it to diagnose composition, not to chase a number. The course cites 200,000–250,000 tokens as a threshold some users avoid; that is practitioner experience, not a product limit [42]. Windows, models and auto-compaction behaviour change, and \texttt{/\allowbreak{}autocompact} sets the auto-compact window explicitly [35], [55].


\FloatBarrier
\setcounter{section}{1}
\section{Four ways to recover focus}

\textbf{1 — Persist durable state.} Findings, plans, decisions and open questions belong in project files. A conversation is a working surface, not a database.


\textbf{2 — Hand off to a clean session.}


\begin{codeblock}{8.0}{9.4}
\cl{Prepare~a~handoff~for~a~clean~session.~Include~the~objective,~completed~work,~approved}
\cl{decisions,~open~issues,~exact~file~references,~verification~already~performed,~and~the~next}
\cl{safe~action.~Exclude~conversational~history~that~does~not~affect~the~task.}
\end{codeblock}

\textbf{3 — \texttt{/\allowbreak{}clear}.} Starts fresh while retaining project instructions and memory. It also clears session-scoped scheduled tasks and any active goal, and it closes the Chrome tab group [21], [24], [74]. Save your handoff first. The previous conversation is not lost: resume it with \texttt{/\allowbreak{}resume}, or reach it from the rewind menu's previous-session entry in the same process [30].


\textbf{4 — \texttt{/\allowbreak{}compact [in\allowbreak{}structions]}.} Summarises to free context. Compaction is lossy; say what must survive:


\begin{codeblock}{8.0}{9.4}
\cl{/compact~Preserve~all~approved~decisions,~unresolved~contradictions,~file~paths,~verification}
\cl{results,~and~the~human~approval~gate.~Remove~discarded~options.}
\end{codeblock}

Read the resulting summary when the task is consequential. Note the cost: \texttt{/\allowbreak{}compact} forces one summarisation request over the full history and then invalidates the prompt cache [60]. Figure 3 sets out what each of these operations preserves and what it discards.


\begin{figure}[tbp]
\begin{codefig}{9.0}{10.6}
\cl{~~~~~~~~~~~~~~~~~~~~~~~~/compact~~~/clear~~~resume~~~/branch~~~fork~~~/rewind}
\cl{~~~\pmboxdrawuni{2500}\pmboxdrawuni{2500}\pmboxdrawuni{2500}\pmboxdrawuni{2500}\pmboxdrawuni{2500}\pmboxdrawuni{2500}\pmboxdrawuni{2500}\pmboxdrawuni{2500}\pmboxdrawuni{2500}\pmboxdrawuni{2500}\pmboxdrawuni{2500}\pmboxdrawuni{2500}\pmboxdrawuni{2500}\pmboxdrawuni{2500}\pmboxdrawuni{2500}\pmboxdrawuni{2500}\pmboxdrawuni{2500}\pmboxdrawuni{2500}\pmboxdrawuni{2500}\pmboxdrawuni{2500}\pmboxdrawuni{2500}\pmboxdrawuni{2500}\pmboxdrawuni{2500}\pmboxdrawuni{2500}\pmboxdrawuni{2500}\pmboxdrawuni{2500}\pmboxdrawuni{2500}\pmboxdrawuni{2500}\pmboxdrawuni{2500}\pmboxdrawuni{2500}\pmboxdrawuni{2500}\pmboxdrawuni{2500}\pmboxdrawuni{2500}\pmboxdrawuni{2500}\pmboxdrawuni{2500}\pmboxdrawuni{2500}\pmboxdrawuni{2500}\pmboxdrawuni{2500}\pmboxdrawuni{2500}\pmboxdrawuni{2500}\pmboxdrawuni{2500}\pmboxdrawuni{2500}\pmboxdrawuni{2500}\pmboxdrawuni{2500}\pmboxdrawuni{2500}\pmboxdrawuni{2500}\pmboxdrawuni{2500}\pmboxdrawuni{2500}\pmboxdrawuni{2500}\pmboxdrawuni{2500}\pmboxdrawuni{2500}\pmboxdrawuni{2500}\pmboxdrawuni{2500}\pmboxdrawuni{2500}\pmboxdrawuni{2500}\pmboxdrawuni{2500}\pmboxdrawuni{2500}\pmboxdrawuni{2500}\pmboxdrawuni{2500}\pmboxdrawuni{2500}\pmboxdrawuni{2500}\pmboxdrawuni{2500}\pmboxdrawuni{2500}\pmboxdrawuni{2500}\pmboxdrawuni{2500}\pmboxdrawuni{2500}\pmboxdrawuni{2500}\pmboxdrawuni{2500}\pmboxdrawuni{2500}\pmboxdrawuni{2500}\pmboxdrawuni{2500}\pmboxdrawuni{2500}\pmboxdrawuni{2500}}
\cl{~~~Conversation~history~~~~summary~~~~gone~~~restored~~~copied~~copied~~cut~back}
\cl{~~~Project~CLAUDE.md~~~~~~re-read~~re-read~~~~re-read~~re-read~re-read~~re-read}
\cl{~~~Path-scoped~rules~~~~~~on~match~on~match~~~on~match~on~match~~same~~on~match}
\cl{~~~Auto~memory~~~~~~~~~~~~~~kept~~~~~kept~~~~~~~kept~~~~~kept~~~~n/a~~~~~kept}
\cl{~~~Active~goal~~~~~~~~~~~~~~kept~~~~~gone~~~restored~~~~~kept~~~n/a~~~~~~kept}
\cl{~~~Scheduled~tasks~~~~~~~~~~kept~~~~~gone~~~restored*~~~~kept~~~n/a~~~~~~kept}
\cl{~~~Session~permission}
\cl{~~~~~grants~~~~~~~~~~~~~~~~~kept~~~~~kept~~~~~~~gone~~~kept**~~~gone~~~~~kept}
\cl{~~~Files~on~disk~~~~~~~~~~~~kept~~~~~kept~~~~~~~kept~~~~~kept~~~kept~~~reverted}
\cl{~~~Git~history~~~~~~~~~~~~~~kept~~~~~kept~~~~~~~kept~~~~~kept~~~kept~~~~~kept}
\cl{~~~\pmboxdrawuni{2500}\pmboxdrawuni{2500}\pmboxdrawuni{2500}\pmboxdrawuni{2500}\pmboxdrawuni{2500}\pmboxdrawuni{2500}\pmboxdrawuni{2500}\pmboxdrawuni{2500}\pmboxdrawuni{2500}\pmboxdrawuni{2500}\pmboxdrawuni{2500}\pmboxdrawuni{2500}\pmboxdrawuni{2500}\pmboxdrawuni{2500}\pmboxdrawuni{2500}\pmboxdrawuni{2500}\pmboxdrawuni{2500}\pmboxdrawuni{2500}\pmboxdrawuni{2500}\pmboxdrawuni{2500}\pmboxdrawuni{2500}\pmboxdrawuni{2500}\pmboxdrawuni{2500}\pmboxdrawuni{2500}\pmboxdrawuni{2500}\pmboxdrawuni{2500}\pmboxdrawuni{2500}\pmboxdrawuni{2500}\pmboxdrawuni{2500}\pmboxdrawuni{2500}\pmboxdrawuni{2500}\pmboxdrawuni{2500}\pmboxdrawuni{2500}\pmboxdrawuni{2500}\pmboxdrawuni{2500}\pmboxdrawuni{2500}\pmboxdrawuni{2500}\pmboxdrawuni{2500}\pmboxdrawuni{2500}\pmboxdrawuni{2500}\pmboxdrawuni{2500}\pmboxdrawuni{2500}\pmboxdrawuni{2500}\pmboxdrawuni{2500}\pmboxdrawuni{2500}\pmboxdrawuni{2500}\pmboxdrawuni{2500}\pmboxdrawuni{2500}\pmboxdrawuni{2500}\pmboxdrawuni{2500}\pmboxdrawuni{2500}\pmboxdrawuni{2500}\pmboxdrawuni{2500}\pmboxdrawuni{2500}\pmboxdrawuni{2500}\pmboxdrawuni{2500}\pmboxdrawuni{2500}\pmboxdrawuni{2500}\pmboxdrawuni{2500}\pmboxdrawuni{2500}\pmboxdrawuni{2500}\pmboxdrawuni{2500}\pmboxdrawuni{2500}\pmboxdrawuni{2500}\pmboxdrawuni{2500}\pmboxdrawuni{2500}\pmboxdrawuni{2500}\pmboxdrawuni{2500}\pmboxdrawuni{2500}\pmboxdrawuni{2500}\pmboxdrawuni{2500}\pmboxdrawuni{2500}\pmboxdrawuni{2500}}
\cl{~~~~*~unexpired~tasks~only~~~~**~kept~by~/branch,~lost~by~--fork-session}
\end{codefig}
\figcaption{Figure 3 — What survives each context operation. The column that matters is the last one: only files and version control survive everything.}
\end{figure}


\FloatBarrier
\setcounter{section}{2}
\section{Rewind, branch, fork, subtask}

These solve different problems and are routinely confused.


\begingroup
\def\tblrows{%
\texttt{/\allowbreak{}rewind} & Returns conversation and/or files to a checkpoint & Later state is removed \\
\texttt{/\allowbreak{}branch [nam\allowbreak{}e]} & Copies the conversation and switches you into the copy; the original is untouched on disk & Same process, so session permission grants carry over \\
\texttt{-\allowbreak{}-\allowbreak{}fork-\allowbreak{}session} & Same copy, in a \textbf{new process} & New process starts without "allow for this session" grants \\
\texttt{/\allowbreak{}fork} & Copies the conversation into a background session & Runs separately \\
\texttt{/\allowbreak{}subtask} & Spawns a \textbf{fork subagent} inheriting the full conversation & Returns a result; tool calls stay out of your context \\
}%
\def\tblbody{\begin{minipage}{\textwidth}\boxcaption{Table 8 — Mechanisms that change the shape of a conversation}
{\footnotesize\begin{tabular}{L{28.7mm}L{62.1mm}L{49.0mm}}
\toprule
\textbf{Mechanism} & \textbf{What it does} & \textbf{Where the work goes} \\
\midrule
\tblrows
\bottomrule\end{tabular}}\end{minipage}}%
\begingroup
\def\sloppy{\tolerance 9999\emergencystretch 3em\hfuzz 200pt\vfuzz 200pt}%
\hbadness=10000\vbadness=10000\hfuzz=200pt\vfuzz=200pt
\global\setbox\tblbox=\hbox{\tblbody}%
\endgroup
\par\addvspace{2.6mm}
\ifdim\dimexpr\ht\tblbox+\dp\tblbox\relax>0.55\textheight
  \typeout{HANDBOOK-TABLE broken \the\dimexpr\ht\tblbox+\dp\tblbox\relax}%
  \tabcaption{Table 8 — Mechanisms that change the shape of a conversation}
  {\footnotesize\begin{longtable}{L{28.7mm}L{62.1mm}L{49.0mm}}
  \toprule
\textbf{Mechanism} & \textbf{What it does} & \textbf{Where the work goes} \\
\midrule\endfirsthead
  \multicolumn{3}{@{}l@{}}{%
  \sffamily\footnotesize\itshape\color{inkgrey}Table 8 — Mechanisms that change the shape of a conversation \textemdash\ continued}\\[1.2mm]
  \toprule
\textbf{Mechanism} & \textbf{What it does} & \textbf{Where the work goes} \\
\midrule\endhead
  \bottomrule\endfoot
  \bottomrule\endlastfoot
  \tblrows
  \end{longtable}}%
\else
  \typeout{HANDBOOK-TABLE atomic \the\dimexpr\ht\tblbox+\dp\tblbox\relax}%
  \noindent\tblbody
\fi
\par\addvspace{2.6mm}
\endgroup

\FloatBarrier
\setcounter{section}{3}
\section{Reduce standing overhead}

\begin{itemize}
\item Keep \texttt{CLAUDE.\allowbreak{}md} concise and rules conditional.
\item Disable unused MCP servers, plugins and agents — after measuring, not before (Chapter 13).
\item Keep large source material in files and reference precise sections.
\item Delegate bounded independent analysis, not every small action.
\item Remember that enabling Chrome by default loads browser tools into every session and increases context usage [21].
\end{itemize}

\subsection*{Exercise}

After drafting your brief, run \texttt{/\allowbreak{}context}. Save a handoff, start a clean session, and ask the new session to continue. Verify that it can name completed work, open issues and the next check without receiving the old transcript.


\emph{A worked solution is given in Appendix M.}


\begingroup
\def\kprows{%
\item Context is what the model has this turn; a file on disk is not known merely because it could be read.
\item Durable state belongs in files. A conversation is a working surface, not a database.
\item \texttt{/\allowbreak{}compact} is lossy and invalidates the prompt cache; say explicitly what must survive.
\item Rewind, branch, fork and subtask solve four different problems and are routinely confused.
}%
\def\kpbody{\begin{minipage}{\textwidth}\subsection*{Key points}\begin{itemize}\kprows\end{itemize}\end{minipage}}%
\begingroup
\def\sloppy{\tolerance 9999\emergencystretch 3em\hfuzz 200pt\vfuzz 200pt}%
\hbadness=10000\vbadness=10000\hfuzz=200pt\vfuzz=200pt
\global\setbox\kpbox=\hbox{\kpbody}%
\endgroup
\par\addvspace{4.2mm}
\ifdim\dimexpr\ht\kpbox+\dp\kpbox\relax>0.30\textheight
  \typeout{HANDBOOK-KEYPOINTS broken \the\dimexpr\ht\kpbox+\dp\kpbox\relax}%
  \subsection*{Key points}
  \begin{itemize}\kprows\end{itemize}
\else
  \typeout{HANDBOOK-KEYPOINTS atomic \the\dimexpr\ht\kpbox+\dp\kpbox\relax}%
  \noindent\kpbody
\fi
\par\addvspace{1.4mm}
\endgroup

\FloatBarrier
\renewcommand{\chaptertitlelabel}{Part III \textperiodcentered\ Chapter 12}
\setcounter{chapter}{11}
\chapter{Sessions: resume, name, retain}

A session is a saved conversation tied to a project directory. The CLI, Desktop, web and VS Code each maintain their own history [30].


\FloatBarrier
\setcounter{section}{0}
\section{Resume}

\begin{codeblock}{9.0}{10.6}
\cl{claude~--continue~~~~~~~~~~~~\#~most~recent~interactive~session~here}
\cl{claude~--resume~~~~~~~~~~~~~~\#~session~picker}
\cl{claude~--resume~<name|id>~~~~\#~directly}
\cl{claude~--from-pr~<number>~~~~\#~picker~filtered~to~sessions~linked~to~that~PR}
\cl{claude~-n~research-brief-source-audit~~~\#~name~at~startup~(--name)}
\end{codeblock}

Inside a session, \texttt{/\allowbreak{}resume} switches conversation and \texttt{/\allowbreak{}rename} renames [30], [35]. In the picker, \texttt{Ctrl+R} renames, \texttt{Ctrl+A} widens to all projects, \texttt{Ctrl+W} to all worktrees of the repository, \texttt{Ctrl+B} filters to the current branch, and pasting a pull-request URL finds the session that created it [30].


\textbf{A resumed session restores} conversation history, model, agent, permission mode, an active goal, and unexpired scheduled tasks [30]. Three of those carry conditions, and the conditions are where the surprises live.


\textbf{The mode} has three documented exceptions, not one. \texttt{plan} and \texttt{bypass\allowbreak{}Permissions} are never restored. \texttt{auto} is restored only while the account still meets the auto-mode requirements. Manual is restored as Manual where a new session would start in \texttt{auto} from the built-in default, but where a \texttt{default\allowbreak{}Mode} in a settings file takes effect the session resumes in \emph{that} mode instead. Pass \texttt{-\allowbreak{}-\allowbreak{}permission-\allowbreak{}mode} to override whatever was restored [30].


\textbf{The model} is \textbf{not} restored when it has been retired or is disallowed by \texttt{available\allowbreak{}Models}, when a \texttt{-\allowbreak{}-\allowbreak{}model} flag or an \texttt{ANTHROPIC\_\allowbreak{}MODEL}-family environment variable picks one at launch, or on providers that use provider-specific deployment IDs [30]. \textbf{Scheduled tasks} are restored only if unexpired, and background Bash and monitor tasks are not restored at all (Section 23.3).


\textbf{It does not restore} \texttt{-\allowbreak{}-\allowbreak{}mcp-\allowbreak{}config}, \texttt{-\allowbreak{}-\allowbreak{}settings}, \texttt{-\allowbreak{}-\allowbreak{}plugin-\allowbreak{}dir}, \texttt{-\allowbreak{}-\allowbreak{}fallback-\allowbreak{}model}, or directories added with \texttt{-\allowbreak{}-\allowbreak{}add-\allowbreak{}dir}; pass those again. Directories added mid-session with \texttt{/\allowbreak{}add-\allowbreak{}dir} are not restored either, though the session picker still uses them to find the session [30].


\textbf{A gotcha worth knowing:} sessions whose \emph{first} prompt was a \texttt{/\allowbreak{}loop} command do not appear in the picker and are skipped by \texttt{-\allowbreak{}-\allowbreak{}continue} [30]. If a scheduled task seems to have vanished, that is usually why.


\FloatBarrier
\setcounter{section}{1}
\section{Name by outcome}

Good names survive a week:


\begin{itemize}
\item \texttt{client-\allowbreak{}x-\allowbreak{}brief-\allowbreak{}source-\allowbreak{}audit}
\item \texttt{accessibili\allowbreak{}ty-\allowbreak{}remediation-\allowbreak{}plan}
\item \texttt{invoice-\allowbreak{}workflow-\allowbreak{}test-\allowbreak{}2026-\allowbreak{}08}
\end{itemize}

Avoid \texttt{test}, \texttt{new chat}, or whatever the first prompt happened to be. Unnamed interactive sessions receive a default display name such as \texttt{my-\allowbreak{}app-\allowbreak{}3f} — which is \textbf{not} a resume handle — plus a generated title from the first prompt, which is [30]. If two live sessions share a name, Claude Code renames the newer one with a two-word suffix and tells you (2.1.232 or later).


\FloatBarrier
\setcounter{section}{2}
\section{Retention, storage and privacy}

Transcripts are stored as \textbf{JSONL} at \texttt{\textasciitilde{}/\allowbreak{}.\allowbreak{}claude/\allowbreak{}projects/\allowbreak{}<project>/\allowbreak{}<session-\allowbreak{}id>.\allowbreak{}jsonl}, where \texttt{<project>} is the working directory path with non-alphanumeric characters replaced by hyphens [30]. The entry format is internal and changes between releases; do not build on it. Use \texttt{/\allowbreak{}export} for a rendered transcript, or the structured interfaces (\texttt{claude -\allowbreak{}p -\allowbreak{}-\allowbreak{}output-\allowbreak{}format json}, the \texttt{transcript\_\allowbreak{}path} supplied to hooks and status lines, or the Agent SDK) for scripts [30].


\begingroup
\def\tblrows{%
Move storage off \texttt{\textasciitilde{}/\allowbreak{}.\allowbreak{}claude} & \texttt{CLAUDE\_\allowbreak{}CONFIG\_\allowbreak{}DIR} & Environment variable \\
Name the project directory yourself (2.1.234+) & \texttt{CLAUDE\_\allowbreak{}CODE\_\allowbreak{}PROJECT\_\allowbreak{}DIR\_\allowbreak{}NAME} & Environment variable, alongside \texttt{CLAUDE\_\allowbreak{}CONFIG\_\allowbreak{}DIR} \\
Change the 30-day retention & \texttt{cleanup\allowbreak{}Period\allowbreak{}Days} & \texttt{settings.\allowbreak{}json} \\
Suppress transcript writes in all modes & \texttt{CLAUDE\_\allowbreak{}CODE\_\allowbreak{}SKIP\_\allowbreak{}PROMPT\_\allowbreak{}HISTORY} & Environment variable \\
Suppress writes for one non-interactive run & \texttt{-\allowbreak{}-\allowbreak{}no-\allowbreak{}session-\allowbreak{}persistence} & CLI flag with \texttt{claude -\allowbreak{}p} \\
}%
\def\tblbody{\begin{minipage}{\textwidth}\boxcaption{Table 9 — Controls over transcript location, retention and writing}
{\footnotesize\begin{tabular}{L{38.3mm}L{63.6mm}L{38.0mm}}
\toprule
\textbf{To} & \textbf{Set} & \textbf{Where} \\
\midrule
\tblrows
\bottomrule\end{tabular}}\end{minipage}}%
\begingroup
\def\sloppy{\tolerance 9999\emergencystretch 3em\hfuzz 200pt\vfuzz 200pt}%
\hbadness=10000\vbadness=10000\hfuzz=200pt\vfuzz=200pt
\global\setbox\tblbox=\hbox{\tblbody}%
\endgroup
\par\addvspace{2.6mm}
\ifdim\dimexpr\ht\tblbox+\dp\tblbox\relax>0.55\textheight
  \typeout{HANDBOOK-TABLE broken \the\dimexpr\ht\tblbox+\dp\tblbox\relax}%
  \tabcaption{Table 9 — Controls over transcript location, retention and writing}
  {\footnotesize\begin{longtable}{L{38.3mm}L{63.6mm}L{38.0mm}}
  \toprule
\textbf{To} & \textbf{Set} & \textbf{Where} \\
\midrule\endfirsthead
  \multicolumn{3}{@{}l@{}}{%
  \sffamily\footnotesize\itshape\color{inkgrey}Table 9 — Controls over transcript location, retention and writing \textemdash\ continued}\\[1.2mm]
  \toprule
\textbf{To} & \textbf{Set} & \textbf{Where} \\
\midrule\endhead
  \bottomrule\endfoot
  \bottomrule\endlastfoot
  \tblrows
  \end{longtable}}%
\else
  \typeout{HANDBOOK-TABLE atomic \the\dimexpr\ht\tblbox+\dp\tblbox\relax}%
  \noindent\tblbody
\fi
\par\addvspace{2.6mm}
\endgroup

\begin{calloutbox}{palebrass}{brassdark}{2.0mm}
\textbf{CAUTION:} transcripts are plaintext on disk. Restrict filesystem access, shorten retention where your obligations require it, and never allow a secret into output on the assumption that history is encrypted. Remember that the auto-memory directory is excluded from the retention sweep (Section 10.2). Data governance is treated properly in Chapter 32.
\end{calloutbox}

\FloatBarrier
\setcounter{section}{3}
\section{When to start another session}

Start clean when the objective changes, when a decision needs independent review, when context is dominated by discarded work, when a bounded subtask can run in isolation, or when you need an alternative path without confusing the approved one.


Do not run parallel sessions that edit the same files. Use worktrees or strict file ownership — Chapter 29.



\begingroup
\def\kprows{%
\item A resumed session restores history, model, agent, goal and unexpired tasks — but never \texttt{plan} or \texttt{bypass\allowbreak{}Permissions}.
\item Configuration passed as flags is not restored; pass \texttt{-\allowbreak{}-\allowbreak{}mcp-\allowbreak{}config}, \texttt{-\allowbreak{}-\allowbreak{}settings}, \texttt{-\allowbreak{}-\allowbreak{}plugin-\allowbreak{}dir} and \texttt{-\allowbreak{}-\allowbreak{}add-\allowbreak{}dir} again.
\item A session whose first prompt was \texttt{/\allowbreak{}loop} is hidden from the picker and skipped by \texttt{-\allowbreak{}-\allowbreak{}continue}.
\item Transcripts are plaintext JSONL on disk; retention, location and whether they are written at all are configurable.
}%
\def\kpbody{\begin{minipage}{\textwidth}\subsection*{Key points}\begin{itemize}\kprows\end{itemize}\end{minipage}}%
\begingroup
\def\sloppy{\tolerance 9999\emergencystretch 3em\hfuzz 200pt\vfuzz 200pt}%
\hbadness=10000\vbadness=10000\hfuzz=200pt\vfuzz=200pt
\global\setbox\kpbox=\hbox{\kpbody}%
\endgroup
\par\addvspace{4.2mm}
\ifdim\dimexpr\ht\kpbox+\dp\kpbox\relax>0.30\textheight
  \typeout{HANDBOOK-KEYPOINTS broken \the\dimexpr\ht\kpbox+\dp\kpbox\relax}%
  \subsection*{Key points}
  \begin{itemize}\kprows\end{itemize}
\else
  \typeout{HANDBOOK-KEYPOINTS atomic \the\dimexpr\ht\kpbox+\dp\kpbox\relax}%
  \noindent\kpbody
\fi
\par\addvspace{1.4mm}
\endgroup

\breakrule

\FloatBarrier
\parttitle{Part IV}{Extending Claude Code}
\addcontentsline{toc}{part}{Part IV \textemdash\ Extending Claude Code}

Every chapter in this part enlarges what Claude Code can reach. Each is therefore also a supply-chain chapter. The organising principle is stated once here and applied throughout:


\begin{calloutbox}{palegrey}{quoterule}{1.2mm}
Anthropic does not security-audit MCP servers, and plugins and marketplaces can execute arbitrary code with your user privileges. Connectors are reviewed against listing criteria before entering the Anthropic Directory, but that is a listing review, not a security audit [38], [53].
\end{calloutbox}

The consequence is that trust in agentic work is not a perimeter property. Figure 4 enumerates the channels through which content the agent did not author reaches the session, and none of them is the network boundary.


\begin{figure}[tbp]
\begin{codefig}{8.0}{9.4}
\cl{~~~~~~~~~~~~~~~~~~~~\pmboxdrawuni{250C}\pmboxdrawuni{2500}\pmboxdrawuni{2500}\pmboxdrawuni{2500}\pmboxdrawuni{2500}\pmboxdrawuni{2500}\pmboxdrawuni{2500}\pmboxdrawuni{2500}\pmboxdrawuni{2500}\pmboxdrawuni{2500}\pmboxdrawuni{2500}\pmboxdrawuni{2500}\pmboxdrawuni{2500}\pmboxdrawuni{2500}\pmboxdrawuni{2500}\pmboxdrawuni{2500}\pmboxdrawuni{2500}\pmboxdrawuni{2500}\pmboxdrawuni{2500}\pmboxdrawuni{2500}\pmboxdrawuni{2500}\pmboxdrawuni{2500}\pmboxdrawuni{2500}\pmboxdrawuni{2500}\pmboxdrawuni{2500}\pmboxdrawuni{2500}\pmboxdrawuni{2500}\pmboxdrawuni{2500}\pmboxdrawuni{2500}\pmboxdrawuni{2500}\pmboxdrawuni{2500}\pmboxdrawuni{2500}\pmboxdrawuni{2500}\pmboxdrawuni{2500}\pmboxdrawuni{2500}\pmboxdrawuni{2500}\pmboxdrawuni{2500}\pmboxdrawuni{2500}\pmboxdrawuni{2500}\pmboxdrawuni{2500}\pmboxdrawuni{2500}\pmboxdrawuni{2500}\pmboxdrawuni{2500}\pmboxdrawuni{2500}\pmboxdrawuni{2510}}
\cl{~~~web~page~\pmboxdrawuni{2500}\pmboxdrawuni{2500}\pmboxdrawuni{2500}\pmboxdrawuni{2500}\pmboxdrawuni{2500}\pmboxdrawuni{2500}\pmboxdrawuni{2500}▶\pmboxdrawuni{2502}~~~~~~~~~~~~~~~~~~~~~~~~~~~~~~~~~~~~~~~~~~~\pmboxdrawuni{2502}}
\cl{~~~tool~result~\pmboxdrawuni{2500}\pmboxdrawuni{2500}\pmboxdrawuni{2500}\pmboxdrawuni{2500}▶\pmboxdrawuni{2502}~~~~~~~~~~~~~~~~~~~~~~~~~~~~~~~~~~~~~~~~~~~\pmboxdrawuni{2502}}
\cl{~~~PR~comment~\pmboxdrawuni{2500}\pmboxdrawuni{2500}\pmboxdrawuni{2500}\pmboxdrawuni{2500}\pmboxdrawuni{2500}▶\pmboxdrawuni{2502}~~~~~~~~THE~SESSION\textquotesingle{}S~CONTEXT~~~~~~~~~~~~~~\pmboxdrawuni{2502}◀\pmboxdrawuni{2500}\pmboxdrawuni{2500}~channel~message}
\cl{~~~fire~payload~\pmboxdrawuni{2500}\pmboxdrawuni{2500}\pmboxdrawuni{2500}▶\pmboxdrawuni{2502}~~~(everything~here~is~read~as~input,~~~~~~\pmboxdrawuni{2502}◀\pmboxdrawuni{2500}\pmboxdrawuni{2500}~peer~session}
\cl{~~~file~content~\pmboxdrawuni{2500}\pmboxdrawuni{2500}\pmboxdrawuni{2500}▶\pmboxdrawuni{2502}~~~~and~some~of~it~will~try~to~instruct)~~~\pmboxdrawuni{2502}◀\pmboxdrawuni{2500}\pmboxdrawuni{2500}~imported~CLAUDE.md}
\cl{~~~search~result~\pmboxdrawuni{2500}\pmboxdrawuni{2500}▶\pmboxdrawuni{2502}~~~~~~~~~~~~~~~~~~~~~~~~~~~~~~~~~~~~~~~~~~~\pmboxdrawuni{2502}}
\cl{~~~~~~~~~~~~~~~~~~~~\pmboxdrawuni{2514}\pmboxdrawuni{2500}\pmboxdrawuni{2500}\pmboxdrawuni{2500}\pmboxdrawuni{2500}\pmboxdrawuni{2500}\pmboxdrawuni{2500}\pmboxdrawuni{2500}\pmboxdrawuni{2500}\pmboxdrawuni{2500}\pmboxdrawuni{2500}\pmboxdrawuni{2500}\pmboxdrawuni{2500}\pmboxdrawuni{2500}\pmboxdrawuni{2500}\pmboxdrawuni{2500}\pmboxdrawuni{2500}\pmboxdrawuni{2500}\pmboxdrawuni{2500}\pmboxdrawuni{2500}\pmboxdrawuni{2500}\pmboxdrawuni{252C}\pmboxdrawuni{2500}\pmboxdrawuni{2500}\pmboxdrawuni{2500}\pmboxdrawuni{2500}\pmboxdrawuni{2500}\pmboxdrawuni{2500}\pmboxdrawuni{2500}\pmboxdrawuni{2500}\pmboxdrawuni{2500}\pmboxdrawuni{2500}\pmboxdrawuni{2500}\pmboxdrawuni{2500}\pmboxdrawuni{2500}\pmboxdrawuni{2500}\pmboxdrawuni{2500}\pmboxdrawuni{2500}\pmboxdrawuni{2500}\pmboxdrawuni{2500}\pmboxdrawuni{2500}\pmboxdrawuni{2500}\pmboxdrawuni{2500}\pmboxdrawuni{2500}\pmboxdrawuni{2518}}
\cl{~~~~~~~~~~~~~~~~~~~~~~~~~~~~~~~~~~~~~~~~~\pmboxdrawuni{2502}~proposes~actions}
\cl{~~~~~~~~~~~~~~~~~~~~~~~~~~~~~~~~~~~~~~~~~▼}
\cl{~~~~~~~~~~~~~~~~~~~~\pmboxdrawuni{250C}\pmboxdrawuni{2500}\pmboxdrawuni{2500}\pmboxdrawuni{2500}\pmboxdrawuni{2500}\pmboxdrawuni{2500}\pmboxdrawuni{2500}\pmboxdrawuni{2500}\pmboxdrawuni{2500}\pmboxdrawuni{2500}\pmboxdrawuni{2500}\pmboxdrawuni{2500}\pmboxdrawuni{2500}\pmboxdrawuni{2500}\pmboxdrawuni{2500}\pmboxdrawuni{2500}\pmboxdrawuni{2500}\pmboxdrawuni{2500}\pmboxdrawuni{2500}\pmboxdrawuni{2500}\pmboxdrawuni{2500}\pmboxdrawuni{2500}\pmboxdrawuni{2500}\pmboxdrawuni{2500}\pmboxdrawuni{2500}\pmboxdrawuni{2500}\pmboxdrawuni{2500}\pmboxdrawuni{2500}\pmboxdrawuni{2500}\pmboxdrawuni{2500}\pmboxdrawuni{2500}\pmboxdrawuni{2500}\pmboxdrawuni{2500}\pmboxdrawuni{2500}\pmboxdrawuni{2500}\pmboxdrawuni{2500}\pmboxdrawuni{2500}\pmboxdrawuni{2500}\pmboxdrawuni{2500}\pmboxdrawuni{2500}\pmboxdrawuni{2500}\pmboxdrawuni{2500}\pmboxdrawuni{2500}\pmboxdrawuni{2500}\pmboxdrawuni{2510}}
\cl{~~~~~~~~~~~~~~~~~~~~\pmboxdrawuni{2502}~~THE~CONTROL~STACK~~(see~Figure~1)~~~~~~~~\pmboxdrawuni{2502}}
\cl{~~~~~~~~~~~~~~~~~~~~\pmboxdrawuni{2514}\pmboxdrawuni{2500}\pmboxdrawuni{2500}\pmboxdrawuni{2500}\pmboxdrawuni{2500}\pmboxdrawuni{2500}\pmboxdrawuni{2500}\pmboxdrawuni{2500}\pmboxdrawuni{2500}\pmboxdrawuni{2500}\pmboxdrawuni{2500}\pmboxdrawuni{2500}\pmboxdrawuni{2500}\pmboxdrawuni{2500}\pmboxdrawuni{2500}\pmboxdrawuni{2500}\pmboxdrawuni{2500}\pmboxdrawuni{2500}\pmboxdrawuni{2500}\pmboxdrawuni{2500}\pmboxdrawuni{2500}\pmboxdrawuni{252C}\pmboxdrawuni{2500}\pmboxdrawuni{2500}\pmboxdrawuni{2500}\pmboxdrawuni{2500}\pmboxdrawuni{2500}\pmboxdrawuni{2500}\pmboxdrawuni{2500}\pmboxdrawuni{2500}\pmboxdrawuni{2500}\pmboxdrawuni{2500}\pmboxdrawuni{2500}\pmboxdrawuni{2500}\pmboxdrawuni{2500}\pmboxdrawuni{2500}\pmboxdrawuni{2500}\pmboxdrawuni{2500}\pmboxdrawuni{2500}\pmboxdrawuni{2500}\pmboxdrawuni{2500}\pmboxdrawuni{2500}\pmboxdrawuni{2500}\pmboxdrawuni{2500}\pmboxdrawuni{2518}}
\cl{~~~~~~~~~~~~~~~~~~~~~~~~~~~~~~~~~~~~~~~~~\pmboxdrawuni{2502}~what~actually~happens}
\cl{~~~~~~~~~~~~~~~~~~~~\pmboxdrawuni{250C}\pmboxdrawuni{2500}\pmboxdrawuni{2500}\pmboxdrawuni{2500}\pmboxdrawuni{2500}\pmboxdrawuni{2500}\pmboxdrawuni{2500}\pmboxdrawuni{2500}\pmboxdrawuni{2500}\pmboxdrawuni{2500}\pmboxdrawuni{2500}\pmboxdrawuni{2500}\pmboxdrawuni{2500}\pmboxdrawuni{2500}\pmboxdrawuni{2500}\pmboxdrawuni{2500}\pmboxdrawuni{2500}\pmboxdrawuni{2500}\pmboxdrawuni{2500}\pmboxdrawuni{2500}\pmboxdrawuni{2500}\pmboxdrawuni{2534}\pmboxdrawuni{2500}\pmboxdrawuni{2500}\pmboxdrawuni{2500}\pmboxdrawuni{2500}\pmboxdrawuni{2500}\pmboxdrawuni{2500}\pmboxdrawuni{2500}\pmboxdrawuni{2500}\pmboxdrawuni{2500}\pmboxdrawuni{2500}\pmboxdrawuni{2500}\pmboxdrawuni{2500}\pmboxdrawuni{2500}\pmboxdrawuni{2500}\pmboxdrawuni{2500}\pmboxdrawuni{2500}\pmboxdrawuni{2500}\pmboxdrawuni{2500}\pmboxdrawuni{2500}\pmboxdrawuni{2500}\pmboxdrawuni{2500}\pmboxdrawuni{2500}\pmboxdrawuni{2510}}
\cl{~~~~~~~~~~~~~~~~~~~~▼~~~~~~~~~~~~~~~~~~~~▼~~~~~~~~~~~~~~~~~~~~~~▼}
\cl{~~~~~~~~~~~~~~your~filesystem~~~~~~your~accounts~~~~~~~~~~published~output}
\cl{~~~~~~~~~~~~~~~~~~~~~~~~~~~~~~~~~~~(MCP,~browser,~~~~~~~~~~(artifacts,~PRs,}
\cl{~~~~~~~~~~~~~~~~~~~~~~~~~~~~~~~~~~~~git,~connectors)~~~~~~~~messages)}
\cl{}
\cl{~~~Product~mitigations,~all~partial:~web~fetches~use~a~separate~context~window;}
\cl{~~~routine~fire~text~arrives~wrapped~and~labelled~untrusted;~inter-agent~messages}
\cl{~~~are~classifier-reviewed~and~cannot~carry~consent;~connectors~are~listing-reviewed}
\cl{~~~but~not~security-audited.}
\end{codefig}
\figcaption{Figure 4 — Where untrusted content enters an agentic session. Every arrow is a documented injection channel; none of them is the network perimeter.}
\end{figure}


Appendix F maps each mechanism in this part to its threat model and the controls available.


\FloatBarrier
\renewcommand{\chaptertitlelabel}{Part IV \textperiodcentered\ Chapter 13}
\setcounter{chapter}{12}
\chapter{Connect external tools with MCP}

The Model Context Protocol is an open standard for connecting AI clients to external tools and data [75]. A server might search a knowledge base, read an issue tracker, query a database, or act on a service. It expands capability and trust simultaneously: a connected tool acts with whatever authority you gave it.


\FloatBarrier
\setcounter{section}{0}
\section{Scopes}

\begingroup
\def\tblrows{%
Local (default) & \texttt{\textasciitilde{}/\allowbreak{}.\allowbreak{}claude.\allowbreak{}json} & One project, current user; private configuration \\
User & \texttt{\textasciitilde{}/\allowbreak{}.\allowbreak{}claude.\allowbreak{}json} & Available across your projects \\
Project & \texttt{.\allowbreak{}mcp.\allowbreak{}json} at the repository root & Shared team configuration, committed \\
}%
\def\tblbody{\begin{minipage}{\textwidth}\boxcaption{Table 10 — MCP configuration scopes}
{\footnotesize\begin{tabular}{L{19.0mm}L{44.0mm}L{76.8mm}}
\toprule
\textbf{Scope} & \textbf{Storage} & \textbf{Use} \\
\midrule
\tblrows
\bottomrule\end{tabular}}\end{minipage}}%
\begingroup
\def\sloppy{\tolerance 9999\emergencystretch 3em\hfuzz 200pt\vfuzz 200pt}%
\hbadness=10000\vbadness=10000\hfuzz=200pt\vfuzz=200pt
\global\setbox\tblbox=\hbox{\tblbody}%
\endgroup
\par\addvspace{2.6mm}
\ifdim\dimexpr\ht\tblbox+\dp\tblbox\relax>0.55\textheight
  \typeout{HANDBOOK-TABLE broken \the\dimexpr\ht\tblbox+\dp\tblbox\relax}%
  \tabcaption{Table 10 — MCP configuration scopes}
  {\footnotesize\begin{longtable}{L{19.0mm}L{44.0mm}L{76.8mm}}
  \toprule
\textbf{Scope} & \textbf{Storage} & \textbf{Use} \\
\midrule\endfirsthead
  \multicolumn{3}{@{}l@{}}{%
  \sffamily\footnotesize\itshape\color{inkgrey}Table 10 — MCP configuration scopes \textemdash\ continued}\\[1.2mm]
  \toprule
\textbf{Scope} & \textbf{Storage} & \textbf{Use} \\
\midrule\endhead
  \bottomrule\endfoot
  \bottomrule\endlastfoot
  \tblrows
  \end{longtable}}%
\else
  \typeout{HANDBOOK-TABLE atomic \the\dimexpr\ht\tblbox+\dp\tblbox\relax}%
  \noindent\tblbody
\fi
\par\addvspace{2.6mm}
\endgroup

Project scope carries an approval mechanism that the course omitted and that the first edition under-described. A project-scoped server from \texttt{.\allowbreak{}mcp.\allowbreak{}json} appears as \texttt{⏸ Pending a\allowbreak{}pproval} until you run \texttt{claude} interactively and approve it; rejected servers are recorded in \texttt{disabled\allowbreak{}Mcpjson\allowbreak{}Servers} [32]. From 2.1.196, approvals are read only from settings files that are \textbf{not} checked into the repository until you trust the workspace, so a cloned repository cannot approve its own servers — \texttt{enable\allowbreak{}All\allowbreak{}Project\allowbreak{}Mcp\allowbreak{}Servers} or \texttt{enabled\allowbreak{}Mcpjson\allowbreak{}Servers} committed to a project's \texttt{.\allowbreak{}claude/\allowbreak{}settings.\allowbreak{}json} is ignored in an untrusted folder [32]. That is the control that stops a malicious pull request from silently connecting a server.


\begin{calloutbox}{palebrass}{brassdark}{2.0mm}
\textbf{CAUTION:} never place a literal credential in \texttt{.\allowbreak{}mcp.\allowbreak{}json}. Use OAuth where the provider supports it, or \texttt{\$\{VAR\}} expansion referencing an environment variable, and use a least-privilege account. See M11 in Table 1.
\end{calloutbox}

\FloatBarrier
\setcounter{section}{1}
\section{Transports}

\begin{codeblock}{9.0}{10.6}
\cl{claude~mcp~add~--transport~http~notion~https://mcp.notion.com/mcp~--scope~project}
\cl{claude~mcp~list}
\cl{claude~mcp~get~notion}
\end{codeblock}

Four transports exist [32]:


\begin{itemize}
\item \textbf{HTTP} — preferred. In JSON, \texttt{type:\allowbreak{} "http"}; \texttt{streamable-\allowbreak{}http} is accepted as an alias so configurations copied from a server's own documentation work unmodified.
\item \textbf{SSE} — \textbf{deprecated}. Use HTTP where available; some services still expose only SSE.
\item \textbf{stdio} — a local process. Everything after \texttt{-\allowbreak{}-\allowbreak{}} is passed to the server untouched, so \texttt{claude mcp \allowbreak{}add -\allowbreak{}-\allowbreak{}env KEY=\allowbreak{}value -\allowbreak{}-\allowbreak{}transport s\allowbreak{}tdio myserv\allowbreak{}er -\allowbreak{}-\allowbreak{} python ser\allowbreak{}ver.\allowbreak{}py -\allowbreak{}-\allowbreak{}port 8080} runs the server with that argument and environment.
\item \textbf{WebSocket} (\texttt{type:\allowbreak{} "ws"}) — persistent bidirectional, for servers that push events. It supports neither OAuth nor the \texttt{-\allowbreak{}-\allowbreak{}transport} flag; authentication is header-only.
\end{itemize}

A JSON entry with a \texttt{url} but no \texttt{type} is a configuration error: Claude Code reads a type-less entry as stdio and skips it with an explicit message [32].


Inside a session, \texttt{/\allowbreak{}mcp} inspects status, completes OAuth, and lists a server's tools. From the shell, \texttt{claude mcp} supports \texttt{add}, \texttt{add-\allowbreak{}json}, \texttt{list}, \texttt{get}, \texttt{remove}, \texttt{login} and \texttt{logout} [32], [62].


\FloatBarrier
\setcounter{section}{2}
\section{A safe installation procedure}

\begin{enumerate}
\item Identify the provider's official repository or documentation. Verify the URL character by character.
\item Define the minimum data and actions the task needs.
\item Use a test account or a least-privilege workspace.
\item Inspect source, publisher, package and network destinations where you can.
\item Choose the narrowest scope that works.
\item Keep secrets out of \texttt{.\allowbreak{}mcp.\allowbreak{}json} and out of version control.
\item Start a new session so configuration reloads.
\item Authenticate and inspect the exposed tools with \texttt{/\allowbreak{}mcp}.
\item Exercise a read-only operation first.
\item Add explicit approval — an \texttt{ask} rule or a human gate — for writes.
\end{enumerate}

\begin{calloutbox}{palebrass}{brassdark}{2.0mm}
\textbf{CAUTION — tool output is untrusted input.} Content returned by an MCP server can contain instructions aimed at the agent, private data, or material designed to redirect the task. This is indirect prompt injection [4], [22], [23]. Treat every tool result as data, constrain what the agent may do next, and require confirmation before any external write [53].
\end{calloutbox}

\FloatBarrier
\setcounter{section}{3}
\section{Context cost, measured rather than assumed}

The course warns that tool definitions consume context [42]. That was true and is now partly mitigated: Tool Search defers many definitions until needed [32], [33]. Measure with \texttt{/\allowbreak{}context} and \texttt{/\allowbreak{}usage} rather than assuming. Remove servers you do not need anyway — the argument is security and clarity, not tokens.


Per-tool result-size overrides and \texttt{always\allowbreak{}Load} are available where a specific server needs different handling [32].


\FloatBarrier
\setcounter{section}{4}
\section{Organisational control}

Administrators can restrict which servers users may add or connect to, using allowlists and denylists in managed configuration [76]. \texttt{denied\allowbreak{}Mcp\allowbreak{}Servers} can block a named server outright — including \texttt{claude-\allowbreak{}in-\allowbreak{}chrome}, which is how an organisation turns off browser integration centrally [21].


\subsection*{Exercise}

Do not install a live server yet. Ask Claude to produce a threat model for a hypothetical Notion connection: data reachable, write operations exposed, credential storage, scope, approval gates, logging, and revocation path. Approve installation only after replacing each assumption with a statement from the provider's documentation.


\emph{A worked solution is given in Appendix M.}


\begingroup
\def\kprows{%
\item A connected server acts with whatever authority you granted it, and Anthropic does not security-audit MCP servers.
\item Project-scoped servers require per-user approval and workspace trust, which is what stops a pull request from connecting one silently.
\item Prefer HTTP transport; SSE is deprecated and a \texttt{url} entry without a \texttt{type} is a configuration error.
\item Treat every tool result as untrusted data, and require confirmation before any external write.
}%
\def\kpbody{\begin{minipage}{\textwidth}\subsection*{Key points}\begin{itemize}\kprows\end{itemize}\end{minipage}}%
\begingroup
\def\sloppy{\tolerance 9999\emergencystretch 3em\hfuzz 200pt\vfuzz 200pt}%
\hbadness=10000\vbadness=10000\hfuzz=200pt\vfuzz=200pt
\global\setbox\kpbox=\hbox{\kpbody}%
\endgroup
\par\addvspace{4.2mm}
\ifdim\dimexpr\ht\kpbox+\dp\kpbox\relax>0.30\textheight
  \typeout{HANDBOOK-KEYPOINTS broken \the\dimexpr\ht\kpbox+\dp\kpbox\relax}%
  \subsection*{Key points}
  \begin{itemize}\kprows\end{itemize}
\else
  \typeout{HANDBOOK-KEYPOINTS atomic \the\dimexpr\ht\kpbox+\dp\kpbox\relax}%
  \noindent\kpbody
\fi
\par\addvspace{1.4mm}
\endgroup

\FloatBarrier
\renewcommand{\chaptertitlelabel}{Part IV \textperiodcentered\ Chapter 14}
\setcounter{chapter}{13}
\chapter{Environment variables and secrets}

A \texttt{.\allowbreak{}env} file is a developer convenience, not a security boundary. Commit only a template:


\begin{codeblock}{9.0}{10.6}
\cl{\#~.env.example~—~names~only,~never~values}
\cl{NOTION\_TOKEN=}
\cl{DATABASE\_URL=}
\end{codeblock}

\FloatBarrier
\setcounter{section}{0}
\section{Why "Claude cannot see it" is unsafe}

If Claude may read \texttt{.\allowbreak{}env}, its contents enter tool output. If it runs a command that prints an environment variable, the value enters the plaintext transcript (Section 12.3). A subprocess may transmit it. A sentence in \texttt{CLAUDE.\allowbreak{}md} saying "do not read secrets" is a preference, not enforcement [26]. This is M04 in Table 1, and it is the most common false belief among new users.


\FloatBarrier
\setcounter{section}{1}
\section{Layered controls}

\begin{enumerate}
\item \textbf{Least privilege.} Issue a credential limited to the service, action and environment the task needs. This is the only control that limits damage after exposure.
\item \textbf{Permission deny rules.} Block reads of secret files and directories (Section 7.2).
\item \textbf{Sandbox credential controls.} \texttt{sandbox.\allowbreak{}credentials.\allowbreak{}files} and \texttt{.\allowbreak{}env\allowbreak{}Vars}, \texttt{deny} or \texttt{mask} (Section 7.4). Denied variables are unset before each sandboxed command runs.
\item \textbf{Network restriction.} Deny \texttt{curl} and \texttt{wget} at the Bash layer and allow specific domains through \texttt{Web\allowbreak{}Fetch(domai\allowbreak{}n:\allowbreak{}…)} instead [53].
\item \textbf{Output discipline.} Never echo, log or interpolate a secret into a command the model will see.
\item \textbf{Git hygiene.} Ignore secret files and scan the staged diff before every commit.
\item \textbf{Rotation.} Revoke immediately if a credential reaches a transcript, a commit, a log, an artifact, or an untrusted service.
\end{enumerate}

Do not paste a security configuration without testing that it blocks. Ask Claude to attempt the denied read and confirm the denial, then check \texttt{/\allowbreak{}permissions} for the rule's source file.


\FloatBarrier
\setcounter{section}{2}
\section{Incident response}

\begin{enumerate}
\item Stop the task.
\item Revoke or rotate the credential. Do this before anything else.
\item Remove it from files, logs and any published artifact.
\item If committed, follow the hosting provider's sensitive-data removal procedure; deleting the latest line is not enough [77].
\item Review the session transcript, connected-system audit logs, and — for organisations — the Claude Code audit log (Chapter 33).
\item Correct the control that permitted exposure, and add a regression test for it.
\end{enumerate}

\begingroup
\def\kprows{%
\item A \texttt{.\allowbreak{}env} file is a convenience, not a boundary. An instruction not to read secrets is not enforcement.
\item Layer the controls: least privilege, deny rules, sandbox credential deny or mask, network restriction, output discipline, git hygiene, rotation.
\item Test that a deny rule actually blocks — do not paste a security configuration and assume it.
\item On exposure, rotate first and investigate second.
}%
\def\kpbody{\begin{minipage}{\textwidth}\subsection*{Key points}\begin{itemize}\kprows\end{itemize}\end{minipage}}%
\begingroup
\def\sloppy{\tolerance 9999\emergencystretch 3em\hfuzz 200pt\vfuzz 200pt}%
\hbadness=10000\vbadness=10000\hfuzz=200pt\vfuzz=200pt
\global\setbox\kpbox=\hbox{\kpbody}%
\endgroup
\par\addvspace{4.2mm}
\ifdim\dimexpr\ht\kpbox+\dp\kpbox\relax>0.30\textheight
  \typeout{HANDBOOK-KEYPOINTS broken \the\dimexpr\ht\kpbox+\dp\kpbox\relax}%
  \subsection*{Key points}
  \begin{itemize}\kprows\end{itemize}
\else
  \typeout{HANDBOOK-KEYPOINTS atomic \the\dimexpr\ht\kpbox+\dp\kpbox\relax}%
  \noindent\kpbody
\fi
\par\addvspace{1.4mm}
\endgroup


\FloatBarrier
\renewcommand{\chaptertitlelabel}{Part IV \textperiodcentered\ Chapter 15}
\setcounter{chapter}{14}
\chapter{Turn repeatable work into skills}

A skill is a reusable package of instructions and optional resources. Claude invokes it when its description matches the task, or you call it as \texttt{/\allowbreak{}skill-\allowbreak{}name}. Unlike \texttt{CLAUDE.\allowbreak{}md}, a skill's body loads only when used, so long reference material costs almost nothing until needed [34].


\begin{codeblock}{9.0}{10.6}
\cl{\textasciitilde{}/.claude/skills/<name>/SKILL.md~~~~~\#~user~scope}
\cl{.claude/skills/<name>/SKILL.md~~~~~~~\#~project~scope}
\end{codeblock}

Custom commands have merged into skills: \texttt{.\allowbreak{}claude/\allowbreak{}commands/\allowbreak{}deploy.\allowbreak{}md} and \texttt{.\allowbreak{}claude/\allowbreak{}skills/\allowbreak{}deploy/\allowbreak{}SKILL.\allowbreak{}md} both create \texttt{/\allowbreak{}deploy}. Existing command files keep working; use skills for anything new [34].


\FloatBarrier
\setcounter{section}{0}
\section{Build one from observed work}

This is the course's strongest contribution and it is worth stating plainly [42]:


\begin{enumerate}
\item Complete the procedure manually, with Claude, once.
\item Identify inputs, tools, decisions, outputs and failure modes.
\item Ask Claude to write a minimal skill from what actually happened.
\item Review every file and every tool grant.
\item Test it on a genuinely different example.
\item Correct the output by hand.
\item Update the skill with the \textbf{general} lesson, not the incidental detail.
\end{enumerate}

\FloatBarrier
\setcounter{section}{1}
\section{Example project skill}

\begin{codeblock}{8.0}{9.4}
\cl{---}
\cl{name:~research-brief}
\cl{description:~Create~an~evidence-led~decision~brief~from~approved~sources~in~this~workspace.}
\cl{~~Use~when~the~user~asks~for~a~research~brief,~evidence~synthesis,~or~source-backed}
\cl{~~recommendation.}
\cl{allowed-tools:~Read,~Grep,~Glob}
\cl{---}
\cl{}
\cl{\#~Research~brief}
\cl{}
\cl{\#\#~Required~input}
\cl{-~Decision~question~and~audience}
\cl{-~Approved~source~scope}
\cl{-~Output~length~and~deadline}
\cl{}
\cl{\#\#~Procedure}
\cl{1.~Read~\textasciigrave{}CLAUDE.md\textasciigrave{}~and~the~approved~source~list.}
\cl{2.~Extract~one~atomic~claim~per~ledger~row.}
\cl{3.~Separate~observed~facts,~source~claims,~and~analyst~inference.}
\cl{4.~Search~for~contradictions~and~missing~evidence.}
\cl{5.~Draft~in~\textasciigrave{}drafts/\textasciigrave{};~never~write~directly~to~\textasciigrave{}deliverables/\textasciigrave{}.}
\cl{6.~Run~the~evidence-reviewer~subagent~if~available.}
\cl{7.~Report~gaps~and~request~human~approval.}
\cl{}
\cl{\#\#~Output}
\cl{Executive~answer~·~scope~and~method~·~findings~with~inline~citations~·~evidence~table~·}
\cl{contradictions~and~gaps~·~source~register.}
\end{codeblock}

\FloatBarrier
\setcounter{section}{2}
\section{Frontmatter that matters for safety}

\begingroup
\def\tblrows{%
\texttt{allowed-\allowbreak{}tools} & Tools Claude may use \textbf{without asking} during the turn that invokes the skill. The grant clears when you send your next message \\
\texttt{disallowed-\allowbreak{}tools} & Tools removed from the pool while the skill is active; also clears on your next message \\
\texttt{disable-\allowbreak{}model-\allowbreak{}invocation:\allowbreak{} true} & Only you can invoke it. Also prevents preloading into subagents, and prevents it running when a scheduled task fires with the skill as its prompt (2.1.196+) \\
\texttt{context:\allowbreak{} fork} & Runs the skill in its own subagent context \\
\texttt{shell} & \texttt{bash} (default) or \texttt{powershell} for inline shell in the skill body \\
\texttt{hooks} & Lifecycle hooks registered when the skill is invoked, active for the rest of the session; \texttt{once:\allowbreak{} true} limits one to a single run \\
}%
\def\tblbody{\begin{minipage}{\textwidth}\boxcaption{Table 11 — Skill frontmatter fields with a security consequence}
{\footnotesize\begin{tabular}{L{51.2mm}L{91.8mm}}
\toprule
\textbf{Field} & \textbf{Effect} \\
\midrule
\tblrows
\bottomrule\end{tabular}}\end{minipage}}%
\begingroup
\def\sloppy{\tolerance 9999\emergencystretch 3em\hfuzz 200pt\vfuzz 200pt}%
\hbadness=10000\vbadness=10000\hfuzz=200pt\vfuzz=200pt
\global\setbox\tblbox=\hbox{\tblbody}%
\endgroup
\par\addvspace{2.6mm}
\ifdim\dimexpr\ht\tblbox+\dp\tblbox\relax>0.55\textheight
  \typeout{HANDBOOK-TABLE broken \the\dimexpr\ht\tblbox+\dp\tblbox\relax}%
  \tabcaption{Table 11 — Skill frontmatter fields with a security consequence}
  {\footnotesize\begin{longtable}{L{51.2mm}L{91.8mm}}
  \toprule
\textbf{Field} & \textbf{Effect} \\
\midrule\endfirsthead
  \multicolumn{2}{@{}l@{}}{%
  \sffamily\footnotesize\itshape\color{inkgrey}Table 11 — Skill frontmatter fields with a security consequence \textemdash\ continued}\\[1.2mm]
  \toprule
\textbf{Field} & \textbf{Effect} \\
\midrule\endhead
  \bottomrule\endfoot
  \bottomrule\endlastfoot
  \tblrows
  \end{longtable}}%
\else
  \typeout{HANDBOOK-TABLE atomic \the\dimexpr\ht\tblbox+\dp\tblbox\relax}%
  \noindent\tblbody
\fi
\par\addvspace{2.6mm}
\endgroup

Two behaviours are easy to miss. \texttt{\$\{CLAUDE\_\allowbreak{}SKILL\_\allowbreak{}DIR\}} and \texttt{\$\{CLAUDE\_\allowbreak{}PROJECT\_\allowbreak{}DIR\}} are substituted both in the body and in Bash rules inside \texttt{allowed-\allowbreak{}tools}, so a skill can run a bundled script without a prompt — \texttt{allowed-\allowbreak{}tools:\allowbreak{} Bash(\$\{C\allowbreak{}LAUDE\_\allowbreak{}SKILL\_\allowbreak{}DIR\}/\allowbreak{}scripts/\allowbreak{}render.\allowbreak{}sh *)} [34]. And the rendered \texttt{SKILL.\allowbreak{}md} enters the conversation once and \textbf{stays for the rest of the session}; Claude Code does not re-read the file on later turns, so write standing guidance rather than one-time steps [34].


Use \texttt{disable-\allowbreak{}model-\allowbreak{}invocation:\allowbreak{} true} for anything with side effects — \texttt{/\allowbreak{}commit}, \texttt{/\allowbreak{}deploy}, \texttt{/\allowbreak{}send-\allowbreak{}message}. You do not want Claude deciding your code looks ready to ship.


\FloatBarrier
\setcounter{section}{3}
\section{Skill security review}

Before installing or running a third-party skill, inspect:


\begin{itemize}
\item \texttt{SKILL.\allowbreak{}md} and every referenced file;
\item \texttt{allowed-\allowbreak{}tools}, \texttt{disallowed-\allowbreak{}tools} and any \texttt{hooks} block;
\item inline shell execution (`\texttt{ !}command\texttt{ }\texttt{ lines and }\texttt{ }`\texttt{! }` blocks);
\item package installation and network calls;
\item paths read or written;
\item credential access;
\item any instruction that discourages review or suppresses output.
\end{itemize}

Organisational and personal controls: \texttt{disable\allowbreak{}Bundled\allowbreak{}Skills} turns off every bundled skill except \texttt{/\allowbreak{}doctor}; \texttt{skill\allowbreak{}Overrides} sets an individual skill to \texttt{"off"} or \texttt{"user-\allowbreak{}invocable-\allowbreak{}only"} without editing its file; a \texttt{Skill} deny rule restricts Claude's access; and \texttt{disable\allowbreak{}Skill\allowbreak{}Shell\allowbreak{}Execution} replaces inline shell command lines with a placeholder [34].


\begin{calloutbox}{palebrass}{brassdark}{2.0mm}
\textbf{CAUTION:} stars, download counts, a video recommendation, or an official-sounding name are not a security review. Pin a reviewed version and repeat the review after every update. This is M12 in Table 1.
\end{calloutbox}

\subsection*{Exercise}

Install the example skill locally, invoke \texttt{/\allowbreak{}research-\allowbreak{}brief} on a second source set, and record every place it made an assumption. Improve the skill only with a rule that generalises.


\emph{A worked solution is given in Appendix M.}


\begingroup
\def\kprows{%
\item Derive a skill from a procedure you have already performed manually, and generalise the lesson rather than the instance.
\item \texttt{allowed-\allowbreak{}tools} pre-approves tools for the invoking turn only, and clears on your next message.
\item Use \texttt{disable-\allowbreak{}model-\allowbreak{}invocation:\allowbreak{} true} for anything with side effects.
\item The rendered skill stays in the conversation for the rest of the session and is not re-read on later turns.
}%
\def\kpbody{\begin{minipage}{\textwidth}\subsection*{Key points}\begin{itemize}\kprows\end{itemize}\end{minipage}}%
\begingroup
\def\sloppy{\tolerance 9999\emergencystretch 3em\hfuzz 200pt\vfuzz 200pt}%
\hbadness=10000\vbadness=10000\hfuzz=200pt\vfuzz=200pt
\global\setbox\kpbox=\hbox{\kpbody}%
\endgroup
\par\addvspace{4.2mm}
\ifdim\dimexpr\ht\kpbox+\dp\kpbox\relax>0.30\textheight
  \typeout{HANDBOOK-KEYPOINTS broken \the\dimexpr\ht\kpbox+\dp\kpbox\relax}%
  \subsection*{Key points}
  \begin{itemize}\kprows\end{itemize}
\else
  \typeout{HANDBOOK-KEYPOINTS atomic \the\dimexpr\ht\kpbox+\dp\kpbox\relax}%
  \noindent\kpbody
\fi
\par\addvspace{1.4mm}
\endgroup


\FloatBarrier
\renewcommand{\chaptertitlelabel}{Part IV \textperiodcentered\ Chapter 16}
\setcounter{chapter}{15}
\chapter{Delegate bounded work to subagents}

A subagent is a specialised instance with its own context window, system prompt, tool access and permissions [25]. Delegate when independence, focus or parallelism improves the result — not to make a small action feel important.


\FloatBarrier
\setcounter{section}{0}
\section{Fresh context, forks, and why it matters}

An ordinary subagent starts fresh: it receives its own system prompt, the delegation message, the \texttt{CLAUDE.\allowbreak{}md} hierarchy, a git-status snapshot and any preloaded skills. It does \textbf{not} receive the parent's conversation history, output style, or auto memory [25].


A \textbf{fork} is different: it inherits the entire conversation, the parent's system prompt, tools and model, and shares the parent's prompt cache. \texttt{/\allowbreak{}subtask} creates one.


\begin{calloutbox}{palebrass}{brassdark}{2.0mm}
\textbf{CAUTION — this is the correction that matters most in this chapter.} Fork mode is \textbf{on by default in interactive sessions}, controlled by the \texttt{CLAUDE\_\allowbreak{}CODE\_\allowbreak{}FORK\_\allowbreak{}SUBAGENT} environment variable (\texttt{1} on, \texttt{0} off) [25]. A "reviewer" you create casually in an interactive session may therefore inherit the reasoning it is supposed to audit. If you need genuine independence — and the capstone of this book does — use an explicit agent definition and confirm the reviewer's context. This is M03 in Table 1.
\end{calloutbox}

\FloatBarrier
\setcounter{section}{1}
\section{Where subagents live, and which one wins}

\begingroup
\def\tblrows{%
Managed settings & Organisation-wide & 1 (highest) \\
\texttt{-\allowbreak{}-\allowbreak{}agents} CLI flag & Current session & 2 \\
\texttt{.\allowbreak{}claude/\allowbreak{}agents/\allowbreak{}} & Current project & 3 \\
\texttt{\textasciitilde{}/\allowbreak{}.\allowbreak{}claude/\allowbreak{}agents/\allowbreak{}} & All your projects & 4 \\
Plugin \texttt{agents/\allowbreak{}} & Where the plugin is enabled & 5 (lowest) \\
}%
\def\tblbody{\begin{minipage}{\textwidth}\boxcaption{Table 12 — Subagent definition precedence}
{\footnotesize\begin{tabular}{L{56.6mm}L{63.9mm}L{19.4mm}}
\toprule
\textbf{Source} & \textbf{Scope} & \textbf{Priority} \\
\midrule
\tblrows
\bottomrule\end{tabular}}\end{minipage}}%
\begingroup
\def\sloppy{\tolerance 9999\emergencystretch 3em\hfuzz 200pt\vfuzz 200pt}%
\hbadness=10000\vbadness=10000\hfuzz=200pt\vfuzz=200pt
\global\setbox\tblbox=\hbox{\tblbody}%
\endgroup
\par\addvspace{2.6mm}
\ifdim\dimexpr\ht\tblbox+\dp\tblbox\relax>0.55\textheight
  \typeout{HANDBOOK-TABLE broken \the\dimexpr\ht\tblbox+\dp\tblbox\relax}%
  \tabcaption{Table 12 — Subagent definition precedence}
  {\footnotesize\begin{longtable}{L{56.6mm}L{63.9mm}L{19.4mm}}
  \toprule
\textbf{Source} & \textbf{Scope} & \textbf{Priority} \\
\midrule\endfirsthead
  \multicolumn{3}{@{}l@{}}{%
  \sffamily\footnotesize\itshape\color{inkgrey}Table 12 — Subagent definition precedence \textemdash\ continued}\\[1.2mm]
  \toprule
\textbf{Source} & \textbf{Scope} & \textbf{Priority} \\
\midrule\endhead
  \bottomrule\endfoot
  \bottomrule\endlastfoot
  \tblrows
  \end{longtable}}%
\else
  \typeout{HANDBOOK-TABLE atomic \the\dimexpr\ht\tblbox+\dp\tblbox\relax}%
  \noindent\tblbody
\fi
\par\addvspace{2.6mm}
\endgroup

\texttt{/\allowbreak{}agents} no longer opens an agent editor \emph{(changed at 2.1.198)}; it directs you to ask Claude to create or update an agent, which you then read as Markdown yourself [35]. This is M13 in Table 1. \texttt{/\allowbreak{}list-\allowbreak{}agents} (alias \texttt{/\allowbreak{}peers}) lists subagents, teammates and other sessions [35].


\FloatBarrier
\setcounter{section}{2}
\section{An independent read-only reviewer}

\texttt{.\allowbreak{}claude/\allowbreak{}agents/\allowbreak{}evidence-\allowbreak{}reviewer.\allowbreak{}md}:


\begin{codeblock}{7.0}{8.3}
\cl{---}
\cl{name:~evidence-reviewer}
\cl{description:~Independently~review~a~draft~for~unsupported,~overstated,~stale,~or~contradictory}
\cl{~~claims.~Use~after~a~research~brief~is~drafted.}
\cl{tools:~Read,~Grep,~Glob}
\cl{permissionMode:~plan}
\cl{model:~inherit}
\cl{maxTurns:~25}
\cl{---}
\cl{}
\cl{You~are~an~independent~evidence~reviewer.~You~did~not~author~the~draft~and~you~have~no~access}
\cl{to~the~reasoning~that~produced~it.}
\cl{}
\cl{For~each~material~claim:}
\cl{1.~Locate~its~cited~source~in~the~workspace.}
\cl{2.~Classify~support~as~supports~/~partially~supports~/~contradicts~/~missing.}
\cl{3.~Check~that~date,~scope~and~attribution~match.}
\cl{4.~Report~the~exact~file~and~passage~requiring~correction.}
\cl{}
\cl{Do~not~edit~files.~Return~a~severity-ordered~review~and~a~final~pass/fail~verdict.}
\end{codeblock}

\begin{calloutbox}{palebrass}{brassdark}{2.0mm}
\textbf{CAUTION:} a subagent's own \texttt{permission\allowbreak{}Mode} does not always apply. Three parent modes defeat it. A parent in \texttt{bypass\allowbreak{}Permissions} or \texttt{accept\allowbreak{}Edits} takes precedence and cannot be overridden. A parent in \texttt{auto} hands the subagent auto mode, and any \texttt{permission\allowbreak{}Mode} in the subagent's frontmatter is ignored: the classifier evaluates its calls under the parent's own block and allow rules. And \texttt{auto} is the mode a session starts in on Pro, Max and Team plans unless settings or an organisation change it, so this is the ordinary case rather than an exotic one [25]. A read-only reviewer is only read-only if its parent is in none of those three modes.
\end{calloutbox}

\FloatBarrier
\setcounter{section}{3}
\section{The frontmatter you should know}

Only \texttt{name} and \texttt{description} are required. The rest change behaviour materially [25]: \texttt{tools} and \texttt{disallowed\allowbreak{}Tools} (both accept MCP patterns such as \texttt{mcp\_\allowbreak{}\_\allowbreak{}server\_\allowbreak{}\_\allowbreak{}*}, and \texttt{Agent(worke\allowbreak{}r,\allowbreak{} researcher\allowbreak{})} restricts which subagents this one may spawn); \texttt{model} (\texttt{sonnet}, \texttt{opus}, \texttt{haiku}, \texttt{fable}, a full ID, or \texttt{inherit}); \texttt{permission\allowbreak{}Mode}; \texttt{max\allowbreak{}Turns}; \texttt{skills} to preload; \texttt{mcp\allowbreak{}Servers}; \texttt{hooks}; \texttt{memory} (\texttt{user}, \texttt{project} or \texttt{local}, each with its own \texttt{agent-\allowbreak{}memory} directory and the same 200-line/25 KB index limit); \texttt{background}; \texttt{effort}; \texttt{isolation:\allowbreak{} worktree} to run in a temporary git worktree; \texttt{color}; and \texttt{initial\allowbreak{}Prompt}.


\FloatBarrier
\setcounter{section}{4}
\section{Limits and defaults}

\begin{itemize}
\item \textbf{Background is the default.} Background subagents run concurrently, surface permission prompts in your main session, carry a smaller built-in tool set, and return results in a later turn [25].
\item \textbf{Nesting is three layers deep} by default; at the limit the \texttt{Agent} tool is withheld. \texttt{CLAUDE\_\allowbreak{}CODE\_\allowbreak{}MAX\_\allowbreak{}SUBAGENT\_\allowbreak{}SPAWN\_\allowbreak{}DEPTH=\allowbreak{}1} disables nesting entirely [25].
\item \textbf{Twenty concurrent subagents} by default, then \texttt{Concurrent \allowbreak{}subagent li\allowbreak{}mit reached}. \texttt{CLAUDE\_\allowbreak{}CODE\_\allowbreak{}MAX\_\allowbreak{}CONCURRENT\_\allowbreak{}SUBAGENTS} changes it [25].
\item \textbf{Built-ins} — \texttt{Explore} (read-only, model capped at Opus), \texttt{Plan} (read-only) and \texttt{General-\allowbreak{}purpose}. \texttt{Explore} and \texttt{Plan} skip \texttt{CLAUDE.\allowbreak{}md} and git status for speed. Disable them with \texttt{permissions.\allowbreak{}deny:\allowbreak{} ["Agent(Ex\allowbreak{}plore)"]} [25].
\end{itemize}

These mechanisms differ mainly in what they inherit. Figure 5 compares the five topologies on that axis, which is also the axis that decides how independent a second opinion really is.


\begin{figure}[tbp]
\begin{codefig}{9.0}{10.6}
\cl{~~~SUBAGENT~(non-fork)~~~~~~~~~~FORK~~(~/subtask~)~~~~~~~~~TEAMMATE}
\cl{~~~\pmboxdrawuni{250C}\pmboxdrawuni{2500}\pmboxdrawuni{2500}\pmboxdrawuni{2500}\pmboxdrawuni{2500}\pmboxdrawuni{2500}\pmboxdrawuni{2500}\pmboxdrawuni{2500}\pmboxdrawuni{2500}\pmboxdrawuni{2500}\pmboxdrawuni{2500}\pmboxdrawuni{2500}\pmboxdrawuni{2500}\pmboxdrawuni{2500}\pmboxdrawuni{2500}\pmboxdrawuni{2510}~~~~~~~~~~~~~\pmboxdrawuni{250C}\pmboxdrawuni{2500}\pmboxdrawuni{2500}\pmboxdrawuni{2500}\pmboxdrawuni{2500}\pmboxdrawuni{2500}\pmboxdrawuni{2500}\pmboxdrawuni{2500}\pmboxdrawuni{2500}\pmboxdrawuni{2500}\pmboxdrawuni{2500}\pmboxdrawuni{2500}\pmboxdrawuni{2500}\pmboxdrawuni{2500}\pmboxdrawuni{2500}\pmboxdrawuni{2510}~~~~~~~~~~~\pmboxdrawuni{250C}\pmboxdrawuni{2500}\pmboxdrawuni{2500}\pmboxdrawuni{2500}\pmboxdrawuni{2500}\pmboxdrawuni{2500}\pmboxdrawuni{2500}\pmboxdrawuni{2500}\pmboxdrawuni{2500}\pmboxdrawuni{2500}\pmboxdrawuni{2500}\pmboxdrawuni{2500}\pmboxdrawuni{2500}\pmboxdrawuni{2500}\pmboxdrawuni{2500}\pmboxdrawuni{2510}}
\cl{~~~\pmboxdrawuni{2502}~own~context~~\pmboxdrawuni{2502}~~~~~~~~~~~~~\pmboxdrawuni{2502}~FULL~parent~~\pmboxdrawuni{2502}~~~~~~~~~~~\pmboxdrawuni{2502}~own~context~~\pmboxdrawuni{2502}}
\cl{~~~\pmboxdrawuni{2502}~own~prompt~~~\pmboxdrawuni{2502}~~~~~~~~~~~~~\pmboxdrawuni{2502}~conversation~\pmboxdrawuni{2502}~~~~~~~~~~~\pmboxdrawuni{2502}~own~session~~\pmboxdrawuni{2502}}
\cl{~~~\pmboxdrawuni{2502}~own~tools~~~~\pmboxdrawuni{2502}~~~~~~~~~~~~~\pmboxdrawuni{2502}~parent~tools~\pmboxdrawuni{2502}~~~~~~~~~~~\pmboxdrawuni{2502}~lead\textquotesingle{}s~perms~\pmboxdrawuni{2502}}
\cl{~~~\pmboxdrawuni{2502}~CLAUDE.md~✓~~\pmboxdrawuni{2502}~~~~~~~~~~~~~\pmboxdrawuni{2502}~same~model~~~\pmboxdrawuni{2502}~~~~~~~~~~~\pmboxdrawuni{2502}~CLAUDE.md~✓~~\pmboxdrawuni{2502}}
\cl{~~~\pmboxdrawuni{2502}~history~~~✗~~\pmboxdrawuni{2502}~~~~~~~~~~~~~\pmboxdrawuni{2502}~shared~cache~\pmboxdrawuni{2502}~~~~~~~~~~~\pmboxdrawuni{2502}~history~~~✗~~\pmboxdrawuni{2502}}
\cl{~~~\pmboxdrawuni{2514}\pmboxdrawuni{2500}\pmboxdrawuni{2500}\pmboxdrawuni{2500}\pmboxdrawuni{2500}\pmboxdrawuni{2500}\pmboxdrawuni{2500}\pmboxdrawuni{252C}\pmboxdrawuni{2500}\pmboxdrawuni{2500}\pmboxdrawuni{2500}\pmboxdrawuni{2500}\pmboxdrawuni{2500}\pmboxdrawuni{2500}\pmboxdrawuni{2500}\pmboxdrawuni{2518}~~~~~~~~~~~~~\pmboxdrawuni{2514}\pmboxdrawuni{2500}\pmboxdrawuni{2500}\pmboxdrawuni{2500}\pmboxdrawuni{2500}\pmboxdrawuni{2500}\pmboxdrawuni{2500}\pmboxdrawuni{252C}\pmboxdrawuni{2500}\pmboxdrawuni{2500}\pmboxdrawuni{2500}\pmboxdrawuni{2500}\pmboxdrawuni{2500}\pmboxdrawuni{2500}\pmboxdrawuni{2500}\pmboxdrawuni{2518}~~~~~~~~~~~\pmboxdrawuni{2514}\pmboxdrawuni{2500}\pmboxdrawuni{2500}\pmboxdrawuni{2500}\pmboxdrawuni{2500}\pmboxdrawuni{2500}\pmboxdrawuni{2500}\pmboxdrawuni{252C}\pmboxdrawuni{2500}\pmboxdrawuni{2500}\pmboxdrawuni{2500}\pmboxdrawuni{2500}\pmboxdrawuni{2500}\pmboxdrawuni{2500}\pmboxdrawuni{2500}\pmboxdrawuni{2518}}
\cl{~~~~~~~~~~\pmboxdrawuni{2502}~returns~a~result~~~~~~~~~~~\pmboxdrawuni{2502}~returns~a~result~~~~~~~~~\pmboxdrawuni{2502}~messages}
\cl{~~~~~~~~~~▼~~~~~~~~~~~~~~~~~~~~~~~~~~~~▼~~~~~~~~~~~~~~~~~~~~~~~~~~▼}
\cl{~~~~~~main~conversation~~~~~~~~~~~main~conversation~~~~~~~~peers~+~task~list}
\cl{}
\cl{~~~WORKFLOW~AGENT~~~~~~~~~~~~~~~~~~~~~~~~~~AGENT~TEAM~LEAD}
\cl{~~~\pmboxdrawuni{250C}\pmboxdrawuni{2500}\pmboxdrawuni{2500}\pmboxdrawuni{2500}\pmboxdrawuni{2500}\pmboxdrawuni{2500}\pmboxdrawuni{2500}\pmboxdrawuni{2500}\pmboxdrawuni{2500}\pmboxdrawuni{2500}\pmboxdrawuni{2500}\pmboxdrawuni{2500}\pmboxdrawuni{2500}\pmboxdrawuni{2500}\pmboxdrawuni{2500}\pmboxdrawuni{2500}\pmboxdrawuni{2500}\pmboxdrawuni{2500}\pmboxdrawuni{2500}\pmboxdrawuni{2500}\pmboxdrawuni{2500}\pmboxdrawuni{2500}\pmboxdrawuni{2500}\pmboxdrawuni{2500}\pmboxdrawuni{2500}\pmboxdrawuni{2500}\pmboxdrawuni{2500}\pmboxdrawuni{2510}~~~~~~~~~~~~\pmboxdrawuni{250C}\pmboxdrawuni{2500}\pmboxdrawuni{2500}\pmboxdrawuni{2500}\pmboxdrawuni{2500}\pmboxdrawuni{2500}\pmboxdrawuni{2500}\pmboxdrawuni{2500}\pmboxdrawuni{2500}\pmboxdrawuni{2500}\pmboxdrawuni{2500}\pmboxdrawuni{2500}\pmboxdrawuni{2500}\pmboxdrawuni{2500}\pmboxdrawuni{2500}\pmboxdrawuni{2500}\pmboxdrawuni{2500}\pmboxdrawuni{2500}\pmboxdrawuni{2500}\pmboxdrawuni{2500}\pmboxdrawuni{2500}\pmboxdrawuni{2500}\pmboxdrawuni{2500}\pmboxdrawuni{2500}\pmboxdrawuni{2500}\pmboxdrawuni{2500}\pmboxdrawuni{2500}\pmboxdrawuni{2510}}
\cl{~~~\pmboxdrawuni{2502}~spawned~by~a~script~~~~~~\pmboxdrawuni{2502}~~~~~~~~~~~~\pmboxdrawuni{2502}~your~session~~~~~~~~~~~~~~\pmboxdrawuni{2502}}
\cl{~~~\pmboxdrawuni{2502}~ALWAYS~acceptEdits~~~~~~~\pmboxdrawuni{2502}~~~~~~~~~~~~\pmboxdrawuni{2502}~approves~teammate~plans~~~\pmboxdrawuni{2502}}
\cl{~~~\pmboxdrawuni{2502}~result~→~script~variable~\pmboxdrawuni{2502}~~~~~~~~~~~~\pmboxdrawuni{2502}~autonomously~~~~~~~~~~~~~~\pmboxdrawuni{2502}}
\cl{~~~\pmboxdrawuni{2502}~up~to~16~at~once,~1,000~~\pmboxdrawuni{2502}~~~~~~~~~~~~\pmboxdrawuni{2502}~one~team~per~session~~~~~~\pmboxdrawuni{2502}}
\cl{~~~\pmboxdrawuni{2502}~per~run~~~~~~~~~~~~~~~~~~\pmboxdrawuni{2502}~~~~~~~~~~~~\pmboxdrawuni{2502}~no~nested~teams~~~~~~~~~~~\pmboxdrawuni{2502}}
\cl{~~~\pmboxdrawuni{2514}\pmboxdrawuni{2500}\pmboxdrawuni{2500}\pmboxdrawuni{2500}\pmboxdrawuni{2500}\pmboxdrawuni{2500}\pmboxdrawuni{2500}\pmboxdrawuni{2500}\pmboxdrawuni{2500}\pmboxdrawuni{2500}\pmboxdrawuni{2500}\pmboxdrawuni{2500}\pmboxdrawuni{2500}\pmboxdrawuni{2500}\pmboxdrawuni{2500}\pmboxdrawuni{2500}\pmboxdrawuni{2500}\pmboxdrawuni{2500}\pmboxdrawuni{2500}\pmboxdrawuni{2500}\pmboxdrawuni{2500}\pmboxdrawuni{2500}\pmboxdrawuni{2500}\pmboxdrawuni{2500}\pmboxdrawuni{2500}\pmboxdrawuni{2500}\pmboxdrawuni{2500}\pmboxdrawuni{2518}~~~~~~~~~~~~\pmboxdrawuni{2514}\pmboxdrawuni{2500}\pmboxdrawuni{2500}\pmboxdrawuni{2500}\pmboxdrawuni{2500}\pmboxdrawuni{2500}\pmboxdrawuni{2500}\pmboxdrawuni{2500}\pmboxdrawuni{2500}\pmboxdrawuni{2500}\pmboxdrawuni{2500}\pmboxdrawuni{2500}\pmboxdrawuni{2500}\pmboxdrawuni{2500}\pmboxdrawuni{2500}\pmboxdrawuni{2500}\pmboxdrawuni{2500}\pmboxdrawuni{2500}\pmboxdrawuni{2500}\pmboxdrawuni{2500}\pmboxdrawuni{2500}\pmboxdrawuni{2500}\pmboxdrawuni{2500}\pmboxdrawuni{2500}\pmboxdrawuni{2500}\pmboxdrawuni{2500}\pmboxdrawuni{2500}\pmboxdrawuni{2518}}
\cl{}
\cl{~~~Independence~for~review~purposes:~~subagent~>~teammate~>~workflow~agent~>~fork}
\cl{~~~Fork~mode~is~the~interactive~DEFAULT.~An~unqualified~request~for~"a~reviewer"}
\cl{~~~in~an~interactive~session~produces~the~least~independent~option~on~this~chart.}
\end{codefig}
\figcaption{Figure 5 — Five delegation topologies and what each one inherits. Independence decreases from left to right along the top row.}
\end{figure}


\FloatBarrier
\setcounter{section}{5}
\section{Skill or subagent?}

\begingroup
\def\tblrows{%
Apply a procedure in the current context & Skill \\
Independent judgement, fresh context & Subagent (non-fork) \\
Continue the current reasoning elsewhere without polluting context & Fork / \texttt{/\allowbreak{}subtask} \\
Parallel bounded analysis & Several subagents \\
Deterministic file or command action & A script or hook — not a model \\
Reusable procedure that coordinates reviewers & Skill that invokes subagents \\
Dozens to hundreds of agents with codified orchestration & Dynamic workflow (Chapter 22) \\
Peers that talk to each other and share a task list & Agent teams (Chapter 30) \\
}%
\def\tblbody{\begin{minipage}{\textwidth}\boxcaption{Table 13 — Choosing a delegation mechanism}
{\footnotesize\begin{tabular}{L{97.6mm}L{45.4mm}}
\toprule
\textbf{Need} & \textbf{Prefer} \\
\midrule
\tblrows
\bottomrule\end{tabular}}\end{minipage}}%
\begingroup
\def\sloppy{\tolerance 9999\emergencystretch 3em\hfuzz 200pt\vfuzz 200pt}%
\hbadness=10000\vbadness=10000\hfuzz=200pt\vfuzz=200pt
\global\setbox\tblbox=\hbox{\tblbody}%
\endgroup
\par\addvspace{2.6mm}
\ifdim\dimexpr\ht\tblbox+\dp\tblbox\relax>0.55\textheight
  \typeout{HANDBOOK-TABLE broken \the\dimexpr\ht\tblbox+\dp\tblbox\relax}%
  \tabcaption{Table 13 — Choosing a delegation mechanism}
  {\footnotesize\begin{longtable}{L{97.6mm}L{45.4mm}}
  \toprule
\textbf{Need} & \textbf{Prefer} \\
\midrule\endfirsthead
  \multicolumn{2}{@{}l@{}}{%
  \sffamily\footnotesize\itshape\color{inkgrey}Table 13 — Choosing a delegation mechanism \textemdash\ continued}\\[1.2mm]
  \toprule
\textbf{Need} & \textbf{Prefer} \\
\midrule\endhead
  \bottomrule\endfoot
  \bottomrule\endlastfoot
  \tblrows
  \end{longtable}}%
\else
  \typeout{HANDBOOK-TABLE atomic \the\dimexpr\ht\tblbox+\dp\tblbox\relax}%
  \noindent\tblbody
\fi
\par\addvspace{2.6mm}
\endgroup

\FloatBarrier
\setcounter{section}{6}
\section{Delegation brief}

Every subagent prompt should state the exact question; the files or sources it may use; the tools and actions allowed; the required return format; the success criteria; and when to stop and report uncertainty. Never delegate "fix everything".


\begingroup
\def\kprows{%
\item An ordinary subagent starts fresh; a fork inherits everything — and fork mode is the interactive default.
\item A parent in \texttt{accept\allowbreak{}Edits}, \texttt{bypass\allowbreak{}Permissions} or \texttt{auto} defeats a subagent's own \texttt{permission\allowbreak{}Mode}; \texttt{auto} is where sessions start on Pro, Max and Team.
\item Background is the default, nesting goes three layers deep, and twenty concurrent subagents is the cap.
\item Delegate a bounded question with named sources, allowed tools, a return format and a stop condition — never 'fix everything'.
}%
\def\kpbody{\begin{minipage}{\textwidth}\subsection*{Key points}\begin{itemize}\kprows\end{itemize}\end{minipage}}%
\begingroup
\def\sloppy{\tolerance 9999\emergencystretch 3em\hfuzz 200pt\vfuzz 200pt}%
\hbadness=10000\vbadness=10000\hfuzz=200pt\vfuzz=200pt
\global\setbox\kpbox=\hbox{\kpbody}%
\endgroup
\par\addvspace{4.2mm}
\ifdim\dimexpr\ht\kpbox+\dp\kpbox\relax>0.30\textheight
  \typeout{HANDBOOK-KEYPOINTS broken \the\dimexpr\ht\kpbox+\dp\kpbox\relax}%
  \subsection*{Key points}
  \begin{itemize}\kprows\end{itemize}
\else
  \typeout{HANDBOOK-KEYPOINTS atomic \the\dimexpr\ht\kpbox+\dp\kpbox\relax}%
  \noindent\kpbody
\fi
\par\addvspace{1.4mm}
\endgroup


\FloatBarrier
\renewcommand{\chaptertitlelabel}{Part IV \textperiodcentered\ Chapter 17}
\setcounter{chapter}{16}
\chapter{Event-driven behaviour with hooks}

A hook runs a configured action at a lifecycle event. Hooks are how you make something happen \emph{because of an event} rather than because the model remembered a sentence. They are the only mechanism in this book that is both deterministic and user-authored.


\FloatBarrier
\setcounter{section}{0}
\section{Events}

Claude Code exposes a large event surface. The ones that carry most of the practical weight [73]:


\begin{itemize}
\item \texttt{Session\allowbreak{}Start}, \texttt{Session\allowbreak{}End} — matchable by start type (\texttt{startup}, \texttt{resume}, \texttt{clear}, \texttt{compact}, \texttt{fork}) and by end reason.
\item \texttt{User\allowbreak{}Prompt\allowbreak{}Submit} — before Claude processes a prompt; exit 2 rejects it.
\item \texttt{Pre\allowbreak{}Tool\allowbreak{}Use} — before a tool call; exit 2 blocks it.
\item \texttt{Permission\allowbreak{}Request} — when a call needs a permission decision. Return a JSON decision rather than exit 2.
\item \texttt{Permission\allowbreak{}Denied} — after auto mode denies a call; may set \texttt{retry}.
\item \texttt{Post\allowbreak{}Tool\allowbreak{}Use}, \texttt{Post\allowbreak{}Tool\allowbreak{}Use\allowbreak{}Failure}, \texttt{Post\allowbreak{}Tool\allowbreak{}Batch} — after the fact; cannot undo an external side effect.
\item \texttt{Stop}, \texttt{Stop\allowbreak{}Failure}, \texttt{Subagent\allowbreak{}Start}, \texttt{Subagent\allowbreak{}Stop}, \texttt{Teammate\allowbreak{}Idle} — turn and agent lifecycle; exit 2 on \texttt{Stop} prevents stopping.
\item \texttt{Task\allowbreak{}Created}, \texttt{Task\allowbreak{}Completed} — shared task list; exit 2 blocks the operation.
\item \texttt{Pre\allowbreak{}Compact}, \texttt{Post\allowbreak{}Compact}, \texttt{Instruction\allowbreak{}s\allowbreak{}Loaded}, \texttt{Config\allowbreak{}Change}, \texttt{Cwd\allowbreak{}Changed}, \texttt{Directory\allowbreak{}Added}, \texttt{File\allowbreak{}Changed}, \texttt{Worktree\allowbreak{}Create}, \texttt{Worktree\allowbreak{}Remove}, \texttt{Notificatio\allowbreak{}n}, \texttt{Message\allowbreak{}Display}, \texttt{Elicitation}, \texttt{Elicitation\allowbreak{}Result}.
\end{itemize}

\texttt{Config\allowbreak{}Change} deserves a specific mention: it lets an organisation audit or block settings changes made during a session [53].


\FloatBarrier
\setcounter{section}{1}
\section{Handler types}

Five handler types are available [73]: \texttt{command} (shell or exec form, with \texttt{async} and \texttt{async\allowbreak{}Rewake}), \texttt{http} (POST, with \texttt{allowed\allowbreak{}Env\allowbreak{}Vars} header interpolation), \texttt{mcp\_\allowbreak{}tool} (call a tool on an already-connected server), \texttt{prompt} (single-turn model evaluation, \texttt{\$ARGUMENTS} placeholder, 30-second default), and \texttt{agent} (spawns a subagent; experimental).


Default timeouts are 600 seconds for \texttt{command}, \texttt{http} and \texttt{mcp\_\allowbreak{}tool} — 30 seconds under \texttt{User\allowbreak{}Prompt\allowbreak{}Submit}, 10 under \texttt{Message\allowbreak{}Display} — 30 for \texttt{prompt} and 60 for \texttt{agent}. \texttt{Session\allowbreak{}End} hooks share a 1.5-second budget by default [73].


\texttt{\$\{CLAUDE\_\allowbreak{}PROJECT\_\allowbreak{}DIR\}}, \texttt{\$\{CLAUDE\_\allowbreak{}PLUGIN\_\allowbreak{}ROOT\}} and \texttt{\$\{CLAUDE\_\allowbreak{}PLUGIN\_\allowbreak{}DATA\}} are substituted in hook commands; \texttt{\$\{CLAUDE\_\allowbreak{}PROJECT\_\allowbreak{}DIR\}} stays at the project root even inside a worktree [73].


\FloatBarrier
\setcounter{section}{2}
\section{A blocking example}

\begin{codeblock}{9.0}{10.6}
\cl{\{}
\cl{~~"hooks":~\{}
\cl{~~~~"PreToolUse":~[}
\cl{~~~~~~\{}
\cl{~~~~~~~~"matcher":~"Bash",}
\cl{~~~~~~~~"hooks":~[}
\cl{~~~~~~~~~~\{}
\cl{~~~~~~~~~~~~"type":~"command",}
\cl{~~~~~~~~~~~~"if":~"Bash(rm~*)",}
\cl{~~~~~~~~~~~~"command":~"\$\{CLAUDE\_PROJECT\_DIR\}/.claude/hooks/guard-rm.sh",}
\cl{~~~~~~~~~~~~"timeout":~20,}
\cl{~~~~~~~~~~~~"statusMessage":~"Checking~destructive~command…"}
\cl{~~~~~~~~~~\}}
\cl{~~~~~~~~]}
\cl{~~~~~~\}}
\cl{~~~~]}
\cl{~~\}}
\cl{\}}
\end{codeblock}

\begin{codeblock}{8.0}{9.4}
\cl{\#!/usr/bin/env~bash}
\cl{\#~.claude/hooks/guard-rm.sh~—~deny~rm~-rf;~stay~silent~otherwise}
\cl{COMMAND=\$(jq~-r~\textquotesingle{}.tool\_input.command\textquotesingle{})}
\cl{if~printf~\textquotesingle{}\%s\textquotesingle{}~"\$COMMAND"~|~grep~-q~\textquotesingle{}rm~-rf\textquotesingle{};~then}
\cl{~~jq~-n~\textquotesingle{}\{hookSpecificOutput:\{hookEventName:"PreToolUse",}
\cl{~~~~~~~~~~permissionDecision:"deny",}
\cl{~~~~~~~~~~permissionDecisionReason:"Destructive~rm~-rf~blocked~by~project~policy"\}\}\textquotesingle{}}
\cl{else}
\cl{~~exit~0}
\cl{fi}
\end{codeblock}

\textbf{Hooks do not bypass permission rules.} Deny and ask rules are evaluated regardless of what a \texttt{Pre\allowbreak{}Tool\allowbreak{}Use} hook returns: a matching deny still blocks, and a matching ask still prompts, even when the hook returned \texttt{allow} [28]. Conversely, a hook that exits 2 stops the call \emph{before} permission rules are evaluated, so it overrides an allow rule. That asymmetry is the basis of the recommended pattern: allow \texttt{Bash} broadly, then block specific commands with a hook [28].


\FloatBarrier
\setcounter{section}{3}
\section{Safer creation workflow}

\begin{enumerate}
\item Define the exact event and the narrowest matcher.
\item Write a standalone script with deterministic input and output.
\item Test it manually against benign fixtures.
\item Ask Claude to propose the configuration as a diff.
\item Confirm paths, quoting, timeout and exit behaviour.
\item Test in the practice repository, including the failure path.
\item Add failure logging and write down how to remove it.
\end{enumerate}

Inspect what is actually configured with \texttt{/\allowbreak{}hooks}, which shows event, matcher, type, source file and full command [35], [72].


\FloatBarrier
\setcounter{section}{4}
\section{Threat model}

\begin{itemize}
\item Can its inputs contain untrusted filenames, prompt text or tool output?
\item Does it build a shell command by string concatenation?
\item Can it approve a permission request? (\texttt{Permission\allowbreak{}Request} hooks can. Treat them as privileged code.)
\item What happens if it runs twice, or in parallel?
\item Does it send data over the network? If so, is the URL on \texttt{allowed\allowbreak{}Http\allowbreak{}Hook\allowbreak{}Urls}?
\item Does a failure block work, or silently allow it?
\item Is a post-action check being mistaken for prevention?
\end{itemize}

Organisational controls: \texttt{disable\allowbreak{}All\allowbreak{}Hooks} turns hooks off subject to settings precedence but cannot disable managed policy hooks; \texttt{allow\allowbreak{}Managed\allowbreak{}Hooks\allowbreak{}Only} in managed settings blocks user, project, local and plugin hooks entirely; \texttt{allowed\allowbreak{}Http\allowbreak{}Hook\allowbreak{}Urls} restricts HTTP hook destinations from every source [73].


\begin{calloutbox}{palebrass}{brassdark}{2.0mm}
\textbf{CAUTION:} a natural-language rule is probabilistic; a hook is executable automation. That makes it more reliable and correspondingly more dangerous. Note also that \texttt{/\allowbreak{}goal} (Chapter 21) is implemented as a session-scoped prompt Stop hook, so \texttt{disable\allowbreak{}All\allowbreak{}Hooks} or \texttt{allow\allowbreak{}Managed\allowbreak{}Hooks\allowbreak{}Only} disables it too [74].
\end{calloutbox}

\subsection*{Exercise}

Create a project-only \texttt{Stop} hook that appends a timestamped, non-sensitive completion record to a local file. It must not make a network call, approve a permission, or include conversation content. Test success, failure and removal.


\emph{A worked solution is given in Appendix M.}


\begingroup
\def\kprows{%
\item A hook is executable automation: more reliable than a written rule, and correspondingly more dangerous.
\item A hook exiting 2 blocks a call before permission rules are evaluated; no hook decision bypasses a matching deny rule.
\item \texttt{Permission\allowbreak{}Request} hooks can influence approval and must be treated as privileged code.
\item A post-action hook cannot undo an external side effect. Prevention happens at \texttt{Pre\allowbreak{}Tool\allowbreak{}Use} or not at all.
}%
\def\kpbody{\begin{minipage}{\textwidth}\subsection*{Key points}\begin{itemize}\kprows\end{itemize}\end{minipage}}%
\begingroup
\def\sloppy{\tolerance 9999\emergencystretch 3em\hfuzz 200pt\vfuzz 200pt}%
\hbadness=10000\vbadness=10000\hfuzz=200pt\vfuzz=200pt
\global\setbox\kpbox=\hbox{\kpbody}%
\endgroup
\par\addvspace{4.2mm}
\ifdim\dimexpr\ht\kpbox+\dp\kpbox\relax>0.30\textheight
  \typeout{HANDBOOK-KEYPOINTS broken \the\dimexpr\ht\kpbox+\dp\kpbox\relax}%
  \subsection*{Key points}
  \begin{itemize}\kprows\end{itemize}
\else
  \typeout{HANDBOOK-KEYPOINTS atomic \the\dimexpr\ht\kpbox+\dp\kpbox\relax}%
  \noindent\kpbody
\fi
\par\addvspace{1.4mm}
\endgroup


\FloatBarrier
\renewcommand{\chaptertitlelabel}{Part IV \textperiodcentered\ Chapter 18}
\setcounter{chapter}{17}
\chapter{Evaluate and install plugins}

A plugin is a distributable bundle. A marketplace is a catalogue pointing Claude Code at plugins; it is not an app-store security review.


\FloatBarrier
\setcounter{section}{0}
\section{What a plugin can contain}

This inventory is the reason plugins deserve a gate [37], [78]:


\begingroup
\def\tblrows{%
\texttt{.\allowbreak{}claude-\allowbreak{}plugin/\allowbreak{}plugin.\allowbreak{}json} & Manifest: name (the skill namespace), description, version, author \\
\texttt{skills/\allowbreak{}} & Skills, invoked as \texttt{/\allowbreak{}plugin-\allowbreak{}name:\allowbreak{}skill-\allowbreak{}name} \\
\texttt{commands/\allowbreak{}} & Legacy flat-file skills \\
\texttt{agents/\allowbreak{}} & Subagent definitions \\
\texttt{hooks/\allowbreak{}hooks.\allowbreak{}json} & Event handlers \\
\texttt{.\allowbreak{}mcp.\allowbreak{}json} & MCP server configurations \\
\texttt{.\allowbreak{}lsp.\allowbreak{}json} & Language-server configurations \\
\texttt{monitors/\allowbreak{}monitors.\allowbreak{}json} & \textbf{Background monitors} started automatically when the plugin is active; each stdout line reaches Claude as a notification \\
\texttt{bin/\allowbreak{}} & \textbf{Executables added to the Bash tool's \texttt{PATH}} while the plugin is enabled \\
\texttt{settings.\allowbreak{}json} & Default settings applied when enabled — including \texttt{agent}, which activates one of the plugin's agents as the \textbf{main thread}, replacing the system prompt, tool restrictions and model \\
\texttt{output-\allowbreak{}styles/\allowbreak{}} & Output styles, which a plugin may force with \texttt{force-\allowbreak{}for-\allowbreak{}plugin} \\
}%
\def\tblbody{\begin{minipage}{\textwidth}\boxcaption{Table 14 — Plugin components, all at the plugin root}
{\footnotesize\begin{tabular}{H{53.3mm}L{89.7mm}}
\toprule
\textbf{Directory or file} & \textbf{Contents} \\
\midrule
\tblrows
\bottomrule\end{tabular}}\end{minipage}}%
\begingroup
\def\sloppy{\tolerance 9999\emergencystretch 3em\hfuzz 200pt\vfuzz 200pt}%
\hbadness=10000\vbadness=10000\hfuzz=200pt\vfuzz=200pt
\global\setbox\tblbox=\hbox{\tblbody}%
\endgroup
\par\addvspace{2.6mm}
\ifdim\dimexpr\ht\tblbox+\dp\tblbox\relax>0.55\textheight
  \typeout{HANDBOOK-TABLE broken \the\dimexpr\ht\tblbox+\dp\tblbox\relax}%
  \tabcaption{Table 14 — Plugin components, all at the plugin root}
  {\footnotesize\begin{longtable}{H{53.3mm}L{89.7mm}}
  \toprule
\textbf{Directory or file} & \textbf{Contents} \\
\midrule\endfirsthead
  \multicolumn{2}{@{}l@{}}{%
  \sffamily\footnotesize\itshape\color{inkgrey}Table 14 — Plugin components, all at the plugin root \textemdash\ continued}\\[1.2mm]
  \toprule
\textbf{Directory or file} & \textbf{Contents} \\
\midrule\endhead
  \bottomrule\endfoot
  \bottomrule\endlastfoot
  \tblrows
  \end{longtable}}%
\else
  \typeout{HANDBOOK-TABLE atomic \the\dimexpr\ht\tblbox+\dp\tblbox\relax}%
  \noindent\tblbody
\fi
\par\addvspace{2.6mm}
\endgroup

Read that table twice. A plugin can place executables on your \texttt{PATH}, start background processes, change the system prompt of your main session, and register hooks — all on enable.


\FloatBarrier
\setcounter{section}{1}
\section{Installation gate}

\begin{codeblock}{9.0}{10.6}
\cl{/plugin~marketplace~add~owner/repository}
\cl{/plugin~install~plugin-name@marketplace-name}
\cl{/reload-plugins}
\end{codeblock}

\begin{codeblock}{9.0}{10.6}
\cl{claude~plugin~validate~./your-plugin}
\cl{claude~--plugin-dir~./my-plugin~~~~~~~~\#~local~test;~also~accepts~a~.zip}
\end{codeblock}

Before adding a marketplace or plugin:


\begin{enumerate}
\item Confirm the exact publisher and repository.
\item Read the manifest and inventory every component against Table 14.
\item Review hooks, executables, install scripts, MCP endpoints, monitors and tool grants.
\item Search for network calls, credential reads, persistence and automatic updates.
\item Pin a reviewed version. A plugin without a \texttt{version} field takes its version from a fallback chain, and a \texttt{command} source updates outside the version gate [37].
\item Test in an isolated profile or disposable environment, using \texttt{-\allowbreak{}-\allowbreak{}plugin-\allowbreak{}dir}.
\item Measure context and usage impact.
\item Record owner, purpose, version, approval date and removal procedure.
\end{enumerate}

Anthropic maintains two public marketplaces: \texttt{claude-\allowbreak{}plugins-\allowbreak{}official}, curated at Anthropic's discretion and registered automatically on first interactive launch, and \texttt{claude-\allowbreak{}community}, where third-party submissions land after review and automated safety screening; approved community plugins are pinned to a commit SHA [37]. Neither is a warranty.


\FloatBarrier
\setcounter{section}{2}
\section{Build instead of bundle}

If you need one procedure from a large plugin, writing a small local skill using the idea is usually safer and simpler than importing opaque code and unused connectors. Preserve licences and attribution when adapting third-party material.


\begin{calloutbox}{palebrass}{brassdark}{2.0mm}
\textbf{CAUTION:} official documentation states that plugins and marketplaces can execute arbitrary code with your user privileges [38]. "Installed" means "trusted to run", not "visible in a menu". See M18 in Table 1.
\end{calloutbox}

\begingroup
\def\kprows{%
\item A plugin may ship executables placed on the shell path, background monitors, hooks, MCP and LSP servers, and a settings key that replaces the main session's agent.
\item 'Installed' means 'trusted to run', not 'visible in a menu'.
\item Pin a reviewed version, test with \texttt{-\allowbreak{}-\allowbreak{}plugin-\allowbreak{}dir} in an isolated profile, and record owner, purpose, version and removal procedure.
\item If you need one procedure from a large plugin, write a small local skill instead.
}%
\def\kpbody{\begin{minipage}{\textwidth}\subsection*{Key points}\begin{itemize}\kprows\end{itemize}\end{minipage}}%
\begingroup
\def\sloppy{\tolerance 9999\emergencystretch 3em\hfuzz 200pt\vfuzz 200pt}%
\hbadness=10000\vbadness=10000\hfuzz=200pt\vfuzz=200pt
\global\setbox\kpbox=\hbox{\kpbody}%
\endgroup
\par\addvspace{4.2mm}
\ifdim\dimexpr\ht\kpbox+\dp\kpbox\relax>0.30\textheight
  \typeout{HANDBOOK-KEYPOINTS broken \the\dimexpr\ht\kpbox+\dp\kpbox\relax}%
  \subsection*{Key points}
  \begin{itemize}\kprows\end{itemize}
\else
  \typeout{HANDBOOK-KEYPOINTS atomic \the\dimexpr\ht\kpbox+\dp\kpbox\relax}%
  \noindent\kpbody
\fi
\par\addvspace{1.4mm}
\endgroup

\FloatBarrier
\renewcommand{\chaptertitlelabel}{Part IV \textperiodcentered\ Chapter 19}
\setcounter{chapter}{18}
\chapter{Output styles: Claude Code beyond software engineering}

Output styles modify the \textbf{system prompt}: they change role, tone and default response format, and — uniquely among the mechanisms in this book — they can remove Claude Code's built-in software-engineering instructions altogether [79]. For a reader using Claude Code for research, operations, analysis or writing, this is the most direct fit between the product and the task, and it was absent from the first edition entirely.


\FloatBarrier
\setcounter{section}{0}
\section{The built-in styles}

\begingroup
\def\tblrows{%
Default & The standard software-engineering system prompt \\
Proactive & Executes immediately, makes reasonable assumptions rather than pausing on routine decisions, prefers action over planning \\
Concise & Leads with the result, skips preamble and narration, keeps responses short — while always preserving error reports, security warnings and destructive-action confirmations in full (2.1.237+) \\
Explanatory & Adds educational "Insights" alongside the work \\
Learning & Collaborative; inserts \texttt{TODO(human)} markers for you to implement \\
}%
\def\tblbody{\begin{minipage}{\textwidth}\boxcaption{Table 15 — Built-in output styles}
{\footnotesize\begin{tabular}{L{22.5mm}L{120.5mm}}
\toprule
\textbf{Style} & \textbf{Behaviour} \\
\midrule
\tblrows
\bottomrule\end{tabular}}\end{minipage}}%
\begingroup
\def\sloppy{\tolerance 9999\emergencystretch 3em\hfuzz 200pt\vfuzz 200pt}%
\hbadness=10000\vbadness=10000\hfuzz=200pt\vfuzz=200pt
\global\setbox\tblbox=\hbox{\tblbody}%
\endgroup
\par\addvspace{2.6mm}
\ifdim\dimexpr\ht\tblbox+\dp\tblbox\relax>0.55\textheight
  \typeout{HANDBOOK-TABLE broken \the\dimexpr\ht\tblbox+\dp\tblbox\relax}%
  \tabcaption{Table 15 — Built-in output styles}
  {\footnotesize\begin{longtable}{L{22.5mm}L{120.5mm}}
  \toprule
\textbf{Style} & \textbf{Behaviour} \\
\midrule\endfirsthead
  \multicolumn{2}{@{}l@{}}{%
  \sffamily\footnotesize\itshape\color{inkgrey}Table 15 — Built-in output styles \textemdash\ continued}\\[1.2mm]
  \toprule
\textbf{Style} & \textbf{Behaviour} \\
\midrule\endhead
  \bottomrule\endfoot
  \bottomrule\endlastfoot
  \tblrows
  \end{longtable}}%
\else
  \typeout{HANDBOOK-TABLE atomic \the\dimexpr\ht\tblbox+\dp\tblbox\relax}%
  \noindent\tblbody
\fi
\par\addvspace{2.6mm}
\endgroup

\begin{calloutbox}{palebrass}{brassdark}{2.0mm}
\textbf{CAUTION:} Proactive is stronger autonomous-execution guidance than auto mode applies, and it works \textbf{without changing your permission mode} [79]. Your permission mode still decides what runs without asking, so Proactive plus Manual mode is a reasonable combination and Proactive plus \texttt{bypass\allowbreak{}Permissions} is not.
\end{calloutbox}

\FloatBarrier
\setcounter{section}{1}
\section{Setting and creating a style}

Run \texttt{/\allowbreak{}config} and select \textbf{Output style}; the selection is saved as \texttt{output\allowbreak{}Style} in \texttt{.\allowbreak{}claude/\allowbreak{}settings.\allowbreak{}local.\allowbreak{}json}. The standalone \texttt{/\allowbreak{}output-\allowbreak{}style} command was deprecated in 2.1.73 and removed in 2.1.91 [79]. Because the style is part of the system prompt, which is read once at session start, a change takes effect after \texttt{/\allowbreak{}clear} or in a new session — and it invalidates the prompt cache [60], [79].


A custom style is a Markdown file in \texttt{\textasciitilde{}/\allowbreak{}.\allowbreak{}claude/\allowbreak{}output-\allowbreak{}styles}, \texttt{.\allowbreak{}claude/\allowbreak{}output-\allowbreak{}styles}, or the managed settings directory:


\begin{codeblock}{7.0}{8.3}
\cl{---}
\cl{name:~Evidence~analyst}
\cl{description:~Research~and~analysis~register~with~explicit~provenance}
\cl{keep-coding-instructions:~false}
\cl{---}
\cl{}
\cl{You~are~a~research~analyst~producing~decision-grade~written~work.~You~are~not~writing~software.}
\cl{}
\cl{-~Lead~with~the~answer~to~the~decision~question,~then~the~evidence.}
\cl{-~Distinguish~observed~fact,~source~claim,~and~your~own~inference,~explicitly,~every~time.}
\cl{-~Never~assert~a~figure,~date~or~attribution~without~naming~where~it~came~from.}
\cl{-~State~what~you~could~not~verify~as~a~numbered~list~of~open~items.~Do~not~smooth~over~a~gap.}
\cl{-~Use~British~English~and~the~register~of~a~professional~briefing~note.}
\end{codeblock}

Set \texttt{keep-\allowbreak{}coding-\allowbreak{}instruction\allowbreak{}s:\allowbreak{} true} when you are changing how Claude communicates but still want it engineering the same way; leave it out when the work is not software at all [79].


\FloatBarrier
\setcounter{section}{2}
\section{Where a style does and does not apply}

An output style applies to the \textbf{main conversation only}. A subagent runs its own system prompt and is unaffected — with the single exception of a fork, which inherits the parent's full system prompt [25], [79]. If you rely on a style for tone or discipline in delegated work, you must put the same instruction in the agent definition.


Compare the mechanisms: output styles modify the system prompt; \texttt{CLAUDE.\allowbreak{}md} adds a user message after it; \texttt{-\allowbreak{}-\allowbreak{}append-\allowbreak{}system-\allowbreak{}prompt} appends for one invocation; agents run a separate prompt entirely; skills load task-specific instructions on demand [79].


\begingroup
\def\kprows{%
\item Output styles modify the system prompt and can remove the built-in software-engineering instructions entirely.
\item The Proactive style is stronger autonomy guidance than auto mode, and it does not change your permission mode.
\item A style takes effect after \texttt{/\allowbreak{}clear} or in a new session, and invalidates the prompt cache.
\item A style applies to the main conversation only — a subagent runs its own system prompt, a fork inherits yours.
}%
\def\kpbody{\begin{minipage}{\textwidth}\subsection*{Key points}\begin{itemize}\kprows\end{itemize}\end{minipage}}%
\begingroup
\def\sloppy{\tolerance 9999\emergencystretch 3em\hfuzz 200pt\vfuzz 200pt}%
\hbadness=10000\vbadness=10000\hfuzz=200pt\vfuzz=200pt
\global\setbox\kpbox=\hbox{\kpbody}%
\endgroup
\par\addvspace{4.2mm}
\ifdim\dimexpr\ht\kpbox+\dp\kpbox\relax>0.30\textheight
  \typeout{HANDBOOK-KEYPOINTS broken \the\dimexpr\ht\kpbox+\dp\kpbox\relax}%
  \subsection*{Key points}
  \begin{itemize}\kprows\end{itemize}
\else
  \typeout{HANDBOOK-KEYPOINTS atomic \the\dimexpr\ht\kpbox+\dp\kpbox\relax}%
  \noindent\kpbody
\fi
\par\addvspace{1.4mm}
\endgroup

\breakrule

\FloatBarrier
\parttitle{Part V}{Browser use, autonomy and orchestration}
\addcontentsline{toc}{part}{Part V \textemdash\ Browser use, autonomy and orchestration}

\FloatBarrier
\renewcommand{\chaptertitlelabel}{Part V \textperiodcentered\ Chapter 20}
\setcounter{chapter}{19}
\chapter{Use Claude in Chrome safely}

Claude in Chrome lets Claude inspect and interact with browser tabs. It reached general availability in week 27 of 2026 [36]. Claude opens new tabs for browser tasks and \textbf{shares your browser's login state}, so it can reach any site you are already signed into [21]. That is the feature and the risk in one sentence.


\FloatBarrier
\setcounter{section}{0}
\section{Prerequisites and connection}

You need Chrome, Edge, or another Chromium browser (Brave, Arc, Vivaldi, Opera are detected); the Claude in Chrome extension version 1.0.36 or later; and a direct Anthropic plan — Pro, Max, Team or Enterprise [21]. Two constraints catch people out: \textbf{WSL is not supported}, and \textbf{API-key or long-lived-token sessions cannot use Chrome integration at all}, even with \texttt{-\allowbreak{}-\allowbreak{}chrome}, because the extension cannot authenticate with those credentials [21]. Sign in with \texttt{/\allowbreak{}login}.


\begin{codeblock}{9.0}{10.6}
\cl{claude~--chrome}
\end{codeblock}

\begin{codeblock}{9.0}{10.6}
\cl{/chrome}
\end{codeblock}

\texttt{/\allowbreak{}chrome} checks connection status, manages permissions, reconnects the extension, and selects among connected browsers. The integration is working when the panel shows \emph{Status: Enabled} and \emph{Extension: Installed}.


\textbf{PRACTITIONER NOTE:} enabling Chrome by default loads browser tools into every session and increases context usage [21]. Prefer \texttt{-\allowbreak{}-\allowbreak{}chrome} when you need it.


\FloatBarrier
\setcounter{section}{1}
\section{Use a dedicated browser profile}

Create a profile for agent work with only the accounts the task requires; no personal password manager unlocked; no unrelated tabs; restricted site permissions; test rather than production accounts where possible; and downloads directed to a controlled folder. An incognito window is not account isolation if the extension and login state are still available.


Site-level permissions are inherited from the Chrome extension and managed in its settings [21].


\FloatBarrier
\setcounter{section}{2}
\section{Escalate capability in stages}

\begin{enumerate}
\item \textbf{Observe} — "Describe this page; take no action."
\item \textbf{Navigate} — "Open the settings page; change nothing."
\item \textbf{Draft} — "Fill the form but do not submit."
\item \textbf{Verify} — "Show the entered values and the validation state."
\item \textbf{Act} — a human approves the final submission.
\end{enumerate}

The staging is enforced for you in Plan mode: read-only browser calls (\texttt{read\_\allowbreak{}page}, \texttt{get\_\allowbreak{}page\_\allowbreak{}text}, \texttt{find}, reading console messages or network requests, taking a screenshot) run without prompting, while clicks, typing, navigation, tab management and GIF recording prompt for approval [21]. An otherwise read-only call also prompts when it sets a state-changing flag, such as \texttt{save\_\allowbreak{}to\_\allowbreak{}disk} on a screenshot, and a \texttt{browser\_\allowbreak{}batch} runs without a prompt only when every action inside it is read-only.


\FloatBarrier
\setcounter{section}{3}
\section{What the browser can carry out}

File uploads are subject to three restrictions worth stating: Claude can upload a file only when the session is permitted to read it, so a \texttt{Read} deny rule also blocks uploading it; a single upload is capped at 10 MB in total; and files with multiple hard links are refused, which is common inside \texttt{node\_\allowbreak{}modules} [21].


\begin{calloutbox}{palebrass}{brassdark}{2.0mm}
\textbf{CAUTION:} screenshots, GIF recordings, downloads and page text may contain sensitive information, and a recording captures everything visible on a logged-in page. These artefacts enter the session transcript or the filesystem [21]. Review before sharing.
\end{calloutbox}

\FloatBarrier
\setcounter{section}{4}
\section{Prompt injection and authenticated action}

A web page is untrusted input. Text on a page can instruct an agent to abandon your task, disclose data, or navigate elsewhere; this is the canonical indirect prompt injection channel [4], [22]. Keep the objective and the allowed domains explicit, restrict network access, and require confirmation before any external write.


\begin{calloutbox}{palebrass}{brassdark}{2.0mm}
\textbf{CAUTION:} do not use browser automation for financial transfers, irreversible account operations, legal acceptance, or safety-critical control without a purpose-built system and an independent human confirmation path.
\end{calloutbox}

\FloatBarrier
\setcounter{section}{5}
\section{Computer use}

Where a task cannot be done in a browser, computer use lets Claude open applications, click, type and see the screen on macOS from the CLI [80]. The same staging discipline applies, with a wider blast radius: a browser profile can be isolated, a desktop session generally cannot.


\subsection*{Exercise}

Use a harmless public page or a local HTML fixture. Ask Claude to extract a small table, cite the page location, and take no action. Then ask it to navigate to one linked page. Inspect the browser history and the transcript to confirm the exact scope of what happened.


\emph{A worked solution is given in Appendix M.}


\begingroup
\def\kprows{%
\item Chrome integration shares your browser's authenticated state, so it can reach anything you are signed into.
\item Escalate in stages: observe, navigate, draft, verify, act — and let a human perform the final submission.
\item Plan mode already enforces the staging: read-only browser calls run, state-changing ones prompt.
\item A web page is untrusted input. Keep the objective and the allowed domains explicit.
}%
\def\kpbody{\begin{minipage}{\textwidth}\subsection*{Key points}\begin{itemize}\kprows\end{itemize}\end{minipage}}%
\begingroup
\def\sloppy{\tolerance 9999\emergencystretch 3em\hfuzz 200pt\vfuzz 200pt}%
\hbadness=10000\vbadness=10000\hfuzz=200pt\vfuzz=200pt
\global\setbox\kpbox=\hbox{\kpbody}%
\endgroup
\par\addvspace{4.2mm}
\ifdim\dimexpr\ht\kpbox+\dp\kpbox\relax>0.30\textheight
  \typeout{HANDBOOK-KEYPOINTS broken \the\dimexpr\ht\kpbox+\dp\kpbox\relax}%
  \subsection*{Key points}
  \begin{itemize}\kprows\end{itemize}
\else
  \typeout{HANDBOOK-KEYPOINTS atomic \the\dimexpr\ht\kpbox+\dp\kpbox\relax}%
  \noindent\kpbody
\fi
\par\addvspace{1.4mm}
\endgroup

\FloatBarrier
\renewcommand{\chaptertitlelabel}{Part V \textperiodcentered\ Chapter 21}
\setcounter{chapter}{20}
\chapter{Keep work moving with goals}

\texttt{/\allowbreak{}goal} sets a completion condition. After each turn a small fast model — Haiku by default on the Anthropic API — checks whether the condition holds; if not, Claude starts another turn instead of returning control to you [74]. One goal is active per session, and the condition may be up to 4,000 characters.


\begin{codeblock}{7.0}{8.3}
\cl{/goal~<condition>~~~~~~\#~set;~starts~a~turn~immediately}
\cl{/goal~~~~~~~~~~~~~~~~~~\#~status:~condition,~duration,~turns~evaluated,~token~spend,~last~reason}
\cl{/goal~clear~~~~~~~~~~~~\#~remove~(aliases:~stop,~off,~reset,~none,~cancel)}
\end{codeblock}

\FloatBarrier
\setcounter{section}{0}
\section{What the evaluator can and cannot see}

\textbf{The evaluator does not run commands or read files.} It judges only what Claude has surfaced in the conversation [74]. This single fact determines what a good condition looks like: "all tests in \texttt{test/\allowbreak{}auth} pass" works because Claude runs the tests and the output lands in the transcript where the evaluator can read it. "Finish the project" does not, because it is satisfiable by an assertion.


The evaluator returns one of three verdicts: \emph{not yet met} (Claude continues, taking the reason as guidance), \emph{met} (goal cleared, achieved entry recorded), or \emph{impossible} (goal cleared, failed entry recorded with the reason). Press \texttt{Ctrl+O} to see the reason behind a verdict [74].


\texttt{/\allowbreak{}goal} is implemented as a session-scoped \textbf{prompt-based Stop hook}, which has two consequences: it is unavailable when \texttt{disable\allowbreak{}All\allowbreak{}Hooks} is true after settings precedence, or when \texttt{allow\allowbreak{}Managed\allowbreak{}Hooks\allowbreak{}Only} is set in managed settings; and it is subject to the same workspace-trust rule as hooks in settings files [73], [74].


\FloatBarrier
\setcounter{section}{1}
\section{Anatomy of a safe goal}

Include a concrete deliverable; the checks required; the evidence to report; explicit exclusions; human gates; a turn, time or resource bound; and a condition for stopping when blocked.


\begin{codeblock}{7.0}{8.3}
\cl{/goal~Create~drafts/brief-v2.md~from~the~approved~sources.~The~goal~is~met~only~when~every}
\cl{material~claim~has~a~source-ledger~row,~the~evidence-reviewer~subagent~reports~no~high-severity}
\cl{unsupported~claim,~the~local~link~check~exits~0~with~its~output~shown,~and~the~final~message}
\cl{lists~the~commands~run~and~the~files~changed~as~evidence.~Do~not~publish,~push,~install}
\cl{software,~or~access~accounts.~Stop~and~ask~if~blocked,~or~after~20~turns.}
\end{codeblock}

\FloatBarrier
\setcounter{section}{2}
\section{Behaviour you should expect}

\begin{itemize}
\item \textbf{A goal does not change your permission mode} [74]. In Manual mode Claude still asks before tool calls your settings do not allow. To let goal turns run unattended, combine it with auto mode — deliberately, and with sandboxing.
\item \textbf{Background work defers evaluation.} If a subagent or background shell command is still running when a turn ends, evaluation is skipped until a turn finishes with nothing running [74].
\item \textbf{Check-ins back off.} Once background work has kept the goal waiting 30 minutes, Claude Code delivers a check-in asking Claude to read the running tasks' output and fix or stop anything stuck; later check-ins double the interval up to four times the first. Set \texttt{CLAUDE\_\allowbreak{}CODE\_\allowbreak{}GOAL\_\allowbreak{}CHECKIN\_\allowbreak{}MINUTES} to change it, or \texttt{0} to turn check-ins off (2.1.234+; idle check-ins 2.1.236+) [74].
\item \textbf{A no-progress guard exists.} If Claude answers the evaluator repeatedly without tool use, Claude Code stops the loop, warns, and returns control with the goal still set [74].
\item \textbf{Four failures clear the goal}: an authentication failure where Claude Code manages its own credentials, an exhausted credit balance, a context overflow auto-compaction could not clear, and an unavailable model. Transient errors such as rate limits leave the goal active [74].
\item \textbf{Goals survive resume}, with turn count, timer and token baseline reset [74].
\end{itemize}

\FloatBarrier
\setcounter{section}{3}
\section{Context is the real multiplier}

The course's long-running agent example works not because \texttt{/\allowbreak{}goal} is magical but because it combines a detailed proposal, tool documentation, learned procedures, a plan, tests and a measurable finished state [42]. A goal cannot repair missing inputs, and it cannot make an unsafe access safe. Clear a stale goal when the objective changes; \texttt{/\allowbreak{}clear} also removes it.


\begingroup
\def\kprows{%
\item The evaluator does not run commands or read files; it judges only what Claude surfaced in the conversation.
\item A good condition names a deliverable, the checks, the evidence to report, the exclusions and a bound.
\item A goal does not change your permission mode, and background work defers evaluation.
\item \texttt{/\allowbreak{}goal} is a prompt Stop hook, so \texttt{disable\allowbreak{}All\allowbreak{}Hooks} or \texttt{allow\allowbreak{}Managed\allowbreak{}Hooks\allowbreak{}Only} disables it.
}%
\def\kpbody{\begin{minipage}{\textwidth}\subsection*{Key points}\begin{itemize}\kprows\end{itemize}\end{minipage}}%
\begingroup
\def\sloppy{\tolerance 9999\emergencystretch 3em\hfuzz 200pt\vfuzz 200pt}%
\hbadness=10000\vbadness=10000\hfuzz=200pt\vfuzz=200pt
\global\setbox\kpbox=\hbox{\kpbody}%
\endgroup
\par\addvspace{4.2mm}
\ifdim\dimexpr\ht\kpbox+\dp\kpbox\relax>0.30\textheight
  \typeout{HANDBOOK-KEYPOINTS broken \the\dimexpr\ht\kpbox+\dp\kpbox\relax}%
  \subsection*{Key points}
  \begin{itemize}\kprows\end{itemize}
\else
  \typeout{HANDBOOK-KEYPOINTS atomic \the\dimexpr\ht\kpbox+\dp\kpbox\relax}%
  \noindent\kpbody
\fi
\par\addvspace{1.4mm}
\endgroup

\FloatBarrier
\renewcommand{\chaptertitlelabel}{Part V \textperiodcentered\ Chapter 22}
\setcounter{chapter}{21}
\chapter{Dynamic workflows and ultracode}

A dynamic workflow is a JavaScript script that orchestrates subagents at scale. Claude writes the script for the task you describe, and a runtime executes it in the background while your session stays responsive [19]. Requires 2.1.154 or later; on Pro, enable it from the Dynamic workflows row in \texttt{/\allowbreak{}config}.


The distinction that matters is \textbf{who holds the plan}. With subagents, skills and agent teams, Claude decides turn by turn what to spawn next and every result lands in a context window. A workflow script holds the loop, the branching and the intermediate results, so Claude's context holds only the final answer [19].


\FloatBarrier
\setcounter{section}{0}
\section{Entry points}

\begin{codeblock}{8.0}{9.4}
\cl{/deep-research~<question>~~~~~\#~bundled~workflow:~fan-out~search,~cross-check,~cited~report}
\cl{/workflows~~~~~~~~~~~~~~~~~~~~\#~progress~view}
\cl{/effort~ultracode~~~~~~~~~~~~~\#~xhigh~effort~plus~automatic~workflow~orchestration}
\end{codeblock}

\begin{codeblock}{9.0}{10.6}
\cl{claude~--effort~ultracode~~~~~\#~requires~2.1.203~or~later}
\end{codeblock}

\texttt{ultracode} is accepted by \texttt{claude -\allowbreak{}-\allowbreak{}effort} and by \texttt{/\allowbreak{}effort} in the terminal, and is absent from the VS Code extension's \texttt{/\allowbreak{}effort} usage string; \texttt{claude -\allowbreak{}-\allowbreak{}help} does not list it either, though the flag accepts it. Table 3 records what each surface did.


Including the keyword \texttt{ultracode} in a prompt you type yourself opts a single task in without changing session effort; asking in your own words works too [19]. The keyword is an opt-in \textbf{only in a prompt you type}: it does not start a workflow from \texttt{-\allowbreak{}p}, from an SDK message not stamped as human input, from a scheduled task prompt, or from a webhook or pull-request comment relayed into the conversation (from 2.1.210). Before 2.1.160 the literal trigger word was \texttt{workflow}.


\FloatBarrier
\setcounter{section}{1}
\section{Limits and permissions}

\begingroup
\def\tblrows{%
Concurrent agents & Up to 16, fewer when Claude Code has fewer CPUs available \\
Total agents per run & 1,000 \\
Mid-run user input & Not possible; only agent permission prompts pause a run \\
Filesystem or shell from the script & None; agents do the work, the script coordinates \\
Module loading & A script containing \texttt{import()} fails before the run starts \\
Prompt-cache stagger & Matching agents start up to \texttt{CLAUDE\_\allowbreak{}CODE\_\allowbreak{}WORKFLOW\_\allowbreak{}PREFIX\_\allowbreak{}STAGGER\_\allowbreak{}MS} ms after the first, default 5000 \\
}%
\def\tblbody{\begin{minipage}{\textwidth}\boxcaption{Table 16 — Dynamic workflow runtime constraints}
{\footnotesize\begin{tabular}{L{41.9mm}L{101.1mm}}
\toprule
\textbf{Constraint} & \textbf{Value} \\
\midrule
\tblrows
\bottomrule\end{tabular}}\end{minipage}}%
\begingroup
\def\sloppy{\tolerance 9999\emergencystretch 3em\hfuzz 200pt\vfuzz 200pt}%
\hbadness=10000\vbadness=10000\hfuzz=200pt\vfuzz=200pt
\global\setbox\tblbox=\hbox{\tblbody}%
\endgroup
\par\addvspace{2.6mm}
\ifdim\dimexpr\ht\tblbox+\dp\tblbox\relax>0.55\textheight
  \typeout{HANDBOOK-TABLE broken \the\dimexpr\ht\tblbox+\dp\tblbox\relax}%
  \tabcaption{Table 16 — Dynamic workflow runtime constraints}
  {\footnotesize\begin{longtable}{L{41.9mm}L{101.1mm}}
  \toprule
\textbf{Constraint} & \textbf{Value} \\
\midrule\endfirsthead
  \multicolumn{2}{@{}l@{}}{%
  \sffamily\footnotesize\itshape\color{inkgrey}Table 16 — Dynamic workflow runtime constraints \textemdash\ continued}\\[1.2mm]
  \toprule
\textbf{Constraint} & \textbf{Value} \\
\midrule\endhead
  \bottomrule\endfoot
  \bottomrule\endlastfoot
  \tblrows
  \end{longtable}}%
\else
  \typeout{HANDBOOK-TABLE atomic \the\dimexpr\ht\tblbox+\dp\tblbox\relax}%
  \noindent\tblbody
\fi
\par\addvspace{2.6mm}
\endgroup

\begin{calloutbox}{palebrass}{brassdark}{2.0mm}
\textbf{CAUTION:} The subagents a workflow spawns always run in \texttt{accept\allowbreak{}Edits} and inherit your tool allowlist, whatever the session's mode [19]. Your mode governs only the launch prompt. File edits are auto-approved for the duration of the run, so commit a clean baseline before starting one.
\end{calloutbox}

Size guidance is separate from the hard caps. \texttt{workflow\allowbreak{}Size\allowbreak{}Guideline} takes \texttt{unrestricte\allowbreak{}d}, \texttt{small} (fewer than 5 agents), \texttt{medium} (fewer than 15, the default), or \texttt{large} (fewer than 50); set it in \texttt{/\allowbreak{}config} or in any settings file from 2.1.219 [19]. Claude Code shows a \texttt{Large workf\allowbreak{}low} warning when a run schedules more than 25 agents or projects more than 1.5 million tokens — advisory only, and suppressed when ultracode is on. \textbf{The 25 is the default guideline's number, not a constant:} choose a guideline yourself and its own agent count replaces the threshold, so \texttt{small} warns above five and \texttt{large} above fifty [19].


\FloatBarrier
\setcounter{section}{2}
\section{Resume semantics, and why they matter for cost}

Resume works only within the same session. Replay follows the order agents \emph{started}: cached results stop at the first agent that did not finish, and every agent that started after it runs again even if it completed [19]. Stopping mid fan-out is therefore expensive, and a workflow that spreads work across many small agents preserves more progress than one with a few long ones. Exiting Claude Code while a workflow runs means the next session starts it fresh.


\FloatBarrier
\setcounter{section}{3}
\section{When a workflow earns its cost}

Use one when the task has genuinely independent lanes: several source families researched in parallel; multiple reviewers testing the same proposal independently; a large change divisible by non-overlapping ownership; research, challenge and synthesis benefiting from separate contexts.


Do not use one to format a file, answer a simple question, or manufacture the appearance of corroboration by asking identical agents the same question from the same evidence.


\FloatBarrier
\setcounter{section}{4}
\section{A research workflow pattern}

The pattern separates gathering from checking and never lets one context do both. Figure 6 shows it end to end, including the release gate that a person, not an agent, has to pass.


\begin{figure}[tbp]
\begin{codefig}{9.0}{10.6}
\cl{~~~~~~~~~~~~~~~~~~~~~\pmboxdrawuni{250C}\pmboxdrawuni{2500}\pmboxdrawuni{2500}\pmboxdrawuni{2500}\pmboxdrawuni{2500}\pmboxdrawuni{2500}\pmboxdrawuni{2500}\pmboxdrawuni{2500}\pmboxdrawuni{2500}\pmboxdrawuni{2500}\pmboxdrawuni{2500}\pmboxdrawuni{2500}\pmboxdrawuni{2500}\pmboxdrawuni{2500}\pmboxdrawuni{2500}\pmboxdrawuni{2510}}
\cl{~~~~~~~~~~~~~~~~~~~~~\pmboxdrawuni{2502}~~Scope~agent~\pmboxdrawuni{2502}~~~defines~the~question,~the~exclusions,}
\cl{~~~~~~~~~~~~~~~~~~~~~\pmboxdrawuni{2514}\pmboxdrawuni{2500}\pmboxdrawuni{2500}\pmboxdrawuni{2500}\pmboxdrawuni{2500}\pmboxdrawuni{2500}\pmboxdrawuni{2500}\pmboxdrawuni{252C}\pmboxdrawuni{2500}\pmboxdrawuni{2500}\pmboxdrawuni{2500}\pmboxdrawuni{2500}\pmboxdrawuni{2500}\pmboxdrawuni{2500}\pmboxdrawuni{2500}\pmboxdrawuni{2518}~~~and~the~source~families}
\cl{~~~~~~~~~~~~\pmboxdrawuni{250C}\pmboxdrawuni{2500}\pmboxdrawuni{2500}\pmboxdrawuni{2500}\pmboxdrawuni{2500}\pmboxdrawuni{2500}\pmboxdrawuni{2500}\pmboxdrawuni{2500}\pmboxdrawuni{2500}\pmboxdrawuni{2500}\pmboxdrawuni{2500}\pmboxdrawuni{2500}\pmboxdrawuni{2500}\pmboxdrawuni{2500}\pmboxdrawuni{2500}\pmboxdrawuni{2500}\pmboxdrawuni{253C}\pmboxdrawuni{2500}\pmboxdrawuni{2500}\pmboxdrawuni{2500}\pmboxdrawuni{2500}\pmboxdrawuni{2500}\pmboxdrawuni{2500}\pmboxdrawuni{2500}\pmboxdrawuni{2500}\pmboxdrawuni{2500}\pmboxdrawuni{2500}\pmboxdrawuni{2500}\pmboxdrawuni{2500}\pmboxdrawuni{2500}\pmboxdrawuni{2500}\pmboxdrawuni{2500}\pmboxdrawuni{2510}}
\cl{~~~~~~~~~~~~▼~~~~~~~~~~~~~~~▼~~~~~~~~~~~~~~~▼}
\cl{~~~~~~\pmboxdrawuni{250C}\pmboxdrawuni{2500}\pmboxdrawuni{2500}\pmboxdrawuni{2500}\pmboxdrawuni{2500}\pmboxdrawuni{2500}\pmboxdrawuni{2500}\pmboxdrawuni{2500}\pmboxdrawuni{2500}\pmboxdrawuni{2500}\pmboxdrawuni{2500}\pmboxdrawuni{2500}\pmboxdrawuni{2510}~~~\pmboxdrawuni{250C}\pmboxdrawuni{2500}\pmboxdrawuni{2500}\pmboxdrawuni{2500}\pmboxdrawuni{2500}\pmboxdrawuni{2500}\pmboxdrawuni{2500}\pmboxdrawuni{2500}\pmboxdrawuni{2500}\pmboxdrawuni{2500}\pmboxdrawuni{2500}\pmboxdrawuni{2500}\pmboxdrawuni{2510}~~~\pmboxdrawuni{250C}\pmboxdrawuni{2500}\pmboxdrawuni{2500}\pmboxdrawuni{2500}\pmboxdrawuni{2500}\pmboxdrawuni{2500}\pmboxdrawuni{2500}\pmboxdrawuni{2500}\pmboxdrawuni{2500}\pmboxdrawuni{2500}\pmboxdrawuni{2500}\pmboxdrawuni{2500}\pmboxdrawuni{2510}}
\cl{~~~~~~\pmboxdrawuni{2502}~~Lane~A~~~\pmboxdrawuni{2502}~~~\pmboxdrawuni{2502}~~Lane~B~~~\pmboxdrawuni{2502}~~~\pmboxdrawuni{2502}~~Lane~C~~~\pmboxdrawuni{2502}~~~independent~source~families}
\cl{~~~~~~\pmboxdrawuni{2514}\pmboxdrawuni{2500}\pmboxdrawuni{2500}\pmboxdrawuni{2500}\pmboxdrawuni{2500}\pmboxdrawuni{2500}\pmboxdrawuni{252C}\pmboxdrawuni{2500}\pmboxdrawuni{2500}\pmboxdrawuni{2500}\pmboxdrawuni{2500}\pmboxdrawuni{2500}\pmboxdrawuni{2518}~~~\pmboxdrawuni{2514}\pmboxdrawuni{2500}\pmboxdrawuni{2500}\pmboxdrawuni{2500}\pmboxdrawuni{2500}\pmboxdrawuni{2500}\pmboxdrawuni{252C}\pmboxdrawuni{2500}\pmboxdrawuni{2500}\pmboxdrawuni{2500}\pmboxdrawuni{2500}\pmboxdrawuni{2500}\pmboxdrawuni{2518}~~~\pmboxdrawuni{2514}\pmboxdrawuni{2500}\pmboxdrawuni{2500}\pmboxdrawuni{2500}\pmboxdrawuni{2500}\pmboxdrawuni{2500}\pmboxdrawuni{252C}\pmboxdrawuni{2500}\pmboxdrawuni{2500}\pmboxdrawuni{2500}\pmboxdrawuni{2500}\pmboxdrawuni{2500}\pmboxdrawuni{2518}}
\cl{~~~~~~~~~~~~\pmboxdrawuni{2514}\pmboxdrawuni{2500}\pmboxdrawuni{2500}\pmboxdrawuni{2500}\pmboxdrawuni{2500}\pmboxdrawuni{2500}\pmboxdrawuni{2500}\pmboxdrawuni{2500}\pmboxdrawuni{2500}\pmboxdrawuni{2500}\pmboxdrawuni{2500}\pmboxdrawuni{2500}\pmboxdrawuni{2500}\pmboxdrawuni{2500}\pmboxdrawuni{2500}\pmboxdrawuni{2500}\pmboxdrawuni{253C}\pmboxdrawuni{2500}\pmboxdrawuni{2500}\pmboxdrawuni{2500}\pmboxdrawuni{2500}\pmboxdrawuni{2500}\pmboxdrawuni{2500}\pmboxdrawuni{2500}\pmboxdrawuni{2500}\pmboxdrawuni{2500}\pmboxdrawuni{2500}\pmboxdrawuni{2500}\pmboxdrawuni{2500}\pmboxdrawuni{2500}\pmboxdrawuni{2500}\pmboxdrawuni{2500}\pmboxdrawuni{2518}}
\cl{~~~~~~~~~~~~~~~~~~~~~~~~~~~~▼}
\cl{~~~~~~~~~~~~~~\pmboxdrawuni{250C}\pmboxdrawuni{2500}\pmboxdrawuni{2500}\pmboxdrawuni{2500}\pmboxdrawuni{2500}\pmboxdrawuni{2500}\pmboxdrawuni{2500}\pmboxdrawuni{2500}\pmboxdrawuni{2500}\pmboxdrawuni{2500}\pmboxdrawuni{2500}\pmboxdrawuni{2500}\pmboxdrawuni{2500}\pmboxdrawuni{2500}\pmboxdrawuni{2500}\pmboxdrawuni{2500}\pmboxdrawuni{2500}\pmboxdrawuni{2500}\pmboxdrawuni{2500}\pmboxdrawuni{2500}\pmboxdrawuni{2500}\pmboxdrawuni{2500}\pmboxdrawuni{2500}\pmboxdrawuni{2500}\pmboxdrawuni{2500}\pmboxdrawuni{2500}\pmboxdrawuni{2500}\pmboxdrawuni{2500}\pmboxdrawuni{2510}}
\cl{~~~~~~~~~~~~~~\pmboxdrawuni{2502}~Contradiction~\&~provenance~\pmboxdrawuni{2502}~~distinguishes~independent~sources}
\cl{~~~~~~~~~~~~~~\pmboxdrawuni{2502}~~~~~~~~~reviewer~~~~~~~~~~~\pmboxdrawuni{2502}~~from~repetitions~of~one~original}
\cl{~~~~~~~~~~~~~~\pmboxdrawuni{2514}\pmboxdrawuni{2500}\pmboxdrawuni{2500}\pmboxdrawuni{2500}\pmboxdrawuni{2500}\pmboxdrawuni{2500}\pmboxdrawuni{2500}\pmboxdrawuni{2500}\pmboxdrawuni{2500}\pmboxdrawuni{2500}\pmboxdrawuni{2500}\pmboxdrawuni{2500}\pmboxdrawuni{2500}\pmboxdrawuni{2500}\pmboxdrawuni{252C}\pmboxdrawuni{2500}\pmboxdrawuni{2500}\pmboxdrawuni{2500}\pmboxdrawuni{2500}\pmboxdrawuni{2500}\pmboxdrawuni{2500}\pmboxdrawuni{2500}\pmboxdrawuni{2500}\pmboxdrawuni{2500}\pmboxdrawuni{2500}\pmboxdrawuni{2500}\pmboxdrawuni{2500}\pmboxdrawuni{2500}\pmboxdrawuni{2518}}
\cl{~~~~~~~~~~~~~~~~~~~~~~~~~~~~▼}
\cl{~~~~~~~~~~~~~~~~~~~\pmboxdrawuni{250C}\pmboxdrawuni{2500}\pmboxdrawuni{2500}\pmboxdrawuni{2500}\pmboxdrawuni{2500}\pmboxdrawuni{2500}\pmboxdrawuni{2500}\pmboxdrawuni{2500}\pmboxdrawuni{2500}\pmboxdrawuni{2500}\pmboxdrawuni{2500}\pmboxdrawuni{2500}\pmboxdrawuni{2500}\pmboxdrawuni{2500}\pmboxdrawuni{2500}\pmboxdrawuni{2500}\pmboxdrawuni{2500}\pmboxdrawuni{2500}\pmboxdrawuni{2510}}
\cl{~~~~~~~~~~~~~~~~~~~\pmboxdrawuni{2502}~Synthesis~agent~\pmboxdrawuni{2502}}
\cl{~~~~~~~~~~~~~~~~~~~\pmboxdrawuni{2514}\pmboxdrawuni{2500}\pmboxdrawuni{2500}\pmboxdrawuni{2500}\pmboxdrawuni{2500}\pmboxdrawuni{2500}\pmboxdrawuni{2500}\pmboxdrawuni{2500}\pmboxdrawuni{2500}\pmboxdrawuni{252C}\pmboxdrawuni{2500}\pmboxdrawuni{2500}\pmboxdrawuni{2500}\pmboxdrawuni{2500}\pmboxdrawuni{2500}\pmboxdrawuni{2500}\pmboxdrawuni{2500}\pmboxdrawuni{2500}\pmboxdrawuni{2518}}
\cl{~~~~~~~~~~~~~~~~~~~~~~~~~~~~▼}
\cl{~~~~~~~~~~~~\pmboxdrawuni{2554}\pmboxdrawuni{2550}\pmboxdrawuni{2550}\pmboxdrawuni{2550}\pmboxdrawuni{2550}\pmboxdrawuni{2550}\pmboxdrawuni{2550}\pmboxdrawuni{2550}\pmboxdrawuni{2550}\pmboxdrawuni{2550}\pmboxdrawuni{2550}\pmboxdrawuni{2550}\pmboxdrawuni{2550}\pmboxdrawuni{2550}\pmboxdrawuni{2550}\pmboxdrawuni{2550}\pmboxdrawuni{2550}\pmboxdrawuni{2550}\pmboxdrawuni{2550}\pmboxdrawuni{2550}\pmboxdrawuni{2550}\pmboxdrawuni{2550}\pmboxdrawuni{2550}\pmboxdrawuni{2550}\pmboxdrawuni{2550}\pmboxdrawuni{2550}\pmboxdrawuni{2550}\pmboxdrawuni{2550}\pmboxdrawuni{2550}\pmboxdrawuni{2550}\pmboxdrawuni{2550}\pmboxdrawuni{2550}\pmboxdrawuni{2557}}
\cl{~~~~~~~~~~~~\pmboxdrawuni{2551}~~HUMAN~REVIEW~AND~RELEASE~GATE~\pmboxdrawuni{2551}}
\cl{~~~~~~~~~~~~\pmboxdrawuni{255A}\pmboxdrawuni{2550}\pmboxdrawuni{2550}\pmboxdrawuni{2550}\pmboxdrawuni{2550}\pmboxdrawuni{2550}\pmboxdrawuni{2550}\pmboxdrawuni{2550}\pmboxdrawuni{2550}\pmboxdrawuni{2550}\pmboxdrawuni{2550}\pmboxdrawuni{2550}\pmboxdrawuni{2550}\pmboxdrawuni{2550}\pmboxdrawuni{2550}\pmboxdrawuni{2550}\pmboxdrawuni{2550}\pmboxdrawuni{2550}\pmboxdrawuni{2550}\pmboxdrawuni{2550}\pmboxdrawuni{2550}\pmboxdrawuni{2550}\pmboxdrawuni{2550}\pmboxdrawuni{2550}\pmboxdrawuni{2550}\pmboxdrawuni{2550}\pmboxdrawuni{2550}\pmboxdrawuni{2550}\pmboxdrawuni{2550}\pmboxdrawuni{2550}\pmboxdrawuni{2550}\pmboxdrawuni{2550}\pmboxdrawuni{255D}}
\end{codefig}
\figcaption{Figure 6 — A cross-checked research workflow with a human release gate}
\end{figure}

The contradiction reviewer is the load-bearing element. Three pages repeating one press release are one source, and a workflow that counts them as three has manufactured confidence rather than evidence. The bundled \texttt{/\allowbreak{}deep-\allowbreak{}research} workflow already votes on each claim and filters out those that do not survive cross-checking, reporting a claim as \emph{unverified} rather than refuted when a verifier could not check it [19].


\FloatBarrier
\setcounter{section}{5}
\section{Saving and distributing a workflow}

Run \texttt{/\allowbreak{}workflows}, select the run, press \texttt{s}, and choose \texttt{.\allowbreak{}claude/\allowbreak{}workflows/\allowbreak{}} (shared with the repository) or \texttt{\textasciitilde{}/\allowbreak{}.\allowbreak{}claude/\allowbreak{}workflows/\allowbreak{}} (yours alone). It then runs as \texttt{/\allowbreak{}<name>}. Claude Code refuses to write through a symlink at either location — before 2.1.216 it followed the link [19]. Distribute across teams by placing the script in a plugin's \texttt{workflows/\allowbreak{}} directory, where it is namespaced as \texttt{/\allowbreak{}plugin-\allowbreak{}name:\allowbreak{}workflow-\allowbreak{}name}.


Turn workflows off with the \texttt{/\allowbreak{}config} toggle, \texttt{"disable\allowbreak{}Workflows":\allowbreak{} true}, or \texttt{CLAUDE\_\allowbreak{}CODE\_\allowbreak{}DISABLE\_\allowbreak{}WORKFLOWS=\allowbreak{}1}; organisations use managed settings [19].


\subsection*{Exercise}

Run \texttt{/\allowbreak{}deep-\allowbreak{}research} on a non-sensitive topic with a \texttt{small} size guideline. Require direct sources, a contradiction pass and a source register. Compare the result with a single-agent attempt, and use \texttt{/\allowbreak{}usage} to decide whether the workflow added decision value or only cost.


\emph{A worked solution is given in Appendix M.}


\begingroup
\def\kprows{%
\item A workflow moves the orchestration into a script, so the conversation holds only the final answer.
\item Workflow subagents always run in \texttt{accept\allowbreak{}Edits} regardless of your session's mode. Commit a clean baseline first.
\item Up to 16 concurrent and 1,000 total agents per run; the size guideline is advice, the caps are not.
\item Use one only where the lanes are genuinely independent — identical agents on identical evidence manufacture confidence, not corroboration.
}%
\def\kpbody{\begin{minipage}{\textwidth}\subsection*{Key points}\begin{itemize}\kprows\end{itemize}\end{minipage}}%
\begingroup
\def\sloppy{\tolerance 9999\emergencystretch 3em\hfuzz 200pt\vfuzz 200pt}%
\hbadness=10000\vbadness=10000\hfuzz=200pt\vfuzz=200pt
\global\setbox\kpbox=\hbox{\kpbody}%
\endgroup
\par\addvspace{4.2mm}
\ifdim\dimexpr\ht\kpbox+\dp\kpbox\relax>0.30\textheight
  \typeout{HANDBOOK-KEYPOINTS broken \the\dimexpr\ht\kpbox+\dp\kpbox\relax}%
  \subsection*{Key points}
  \begin{itemize}\kprows\end{itemize}
\else
  \typeout{HANDBOOK-KEYPOINTS atomic \the\dimexpr\ht\kpbox+\dp\kpbox\relax}%
  \noindent\kpbody
\fi
\par\addvspace{1.4mm}
\endgroup

\FloatBarrier
\renewcommand{\chaptertitlelabel}{Part V \textperiodcentered\ Chapter 23}
\setcounter{chapter}{22}
\chapter{Session-scoped scheduling with \texttt{/\allowbreak{}loop}}

\texttt{/\allowbreak{}loop} is a bundled skill (alias \texttt{/\allowbreak{}proactive}) that re-runs a prompt while the session stays open [24], [35]. It is not cron. This chapter carries three of the six first-edition corrections, so read it even if you know the feature.


\FloatBarrier
\setcounter{section}{0}
\section{Three forms, three behaviours}

\begingroup
\def\tblrows{%
Interval and prompt & \texttt{/\allowbreak{}loop 5m che\allowbreak{}ck the depl\allowbreak{}oy} & Runs on a fixed cron schedule \\
Prompt only & \texttt{/\allowbreak{}loop check \allowbreak{}the deploy} & Runs at an interval Claude chooses each iteration, between one minute and one hour, printing the delay and the reason \\
Interval only, or nothing & \texttt{/\allowbreak{}loop} & Runs the \textbf{built-in maintenance prompt}, or your \texttt{loop.\allowbreak{}md} if one exists \\
}%
\def\tblbody{\begin{minipage}{\textwidth}\boxcaption{Table 17 — What \texttt{/\allowbreak{}loop} does with each combination of arguments}
{\footnotesize\begin{tabular}{L{26.4mm}L{23.6mm}L{89.9mm}}
\toprule
\textbf{You provide} & \textbf{Example} & \textbf{What happens} \\
\midrule
\tblrows
\bottomrule\end{tabular}}\end{minipage}}%
\begingroup
\def\sloppy{\tolerance 9999\emergencystretch 3em\hfuzz 200pt\vfuzz 200pt}%
\hbadness=10000\vbadness=10000\hfuzz=200pt\vfuzz=200pt
\global\setbox\tblbox=\hbox{\tblbody}%
\endgroup
\par\addvspace{2.6mm}
\ifdim\dimexpr\ht\tblbox+\dp\tblbox\relax>0.55\textheight
  \typeout{HANDBOOK-TABLE broken \the\dimexpr\ht\tblbox+\dp\tblbox\relax}%
  \tabcaption{Table 17 — What \texttt{/\allowbreak{}loop} does with each combination of arguments}
  {\footnotesize\begin{longtable}{L{26.4mm}L{23.6mm}L{89.9mm}}
  \toprule
\textbf{You provide} & \textbf{Example} & \textbf{What happens} \\
\midrule\endfirsthead
  \multicolumn{3}{@{}l@{}}{%
  \sffamily\footnotesize\itshape\color{inkgrey}Table 17 — What \texttt{/\allowbreak{}loop} does with each combination of arguments \textemdash\ continued}\\[1.2mm]
  \toprule
\textbf{You provide} & \textbf{Example} & \textbf{What happens} \\
\midrule\endhead
  \bottomrule\endfoot
  \bottomrule\endlastfoot
  \tblrows
  \end{longtable}}%
\else
  \typeout{HANDBOOK-TABLE atomic \the\dimexpr\ht\tblbox+\dp\tblbox\relax}%
  \noindent\tblbody
\fi
\par\addvspace{2.6mm}
\endgroup

\begin{calloutbox}{palebrass}{brassdark}{2.0mm}
\textbf{CAUTION — M01.} A bare \texttt{/\allowbreak{}loop} does \textbf{not} list your scheduled tasks. It starts an autonomous maintenance loop that continues unfinished work, tends the current branch's pull request — review comments, failed CI, merge conflicts — and runs cleanup passes such as bug hunts when nothing else is pending [24]. Claude does not start new initiatives outside that scope, and irreversible actions proceed only when they continue something the transcript already authorised, but this is still autonomous execution you did not ask for. To list tasks, ask Claude in natural language, or use \texttt{Cron\allowbreak{}List}.
\end{calloutbox}

Replace the built-in prompt with your own by creating \texttt{.\allowbreak{}claude/\allowbreak{}loop.\allowbreak{}md} (project, takes precedence) or \texttt{\textasciitilde{}/\allowbreak{}.\allowbreak{}claude/\allowbreak{}loop.\allowbreak{}md} (user). Edits take effect on the next iteration; content beyond 25,000 bytes is truncated [24].


\FloatBarrier
\setcounter{section}{1}
\section{Operational constraints}

\begin{itemize}
\item Claude Code must be \textbf{running and idle}. A scheduled prompt fires between turns, never mid-response.
\item Minimum interval one minute. Units are \texttt{s}, \texttt{m}, \texttt{h}, \texttt{d}; seconds round up to the nearest minute; intervals that do not map to a clean cron step, such as \texttt{7m} or \texttt{90m}, are rounded and Claude tells you what it chose.
\item \textbf{Jitter is specified, not random.} Recurring tasks fire up to 30 minutes after the scheduled time — or up to half the interval for sub-hourly tasks — and one-shots scheduled for \texttt{:\allowbreak{}00} or \texttt{:\allowbreak{}30} fire up to 90 seconds early. The offset derives from the task ID, so it is stable. \emph{Mitigation:} schedule at a minute that is neither \texttt{:\allowbreak{}00} nor \texttt{:\allowbreak{}30}, for example \texttt{3 9 * * *}.
\item \textbf{No catch-up.} A task whose time passes while Claude is busy fires once when Claude becomes idle, not once per missed interval.
\item \textbf{Seven-day expiry.} Recurring tasks expire seven days after creation, firing one final time.
\item \textbf{Fifty tasks} per session, each with an eight-character ID.
\item Times are interpreted in \textbf{your local timezone}.
\item Scheduling fails if the project \texttt{.\allowbreak{}claude} directory or the task file inside it is a symlink.
\end{itemize}

All of the above from [24].


\FloatBarrier
\setcounter{section}{2}
\section{Lifecycle, and what actually stops a loop}

This is E2. Session-scoped does not mean session-bound:


\begin{itemize}
\item \texttt{Esc} while a \texttt{/\allowbreak{}loop} is waiting clears the pending wakeup. Tasks you scheduled by asking Claude directly are unaffected and stay until deleted.
\item \textbf{Backgrounding the session carries \texttt{/\allowbreak{}loop} tasks into a background session that keeps running without a terminal.}
\item Starting a fresh conversation with \texttt{/\allowbreak{}clear} clears all session-scoped tasks.
\item \texttt{claude -\allowbreak{}-\allowbreak{}resume} or \texttt{-\allowbreak{}-\allowbreak{}continue} \textbf{restores} recurring tasks that have not expired and one-shot tasks whose time has not passed. Background Bash and monitor tasks are never restored.
\item In self-paced mode Claude can end the loop itself by calling \texttt{Schedule\allowbreak{}Wakeup} with \texttt{stop:\allowbreak{} true}. If an iteration ends without rescheduling or stopping, Claude Code schedules one fallback wakeup about 20 minutes later and ends the loop if that iteration does not reschedule either.
\end{itemize}

Underneath, Claude uses \texttt{Cron\allowbreak{}Create}, \texttt{Cron\allowbreak{}List} and \texttt{Cron\allowbreak{}Delete}, which accept standard five-field cron expressions; extended syntax such as \texttt{L}, \texttt{W}, \texttt{?} and name aliases is not supported [24]. \texttt{CLAUDE\_\allowbreak{}CODE\_\allowbreak{}DISABLE\_\allowbreak{}CRON=\allowbreak{}1} disables the scheduler entirely, and \texttt{/\allowbreak{}loop} with it.


\textbf{Provider differences.} On Amazon Bedrock, Claude Platform on AWS, Google Cloud's Agent Platform and Microsoft Foundry — and wherever feature-flag fetching is turned off — a prompt with no interval runs on a fixed ten-minute schedule instead of a Claude-chosen one, and a bare \texttt{/\allowbreak{}loop} prints the usage message instead of running the maintenance prompt [24].


\textbf{Skills on a scheduled fire.} Only skills Claude may invoke on its own actually execute. Built-in commands, skills marked \texttt{disable-\allowbreak{}model-\allowbreak{}invocation:\allowbreak{} true} — including the bundled \texttt{/\allowbreak{}verify} — skills withheld by \texttt{skill\allowbreak{}Overrides} or a \texttt{Skill} deny rule, and MCP prompts all arrive as plain text instead [24].


\FloatBarrier
\setcounter{section}{3}
\section{The Monitor alternative}

When you ask for a dynamic \texttt{/\allowbreak{}loop} schedule, Claude may use the \textbf{Monitor tool} instead. Monitor runs a background script and streams each output line back, avoiding polling altogether; it is usually more token-efficient and more responsive than re-running a prompt on an interval [24], [81]. If you find yourself polling a log, a build, or a queue, ask for a monitor rather than a loop.


\FloatBarrier
\setcounter{section}{4}
\section{Suitable and unsuitable uses}

Suitable: checking a long build; re-running a read-only local verification; reminding the active session to inspect a queue; reporting a non-critical condition during an attended work period.


Unsuitable: anything with a service-level commitment; security, medical or emergency monitoring; payments or trading; anything needing exact timing or catch-up; unattended destructive repair.


Always state what to do when the condition is missing, ambiguous or unsafe. "Report only" is a good default.


\begingroup
\def\kprows{%
\item A bare \texttt{/\allowbreak{}loop} starts an autonomous maintenance loop; it does not list your tasks.
\item Backgrounding carries loops into a background session, and resume restores unexpired ones.
\item Jitter is specified and deterministic — schedule at a minute that is neither \texttt{:\allowbreak{}00} nor \texttt{:\allowbreak{}30}.
\item For a stream to watch rather than a state to poll, ask for the Monitor tool instead.
}%
\def\kpbody{\begin{minipage}{\textwidth}\subsection*{Key points}\begin{itemize}\kprows\end{itemize}\end{minipage}}%
\begingroup
\def\sloppy{\tolerance 9999\emergencystretch 3em\hfuzz 200pt\vfuzz 200pt}%
\hbadness=10000\vbadness=10000\hfuzz=200pt\vfuzz=200pt
\global\setbox\kpbox=\hbox{\kpbody}%
\endgroup
\par\addvspace{4.2mm}
\ifdim\dimexpr\ht\kpbox+\dp\kpbox\relax>0.30\textheight
  \typeout{HANDBOOK-KEYPOINTS broken \the\dimexpr\ht\kpbox+\dp\kpbox\relax}%
  \subsection*{Key points}
  \begin{itemize}\kprows\end{itemize}
\else
  \typeout{HANDBOOK-KEYPOINTS atomic \the\dimexpr\ht\kpbox+\dp\kpbox\relax}%
  \noindent\kpbody
\fi
\par\addvspace{1.4mm}
\endgroup

\FloatBarrier
\renewcommand{\chaptertitlelabel}{Part V \textperiodcentered\ Chapter 24}
\setcounter{chapter}{23}
\chapter{Cloud routines and desktop scheduled tasks}

Three scheduling mechanisms exist and they are not interchangeable. Choose by what must be true when the work runs.


\begingroup
\def\tblrows{%
Runs on & Anthropic cloud (or your self-hosted environment) & Your machine & Your machine \\
Requires machine on & No & Yes & Yes \\
Requires open session & No & No & Yes \\
Persists across restarts & Yes & Yes & Restored on \texttt{-\allowbreak{}-\allowbreak{}resume} if unexpired \\
Access to local files & No — a fresh clone & Yes & Yes \\
MCP & Connectors configured per task & Config files and connectors & Inherits from the session \\
Permission prompts & \textbf{None — runs autonomously} & Configurable per task & Inherits from the session \\
Minimum interval & \textbf{1 hour} & 1 minute & 1 minute \\
}%
\def\tblbody{\begin{minipage}{\textwidth}\boxcaption{Table 18 — Comparing the three scheduling mechanisms}
{\footnotesize\begin{tabular}{L{35.8mm}L{40.3mm}L{28.0mm}L{32.5mm}}
\toprule
\textbf{} & \textbf{Cloud routine} & \textbf{Desktop scheduled task} & \textbf{\texttt{/\allowbreak{}loop}} \\
\midrule
\tblrows
\bottomrule\end{tabular}}\end{minipage}}%
\begingroup
\def\sloppy{\tolerance 9999\emergencystretch 3em\hfuzz 200pt\vfuzz 200pt}%
\hbadness=10000\vbadness=10000\hfuzz=200pt\vfuzz=200pt
\global\setbox\tblbox=\hbox{\tblbody}%
\endgroup
\par\addvspace{2.6mm}
\ifdim\dimexpr\ht\tblbox+\dp\tblbox\relax>0.55\textheight
  \typeout{HANDBOOK-TABLE broken \the\dimexpr\ht\tblbox+\dp\tblbox\relax}%
  \tabcaption{Table 18 — Comparing the three scheduling mechanisms}
  {\footnotesize\begin{longtable}{L{35.8mm}L{40.3mm}L{28.0mm}L{32.5mm}}
  \toprule
\textbf{} & \textbf{Cloud routine} & \textbf{Desktop scheduled task} & \textbf{\texttt{/\allowbreak{}loop}} \\
\midrule\endfirsthead
  \multicolumn{4}{@{}l@{}}{%
  \sffamily\footnotesize\itshape\color{inkgrey}Table 18 — Comparing the three scheduling mechanisms \textemdash\ continued}\\[1.2mm]
  \toprule
\textbf{} & \textbf{Cloud routine} & \textbf{Desktop scheduled task} & \textbf{\texttt{/\allowbreak{}loop}} \\
\midrule\endhead
  \bottomrule\endfoot
  \bottomrule\endlastfoot
  \tblrows
  \end{longtable}}%
\else
  \typeout{HANDBOOK-TABLE atomic \the\dimexpr\ht\tblbox+\dp\tblbox\relax}%
  \noindent\tblbody
\fi
\par\addvspace{2.6mm}
\endgroup

Source: [20], [24], [82].


Figure 7 reduces the choice to the question that actually settles it: what has to be true already, at the moment the work runs, with nobody watching.


\begin{figure}[tbp]
\begin{codefig}{8.0}{9.4}
\cl{~~~Must~it~run~when~your~machine~is~off?}
\cl{~~~~~~\pmboxdrawuni{2502}}
\cl{~~~~~~\pmboxdrawuni{251C}\pmboxdrawuni{2500}\pmboxdrawuni{2500}~yes~\pmboxdrawuni{2500}\pmboxdrawuni{2500}▶~Does~it~need~your~local~files~or~desktop~software?}
\cl{~~~~~~\pmboxdrawuni{2502}~~~~~~~~~~~~~~\pmboxdrawuni{251C}\pmboxdrawuni{2500}\pmboxdrawuni{2500}~yes~\pmboxdrawuni{2500}\pmboxdrawuni{2500}▶~An~always-on~controlled~host~running}
\cl{~~~~~~\pmboxdrawuni{2502}~~~~~~~~~~~~~~\pmboxdrawuni{2502}~~~~~~~~~~~headless~Claude~Code,~or~CI}
\cl{~~~~~~\pmboxdrawuni{2502}~~~~~~~~~~~~~~\pmboxdrawuni{2514}\pmboxdrawuni{2500}\pmboxdrawuni{2500}~no~~\pmboxdrawuni{2500}\pmboxdrawuni{2500}▶~CLOUD~ROUTINE~~~~~~~~(min.~interval~1~hour,}
\cl{~~~~~~\pmboxdrawuni{2502}~~~~~~~~~~~~~~~~~~~~~~~~~~runs~autonomously,~all~connectors~by~default)}
\cl{~~~~~~\pmboxdrawuni{2502}}
\cl{~~~~~~\pmboxdrawuni{2514}\pmboxdrawuni{2500}\pmboxdrawuni{2500}~no~~\pmboxdrawuni{2500}\pmboxdrawuni{2500}▶~Must~it~survive~a~closed~terminal?}
\cl{~~~~~~~~~~~~~~~~~~~~~\pmboxdrawuni{251C}\pmboxdrawuni{2500}\pmboxdrawuni{2500}~yes~\pmboxdrawuni{2500}\pmboxdrawuni{2500}▶~DESKTOP~SCHEDULED~TASK~~~(min.~1~minute,}
\cl{~~~~~~~~~~~~~~~~~~~~~\pmboxdrawuni{2502}~~~~~~~~~~~local~files,~per-task~permissions)}
\cl{~~~~~~~~~~~~~~~~~~~~~\pmboxdrawuni{2514}\pmboxdrawuni{2500}\pmboxdrawuni{2500}~no~~\pmboxdrawuni{2500}\pmboxdrawuni{2500}▶~Is~the~trigger~a~time,~or~an~event?}
\cl{~~~~~~~~~~~~~~~~~~~~~~~~~~~~~~~~~~~~\pmboxdrawuni{251C}\pmboxdrawuni{2500}\pmboxdrawuni{2500}~time~~\pmboxdrawuni{2500}\pmboxdrawuni{2500}▶~/loop~~~~~~(min.~1~minute,}
\cl{~~~~~~~~~~~~~~~~~~~~~~~~~~~~~~~~~~~~\pmboxdrawuni{2502}~~~~~~~~~~~~~session-scoped,~7-day~expiry)}
\cl{~~~~~~~~~~~~~~~~~~~~~~~~~~~~~~~~~~~~\pmboxdrawuni{251C}\pmboxdrawuni{2500}\pmboxdrawuni{2500}~event~\pmboxdrawuni{2500}\pmboxdrawuni{2500}▶~CHANNEL~~~~(pushed~in)}
\cl{~~~~~~~~~~~~~~~~~~~~~~~~~~~~~~~~~~~~\pmboxdrawuni{2514}\pmboxdrawuni{2500}\pmboxdrawuni{2500}~a~stream~to~watch~\pmboxdrawuni{2500}\pmboxdrawuni{2500}▶~MONITOR~tool}
\cl{}
\cl{~~~Needs~exact~timing,~catch-up,~retries~or~a~service-level~commitment?}
\cl{~~~~~~\pmboxdrawuni{2500}\pmboxdrawuni{2500}▶~None~of~the~above.~Use~a~purpose-built~scheduler~and~call~Claude~Code}
\cl{~~~~~~~~~~from~it,~securing~the~account,~directory,~permissions~and~logging~yourself.}
\end{codefig}
\figcaption{Figure 7 — Choosing a scheduling mechanism. Start from what must be true when the work runs, not from how often it runs.}
\end{figure}


\FloatBarrier
\setcounter{section}{0}
\section{What a routine actually is}

\begin{calloutbox}{palegrey}{quoterule}{1.2mm}
\textbf{Routines are in research preview.} Behaviour, limits and the API surface may change [20].
\end{calloutbox}

A routine is a saved configuration — a prompt, one or more repositories, an environment, and a set of connectors — with one or more triggers attached [20]:


\begin{itemize}
\item \textbf{Scheduled} — hourly, daily, weekdays, weekly, or a single future timestamp.
\item \textbf{API} — an HTTP POST to a per-routine \texttt{/\allowbreak{}fire} endpoint with a bearer token.
\item \textbf{GitHub} — pull-request and release events, with filters on author, title, body, base and head branch, labels, draft state and merged state.
\end{itemize}

A single routine can combine all three. This is M19 in Table 1: routines are widely described as scheduled sessions, which is a third of the feature.


Create them at \texttt{claude.\allowbreak{}ai/\allowbreak{}code/\allowbreak{}routines}, in the Desktop app's \textbf{Routines} sidebar (choosing \textbf{Cloud}, not \textbf{Local}), or from the CLI with \texttt{/\allowbreak{}schedule} (alias \texttt{/\allowbreak{}routines}), which also supports \texttt{/\allowbreak{}schedule li\allowbreak{}st}, \texttt{/\allowbreak{}schedule up\allowbreak{}date} and \texttt{/\allowbreak{}schedule ru\allowbreak{}n} [20]. Custom cron intervals are set through \texttt{/\allowbreak{}schedule up\allowbreak{}date}; expressions more frequent than hourly are rejected.


\FloatBarrier
\setcounter{section}{1}
\section{Why a routine deserves a design review}

\textbf{Routine runs are autonomous. There is no permission-mode picker and no approval prompts} [20]. Four consequences follow, and all four are easy to miss:


\begin{enumerate}
\item \textbf{All of your connected connectors are included by default.} Claude can use every tool from an included connector, including writes, without asking. Remove everything the routine does not need.
\item \textbf{It acts as you.} Commits and pull requests carry your GitHub user; Slack messages, tickets and other connector actions use your linked accounts.
\item \textbf{Branch protection is partial.} Claude pushes to \texttt{claude/\allowbreak{}}-prefixed branches, which are always accepted. A push to another branch is rejected if the branch is protected, if someone else has an open pull request from it, or if it carries commits authored by someone else.
\item \textbf{Environment variables are visible to anyone who uses the environment} [83]. Scope credentials accordingly.
\end{enumerate}

The \textbf{Default} environment uses \emph{Trusted} network access, allowing only a default allowlist of package registries and common development domains; requests outside it fail with \texttt{403} and \texttt{x-\allowbreak{}deny-\allowbreak{}reason:\allowbreak{} host\_\allowbreak{}not\_\allowbreak{}allowed}. Connector traffic routes through Anthropic's servers and does not need allowlist changes [20], [83].


\FloatBarrier
\setcounter{section}{2}
\section{Fire-payload text is untrusted by design}

Text supplied with an API trigger or with \textbf{Run now} arrives wrapped in a \texttt{<routine-\allowbreak{}fire-\allowbreak{}payload>} block that labels it as untrusted data and tells Claude not to follow instructions inside it unless the routine's own prompt says to [20]. A routine must therefore \emph{opt in} — "investigate the alert described in the routine-fire-payload block" — or the text is inert context. This is a deliberate injection control: anyone holding the bearer token can send \texttt{text}, and the wrapper ensures a leaked token delivers data rather than instructions.


The \texttt{/\allowbreak{}fire} endpoint ships under the \texttt{experimenta\allowbreak{}l-\allowbreak{}cc-\allowbreak{}routine-\allowbreak{}2026-\allowbreak{}04-\allowbreak{}01} beta header; request and response shapes may change [20].


\FloatBarrier
\setcounter{section}{3}
\section{Before you activate one}

\begin{enumerate}
\item Remove every connector the task does not need.
\item Use read-only or least-privilege service accounts.
\item Specify the repository and the branch explicitly.
\item Prohibit direct production push and external publication in the prompt.
\item Define the safe failure state.
\item Configure notifications.
\item Run the prompt interactively, against test data, first.
\item Inspect the first several scheduled runs by opening them.
\end{enumerate}

\begin{calloutbox}{palebrass}{brassdark}{2.0mm}
\textbf{CAUTION:} a green status in the run list means the session started and exited without an \emph{infrastructure} error. It does not mean the task succeeded [20]. Blocked network requests, missing connector tools and task-level failures all appear inside the transcript, not in the status indicator. Never treat the indicator as assurance.
\end{calloutbox}

Routines draw down subscription usage and are additionally subject to a daily per-account run cap; one-off runs do not count against it [20]. Team and Enterprise Owners can disable routines organisation-wide, which also hides \texttt{/\allowbreak{}schedule} in the CLI from 2.1.227 [20].


\FloatBarrier
\setcounter{section}{4}
\section{Desktop scheduled tasks and headless automation}

A Desktop scheduled task runs on your machine with access to local files and tools, at intervals down to one minute, with per-task permission configuration [82]. It is the right choice when the work needs local resources and must survive a closed terminal.


A separate architecture is a local operating-system scheduler invoking headless Claude Code (\texttt{claude -\allowbreak{}p}) [84]. That is a legitimate design, and everything about it — the account it uses, its working directory, its permission rules, its logging, its retries, its cost, and the fact that trust verification is disabled under \texttt{-\allowbreak{}p} [53] — must be secured explicitly. It is not a scheduling feature; it is a system you are building.


\begin{calloutbox}{palebrass}{brassdark}{2.0mm}
\textbf{CAUTION:} never make a research-preview routine the only control for business-critical or safety-critical work.
\end{calloutbox}


\begingroup
\def\kprows{%
\item Routines are a research preview that trigger on schedules, API calls and GitHub events.
\item They run autonomously with every connected connector by default, under your identity.
\item A green run status means no infrastructure error, not that the task succeeded.
\item Minimum interval is one hour, and fire-payload text arrives explicitly labelled as untrusted.
}%
\def\kpbody{\begin{minipage}{\textwidth}\subsection*{Key points}\begin{itemize}\kprows\end{itemize}\end{minipage}}%
\begingroup
\def\sloppy{\tolerance 9999\emergencystretch 3em\hfuzz 200pt\vfuzz 200pt}%
\hbadness=10000\vbadness=10000\hfuzz=200pt\vfuzz=200pt
\global\setbox\kpbox=\hbox{\kpbody}%
\endgroup
\par\addvspace{4.2mm}
\ifdim\dimexpr\ht\kpbox+\dp\kpbox\relax>0.30\textheight
  \typeout{HANDBOOK-KEYPOINTS broken \the\dimexpr\ht\kpbox+\dp\kpbox\relax}%
  \subsection*{Key points}
  \begin{itemize}\kprows\end{itemize}
\else
  \typeout{HANDBOOK-KEYPOINTS atomic \the\dimexpr\ht\kpbox+\dp\kpbox\relax}%
  \noindent\kpbody
\fi
\par\addvspace{1.4mm}
\endgroup

\breakrule

\FloatBarrier
\parttitle{Part VI}{Distributed and cloud operation}
\addcontentsline{toc}{part}{Part VI \textemdash\ Distributed and cloud operation}

\FloatBarrier
\renewcommand{\chaptertitlelabel}{Part VI \textperiodcentered\ Chapter 25}
\setcounter{chapter}{24}
\chapter{Continue local work with Remote Control}

Remote Control lets another signed-in device supervise a Claude Code process running on your own machine. \textbf{Execution stays local.} Code, files and commands never leave your computer; what travels is the conversation, which is relayed through the Anthropic API over TLS and stored on Anthropic servers while connected so the conversation syncs across devices [39], [53].


\begin{codeblock}{9.0}{10.6}
\cl{/remote-control~~~~~~~~~~\#~alias~/rc}
\end{codeblock}

\begin{codeblock}{9.0}{10.6}
\cl{claude~--remote-control~research-brief~~~~\#~alias~--rc}
\cl{claude~remote-control~~~~~~~~~~~~~~~~~~~~~\#~server~mode}
\end{codeblock}

Server mode supports multiple or on-demand sessions subject to account and platform availability. Name sessions clearly so a phone cannot confuse two projects; a session renamed from claude.ai or the Claude app is renamed in the CLI too, from 2.1.221 [30].


The connection uses multiple short-lived, narrowly scoped credentials, each limited to a purpose and expiring independently, to limit the blast radius of any single compromise [53].


\FloatBarrier
\setcounter{section}{0}
\section{Operating rules}

\begin{itemize}
\item Start from a narrow project directory, never a home directory.
\item Keep your normal permission mode and sandboxing. Remoteness is not a reason to widen access.
\item Do not expose a broad shell merely because you are away from the desk.
\item Confirm which machine, repository, branch and account the session controls before you approve anything.
\item Close remote control when supervision is over.
\item Treat the relayed transcript as cloud-stored data for as long as you are connected (Chapter 32).
\end{itemize}

Remote Control is good for answering a question, unblocking a decision, or approving a safe local edit. It is poor for inspecting large diffs, reviewing sensitive commands, or making high-impact approvals on a small screen. Defer those until you can review properly — the small screen is not a reason to lower the bar, it is a reason to wait.


The local Claude Code process must remain running. Closing it ends the execution path even though the mobile or web client stays open.


\FloatBarrier
\setcounter{section}{1}
\section{Mobile}

The Claude app for iOS and Android starts, monitors and steers Claude Code tasks, including starting a session on your machine from your phone [36], [48]. \texttt{/\allowbreak{}mobile} (aliases \texttt{/\allowbreak{}ios}, \texttt{/\allowbreak{}android}) shows a download QR code [35]. The same operating rules apply, with the review caveat above sharpened: a phone is an excellent notification surface and a poor audit surface.


\subsection*{Exercise}

Start Remote Control on the practice workspace. From a second device, request a read-only status report. Confirm the output appears locally, then disconnect. Grant no new permissions during the exercise.


\emph{A worked solution is given in Appendix M.}


\begingroup
\def\kprows{%
\item Execution stays local; the transcript is relayed and stored server-side while connected.
\item Keep your normal permission mode — remoteness is not a reason to widen access.
\item A phone is an excellent notification surface and a poor audit surface. Defer high-impact approvals.
\item Closing the local process ends the execution path even though the client stays open.
}%
\def\kpbody{\begin{minipage}{\textwidth}\subsection*{Key points}\begin{itemize}\kprows\end{itemize}\end{minipage}}%
\begingroup
\def\sloppy{\tolerance 9999\emergencystretch 3em\hfuzz 200pt\vfuzz 200pt}%
\hbadness=10000\vbadness=10000\hfuzz=200pt\vfuzz=200pt
\global\setbox\kpbox=\hbox{\kpbody}%
\endgroup
\par\addvspace{4.2mm}
\ifdim\dimexpr\ht\kpbox+\dp\kpbox\relax>0.30\textheight
  \typeout{HANDBOOK-KEYPOINTS broken \the\dimexpr\ht\kpbox+\dp\kpbox\relax}%
  \subsection*{Key points}
  \begin{itemize}\kprows\end{itemize}
\else
  \typeout{HANDBOOK-KEYPOINTS atomic \the\dimexpr\ht\kpbox+\dp\kpbox\relax}%
  \noindent\kpbody
\fi
\par\addvspace{1.4mm}
\endgroup

\FloatBarrier
\renewcommand{\chaptertitlelabel}{Part VI \textperiodcentered\ Chapter 26}
\setcounter{chapter}{25}
\chapter{Cloud sessions, \texttt{-\allowbreak{}-\allowbreak{}cloud} and \texttt{-\allowbreak{}-\allowbreak{}teleport}}

\begin{calloutbox}{palegrey}{quoterule}{1.2mm}
Claude Code on the web is in research preview for Pro, Max and Team users, and for Enterprise users with premium or Chat + Claude Code seats [47].
\end{calloutbox}

A cloud session runs on Anthropic-managed infrastructure at \texttt{claude.\allowbreak{}ai/\allowbreak{}code}, or on your organisation's self-hosted environment when routed there. Sessions persist when you close the browser and can be monitored from the mobile app [47]. This is a different mechanism from Remote Control, and confusing the two is the most common error in this area: \texttt{-\allowbreak{}-\allowbreak{}cloud} creates a cloud session that runs elsewhere; \texttt{-\allowbreak{}-\allowbreak{}remote-\allowbreak{}control} exposes a \textbf{local} session for monitoring [47].


\FloatBarrier
\setcounter{section}{0}
\section{From terminal to cloud}

\begin{codeblock}{9.0}{10.6}
\cl{claude~--cloud~"Fix~the~authentication~bug~in~src/auth/login.ts"}
\end{codeblock}

The cloud VM clones your current directory's \textbf{GitHub remote at your current branch}, not your local checkout — push first if you have local commits [47]. One repository at a time. From a repository not connected to GitHub, Claude Code bundles the local repository and uploads it: full history across branches plus uncommitted changes to tracked files, under 100 MB, untracked files excluded, and no ability to push back unless GitHub authentication is also configured. Force bundling with \texttt{CCR\_\allowbreak{}FORCE\_\allowbreak{}BUNDLE=\allowbreak{}1} [47].


Send a follow-up to a running cloud session from any machine where you are logged in:


\begin{codeblock}{8.0}{9.4}
\cl{claude~-p~"your~message"~--cloud~<session-id>~~~~~\#~queue~and~exit}
\cl{claude~--cloud~<session-id>~~~~~~~~~~~~~~~~~~~~~~~\#~attach~interactively~(rolling~out)}
\end{codeblock}

\FloatBarrier
\setcounter{section}{1}
\section{From cloud to terminal}

\begin{codeblock}{9.0}{10.6}
\cl{claude~--teleport~~~~~~~~~~~~~~~\#~interactive~picker}
\cl{claude~--teleport~<session-id>~~\#~directly}
\end{codeblock}

\begin{codeblock}{9.0}{10.6}
\cl{/teleport~~~~~~\#~alias~/tp,~from~inside~a~session}
\end{codeblock}

\texttt{-\allowbreak{}-\allowbreak{}teleport} fetches and checks out the cloud session's branch and loads the full conversation history locally. It is \textbf{distinct from \texttt{-\allowbreak{}-\allowbreak{}resume}}: \texttt{-\allowbreak{}-\allowbreak{}resume} reopens a conversation from this machine's local history, \texttt{-\allowbreak{}-\allowbreak{}teleport} pulls a cloud session and its branch [47]. The terminal gets its own copy: new work stays local and does not appear back in the cloud session. To keep steering from a phone afterwards, start \texttt{/\allowbreak{}remote-\allowbreak{}control} in the local session.


Teleport requires a clean git state, a checkout of the same repository (not a fork), the branch pushed to the remote, and the same claude.ai account [47].


\FloatBarrier
\setcounter{section}{2}
\section{GitHub access, and what it actually grants}

Two methods: authorising the Claude GitHub App during web onboarding, or \texttt{/\allowbreak{}web-\allowbreak{}setup}, which syncs your local \texttt{gh} CLI token to your Claude account [47].


\begin{calloutbox}{palebrass}{brassdark}{2.0mm}
\textbf{CAUTION:} with either method, \textbf{a cloud session can access any repository the connecting GitHub account can see} — not only the repositories the App is installed on. App installation enables pull-request webhooks for auto-fix; it is not a session-level access control. To restrict what a team can reach from cloud sessions, restrict access on GitHub itself [47].
\end{calloutbox}

Team and Enterprise Owners can disable \texttt{/\allowbreak{}web-\allowbreak{}setup}. Organisations with Zero Data Retention cannot use \texttt{/\allowbreak{}web-\allowbreak{}setup} or other cloud session features at all [40], [47].


\FloatBarrier
\setcounter{section}{3}
\section{Auto-fix, and a real hazard}

Claude can watch a pull request and respond automatically to CI failures and review comments; turn it on from the CI status bar, with \texttt{/\allowbreak{}autofix-\allowbreak{}pr} from the branch, from mobile, or by pasting a PR URL [47]. Claude makes a change when it is confident and the change does not conflict with earlier instructions, asks when a comment is ambiguous or architecturally significant, and notes duplicates.


\begin{calloutbox}{palebrass}{brassdark}{2.0mm}
\textbf{CAUTION:} Claude may reply to review-comment threads on GitHub using \textbf{your} account, labelled as coming from Claude Code. If your repository uses comment-triggered automation — Atlantis, Terraform Cloud, or custom Actions running on \texttt{issue\_\allowbreak{}comment} — Claude replying on your behalf can trigger those workflows. Review your repository automation before enabling auto-fix, and disable it where a comment can deploy infrastructure [47].
\end{calloutbox}

GitHub emits no webhook when the base branch advances into a merge conflict, so auto-fix cannot react to conflicts; open the session and ask for a rebase [47].


\FloatBarrier
\setcounter{section}{4}
\section{Sharing a cloud session}

On Enterprise and Team accounts the options are \textbf{Private} and \textbf{Team}, with repository-access verification enabled by default. On Max and Pro they are \textbf{Private} and \textbf{Public}, where public means visible to any signed-in claude.ai user and \textbf{repository access verification is not enabled by default} [47].


\begin{calloutbox}{palebrass}{brassdark}{2.0mm}
\textbf{CAUTION:} a session may contain code and credentials from private repositories. Check before sharing. Sharing settings, including requiring repository access and hiding your name, are under Settings → Claude Code → Sharing settings [47].
\end{calloutbox}

\FloatBarrier
\setcounter{section}{5}
\section{Isolation and limits}

Each Anthropic-hosted session runs in an isolated VM; network access is limited by default and configurable; sensitive credentials such as git credentials and signing keys are never inside the sandbox, authentication going through a secure proxy with scoped credentials; git push is restricted to the current working branch; and all operations are logged for audit [47], [53]. Self-hosted environments move all of that responsibility to your deployment.


Two limits deserve planning attention. \textbf{When network access is disabled, Claude Code can still reach the Anthropic API, which may allow data to exit the VM} [47]. And if your organisation enforces IP allowlisting, every Anthropic-hosted cloud session fails authentication, as do Code Review and cloud-hosted routines; an exemption must be arranged [47].


\begingroup
\def\kprows{%
\item \texttt{-\allowbreak{}-\allowbreak{}cloud} creates a session that runs elsewhere; \texttt{-\allowbreak{}-\allowbreak{}remote-\allowbreak{}control} exposes a local one. They are unrelated.
\item A cloud session can reach any repository the connecting GitHub account can see, not only where the App is installed.
\item \texttt{-\allowbreak{}-\allowbreak{}teleport} pulls a cloud session and its branch into the terminal; \texttt{-\allowbreak{}-\allowbreak{}resume} does not list cloud sessions.
\item Auto-fix can reply to review threads using your GitHub account, which can trigger comment-driven CI automation.
}%
\def\kpbody{\begin{minipage}{\textwidth}\subsection*{Key points}\begin{itemize}\kprows\end{itemize}\end{minipage}}%
\begingroup
\def\sloppy{\tolerance 9999\emergencystretch 3em\hfuzz 200pt\vfuzz 200pt}%
\hbadness=10000\vbadness=10000\hfuzz=200pt\vfuzz=200pt
\global\setbox\kpbox=\hbox{\kpbody}%
\endgroup
\par\addvspace{4.2mm}
\ifdim\dimexpr\ht\kpbox+\dp\kpbox\relax>0.30\textheight
  \typeout{HANDBOOK-KEYPOINTS broken \the\dimexpr\ht\kpbox+\dp\kpbox\relax}%
  \subsection*{Key points}
  \begin{itemize}\kprows\end{itemize}
\else
  \typeout{HANDBOOK-KEYPOINTS atomic \the\dimexpr\ht\kpbox+\dp\kpbox\relax}%
  \noindent\kpbody
\fi
\par\addvspace{1.4mm}
\endgroup

\FloatBarrier
\renewcommand{\chaptertitlelabel}{Part VI \textperiodcentered\ Chapter 27}
\setcounter{chapter}{26}
\chapter{Push events into a session with channels}

\begin{calloutbox}{palegrey}{quoterule}{1.2mm}
Channels are a research preview. They require Anthropic authentication through claude.ai or a Console API key, and are unavailable on Amazon Bedrock, Google Cloud's Agent Platform and Microsoft Foundry. Team and Enterprise organisations must explicitly enable them [51].
\end{calloutbox}

A channel is an MCP server that \textbf{pushes} events into a running session, rather than waiting to be queried [51]. Where a standard connector gives Claude on-demand access to a system, a channel lets a system reach into the session you already have open — a CI failure, a monitoring alert, a chat message — while Claude still has your files open and remembers what you were debugging.


Events arrive only while the session is open, so an always-on setup means a background process or a persistent terminal.


\FloatBarrier
\setcounter{section}{0}
\section{Enabling one}

Channels install as plugins and require Bun. Telegram, Discord and iMessage are included in the preview, with \texttt{fakechat} as a local demo [51].


\begin{codeblock}{9.0}{10.6}
\cl{/plugin~install~telegram@claude-plugins-official}
\cl{/telegram:configure~<token>}
\end{codeblock}

\begin{codeblock}{9.0}{10.6}
\cl{claude~--channels~plugin:telegram@claude-plugins-official}
\end{codeblock}

\textbf{Being present in \texttt{.\allowbreak{}mcp.\allowbreak{}json} is not enough}: a server must also be named in \texttt{-\allowbreak{}-\allowbreak{}channels} for that session [51]. Neither \texttt{-\allowbreak{}-\allowbreak{}channels} nor \texttt{-\allowbreak{}-\allowbreak{}dangerously-\allowbreak{}load-\allowbreak{}development-\allowbreak{}channels} appears in \texttt{claude -\allowbreak{}-\allowbreak{}help} while the feature is in preview; the flags work regardless.


\FloatBarrier
\setcounter{section}{1}
\section{The security model, and the part that matters}

Every approved channel plugin maintains a \textbf{sender allowlist}; messages from anyone else are silently dropped. Telegram and Discord bootstrap it by pairing — message the bot, receive a code, approve it in session, then lock down with \texttt{policy allo\allowbreak{}wlist}. iMessage lets your own messages through automatically and takes other contacts by handle [51].


\begin{calloutbox}{palebrass}{brassdark}{2.0mm}
\textbf{CAUTION — read this before enabling a channel.} The allowlist also gates \textbf{permission relay} where a channel declares it. \textbf{Anyone who can reply through the channel can approve or deny tool use in your session} [51]. Add to a channel allowlist only senders you would trust with that authority, and understand that a compromised messaging account becomes a compromised permission gate.
\end{calloutbox}

Organisational controls are two managed settings that users cannot override: \texttt{channels\allowbreak{}Enabled}, the master switch, and \texttt{allowed\allowbreak{}Channel\allowbreak{}Plugins}, which replaces the Anthropic-maintained plugin allowlist with your own [51]. claude.ai Team and Enterprise organisations are blocked until an Owner enables channels; Console organisations with API-key authentication are permitted by default unless managed settings are deployed. An empty \texttt{allowed\allowbreak{}Channel\allowbreak{}Plugins} array blocks allowlisted plugins but the development flag can still bypass it — to block channels entirely, leave \texttt{channels\allowbreak{}Enabled} unset.


In non-interactive mode with \texttt{-\allowbreak{}p}, tools that need terminal input — multiple-choice questions and plan approval — are disabled so the session never stalls [51].


\FloatBarrier
\setcounter{section}{2}
\section{When a channel is the right shape}

\begingroup
\def\tblrows{%
Claude Code on the web & Runs a task in a fresh cloud sandbox cloned from GitHub \\
Claude in Slack & Spawns a web session from an \texttt{@Claude} mention \\
Standard MCP server & Claude queries it during a task; nothing is pushed \\
Remote Control & You drive your local session from claude.ai or mobile \\
\textbf{Channels} & An external system pushes an event into your already-running local session \\
}%
\def\tblbody{\begin{minipage}{\textwidth}\boxcaption{Table 19 — Where channels sit among the ways to reach a session}
{\footnotesize\begin{tabular}{L{33.2mm}L{109.8mm}}
\toprule
\textbf{Mechanism} & \textbf{What it does} \\
\midrule
\tblrows
\bottomrule\end{tabular}}\end{minipage}}%
\begingroup
\def\sloppy{\tolerance 9999\emergencystretch 3em\hfuzz 200pt\vfuzz 200pt}%
\hbadness=10000\vbadness=10000\hfuzz=200pt\vfuzz=200pt
\global\setbox\tblbox=\hbox{\tblbody}%
\endgroup
\par\addvspace{2.6mm}
\ifdim\dimexpr\ht\tblbox+\dp\tblbox\relax>0.55\textheight
  \typeout{HANDBOOK-TABLE broken \the\dimexpr\ht\tblbox+\dp\tblbox\relax}%
  \tabcaption{Table 19 — Where channels sit among the ways to reach a session}
  {\footnotesize\begin{longtable}{L{33.2mm}L{109.8mm}}
  \toprule
\textbf{Mechanism} & \textbf{What it does} \\
\midrule\endfirsthead
  \multicolumn{2}{@{}l@{}}{%
  \sffamily\footnotesize\itshape\color{inkgrey}Table 19 — Where channels sit among the ways to reach a session \textemdash\ continued}\\[1.2mm]
  \toprule
\textbf{Mechanism} & \textbf{What it does} \\
\midrule\endhead
  \bottomrule\endfoot
  \bottomrule\endlastfoot
  \tblrows
  \end{longtable}}%
\else
  \typeout{HANDBOOK-TABLE atomic \the\dimexpr\ht\tblbox+\dp\tblbox\relax}%
  \noindent\tblbody
\fi
\par\addvspace{2.6mm}
\endgroup

Source: [51].


Two shapes recur: a \textbf{chat bridge}, where you ask a question from your phone and the work runs on your machine against real files; and a \textbf{webhook receiver}, where CI, an error tracker or a deploy pipeline delivers an event to a session that already has context.


\begingroup
\def\kprows{%
\item A channel is an MCP server that pushes events into a session you already have open.
\item A server must be named in \texttt{-\allowbreak{}-\allowbreak{}channels} for that session; being in \texttt{.\allowbreak{}mcp.\allowbreak{}json} is not enough.
\item The sender allowlist also gates permission relay: anyone who can reply through the channel can approve tool use in your session.
\item Organisations control availability with \texttt{channels\allowbreak{}Enabled} and \texttt{allowed\allowbreak{}Channel\allowbreak{}Plugins} in managed settings.
}%
\def\kpbody{\begin{minipage}{\textwidth}\subsection*{Key points}\begin{itemize}\kprows\end{itemize}\end{minipage}}%
\begingroup
\def\sloppy{\tolerance 9999\emergencystretch 3em\hfuzz 200pt\vfuzz 200pt}%
\hbadness=10000\vbadness=10000\hfuzz=200pt\vfuzz=200pt
\global\setbox\kpbox=\hbox{\kpbody}%
\endgroup
\par\addvspace{4.2mm}
\ifdim\dimexpr\ht\kpbox+\dp\kpbox\relax>0.30\textheight
  \typeout{HANDBOOK-KEYPOINTS broken \the\dimexpr\ht\kpbox+\dp\kpbox\relax}%
  \subsection*{Key points}
  \begin{itemize}\kprows\end{itemize}
\else
  \typeout{HANDBOOK-KEYPOINTS atomic \the\dimexpr\ht\kpbox+\dp\kpbox\relax}%
  \noindent\kpbody
\fi
\par\addvspace{1.4mm}
\endgroup

\FloatBarrier
\renewcommand{\chaptertitlelabel}{Part VI \textperiodcentered\ Chapter 28}
\setcounter{chapter}{27}
\chapter{Publish work as artifacts}

An artifact is a live, interactive web page published from a session to a private URL on claude.ai, updating in place as the session continues [50]. For the research workflow of this book it is the natural delivery surface for anything easier to look at than to read: an annotated diff, a dashboard, a comparison of options, an investigation timeline that fills in while a long task runs.


Claude may publish one on its own when the output suits a page, or you can ask. Claude writes an HTML or Markdown file in your project and publishes it; the first publish of a new artifact asks permission, republishing an approved artifact does not [50]. \texttt{Ctrl+]} reopens the most recent artifact; \texttt{/\allowbreak{}artifacts} lists every artifact you own and every one shared with you, with \texttt{o} to open, \texttt{c} to copy the link and \texttt{Enter} to attach it to the session (2.1.208+).


To update an artifact from a different session, give Claude its URL or attach it with \texttt{/\allowbreak{}artifacts}. Without either, a new session creates a \textbf{new} artifact instead of updating the existing one [50].


\FloatBarrier
\setcounter{section}{0}
\section{What an artifact is not}

It is a capture of work, not an application: one self-contained page, no backend, no form storage, no routes [50].


\begingroup
\def\tblrows{%
External requests & A strict Content Security Policy blocks scripts, stylesheets, fonts and images from other hosts, and \texttt{fetch}, XHR and WebSocket calls. Google Fonts is the one exception; MCP connector calls go through claude.ai \\
No backend & Cannot store form input or authenticate viewers \\
Single page & Relative links do not resolve; use in-page anchors \\
File types & \texttt{.\allowbreak{}html}, \texttt{.\allowbreak{}htm} or \texttt{.\allowbreak{}md} \\
Rendered size & 16 MiB or smaller \\
}%
\def\tblbody{\begin{minipage}{\textwidth}\boxcaption{Table 20 — Artifact page constraints}
{\footnotesize\begin{tabular}{L{24.3mm}L{118.7mm}}
\toprule
\textbf{Constraint} & \textbf{Effect} \\
\midrule
\tblrows
\bottomrule\end{tabular}}\end{minipage}}%
\begingroup
\def\sloppy{\tolerance 9999\emergencystretch 3em\hfuzz 200pt\vfuzz 200pt}%
\hbadness=10000\vbadness=10000\hfuzz=200pt\vfuzz=200pt
\global\setbox\tblbox=\hbox{\tblbody}%
\endgroup
\par\addvspace{2.6mm}
\ifdim\dimexpr\ht\tblbox+\dp\tblbox\relax>0.55\textheight
  \typeout{HANDBOOK-TABLE broken \the\dimexpr\ht\tblbox+\dp\tblbox\relax}%
  \tabcaption{Table 20 — Artifact page constraints}
  {\footnotesize\begin{longtable}{L{24.3mm}L{118.7mm}}
  \toprule
\textbf{Constraint} & \textbf{Effect} \\
\midrule\endfirsthead
  \multicolumn{2}{@{}l@{}}{%
  \sffamily\footnotesize\itshape\color{inkgrey}Table 20 — Artifact page constraints \textemdash\ continued}\\[1.2mm]
  \toprule
\textbf{Constraint} & \textbf{Effect} \\
\midrule\endhead
  \bottomrule\endfoot
  \bottomrule\endlastfoot
  \tblrows
  \end{longtable}}%
\else
  \typeout{HANDBOOK-TABLE atomic \the\dimexpr\ht\tblbox+\dp\tblbox\relax}%
  \noindent\tblbody
\fi
\par\addvspace{2.6mm}
\endgroup

\FloatBarrier
\setcounter{section}{1}
\section{Sharing, and who sees what}

A new artifact is visible only to you. On Team and Enterprise plans you can grant access to specific people or to everyone in the organisation, and make someone an \textbf{editor} who can publish new versions. On Pro and Max, a public link — open to anyone on the internet, no sign-in — is the only sharing option; on Team and Enterprise, public sharing is off until an Owner enables \textbf{External sharing} [50].


Each publish becomes a version, and the Share control chooses which version viewers see. Organisation viewers see your name; a viewer opening a public link from outside your organisation sees \texttt{Content is \allowbreak{}user-\allowbreak{}generated a\allowbreak{}nd unverifi\allowbreak{}ed.\allowbreak{}} instead [50].


\FloatBarrier
\setcounter{section}{2}
\section{Live data through connectors}

A page can call MCP connectors each time someone views it, so it shows current data rather than a snapshot (2.1.209+). Claude declares which connectors the page may call, and the page cannot call others [50]. The access model is the important part:


\begin{itemize}
\item \textbf{Each viewer uses their own connectors.} Calls go through the viewing account's connections, so two people opening the same dashboard can see different data. The page never sees anyone's credentials.
\item \textbf{Viewers approve access first}, and a decline lasts for that page load.
\item \textbf{Actions use the viewer's account too.} A control that posts a message or updates an issue acts as whoever clicked it.
\end{itemize}

An artifact that calls connectors \textbf{cannot be shared to a public link on any plan} [50]. Ask Claude to include a fallback message in each live section naming the connector it needs, so a viewer without the connection sees what to connect rather than an empty panel.


\FloatBarrier
\setcounter{section}{3}
\section{Governance}

For a regulated organisation, the controls are these [50]: artifacts are off by default on Enterprise and enabled by an Owner; connector calls from artifacts have a \textbf{separate} toggle; public sharing is separately controlled and turning it off blocks existing public links without changing each artifact's audience; retention periods can be set independently for private and shared artifacts; publishing, sharing and deleting appear in the audit log under \texttt{claude\_\allowbreak{}artifact\_\allowbreak{}*} event types; and the Compliance API can list, retrieve and delete artifacts organisation-wide.


Artifacts are unavailable where the organisation has customer-managed encryption keys, HIPAA configuration, or Zero Data Retention enabled, and on third-party model providers [50]. The viewer loads pages from a sandboxed \texttt{*.\allowbreak{}claudeuserc\allowbreak{}ontent.\allowbreak{}com} origin, which must be allowlisted if you restrict outbound access [50], [85].


Disable artifacts for your own sessions with \texttt{"disable\allowbreak{}Artifact":\allowbreak{} true}, \texttt{CLAUDE\_\allowbreak{}CODE\_\allowbreak{}DISABLE\_\allowbreak{}ARTIFACT=\allowbreak{}1}, or by adding \texttt{Artifact} to \texttt{permissions.\allowbreak{}deny} [50].


\begin{calloutbox}{palebrass}{brassdark}{2.0mm}
\textbf{CAUTION:} publishing is distribution. An artifact built from a session that touched private repositories, customer data or internal figures carries that content to a hosted page. Review the page, not the prompt, before sharing — and remember that an editor on Team or Enterprise can publish a new version you have not seen.
\end{calloutbox}


\begingroup
\def\kprows{%
\item An artifact is a hosted page, so publishing is distribution — review the page, not the prompt.
\item Connector-backed pages call connectors as the \emph{viewer}, so two people can see different data.
\item A connector-backed artifact cannot be shared to a public link on any plan.
\item Publishing, sharing and deletion are audit-logged, and retention is separately configurable for private and shared artifacts.
}%
\def\kpbody{\begin{minipage}{\textwidth}\subsection*{Key points}\begin{itemize}\kprows\end{itemize}\end{minipage}}%
\begingroup
\def\sloppy{\tolerance 9999\emergencystretch 3em\hfuzz 200pt\vfuzz 200pt}%
\hbadness=10000\vbadness=10000\hfuzz=200pt\vfuzz=200pt
\global\setbox\kpbox=\hbox{\kpbody}%
\endgroup
\par\addvspace{4.2mm}
\ifdim\dimexpr\ht\kpbox+\dp\kpbox\relax>0.30\textheight
  \typeout{HANDBOOK-KEYPOINTS broken \the\dimexpr\ht\kpbox+\dp\kpbox\relax}%
  \subsection*{Key points}
  \begin{itemize}\kprows\end{itemize}
\else
  \typeout{HANDBOOK-KEYPOINTS atomic \the\dimexpr\ht\kpbox+\dp\kpbox\relax}%
  \noindent\kpbody
\fi
\par\addvspace{1.4mm}
\endgroup

\breakrule

\FloatBarrier
\parttitle{Part VII}{Operating at scale}
\addcontentsline{toc}{part}{Part VII \textemdash\ Operating at scale}

\FloatBarrier
\renewcommand{\chaptertitlelabel}{Part VII \textperiodcentered\ Chapter 29}
\setcounter{chapter}{28}
\chapter{Parallel sessions without collisions}

Parallel agents reduce elapsed time and increase token use, coordination overhead and — the binding constraint — human decision load. The course's warning is the right one: three interactive sessions asking frequent questions feel like three colleagues interrupting at once [42].


\FloatBarrier
\setcounter{section}{0}
\section{The mechanisms}

\begingroup
\def\tblrows{%
Background the current session & \texttt{/\allowbreak{}background}, \texttt{/\allowbreak{}bg} & Detach work and free the terminal; carries \texttt{/\allowbreak{}loop} tasks with it \\
New background session & \texttt{claude -\allowbreak{}-\allowbreak{}bg "<prompt\allowbreak{}>"} & Start bounded work separately \\
Fork the conversation & \texttt{/\allowbreak{}fork} & Copy context into another background session \\
Fork subagent & \texttt{/\allowbreak{}subtask} & Bounded result returned to the parent, full context inherited \\
Subagent & Delegation & Fresh-context worker (Chapter 16) \\
Agent view & \texttt{claude agen\allowbreak{}ts} & Monitor and dispatch many sessions from one screen \\
Worktree & \texttt{claude -\allowbreak{}-\allowbreak{}worktree <n\allowbreak{}ame>} & Isolated git working tree and branch \\
Conversation branch & \texttt{/\allowbreak{}branch} & Explore an alternative in the current interface \\
Agent team & See Chapter 30 & Peer sessions with shared tasks and messaging \\
Dynamic workflow & See Chapter 22 & Codified orchestration of dozens to hundreds of agents \\
}%
\def\tblbody{\begin{minipage}{\textwidth}\boxcaption{Table 21 — Ways to run more than one thing at a time}
{\footnotesize\begin{tabular}{L{30.1mm}L{26.1mm}L{83.7mm}}
\toprule
\textbf{Mechanism} & \textbf{Entry point} & \textbf{Use} \\
\midrule
\tblrows
\bottomrule\end{tabular}}\end{minipage}}%
\begingroup
\def\sloppy{\tolerance 9999\emergencystretch 3em\hfuzz 200pt\vfuzz 200pt}%
\hbadness=10000\vbadness=10000\hfuzz=200pt\vfuzz=200pt
\global\setbox\tblbox=\hbox{\tblbody}%
\endgroup
\par\addvspace{2.6mm}
\ifdim\dimexpr\ht\tblbox+\dp\tblbox\relax>0.55\textheight
  \typeout{HANDBOOK-TABLE broken \the\dimexpr\ht\tblbox+\dp\tblbox\relax}%
  \tabcaption{Table 21 — Ways to run more than one thing at a time}
  {\footnotesize\begin{longtable}{L{30.1mm}L{26.1mm}L{83.7mm}}
  \toprule
\textbf{Mechanism} & \textbf{Entry point} & \textbf{Use} \\
\midrule\endfirsthead
  \multicolumn{3}{@{}l@{}}{%
  \sffamily\footnotesize\itshape\color{inkgrey}Table 21 — Ways to run more than one thing at a time \textemdash\ continued}\\[1.2mm]
  \toprule
\textbf{Mechanism} & \textbf{Entry point} & \textbf{Use} \\
\midrule\endhead
  \bottomrule\endfoot
  \bottomrule\endlastfoot
  \tblrows
  \end{longtable}}%
\else
  \typeout{HANDBOOK-TABLE atomic \the\dimexpr\ht\tblbox+\dp\tblbox\relax}%
  \noindent\tblbody
\fi
\par\addvspace{2.6mm}
\endgroup

Sources: [25], [35], [86]–[88].


\FloatBarrier
\setcounter{section}{1}
\section{Prevent file collisions}

Two sessions editing the same checkout can overwrite each other, conflict, or validate against a moving target. Choose one pattern and state it in \texttt{CLAUDE.\allowbreak{}md}:


\begin{enumerate}
\item \textbf{Single writer} — one session edits; the others only review.
\item \textbf{File ownership} — each session owns explicitly non-overlapping paths.
\item \textbf{Git worktrees} — each implementation session gets its own branch and working tree [88]. \texttt{.\allowbreak{}worktreeinc\allowbreak{}lude} controls what is carried into a new worktree, and a subagent can be given \texttt{isolation:\allowbreak{} worktree} to work in a temporary one [25], [88].
\item \textbf{Sequential integration} — a separate session or a human reviews and merges completed units.
\end{enumerate}

Worktrees are the safest default for concurrent change. Each branch runs its own verification before integration, and integration tests run after the merge — not instead of it.


Note the interaction with hooks: \texttt{\$\{CLAUDE\_\allowbreak{}PROJECT\_\allowbreak{}DIR\}} stays at the project root inside a worktree while the hook input's \texttt{cwd} shows the worktree path, so a hook script must decide which it means [73].


\FloatBarrier
\setcounter{section}{2}
\section{Cross-session messaging}

Claude can list and message your other sessions on this machine, and reach sessions on other machines or on the web [89]. \texttt{/\allowbreak{}list-\allowbreak{}agents} (alias \texttt{/\allowbreak{}peers}) enumerates them [35].


\begin{calloutbox}{palebrass}{brassdark}{2.0mm}
\textbf{CAUTION:} an incoming message is told it came from another Claude session, not from you. A peer session \textbf{cannot approve a permission prompt or supply consent on your behalf}, and a session denied an action cannot relay it to another session to bypass the check. In auto mode the classifier treats a relayed approval claim as untrusted input and reviews each message before delivery [29], [89]. Treat messages as coordination input, never as privileged control.
\end{calloutbox}

\FloatBarrier
\setcounter{section}{3}
\section{Protect your attention}

\begin{itemize}
\item Keep at most two or three interactive sessions until you know your own limit.
\item Use background work for anything that does not need frequent decisions.
\item Name sessions by outcome (Section 12.2).
\item Maintain a small dashboard: owner, branch, task, status, next human decision.
\item Batch reviews rather than checking every minute.
\item Stop starting work when integration becomes the bottleneck. This is the real ceiling, and it arrives sooner than the token limit.
\end{itemize}

\subsection*{A parallel task brief \{\#sec:parallel-brief\}}

\begin{codeblock}{7.0}{8.3}
\cl{Task:~~~~~~Review~only~sources/source-set-b/~for~contradictory~claims.}
\cl{Ownership:~Read~source-set-b.~Write~notes/review-b.md~and~nothing~else.}
\cl{Do~not:~~~~Edit~drafts,~change~configuration,~message~external~systems,~or~push.}
\cl{Return:~~~~Claim~IDs,~the~contradiction,~source~passages,~confidence,~open~questions.}
\cl{Done:~~~~~~Every~file~in~source-set-b~is~accounted~for~and~the~output~passes~its~schema~check.}
\end{codeblock}

\begingroup
\def\kprows{%
\item Parallelism is bounded by your decision load, not by tokens; integration becomes the bottleneck first.
\item Choose one collision pattern — single writer, file ownership, worktrees or sequential integration — and write it into \texttt{CLAUDE.\allowbreak{}md}.
\item Worktrees are the safest default for concurrent change.
\item A peer session cannot approve a permission or supply consent on your behalf.
}%
\def\kpbody{\begin{minipage}{\textwidth}\subsection*{Key points}\begin{itemize}\kprows\end{itemize}\end{minipage}}%
\begingroup
\def\sloppy{\tolerance 9999\emergencystretch 3em\hfuzz 200pt\vfuzz 200pt}%
\hbadness=10000\vbadness=10000\hfuzz=200pt\vfuzz=200pt
\global\setbox\kpbox=\hbox{\kpbody}%
\endgroup
\par\addvspace{4.2mm}
\ifdim\dimexpr\ht\kpbox+\dp\kpbox\relax>0.30\textheight
  \typeout{HANDBOOK-KEYPOINTS broken \the\dimexpr\ht\kpbox+\dp\kpbox\relax}%
  \subsection*{Key points}
  \begin{itemize}\kprows\end{itemize}
\else
  \typeout{HANDBOOK-KEYPOINTS atomic \the\dimexpr\ht\kpbox+\dp\kpbox\relax}%
  \noindent\kpbody
\fi
\par\addvspace{1.4mm}
\endgroup

\FloatBarrier
\renewcommand{\chaptertitlelabel}{Part VII \textperiodcentered\ Chapter 30}
\setcounter{chapter}{29}
\chapter{Agent teams}

\begin{calloutbox}{palegrey}{quoterule}{1.2mm}
\textbf{Agent teams are experimental and disabled by default.} Enable them by setting \texttt{CLAUDE\_\allowbreak{}CODE\_\allowbreak{}EXPERIMENTA\allowbreak{}L\_\allowbreak{}AGENT\_\allowbreak{}TEAMS=\allowbreak{}1} in \texttt{settings.\allowbreak{}json} or the environment. Without it no team is set up, no team directories are written, and Claude does not spawn or propose teammates [29].
\end{calloutbox}

An agent team is a lead session coordinating peer sessions. Teammates work independently, each in its own context window, share a task list, and message each other directly. Unlike subagents, you can interact with an individual teammate without going through the lead [29].


\FloatBarrier
\setcounter{section}{0}
\section{The behaviour change you must understand before enabling}

\begin{calloutbox}{palebrass}{brassdark}{2.0mm}
\textbf{CAUTION:} enabling agent teams changes ordinary delegation. Claude names subagents on its own so it can message them later, and \textbf{while agent teams are enabled a named subagent launches as a teammate} — so teams can form during delegation you never framed as team work [29]. Subagents and teammates report back differently: a subagent's result returns to Claude when it completes, whereas a teammate's idle notification reports only that it stopped, without its output. An orchestration flow that waits on subagent results can stall.
\end{calloutbox}

Set the variable to \texttt{0} to restore subagent behaviour; no restart is needed, because settings-file \texttt{env} values are reapplied to the running session on save [29]. Note the precedence trap: \texttt{0} in your user settings is overridden by \texttt{1} in project settings, local settings, a \texttt{-\allowbreak{}-\allowbreak{}settings} payload, or managed settings.


\FloatBarrier
\setcounter{section}{1}
\section{Architecture}

\begingroup
\def\tblrows{%
Team lead & The main session; spawns teammates and coordinates & The session you are in \\
Teammates & Separate Claude Code instances & Own context windows \\
Team config & Runtime state — session IDs, pane IDs & \texttt{\textasciitilde{}/\allowbreak{}.\allowbreak{}claude/\allowbreak{}teams/\allowbreak{}\{team\}/\allowbreak{}config.\allowbreak{}json} \\
Mailbox & Per-agent message queue & \texttt{\textasciitilde{}/\allowbreak{}.\allowbreak{}claude/\allowbreak{}teams/\allowbreak{}\{team\}/\allowbreak{}inboxes/\allowbreak{}\{agent\}.\allowbreak{}json} \\
Task list & Shared work items & \texttt{\textasciitilde{}/\allowbreak{}.\allowbreak{}claude/\allowbreak{}tasks/\allowbreak{}\{team\}/\allowbreak{}} \\
}%
\def\tblbody{\begin{minipage}{\textwidth}
{\footnotesize\begin{tabular}{L{21.8mm}L{41.1mm}H{76.9mm}}
\toprule
\textbf{Component} & \textbf{Role} & \textbf{Location} \\
\midrule
\tblrows
\bottomrule\end{tabular}}\end{minipage}}%
\begingroup
\def\sloppy{\tolerance 9999\emergencystretch 3em\hfuzz 200pt\vfuzz 200pt}%
\hbadness=10000\vbadness=10000\hfuzz=200pt\vfuzz=200pt
\global\setbox\tblbox=\hbox{\tblbody}%
\endgroup
\par\addvspace{2.6mm}
\ifdim\dimexpr\ht\tblbox+\dp\tblbox\relax>0.55\textheight
  \typeout{HANDBOOK-TABLE broken \the\dimexpr\ht\tblbox+\dp\tblbox\relax}%
  
  {\footnotesize\begin{longtable}{L{21.8mm}L{41.1mm}H{76.9mm}}
  \toprule
\textbf{Component} & \textbf{Role} & \textbf{Location} \\
\midrule\endfirsthead
  \multicolumn{3}{@{}l@{}}{%
  \sffamily\footnotesize\itshape\color{inkgrey}Continued}\\[1.2mm]
  \toprule
\textbf{Component} & \textbf{Role} & \textbf{Location} \\
\midrule\endhead
  \bottomrule\endfoot
  \bottomrule\endlastfoot
  \tblrows
  \end{longtable}}%
\else
  \typeout{HANDBOOK-TABLE atomic \the\dimexpr\ht\tblbox+\dp\tblbox\relax}%
  \noindent\tblbody
\fi
\par\addvspace{2.6mm}
\endgroup

The team name is \texttt{session-\allowbreak{}} plus the first eight characters of the session ID. The config directory is removed when the session ends; the task list persists locally, is never uploaded, and is governed by the same \texttt{cleanup\allowbreak{}Period\allowbreak{}Days} retention as transcripts [29]. Do not hand-edit the config: it is overwritten on the next state update, and there is no project-level equivalent.


\FloatBarrier
\setcounter{section}{2}
\section{Control}

Describe the task and the teammates in natural language. Display mode defaults to \texttt{in-\allowbreak{}process}, where all teammates run in your terminal and the agent panel below the prompt selects among them: arrow keys to select, \texttt{Enter} to open a teammate's transcript and message it, \texttt{Esc} to interrupt it, \texttt{x} to stop it, \texttt{Ctrl+T} to toggle the task list [29]. Split panes require tmux or iTerm2 with the \texttt{it2} CLI and are unsupported in the VS Code integrated terminal, Windows Terminal and Ghostty; set \texttt{teammate\allowbreak{}Mode} or pass \texttt{-\allowbreak{}-\allowbreak{}teammate-\allowbreak{}mode}.


Teammates inherit the lead's model unless your prompt names one or \texttt{CLAUDE\_\allowbreak{}CODE\_\allowbreak{}SUBAGENT\_\allowbreak{}MODEL} is set — \texttt{teammate\allowbreak{}Default\allowbreak{}Model} was removed in 2.1.234 — and inherit the lead's effort level [29]. A teammate can be spawned from an existing subagent definition, honouring its \texttt{tools} allowlist and \texttt{model}, with the definition's body appended to its system prompt; note that the \texttt{skills} and \texttt{mcp\allowbreak{}Servers} frontmatter fields are \textbf{not} applied in that mode [29].


For risky work, require plan approval: the teammate works in read-only plan mode until the lead approves. \textbf{The lead approves autonomously}, so give it criteria in your prompt — "only approve plans that include test coverage", "reject plans that modify the database schema" [29].


\FloatBarrier
\setcounter{section}{3}
\section{Permissions}

\begin{calloutbox}{palebrass}{brassdark}{2.0mm}
\textbf{CAUTION:} teammates start with the lead's permission settings, and per-teammate modes cannot be set at spawn time. If the lead runs with \texttt{-\allowbreak{}-\allowbreak{}dangerously-\allowbreak{}skip-\allowbreak{}permissions}, \textbf{every teammate does too} [29]. Teammate permission prompts surface in the lead session, which becomes a bottleneck at team scale; pre-approve genuinely safe operations before spawning rather than approving under pressure afterwards.
\end{calloutbox}

The inter-agent messaging rules of Section 29.3 apply in full: a teammate cannot approve a prompt or supply consent for you, and the auto-mode classifier reviews every message before delivery, including structured protocol messages such as shutdown requests and plan-approval responses.


\FloatBarrier
\setcounter{section}{4}
\section{Quality gates}

Three hooks are specific to teams and are the mechanism for enforcing standards across peers [29], [73]:


\begin{itemize}
\item \texttt{Teammate\allowbreak{}Idle} — runs when a teammate is about to go idle; exit 2 sends feedback and keeps it working.
\item \texttt{Task\allowbreak{}Created} — exit 2 prevents creation and returns feedback.
\item \texttt{Task\allowbreak{}Completed} — exit 2 prevents completion and returns feedback.
\end{itemize}

A \texttt{Task\allowbreak{}Completed} hook that requires evidence before a task may be closed is the team-scale equivalent of the release gate in Chapter 34.


\FloatBarrier
\setcounter{section}{5}
\section{Limitations to plan around}

From [29]: \texttt{/\allowbreak{}resume} and \texttt{/\allowbreak{}rewind} do not restore in-process teammates, and the lead may try to message teammates that no longer exist; task status can lag, blocking dependents; shutdown waits for the current tool call; there is exactly one team per session and no nested teams; the lead is fixed for the session's lifetime; and an in-process teammate's own subagents run in the foreground, so a definition with \texttt{background:\allowbreak{} true} returns an error.


Start with three to five teammates. Token cost scales linearly with teammates, coordination overhead scales worse, and three focused teammates routinely outperform five scattered ones [29], [59].


\FloatBarrier
\setcounter{section}{6}
\section{When a team beats the alternatives}

Use a team when teammates need to \textbf{share findings and challenge each other}: parallel review with distinct lenses, competing hypotheses in a debugging investigation, cross-layer work with clear ownership. Use subagents when only the result matters. Use a workflow when the orchestration itself should be repeatable (Chapter 22). Use worktrees when the issue is file isolation rather than reasoning (Section 29.2).


The adversarial pattern is the one that earns its cost:


\begin{codeblock}{7.0}{8.3}
\cl{Users~report~the~app~exits~after~one~message~instead~of~staying~connected.~Spawn~five~teammates}
\cl{to~investigate~different~hypotheses.~Have~them~message~each~other~to~try~to~disprove~each}
\cl{other\textquotesingle{}s~theories,~like~a~scientific~debate.~Update~the~findings~document~with~the~consensus}
\cl{that~emerges,~and~record~which~theories~were~eliminated~and~by~what~evidence.}
\end{codeblock}

Sequential investigation anchors on the first plausible explanation. Independent investigators actively trying to falsify each other produce a surviving theory that is much more likely to be the actual cause — the same epistemics that make the contradiction reviewer load-bearing in Section 22.5.



\begingroup
\def\kprows{%
\item Agent teams are experimental and off by default; enabling them changes ordinary delegation.
\item While enabled, a subagent Claude names launches as a teammate, so teams form unasked and results report back differently.
\item Teammates inherit the lead's permission settings and cannot be given their own at spawn.
\item Three to five teammates; the adversarial pattern — teammates trying to falsify each other — is what earns the cost.
}%
\def\kpbody{\begin{minipage}{\textwidth}\subsection*{Key points}\begin{itemize}\kprows\end{itemize}\end{minipage}}%
\begingroup
\def\sloppy{\tolerance 9999\emergencystretch 3em\hfuzz 200pt\vfuzz 200pt}%
\hbadness=10000\vbadness=10000\hfuzz=200pt\vfuzz=200pt
\global\setbox\kpbox=\hbox{\kpbody}%
\endgroup
\par\addvspace{4.2mm}
\ifdim\dimexpr\ht\kpbox+\dp\kpbox\relax>0.30\textheight
  \typeout{HANDBOOK-KEYPOINTS broken \the\dimexpr\ht\kpbox+\dp\kpbox\relax}%
  \subsection*{Key points}
  \begin{itemize}\kprows\end{itemize}
\else
  \typeout{HANDBOOK-KEYPOINTS atomic \the\dimexpr\ht\kpbox+\dp\kpbox\relax}%
  \noindent\kpbody
\fi
\par\addvspace{1.4mm}
\endgroup

\breakrule

\FloatBarrier
\parttitle{Part VIII}{Governance, assurance and compliance}
\addcontentsline{toc}{part}{Part VIII \textemdash\ Governance, assurance and compliance}

This part is written for readers who must answer to somebody else: an audit committee, a regulator, a customer's security questionnaire, or a board. It assumes the operating discipline of Parts I–III and asks a different question: what can be \emph{enforced}, what can be \emph{evidenced}, and what remains a matter of trust.


\FloatBarrier
\renewcommand{\chaptertitlelabel}{Part VIII \textperiodcentered\ Chapter 31}
\setcounter{chapter}{30}
\chapter{Enterprise configuration and managed policy}

Everything in Chapters 7, 13, 15, 17 and 18 is a control a \emph{user} can set and therefore a control a user can unset. Managed settings are the tier above: Claude Code applies them above every other level, and no user, project, local or \texttt{-\allowbreak{}-\allowbreak{}settings} value overrides them, apart from a small set of exceptions where a \textbf{stricter} value from a lower level still counts [63], [64].


\FloatBarrier
\setcounter{section}{0}
\section{Delivery mechanisms}

\begingroup
\def\tblrows{%
Server-managed & claude.ai admin console, or a self-hosted Claude apps gateway & Fetched at startup, polled hourly & One place to change policy for a claude.ai organisation \\
MDM or OS policy & macOS configuration profile (\texttt{com.\allowbreak{}anthropic.\allowbreak{}claudecode}) or Windows \texttt{HKLM\textbackslash{}SOFTWARE\textbackslash{}Policies\textbackslash{}Claude\allowbreak{}Code} & Startup, re-checked every 30 minutes & You already manage devices with Jamf, Intune or Group Policy \\
File-based & \texttt{managed-\allowbreak{}settings.\allowbreak{}json} in a system directory & Startup, reloaded on change & Machines without MDM, Linux hosts, images you build \\
\texttt{HKCU} registry & \texttt{HKCU\textbackslash{}SOFTWARE\textbackslash{}Policies\textbackslash{}Claude\allowbreak{}Code} & Startup, re-checked every 30 minutes & You cannot write the machine-level key \\
}%
\def\tblbody{\begin{minipage}{\textwidth}\boxcaption{Table 22 — How a managed policy reaches a machine}
{\scriptsize\begin{tabular}{L{26.6mm}L{62.7mm}L{20.8mm}L{26.5mm}}
\toprule
\textbf{Mechanism} & \textbf{Delivered as} & \textbf{Read when} & \textbf{Use when} \\
\midrule
\tblrows
\bottomrule\end{tabular}}\end{minipage}}%
\begingroup
\def\sloppy{\tolerance 9999\emergencystretch 3em\hfuzz 200pt\vfuzz 200pt}%
\hbadness=10000\vbadness=10000\hfuzz=200pt\vfuzz=200pt
\global\setbox\tblbox=\hbox{\tblbody}%
\endgroup
\par\addvspace{2.6mm}
\ifdim\dimexpr\ht\tblbox+\dp\tblbox\relax>0.55\textheight
  \typeout{HANDBOOK-TABLE broken \the\dimexpr\ht\tblbox+\dp\tblbox\relax}%
  \tabcaption{Table 22 — How a managed policy reaches a machine}
  {\scriptsize\begin{longtable}{L{26.6mm}L{62.7mm}L{20.8mm}L{26.5mm}}
  \toprule
\textbf{Mechanism} & \textbf{Delivered as} & \textbf{Read when} & \textbf{Use when} \\
\midrule\endfirsthead
  \multicolumn{4}{@{}l@{}}{%
  \sffamily\footnotesize\itshape\color{inkgrey}Table 22 — How a managed policy reaches a machine \textemdash\ continued}\\[1.2mm]
  \toprule
\textbf{Mechanism} & \textbf{Delivered as} & \textbf{Read when} & \textbf{Use when} \\
\midrule\endhead
  \bottomrule\endfoot
  \bottomrule\endlastfoot
  \tblrows
  \end{longtable}}%
\else
  \typeout{HANDBOOK-TABLE atomic \the\dimexpr\ht\tblbox+\dp\tblbox\relax}%
  \noindent\tblbody
\fi
\par\addvspace{2.6mm}
\endgroup

Source: [64].


File paths: \texttt{/\allowbreak{}Library/\allowbreak{}Application\allowbreak{} Support/\allowbreak{}Claude\allowbreak{}Code/\allowbreak{}} on macOS, \texttt{/\allowbreak{}etc/\allowbreak{}claude-\allowbreak{}code/\allowbreak{}} on Linux and WSL, \texttt{C:\allowbreak{}\textbackslash{}Program Fil\allowbreak{}es\textbackslash{}Claude\allowbreak{}Code\textbackslash{}} on Windows. The legacy Windows path \texttt{C:\allowbreak{}\textbackslash{}Program\allowbreak{}Data\textbackslash{}Claude\allowbreak{}Code\textbackslash{}managed-\allowbreak{}settings.\allowbreak{}json} is \textbf{not} read [64]. A \texttt{managed-\allowbreak{}settings.\allowbreak{}d/\allowbreak{}} directory alongside the main file lets several teams own parts of one policy; files merge alphabetically, single values replace, lists combine with duplicates removed, and nested blocks merge key by key.


\begin{calloutbox}{palebrass}{brassdark}{2.0mm}
\textbf{CAUTION — the precedence rule that surprises administrators.} When more than one managed source delivers a policy, Claude Code uses \textbf{the first source that delivers at least one policy key and ignores the rest} — it does not merge them, and it shows no warning for the sources it skipped [64]. The order is remote settings, then MDM or OS policy, then managed settings files, then \texttt{HKCU}. A small number of cross-source keys are read from every admin source, including the sandbox locks, \texttt{allow\allowbreak{}All\allowbreak{}Claude\allowbreak{}Ai\allowbreak{}Mcps}, \texttt{force\allowbreak{}Remote\allowbreak{}Settings\allowbreak{}Refresh}, and — from 2.1.223 — \texttt{env}, merged per variable. Run \texttt{/\allowbreak{}status} on the machine to see which source was actually selected.
\end{calloutbox}

Cloud sessions do not read a device's MDM profile or file at all: policy for an Anthropic-hosted session must come from server-managed settings [64].


\FloatBarrier
\setcounter{section}{1}
\section{A minimum viable enterprise policy}

\begin{codeblock}{9.0}{10.6}
\cl{\{}
\cl{~~"permissions":~\{}
\cl{~~~~"deny":~[}
\cl{~~~~~~"Read(./.env)",}
\cl{~~~~~~"Read(./.env.*)",}
\cl{~~~~~~"Read(./secrets/**)",}
\cl{~~~~~~"Read(\textasciitilde{}/.ssh/**)",}
\cl{~~~~~~"Read(\textasciitilde{}/.aws/**)"}
\cl{~~~~],}
\cl{~~~~"disableBypassPermissionsMode":~"disable"}
\cl{~~\},}
\cl{~~"allowManagedPermissionRulesOnly":~true,}
\cl{~~"allowManagedHooksOnly":~true,}
\cl{~~"allowManagedMcpServersOnly":~true,}
\cl{~~"sandbox":~\{}
\cl{~~~~"enabled":~true,}
\cl{~~~~"network":~\{~"allowManagedDomainsOnly":~true~\},}
\cl{~~~~"credentials":~\{}
\cl{~~~~~~"files":~~~[\{~"path":~"\textasciitilde{}/.aws/credentials",~"mode":~"deny"~\}],}
\cl{~~~~~~"envVars":~[\{~"name":~"GITHUB\_TOKEN",~"mode":~"deny"~\}]}
\cl{~~~~\}}
\cl{~~\},}
\cl{~~"strictKnownMarketplaces":~true,}
\cl{~~"disableSideloadFlags":~true,}
\cl{~~"channelsEnabled":~false,}
\cl{~~"forceLoginMethod":~"claudeai",}
\cl{~~"forceLoginOrgUUID":~"<your-organization-uuid>",}
\cl{~~"cleanupPeriodDays":~7,}
\cl{~~"requiredMinimumVersion":~"2.1.241"}
\cl{\}}
\end{codeblock}

Each key answers a threat from an earlier chapter. \texttt{allow\allowbreak{}Managed\allowbreak{}Permission\allowbreak{}Rules\allowbreak{}Only} closes the gap that a developer can widen their own permission rules (Chapter 7). \texttt{allow\allowbreak{}Managed\allowbreak{}Hooks\allowbreak{}Only} blocks user, project, local and plugin hooks — and, as noted in Chapter 21, disables \texttt{/\allowbreak{}goal} as a side effect, which is a trade-off to make consciously. \texttt{disable\allowbreak{}Sideload\allowbreak{}Flags} rejects \texttt{-\allowbreak{}-\allowbreak{}plugin-\allowbreak{}dir}, \texttt{-\allowbreak{}-\allowbreak{}plugin-\allowbreak{}url}, \texttt{-\allowbreak{}-\allowbreak{}agents} and \texttt{-\allowbreak{}-\allowbreak{}mcp-\allowbreak{}config} at startup, closing the local-testing path of Chapter 18 (2.1.193+). \texttt{strict\allowbreak{}Known\allowbreak{}Marketplace\allowbreak{}s} and \texttt{blocked\allowbreak{}Marketplace\allowbreak{}s} control which plugin sources exist at all, with blocked sources checked before download so they never touch the filesystem. \texttt{force\allowbreak{}Login\allowbreak{}Method} and \texttt{force\allowbreak{}Login\allowbreak{}Org\allowbreak{}UUID} are what route traffic into a ZDR organisation (Chapter 32) [52], [64].


Other managed-only keys worth knowing: \texttt{strict\allowbreak{}Plugin\allowbreak{}Only\allowbreak{}Customizati\allowbreak{}on} blocks skills, agents, hooks and MCP servers from user and project sources; \texttt{plugin\allowbreak{}Trust\allowbreak{}Message} appends your own wording to the plugin trust warning; \texttt{allowed\allowbreak{}Channel\allowbreak{}Plugins} replaces the Anthropic channel allowlist; \texttt{available\allowbreak{}Models} and \texttt{enforce\allowbreak{}Available\allowbreak{}Models} restrict model choice — necessary because \textbf{a managed \texttt{model} value is a default, not a lock}, and \texttt{-\allowbreak{}-\allowbreak{}model} or \texttt{ANTHROPIC\_\allowbreak{}MODEL} still picks the model for a session [64].


\FloatBarrier
\setcounter{section}{2}
\section{Validation behaviour, and which keys fail closed}

When a managed payload fails schema validation, Claude Code repairs what it can, drops top-level keys that still fail, and keeps enforcing the rest [64]. A handful of enforcement keys deliberately \textbf{fail closed} rather than being dropped:


\begin{itemize}
\item \texttt{allowed\allowbreak{}Mcp\allowbreak{}Servers} → enforced as an empty allowlist, admitting no servers;
\item \texttt{allow\allowbreak{}Managed\allowbreak{}Hooks\allowbreak{}Only}, \texttt{allow\allowbreak{}Managed\allowbreak{}Mcp\allowbreak{}Servers\allowbreak{}Only}, \texttt{disable\allowbreak{}Command\allowbreak{}Plugin\allowbreak{}Sources}, \texttt{enforce\allowbreak{}Available\allowbreak{}Models} → treated as \texttt{true};
\item \texttt{available\allowbreak{}Models} → empty allowlist, leaving only the default model;
\item \texttt{force\allowbreak{}Login\allowbreak{}Org\allowbreak{}UUID} → no organisation may log in;
\item \texttt{sandbox.\allowbreak{}credentials} → a recoverable invalid entry degrades to \texttt{mode:\allowbreak{} "deny"}.
\end{itemize}

\texttt{required\allowbreak{}Minimum\allowbreak{}Version} and \texttt{required\allowbreak{}Maximum\allowbreak{}Version} fail \textbf{open} by design, so a bad policy push cannot prevent Claude Code from starting [64]. Find dropped entries in the startup dialog, on stderr under \texttt{-\allowbreak{}p}, or with \texttt{claude doct\allowbreak{}or}.


\FloatBarrier
\setcounter{section}{3}
\section{What a managed policy does not reach}

Be explicit with stakeholders about the boundary [64]:


\begin{itemize}
\item \textbf{A local administrator can edit the managed source itself.} MDM redeployment on a schedule is the mitigation, not an assumption of integrity.
\item \textbf{The server-managed cache can be edited locally}, and the edit lasts until the next successful fetch. \texttt{force\allowbreak{}Remote\allowbreak{}Settings\allowbreak{}Refresh} blocks startup until a fresh fetch succeeds.
\item \textbf{Managed settings bind Claude Code only.} A developer calling the API from another tool is not under them.
\item \textbf{Model selection for a session} escapes a managed \texttt{model} default unless \texttt{available\allowbreak{}Models} is deployed.
\end{itemize}

Two organisation-level toggles live outside settings files entirely: Remote Control and web sessions are enabled or disabled organisation-wide in the claude.ai Claude Code admin settings, with \texttt{disable\allowbreak{}Remote\allowbreak{}Control} available per device and no per-device key for web sessions [64].


\begingroup
\def\kprows{%
\item Managed settings are the only tier a user cannot unset, and no other level overrides them.
\item When several managed sources exist, the first with a policy key wins and the rest are ignored with no warning. Check \texttt{/\allowbreak{}status}.
\item A managed \texttt{model} is a default, not a lock; deploy \texttt{available\allowbreak{}Models} to restrict choice.
\item Several enforcement keys deliberately fail closed when invalid, and version floors deliberately fail open.
}%
\def\kpbody{\begin{minipage}{\textwidth}\subsection*{Key points}\begin{itemize}\kprows\end{itemize}\end{minipage}}%
\begingroup
\def\sloppy{\tolerance 9999\emergencystretch 3em\hfuzz 200pt\vfuzz 200pt}%
\hbadness=10000\vbadness=10000\hfuzz=200pt\vfuzz=200pt
\global\setbox\kpbox=\hbox{\kpbody}%
\endgroup
\par\addvspace{4.2mm}
\ifdim\dimexpr\ht\kpbox+\dp\kpbox\relax>0.30\textheight
  \typeout{HANDBOOK-KEYPOINTS broken \the\dimexpr\ht\kpbox+\dp\kpbox\relax}%
  \subsection*{Key points}
  \begin{itemize}\kprows\end{itemize}
\else
  \typeout{HANDBOOK-KEYPOINTS atomic \the\dimexpr\ht\kpbox+\dp\kpbox\relax}%
  \noindent\kpbody
\fi
\par\addvspace{1.4mm}
\endgroup

\FloatBarrier
\renewcommand{\chaptertitlelabel}{Part VIII \textperiodcentered\ Chapter 32}
\setcounter{chapter}{31}
\chapter{Data governance, residency and retention}

This chapter states what leaves the machine, what is retained and for how long, what an organisation can turn off, and where the boundaries of Zero Data Retention actually fall. For a reader operating across Australia, the European Union and Asia-Pacific, the mapping to obligations under the Privacy Act, the GDPR and the EU AI Act is a matter of applying these mechanics to your own processing purposes [17], [90], [91]; this chapter supplies the mechanics, not the legal conclusion.


\FloatBarrier
\setcounter{section}{0}
\section{What leaves the machine}

Claude Code runs locally and sends prompts and model outputs over TLS 1.2 or later to the configured provider [41]. Beyond inference, four other flows exist and each has its own control:


\begingroup
\def\tblrows{%
Metrics & Latency, reliability, usage patterns. \textbf{Never} code, prompts or file paths & \texttt{DISABLE\_\allowbreak{}TELEMETRY=\allowbreak{}1} \\
Error reports & Error messages and stack traces from Claude Code internals, with known secret patterns, file paths, email addresses and personal information redacted before leaving the machine & \texttt{DISABLE\_\allowbreak{}ERROR\_\allowbreak{}REPORTING=\allowbreak{}1} \\
\texttt{/\allowbreak{}feedback}, \texttt{/\allowbreak{}bug}, \texttt{/\allowbreak{}share} & A copy of the conversation \textbf{including code}; you choose current session, or the project's last 24 hours or 7 days & \texttt{DISABLE\_\allowbreak{}FEEDBACK\_\allowbreak{}COMMAND=\allowbreak{}1} \\
Session quality survey & The rating only. A separate opt-in follow-up uploads the transcript, subagent transcripts and raw session log, with known API-key and token patterns redacted & \texttt{CLAUDE\_\allowbreak{}CODE\_\allowbreak{}DISABLE\_\allowbreak{}FEEDBACK\_\allowbreak{}SURVEY=\allowbreak{}1}, or \texttt{feedback\allowbreak{}Survey\allowbreak{}Rate} \\
WebFetch domain safety check & The \textbf{hostname only}, to \texttt{api.\allowbreak{}anthropic.\allowbreak{}com}, checked against a blocklist; cached five minutes & \texttt{skip\allowbreak{}Web\allowbreak{}Fetch\allowbreak{}Preflight:\allowbreak{} true} \\
}%
\def\tblbody{\begin{minipage}{\textwidth}\boxcaption{Table 23 — Non-inference data flows and their opt-outs}
{\footnotesize\begin{tabular}{L{20.5mm}L{42.4mm}H{76.9mm}}
\toprule
\textbf{Flow} & \textbf{Contents} & \textbf{Opt-out} \\
\midrule
\tblrows
\bottomrule\end{tabular}}\end{minipage}}%
\begingroup
\def\sloppy{\tolerance 9999\emergencystretch 3em\hfuzz 200pt\vfuzz 200pt}%
\hbadness=10000\vbadness=10000\hfuzz=200pt\vfuzz=200pt
\global\setbox\tblbox=\hbox{\tblbody}%
\endgroup
\par\addvspace{2.6mm}
\ifdim\dimexpr\ht\tblbox+\dp\tblbox\relax>0.55\textheight
  \typeout{HANDBOOK-TABLE broken \the\dimexpr\ht\tblbox+\dp\tblbox\relax}%
  \tabcaption{Table 23 — Non-inference data flows and their opt-outs}
  {\footnotesize\begin{longtable}{L{20.5mm}L{42.4mm}H{76.9mm}}
  \toprule
\textbf{Flow} & \textbf{Contents} & \textbf{Opt-out} \\
\midrule\endfirsthead
  \multicolumn{3}{@{}l@{}}{%
  \sffamily\footnotesize\itshape\color{inkgrey}Table 23 — Non-inference data flows and their opt-outs \textemdash\ continued}\\[1.2mm]
  \toprule
\textbf{Flow} & \textbf{Contents} & \textbf{Opt-out} \\
\midrule\endhead
  \bottomrule\endfoot
  \bottomrule\endlastfoot
  \tblrows
  \end{longtable}}%
\else
  \typeout{HANDBOOK-TABLE atomic \the\dimexpr\ht\tblbox+\dp\tblbox\relax}%
  \noindent\tblbody
\fi
\par\addvspace{2.6mm}
\endgroup

Source: [41].


\texttt{CLAUDE\_\allowbreak{}CODE\_\allowbreak{}DISABLE\_\allowbreak{}NONESSENTIA\allowbreak{}L\_\allowbreak{}TRAFFIC} disables all of the above at once \textbf{except} the WebFetch check and official marketplace auto-install, each of which has its own opt-out [41]. Note the coupling: setting \texttt{DISABLE\_\allowbreak{}TELEMETRY} or \texttt{CLAUDE\_\allowbreak{}CODE\_\allowbreak{}DISABLE\_\allowbreak{}NONESSENTIA\allowbreak{}L\_\allowbreak{}TRAFFIC} also disables the feature-flag evaluation that Remote Control depends on, and turns off \texttt{/\allowbreak{}schedule} [20], [41].


\begin{calloutbox}{palebrass}{brassdark}{2.0mm}
\textbf{CAUTION:} if you disable the WebFetch preflight, combine it with \texttt{Web\allowbreak{}Fetch(domai\allowbreak{}n:\allowbreak{}…)} permission rules. Disabling the check means WebFetch will attempt any URL without consulting the blocklist [41].
\end{calloutbox}

\FloatBarrier
\setcounter{section}{1}
\section{Retention}

\begingroup
\def\tblrows{%
Free, Pro, Max — data use for model improvement \textbf{on} & 5 years \\
Free, Pro, Max — data use \textbf{off} & 30 days \\
Team, Enterprise, API — standard & 30 days \\
Claude for Enterprise with ZDR & Not stored after the response is returned, except as required by law or to address misuse \\
Transcripts shared via \texttt{/\allowbreak{}feedback}, \texttt{/\allowbreak{}bug}, \texttt{/\allowbreak{}share} & 5 years \\
Transcripts shared through the session-quality follow-up & Up to 6 months \\
Sessions flagged for a policy violation, even under ZDR & Up to 2 years \\
}%
\def\tblbody{\begin{minipage}{\textwidth}\boxcaption{Table 24 — Retention by account type}
{\footnotesize\begin{tabular}{L{92.8mm}L{50.2mm}}
\toprule
\textbf{Account} & \textbf{Server-side retention} \\
\midrule
\tblrows
\bottomrule\end{tabular}}\end{minipage}}%
\begingroup
\def\sloppy{\tolerance 9999\emergencystretch 3em\hfuzz 200pt\vfuzz 200pt}%
\hbadness=10000\vbadness=10000\hfuzz=200pt\vfuzz=200pt
\global\setbox\tblbox=\hbox{\tblbody}%
\endgroup
\par\addvspace{2.6mm}
\ifdim\dimexpr\ht\tblbox+\dp\tblbox\relax>0.55\textheight
  \typeout{HANDBOOK-TABLE broken \the\dimexpr\ht\tblbox+\dp\tblbox\relax}%
  \tabcaption{Table 24 — Retention by account type}
  {\footnotesize\begin{longtable}{L{92.8mm}L{50.2mm}}
  \toprule
\textbf{Account} & \textbf{Server-side retention} \\
\midrule\endfirsthead
  \multicolumn{2}{@{}l@{}}{%
  \sffamily\footnotesize\itshape\color{inkgrey}Table 24 — Retention by account type \textemdash\ continued}\\[1.2mm]
  \toprule
\textbf{Account} & \textbf{Server-side retention} \\
\midrule\endhead
  \bottomrule\endfoot
  \bottomrule\endlastfoot
  \tblrows
  \end{longtable}}%
\else
  \typeout{HANDBOOK-TABLE atomic \the\dimexpr\ht\tblbox+\dp\tblbox\relax}%
  \noindent\tblbody
\fi
\par\addvspace{2.6mm}
\endgroup

Source: [40], [41].


\textbf{Training.} Anthropic does not train generative models on code or prompts sent to Claude Code under commercial terms unless the customer opts in, for example through the Development Partner Program. Consumer accounts have a choice, and the setting governs Claude Code usage from those accounts too [41].


\textbf{Local caching} is separate from all of the above: transcripts sit in plaintext under \texttt{\textasciitilde{}/\allowbreak{}.\allowbreak{}claude/\allowbreak{}projects/\allowbreak{}} for 30 days by default, adjustable with \texttt{cleanup\allowbreak{}Period\allowbreak{}Days} [41], [92]. Remember that auto memory is excluded from that sweep (Section 10.2) and that agent-team task lists are included in it (Section 30.2).


\textbf{Encryption at rest depends on the provider} [41]: infrastructure-level AES-256 on the Anthropic API; AES-256 with AWS-managed or customer-managed KMS keys on Amazon Bedrock; Google-managed keys with CMEK available on Google Cloud's Agent Platform; and on Microsoft Foundry it depends on the hosting option — for \emph{Hosted on Azure}, prompts and completions remain within Azure and only usage metadata and safety-flagged content egress to Anthropic.


\FloatBarrier
\setcounter{section}{2}
\section{Zero Data Retention, and its exact boundary}

ZDR for Claude Code is available to qualified accounts on Claude for Enterprise. It is \textbf{not} part of the standard Enterprise plan, cannot be enabled from admin settings, and is enabled per organisation by Anthropic after eligibility review — including for each new organisation created under the same account [40].


\textbf{Covered:} model inference calls made through Claude Code on Claude for Enterprise.


\textbf{Not covered, even for a ZDR organisation} [40]: chat on claude.ai; Cowork sessions; Claude Code Analytics, which stores no prompts or responses but does collect productivity metadata such as account emails and usage statistics; user and seat management data; and anything processed by third-party tools, MCP servers or external integrations.


\textbf{Disabled under ZDR}, at the backend and regardless of what the client displays [40]:


\begin{itemize}
\item Claude Code on the web;
\item cloud sessions from the Desktop app;
\item artifacts;
\item feedback submission (\texttt{/\allowbreak{}feedback}, \texttt{/\allowbreak{}bug}, \texttt{/\allowbreak{}share});
\item \textbf{Remote Control}, because it stores the transcript on Anthropic servers to sync across devices.
\end{itemize}

Claude Fable 5 is unavailable under ZDR because that model class requires data retention; the \texttt{best} alias resolves to Opus for ZDR organisations instead [40].


\begin{calloutbox}{palebrass}{brassdark}{2.0mm}
\textbf{CAUTION — the routing gap.} ZDR applies to requests that authenticate \textbf{into} the ZDR-enabled organisation. A developer signing in with a personal account, or with an API key from a different organisation, is not covered. Deploy \texttt{force\allowbreak{}Login\allowbreak{}Method} and \texttt{force\allowbreak{}Login\allowbreak{}Org\allowbreak{}UUID} through managed settings to close it [40]. An organisation that believes it has ZDR but has not deployed those keys has a policy, not a control.
\end{calloutbox}

\FloatBarrier
\setcounter{section}{3}
\section{Cloud execution}

In an Anthropic-hosted cloud session the repository is cloned to an isolated VM; GitHub authentication goes through a secure proxy so credentials never enter the sandbox; \textbf{all outbound traffic passes through a security proxy for audit logging}; and session data follows the retention policy for your account type [41], [53]. A self-hosted environment moves isolation, egress control and git credentials to your own deployment.


Remote Control follows the \emph{local} data flow, because execution is local — with the qualification that while connected the transcript is stored on Anthropic servers to sync across devices [39], [41]. That single sentence is why Remote Control is disabled under ZDR, and it is the fact most often missed when a team assumes "execution is local" means "nothing leaves".


\FloatBarrier
\setcounter{section}{4}
\section{A data-governance questionnaire}

Answer these before an organisational rollout, and record the answers:


\begin{enumerate}
\item Which plan and provider is each cohort of developers on, and therefore which retention row of Table 24 applies?
\item Is training opt-out contractual (commercial terms) or a user setting (consumer)?
\item Is \texttt{force\allowbreak{}Login\allowbreak{}Org\allowbreak{}UUID} deployed, so that sessions actually authenticate into the intended organisation?
\item Which of the flows in Table 23 are disabled, and by which delivery mechanism?
\item What is \texttt{cleanup\allowbreak{}Period\allowbreak{}Days}, and does anyone rely on the sweep covering auto memory? (It does not.)
\item Are artifacts, channels, Remote Control and web sessions enabled, and who approved each?
\item Where do published artifacts live, what is their retention, and who can share them publicly?
\item Which MCP servers and connectors are permitted, and what data can each reach?
\item Is \texttt{claude\_\allowbreak{}artifact\_\allowbreak{}*} and Claude Code audit-log data being collected and reviewed?
\item What is the incident procedure when a credential reaches a transcript (Section 14.3)?
\end{enumerate}

\begingroup
\def\kprows{%
\item Five non-inference flows leave the machine, each with its own opt-out; one variable disables all but two of them.
\item Retention differs by plan, by feature and by consent, and local caching is separate from all of it.
\item ZDR covers inference on Claude for Enterprise and disables web sessions, cloud sessions, artifacts, feedback and Remote Control.
\item ZDR applies to requests authenticating into the ZDR organisation — without \texttt{force\allowbreak{}Login\allowbreak{}Org\allowbreak{}UUID} you have a policy, not a control.
}%
\def\kpbody{\begin{minipage}{\textwidth}\subsection*{Key points}\begin{itemize}\kprows\end{itemize}\end{minipage}}%
\begingroup
\def\sloppy{\tolerance 9999\emergencystretch 3em\hfuzz 200pt\vfuzz 200pt}%
\hbadness=10000\vbadness=10000\hfuzz=200pt\vfuzz=200pt
\global\setbox\kpbox=\hbox{\kpbody}%
\endgroup
\par\addvspace{4.2mm}
\ifdim\dimexpr\ht\kpbox+\dp\kpbox\relax>0.30\textheight
  \typeout{HANDBOOK-KEYPOINTS broken \the\dimexpr\ht\kpbox+\dp\kpbox\relax}%
  \subsection*{Key points}
  \begin{itemize}\kprows\end{itemize}
\else
  \typeout{HANDBOOK-KEYPOINTS atomic \the\dimexpr\ht\kpbox+\dp\kpbox\relax}%
  \noindent\kpbody
\fi
\par\addvspace{1.4mm}
\endgroup

\FloatBarrier
\renewcommand{\chaptertitlelabel}{Part VIII \textperiodcentered\ Chapter 33}
\setcounter{chapter}{32}
\chapter{Observability, cost control and accessibility}

\FloatBarrier
\setcounter{section}{0}
\section{OpenTelemetry}

Claude Code exports metrics, logs and — in beta — traces through OpenTelemetry [93]. This is the mechanism that turns "we use an AI agent" into something an assurance function can examine.


\begin{codeblock}{9.0}{10.6}
\cl{export~CLAUDE\_CODE\_ENABLE\_TELEMETRY=1}
\cl{export~OTEL\_METRICS\_EXPORTER=otlp~~~~~~~~~~\#~otlp~|~prometheus~|~console~|~none}
\cl{export~OTEL\_LOGS\_EXPORTER=otlp~~~~~~~~~~~~~\#~otlp~|~console~|~none}
\cl{export~OTEL\_EXPORTER\_OTLP\_PROTOCOL=grpc~~~~\#~grpc~|~http/json~|~http/protobuf}
\cl{export~OTEL\_EXPORTER\_OTLP\_ENDPOINT=http://collector.example.com:4317}
\cl{export~OTEL\_EXPORTER\_OTLP\_HEADERS="Authorization=Bearer~<token>"}
\end{codeblock}

Deploy it fleet-wide through the \texttt{env} block of managed settings (Chapter 31).


\begingroup
\def\tblrows{%
\texttt{claude\_\allowbreak{}code.\allowbreak{}session.\allowbreak{}count} & CLI sessions started & — \\
\texttt{claude\_\allowbreak{}code.\allowbreak{}lines\_\allowbreak{}of\_\allowbreak{}code.\allowbreak{}count} & Lines of code modified & — \\
\texttt{claude\_\allowbreak{}code.\allowbreak{}pull\_\allowbreak{}request.\allowbreak{}count} & Pull requests created & — \\
\texttt{claude\_\allowbreak{}code.\allowbreak{}commit.\allowbreak{}count} & Commits created & — \\
\texttt{claude\_\allowbreak{}code.\allowbreak{}cost.\allowbreak{}usage} & Session cost & USD \\
\texttt{claude\_\allowbreak{}code.\allowbreak{}token.\allowbreak{}usage} & Tokens used & tokens \\
\texttt{claude\_\allowbreak{}code.\allowbreak{}code\_\allowbreak{}edit\_\allowbreak{}tool.\allowbreak{}decision} & Code-edit permission decisions & — \\
\texttt{claude\_\allowbreak{}code.\allowbreak{}active\_\allowbreak{}time.\allowbreak{}total} & Active time & s \\
}%
\def\tblbody{\begin{minipage}{\textwidth}\boxcaption{Table 25 — Exported metrics}
{\footnotesize\begin{tabular}{L{76.7mm}L{50.8mm}L{12.3mm}}
\toprule
\textbf{Metric} & \textbf{Meaning} & \textbf{Unit} \\
\midrule
\tblrows
\bottomrule\end{tabular}}\end{minipage}}%
\begingroup
\def\sloppy{\tolerance 9999\emergencystretch 3em\hfuzz 200pt\vfuzz 200pt}%
\hbadness=10000\vbadness=10000\hfuzz=200pt\vfuzz=200pt
\global\setbox\tblbox=\hbox{\tblbody}%
\endgroup
\par\addvspace{2.6mm}
\ifdim\dimexpr\ht\tblbox+\dp\tblbox\relax>0.55\textheight
  \typeout{HANDBOOK-TABLE broken \the\dimexpr\ht\tblbox+\dp\tblbox\relax}%
  \tabcaption{Table 25 — Exported metrics}
  {\footnotesize\begin{longtable}{L{76.7mm}L{50.8mm}L{12.3mm}}
  \toprule
\textbf{Metric} & \textbf{Meaning} & \textbf{Unit} \\
\midrule\endfirsthead
  \multicolumn{3}{@{}l@{}}{%
  \sffamily\footnotesize\itshape\color{inkgrey}Table 25 — Exported metrics \textemdash\ continued}\\[1.2mm]
  \toprule
\textbf{Metric} & \textbf{Meaning} & \textbf{Unit} \\
\midrule\endhead
  \bottomrule\endfoot
  \bottomrule\endlastfoot
  \tblrows
  \end{longtable}}%
\else
  \typeout{HANDBOOK-TABLE atomic \the\dimexpr\ht\tblbox+\dp\tblbox\relax}%
  \noindent\tblbody
\fi
\par\addvspace{2.6mm}
\endgroup

Events carry the audit signal. \texttt{claude\_\allowbreak{}code.\allowbreak{}tool\_\allowbreak{}decision} records a permission decision with its \texttt{source}; \texttt{claude\_\allowbreak{}code.\allowbreak{}permission\_\allowbreak{}mode\_\allowbreak{}changed} records \texttt{from\_\allowbreak{}mode}, \texttt{to\_\allowbreak{}mode} and \texttt{trigger}; \texttt{claude\_\allowbreak{}code.\allowbreak{}mcp\_\allowbreak{}server\_\allowbreak{}connection} records connection status and transport; \texttt{claude\_\allowbreak{}code.\allowbreak{}plugin\_\allowbreak{}installed} and \texttt{claude\_\allowbreak{}code.\allowbreak{}plugin\_\allowbreak{}loaded} record plugin name and marketplace; \texttt{claude\_\allowbreak{}code.\allowbreak{}auth} records login and logout; \texttt{claude\_\allowbreak{}code.\allowbreak{}api\_\allowbreak{}refusal} records model refusals [93]. All events share \texttt{prompt.\allowbreak{}id}, so filtering by it traces every action triggered by a single prompt — which is precisely the query an incident review needs.


\begin{calloutbox}{palebrass}{brassdark}{2.0mm}
\textbf{CAUTION — content logging is off by default and should usually stay off.} \texttt{OTEL\_\allowbreak{}LOG\_\allowbreak{}USER\_\allowbreak{}PROMPTS}, \texttt{OTEL\_\allowbreak{}LOG\_\allowbreak{}ASSISTANT\_\allowbreak{}RESPONSES}, \texttt{OTEL\_\allowbreak{}LOG\_\allowbreak{}TOOL\_\allowbreak{}DETAILS}, \texttt{OTEL\_\allowbreak{}LOG\_\allowbreak{}TOOL\_\allowbreak{}CONTENT} and \texttt{OTEL\_\allowbreak{}LOG\_\allowbreak{}RAW\_\allowbreak{}API\_\allowbreak{}BODIES} are all disabled by default [93]. Enabling any of them moves prompt and code content into your telemetry pipeline, which changes that pipeline's classification, retention obligations and access controls. Decide that deliberately, with your privacy function, not as a debugging convenience.
\end{calloutbox}

Cardinality and identity controls matter for both cost and privacy: \texttt{OTEL\_\allowbreak{}METRICS\_\allowbreak{}INCLUDE\_\allowbreak{}SESSION\_\allowbreak{}ID} (default \texttt{true}), \texttt{OTEL\_\allowbreak{}METRICS\_\allowbreak{}INCLUDE\_\allowbreak{}ACCOUNT\_\allowbreak{}UUID} (default \texttt{true}), \texttt{OTEL\_\allowbreak{}METRICS\_\allowbreak{}INCLUDE\_\allowbreak{}VERSION} (default \texttt{false}) and \texttt{OTEL\_\allowbreak{}METRICS\_\allowbreak{}INCLUDE\_\allowbreak{}ENTRYPOINT} (default \texttt{false}) [93]. Standard attributes include \texttt{user.\allowbreak{}email} and \texttt{organizatio\allowbreak{}n.\allowbreak{}id} when authenticated — personal data by most definitions.


For organisations on Claude for Enterprise, the analytics dashboard provides usage metrics and adoption tracking; under ZDR it shows usage metrics only, with contribution metrics unavailable [40], [94].


\FloatBarrier
\setcounter{section}{1}
\section{Cost control}

Three levers, in order of effect [59]:


\begin{enumerate}
\item \textbf{Context discipline.} Usage climbs superlinearly in a long session because every request carries the accumulated history. Handoffs and clean sessions are a cost control, not only a quality one (Chapter 11).
\item \textbf{Architecture.} An agent team costs roughly linearly in teammates; a dynamic workflow can spawn up to 1,000 agents; a fork shares the parent's prompt cache and is cheaper than a fresh subagent (Chapter 30;Chapter 22;Section 16.1). Choose the cheapest structure that answers the question.
\item \textbf{Model and effort.} \texttt{CLAUDE\_\allowbreak{}CODE\_\allowbreak{}SUBAGENT\_\allowbreak{}MODEL} routes delegated work to a smaller model; \texttt{workflow\allowbreak{}Size\allowbreak{}Guideline} bounds fan-out; \texttt{/\allowbreak{}model} before a large run is worth the two seconds.
\end{enumerate}

Watch the \texttt{Large workf\allowbreak{}low} warning, which appears when a run schedules more than 25 agents or projects more than 1.5 million tokens, and remember it is advisory, suppressed under ultracode, and pegged to the size guideline in force rather than to 25 — set one and its agent count becomes the threshold [19]. \texttt{/\allowbreak{}usage} and \texttt{/\allowbreak{}insights} are the in-session views; the analytics dashboard and OTel metrics are the organisational ones [35], [93], [94].


\FloatBarrier
\setcounter{section}{2}
\section{Accessibility}

Claude Code supports screen readers including VoiceOver and NVDA, with settings for screen magnifiers, reduced motion and colourblind-friendly themes [95]. Related controls: \texttt{/\allowbreak{}keybindings} opens the keyboard-shortcuts file [96]; \texttt{/\allowbreak{}statusline} configures a custom status bar showing context usage, cost and git state [97]; terminal configuration covers \texttt{Shift+Enter} for newlines, a terminal bell when Claude finishes, tmux behaviour and theme matching [98]; and the Concise output style materially reduces the volume of text a screen-reader user must traverse while preserving error reports and destructive-action confirmations in full (Chapter 19) [79].


An accessibility statement belongs in any organisational rollout plan. A tool that a colleague cannot use is not deployed; it is deployed to some people.



\begingroup
\def\kprows{%
\item OpenTelemetry turns 'we use an agent' into something an assurance function can examine.
\item Every event carries a prompt identifier, so one instruction's full activity can be reconstructed.
\item Content logging is off by default; enabling it moves prompts and code into your telemetry pipeline and changes its classification.
\item Context discipline, then architecture, then model choice — in that order of effect on cost.
}%
\def\kpbody{\begin{minipage}{\textwidth}\subsection*{Key points}\begin{itemize}\kprows\end{itemize}\end{minipage}}%
\begingroup
\def\sloppy{\tolerance 9999\emergencystretch 3em\hfuzz 200pt\vfuzz 200pt}%
\hbadness=10000\vbadness=10000\hfuzz=200pt\vfuzz=200pt
\global\setbox\kpbox=\hbox{\kpbody}%
\endgroup
\par\addvspace{4.2mm}
\ifdim\dimexpr\ht\kpbox+\dp\kpbox\relax>0.30\textheight
  \typeout{HANDBOOK-KEYPOINTS broken \the\dimexpr\ht\kpbox+\dp\kpbox\relax}%
  \subsection*{Key points}
  \begin{itemize}\kprows\end{itemize}
\else
  \typeout{HANDBOOK-KEYPOINTS atomic \the\dimexpr\ht\kpbox+\dp\kpbox\relax}%
  \noindent\kpbody
\fi
\par\addvspace{1.4mm}
\endgroup

\breakrule

\FloatBarrier
\parttitle{Part IX}{Capstone}
\addcontentsline{toc}{part}{Part IX \textemdash\ Capstone}

\FloatBarrier
\renewcommand{\chaptertitlelabel}{Part IX \textperiodcentered\ Chapter 34}
\setcounter{chapter}{33}
\chapter{Build a verified research-brief agent}

This capstone combines every essential mechanism without requiring production access. Complete it in the practice repository of “The running project”. The measure of success is not that Claude says it is finished; it is that the evidence exists.


\FloatBarrier
\setcounter{section}{0}
\section{Target outcome}

Given an approved question and source set, the workspace produces a draft decision brief, an \textbf{independent} evidence review, a rendered deliverable, and a verification report. It never publishes automatically.


\FloatBarrier
\setcounter{section}{1}
\section{Stage 1 — Establish control}

\begin{enumerate}
\item Run \texttt{claude doct\allowbreak{}or}; record the version.
\item Confirm the project directory, git branch and account with \texttt{/\allowbreak{}status}.
\item Review \texttt{.\allowbreak{}gitignore} and permission rules; run \texttt{/\allowbreak{}permissions} and note the source file of each rule.
\item Enable sandboxing with \texttt{/\allowbreak{}sandbox} where supported, including credential deny entries.
\item Commit a clean baseline.
\end{enumerate}

\textbf{Evidence:} version string, \texttt{git status} output, permission summary, sandbox status, baseline commit hash.


\FloatBarrier
\setcounter{section}{2}
\section{Stage 2 — Create project knowledge}

\begin{enumerate}
\item Finalise \texttt{CLAUDE.\allowbreak{}md} under 200 lines (Chapter 9).
\item Add \texttt{references/\allowbreak{}source-\allowbreak{}policy.\allowbreak{}md} with a retrieval rule rather than inlining it.
\item Add a path-scoped rule for \texttt{sources/\allowbreak{}**} that forbids modification.
\item Add \texttt{templates/\allowbreak{}claim-\allowbreak{}ledger.\allowbreak{}md} and \texttt{templates/\allowbreak{}brief.\allowbreak{}md}.
\item Inspect auto memory with \texttt{/\allowbreak{}memory} and remove anything inaccurate or over-generalised.
\end{enumerate}

\textbf{Evidence:} \texttt{/\allowbreak{}context} shows the instruction files you expect and no others; the instruction files do not contradict each other.


\FloatBarrier
\setcounter{section}{3}
\section{Stage 3 — Choose the register}

Create the \emph{Evidence analyst} output style of Section 19.2 and select it with \texttt{/\allowbreak{}config}. Start a new session so the system prompt is rebuilt.


\textbf{Evidence:} a new session's first response distinguishes source claim from inference without being asked.


\FloatBarrier
\setcounter{section}{4}
\section{Stage 4 — Make the reusable procedure}

\begin{enumerate}
\item Complete one brief manually, with Claude.
\item Create the \texttt{research-\allowbreak{}brief} skill from the observed procedure (Section 15.1).
\item Review every file and every tool grant. Set \texttt{allowed-\allowbreak{}tools} to the minimum.
\item Test it on a genuinely different source set.
\end{enumerate}

\textbf{Evidence:} the second run follows the same output contract without importing topic-specific assumptions from the first.


\FloatBarrier
\setcounter{section}{5}
\section{Stage 5 — Add an independent reviewer}

\begin{enumerate}
\item Create the read-only \texttt{evidence-\allowbreak{}reviewer} subagent of Section 16.3.
\item \textbf{Confirm it is not a fork.} Fork mode is the interactive default (Section 16.1); invoke the agent by name and verify from its opening output that it has no knowledge of the drafting conversation.
\item Confirm the parent session is not in \texttt{accept\allowbreak{}Edits}, \texttt{bypass\allowbreak{}Permissions} or \texttt{auto}, each of which overrides the reviewer's \texttt{permission\allowbreak{}Mode}. On Pro, Max and Team plans \texttt{auto} is where a session starts, so this step is not a formality.
\item Pass only the draft, the ledger and the source files.
\item Require a severity-ordered report and a pass/fail verdict.
\item Correct the draft and re-run the reviewer.
\end{enumerate}

\textbf{Evidence:} no unresolved high-severity unsupported claim; the reviewer's disagreements remain visible in the record rather than being silently reconciled.


\FloatBarrier
\setcounter{section}{6}
\section{Stage 6 — Add exactly one piece of bounded automation}

Choose \textbf{one}:


\begin{itemize}
\item a \texttt{Stop} hook that appends a timestamped, non-sensitive completion record (Chapter 17);
\item a \texttt{/\allowbreak{}loop} running a read-only local check during an attended session (Chapter 23);
\item a dynamic workflow with three genuinely independent research lanes (Chapter 22).
\end{itemize}

Do not add a plugin, a live MCP server, a cloud routine and browser access in order to "use every feature". Capability must answer a need. Every mechanism you add is one you must also review, monitor and be able to remove.


\textbf{Evidence:} a written threat model, a test result including the failure path, a removal procedure, and the measured \texttt{/\allowbreak{}usage} impact.


\FloatBarrier
\setcounter{section}{7}
\section{Stage 7 — Render and verify}

Generate the final draft in an editable source format, render it, then run:


\begin{itemize}
\item \textbf{structural checks} — required sections present, heading order correct, tables complete;
\item \textbf{source checks} — every material claim maps to a ledger row and a named source;
\item \textbf{link checks} — each URL resolves or is explicitly marked inaccessible;
\item \textbf{content checks} — contradictions, dates, scope and attribution;
\item \textbf{file checks} — opens successfully, expected page count, no clipping or overflow;
\item \textbf{security checks} — no secret in files, diff, logs, artifacts or quoted transcript excerpts;
\item \textbf{git checks} — reviewed diff containing only intended files, and a clean commit.
\end{itemize}

Run the checks as commands whose output lands in the transcript, so that a goal evaluator can see them (Section 21.1) and so that you can archive them.


\FloatBarrier
\setcounter{section}{8}
\section{Stage 8 — Human release gate}

Promotion from \texttt{drafts/\allowbreak{}} to \texttt{deliverable\allowbreak{}s/\allowbreak{}} happens only after a human confirms:


\begin{codeblock}{8.0}{9.4}
\cl{[~]~The~decision~question~is~answered.}
\cl{[~]~All~material~claims~are~supported,~or~visibly~qualified.}
\cl{[~]~The~independent~review~passed,~and~it~was~genuinely~independent.}
\cl{[~]~Commands~and~links~were~actually~executed,~and~their~output~was~inspected.}
\cl{[~]~External~side~effects~are~enumerated.}
\cl{[~]~No~credential~or~private~data~appears~in~any~output,~artifact~or~transcript~excerpt.}
\cl{[~]~The~rendered~document~was~opened~and~read.}
\cl{[~]~Publication,~sending~or~sharing~is~separately~approved.}
\end{codeblock}

\FloatBarrier
\setcounter{section}{9}
\section{A final goal}

\begin{codeblock}{7.0}{8.3}
\cl{/goal~Produce~drafts/final-brief.md~and~drafts/final-brief.pdf~from~the~approved~source~set.}
\cl{The~goal~is~met~only~when:~the~evidence-reviewer~subagent~reports~no~unresolved~high-severity}
\cl{finding;~the~structure,~link~and~security~checks~each~exit~0~with~their~output~shown~in~this}
\cl{conversation;~a~render~inspection~report~exists~at~notes/render-check.md;~git~diff~contains~only}
\cl{the~intended~files;~and~the~completion~message~lists~the~exact~commands~run~and~the~files}
\cl{changed.~Do~not~publish,~push,~use~external~accounts,~install~software,~or~promote~anything~to}
\cl{deliverables/.~Stop~and~ask~if~blocked,~on~any~missing~permission,~source~or~credential,~or}
\cl{after~25~turns.}
\end{codeblock}

Note what that condition does: every clause is something Claude must \textbf{surface in the conversation}, because the evaluator cannot run commands or read files (Section 21.1). A condition the evaluator cannot see is a condition it will accept on assertion.


The capstone is complete when the evidence exists — not when Claude says it is complete.


\begingroup
\def\kprows{%
\item Evidence-Gated Delivery: control, knowledge, register, procedure, independent review, bounded automation, verification, release.
\item Confirm the reviewer is genuinely independent and that the parent mode does not override it.
\item Add exactly one automation mechanism, with a threat model, a tested failure path and a removal procedure.
\item Run checks as commands whose output lands in the transcript, so the evidence is both evaluable and archivable.
}%
\def\kpbody{\begin{minipage}{\textwidth}\subsection*{Key points}\begin{itemize}\kprows\end{itemize}\end{minipage}}%
\begingroup
\def\sloppy{\tolerance 9999\emergencystretch 3em\hfuzz 200pt\vfuzz 200pt}%
\hbadness=10000\vbadness=10000\hfuzz=200pt\vfuzz=200pt
\global\setbox\kpbox=\hbox{\kpbody}%
\endgroup
\par\addvspace{4.2mm}
\ifdim\dimexpr\ht\kpbox+\dp\kpbox\relax>0.30\textheight
  \typeout{HANDBOOK-KEYPOINTS broken \the\dimexpr\ht\kpbox+\dp\kpbox\relax}%
  \subsection*{Key points}
  \begin{itemize}\kprows\end{itemize}
\else
  \typeout{HANDBOOK-KEYPOINTS atomic \the\dimexpr\ht\kpbox+\dp\kpbox\relax}%
  \noindent\kpbody
\fi
\par\addvspace{1.4mm}
\endgroup

\breakrule

\FloatBarrier
\frontchapter{Closing perspective}

The most capable Claude Code configuration is not the one with the most plugins, connectors, agents and permissions. It is the one that repeatedly produces an inspectable result from a clear objective, trusted context, narrow access and observed evidence.


Three propositions are worth carrying out of this book.


\textbf{Capability is not the constraint; verification is.} Every part of this system has grown faster than the human capacity to supervise it. Sixteen concurrent agents, a thousand per run, routines that fire on webhooks, channels that push events into a live session, teams that spawn without being asked. The scarce resource is not agent time. It is the number of things a person can meaningfully check, and every design decision should be read as a claim on that budget.


\textbf{Trust boundaries have moved inward.} The interesting attack surface is no longer the network perimeter but the content the agent reads: a web page, a tool result, a pull-request comment, a message from a peer session, a fire payload from a leaked token. Anthropic's own controls reflect this — untrusted-input wrappers, isolated context windows for web fetches, classifier review of inter-agent messages, the rule that a teammate cannot consent on your behalf. Design your own workflows on the same assumption, and treat zero trust as the operating model rather than a slogan [4], [8].


\textbf{Governance is now a build-time concern.} Managed policy, retention, residency, audit and accessibility are not paperwork applied after a tool is adopted. They are settings, and settings have precedence orders, failure modes and gaps — a ZDR organisation without \texttt{force\allowbreak{}Login\allowbreak{}Org\allowbreak{}UUID}, a managed file silently skipped because a higher-priority source exists, an auto-memory directory outside the retention sweep. The frameworks are converging on the same demand: document the system, evidence the controls, and name the human who is accountable [6], [11].


Start small. Preserve what works in reviewed files. Add autonomy only where you can define both completion and failure. Make external actions deliberate. And verify the artefact, not the confidence of the final message.


\breakrule

\FloatBarrier
\appendix
\setcounter{chapter}{0}
\setcounter{section}{0}
\appendixchapter{A}{In-session command reference}
Commands are recognised at the start of a message. Type \texttt{/\allowbreak{}} to see what is available for your plan, platform and installed release. Entries marked \textbf{[S]} are bundled skills rather than built-in commands; a scheduled task can only fire a skill Claude is permitted to invoke on its own (Section 23.3). Verified against [35] on 23 August 2026 and re-adjudicated against it on 26 August 2026.


\begingroup
\def\tblrows{%
\texttt{/\allowbreak{}help} & Show help and available commands & First stop when the interface differs from this book \\
\texttt{/\allowbreak{}status} & Session, account and setting-source status & Shows which managed source was selected (Chapter 31) \\
\texttt{/\allowbreak{}doctor} \textbf{[S]} & Setup checkup; diagnose and fix & Alias \texttt{/\allowbreak{}checkup}; CLI equivalent \texttt{claude doct\allowbreak{}or} \\
\texttt{/\allowbreak{}config} & Open settings or set an individual key & Alias \texttt{/\allowbreak{}settings}. Home of output style, workflows, ultracode keyword \\
\texttt{/\allowbreak{}login} · \texttt{/\allowbreak{}logout} & Authenticate, switch account, sign out & Several features need a claude.ai login, not an API key \\
\texttt{/\allowbreak{}privacy-\allowbreak{}settings} & View and update privacy settings & See Chapter 32 \\
\texttt{/\allowbreak{}model} & Switch model and save as default & Availability changes \\
\texttt{/\allowbreak{}effort} & Set effort: \texttt{low}, \texttt{medium}, \texttt{high}, \texttt{xhigh}, \texttt{max}, \texttt{auto} and \texttt{ultracode} & Availability differs by surface — see Table 3. \texttt{ultracode} also enables workflow orchestration \\
\texttt{/\allowbreak{}fast} & Toggle fast mode & Model-dependent \\
\texttt{/\allowbreak{}advisor} & Enable a stronger advisor model & Consulted at key moments \\
\texttt{/\allowbreak{}usage} & Cost, limits and activity & \texttt{/\allowbreak{}cost} and \texttt{/\allowbreak{}stats} are both documented aliases; \texttt{/\allowbreak{}stats} opens on the Stats tab \\
\texttt{/\allowbreak{}insights} & HTML report analysing recent sessions &  \\
\texttt{/\allowbreak{}context} & Context composition & Diagnose; do not chase a threshold \\
\texttt{/\allowbreak{}autocompact} & Set the auto-compact window &  \\
\texttt{/\allowbreak{}compact [in\allowbreak{}structions]} & Summarise the conversation & Lossy; say what must survive \\
\texttt{/\allowbreak{}clear} & Empty context, new conversation & Aliases \texttt{/\allowbreak{}reset}, \texttt{/\allowbreak{}new}. Clears goal and scheduled tasks \\
\texttt{/\allowbreak{}rewind} & Roll code and conversation to a checkpoint & Cannot undo external side effects \\
\texttt{/\allowbreak{}plan} & Enter plan mode & Optional task description \\
\texttt{/\allowbreak{}permissions} & Manage allow/ask/deny rules & Alias \texttt{/\allowbreak{}allowed-\allowbreak{}tools}. Deny → ask → allow \\
\texttt{/\allowbreak{}sandbox} & Sandbox panel & Not on native Windows \\
\texttt{/\allowbreak{}hooks} & View configured hooks & Read-only; shows source file \\
\texttt{/\allowbreak{}init} & Draft a starter \texttt{CLAUDE.\allowbreak{}md} & Review and shorten \\
\texttt{/\allowbreak{}memory} & Edit instruction files; toggle auto memory & Writes \texttt{auto\allowbreak{}Memory\allowbreak{}Enabled} \\
\texttt{/\allowbreak{}add-\allowbreak{}dir} & Add a working directory & Confirm the target first \\
\texttt{/\allowbreak{}cd} & Move the session to a new directory & Relocates session storage \\
\texttt{/\allowbreak{}diff} & Interactive diff viewer & Review before commit \\
\texttt{/\allowbreak{}export} & Export the conversation as plain text & The supported audit artefact \\
\texttt{/\allowbreak{}recap} & One-line session summary &  \\
\texttt{/\allowbreak{}btw} & Side question, not added to history &  \\
\texttt{/\allowbreak{}focus} & Show only the last prompt and response &  \\
\texttt{/\allowbreak{}resume} · \texttt{/\allowbreak{}rename} · \texttt{/\allowbreak{}branch} & Session control & See Chapter 12 \\
\texttt{/\allowbreak{}fork} · \texttt{/\allowbreak{}subtask} & Copy to a background session; fork a subagent & Forks inherit the conversation \\
\texttt{/\allowbreak{}background} & Detach the session & Alias \texttt{/\allowbreak{}bg}; carries \texttt{/\allowbreak{}loop} tasks with it \\
\texttt{/\allowbreak{}tasks} & Background work and watched artifacts &  \\
\texttt{/\allowbreak{}list-\allowbreak{}agents} & Subagents, teammates, other sessions & Alias \texttt{/\allowbreak{}peers} \\
\texttt{/\allowbreak{}agents} & Guidance for agent files & No longer the editor shown in the recording \\
\texttt{/\allowbreak{}mcp} & Inspect and authenticate MCP servers & Review tools and scopes \\
\texttt{/\allowbreak{}plugin} · \texttt{/\allowbreak{}reload-\allowbreak{}plugins} & Plugin management & Review arbitrary-code risk first \\
\texttt{/\allowbreak{}import} & Bring another agent's configuration in & Requires 2.1.213+ \\
\texttt{/\allowbreak{}chrome} & Configure the browser connection & Shares authenticated browser state \\
\texttt{/\allowbreak{}artifacts} & List, attach, open or copy artifacts & Attach before updating from a new session \\
\texttt{/\allowbreak{}goal} · \texttt{/\allowbreak{}goal clear} & Persistent completion condition & Implemented as a prompt Stop hook \\
\texttt{/\allowbreak{}loop [inter\allowbreak{}val] [promp\allowbreak{}t]} \textbf{[S]} & Session-scoped recurring prompt & Alias \texttt{/\allowbreak{}proactive}. \textbf{Bare \texttt{/\allowbreak{}loop} runs a maintenance prompt} \\
\texttt{/\allowbreak{}schedule} & Create and manage cloud routines & Alias \texttt{/\allowbreak{}routines}; research preview \\
\texttt{/\allowbreak{}remote-\allowbreak{}control} & Attach remote access & Alias \texttt{/\allowbreak{}rc} \\
\texttt{/\allowbreak{}desktop} & Continue the session in Desktop & Alias \texttt{/\allowbreak{}app} \\
\texttt{/\allowbreak{}teleport} & Pull a cloud session into the terminal & Alias \texttt{/\allowbreak{}tp} \\
\texttt{/\allowbreak{}mobile} & QR code for the mobile app & Aliases \texttt{/\allowbreak{}ios}, \texttt{/\allowbreak{}android} \\
\texttt{/\allowbreak{}workflows} & Workflow progress view & \texttt{s} saves a run as a command \\
\texttt{/\allowbreak{}deep-\allowbreak{}research <q\allowbreak{}uestion>} & Bundled research workflow & Requires WebSearch \\
\texttt{/\allowbreak{}batch} \textbf{[S]} & Large-scale parallel changes &  \\
\texttt{/\allowbreak{}code-\allowbreak{}review} \textbf{[S]} & Review a diff or PR & Alias \texttt{/\allowbreak{}review}; \texttt{ultra} for the cloud pass \\
\texttt{/\allowbreak{}security-\allowbreak{}review} \textbf{[S]} & Security pass over the branch & Not a substitute for specialist review \\
\texttt{/\allowbreak{}verify} \textbf{[S]} & Run and observe verification & \texttt{disable-\allowbreak{}model-\allowbreak{}invocation:\allowbreak{} true} \\
\texttt{/\allowbreak{}fewer-\allowbreak{}permission-\allowbreak{}prompts} \textbf{[S]} & Propose an allowlist from transcripts & Review before applying \\
\texttt{/\allowbreak{}auto-\allowbreak{}mode-\allowbreak{}setup} & Draft \texttt{auto\allowbreak{}Mode.\allowbreak{}environment} entries & See [99] \\
\texttt{/\allowbreak{}debug} \textbf{[S]} · \texttt{/\allowbreak{}heapdump} & Diagnostics &  \\
\texttt{/\allowbreak{}voice} & Voice dictation & Proofread negations and numbers \\
\texttt{/\allowbreak{}keybindings} · \texttt{/\allowbreak{}color} · \texttt{/\allowbreak{}theme} & Interface configuration & See [96] \\
\texttt{/\allowbreak{}feedback} · \texttt{/\allowbreak{}bug} · \texttt{/\allowbreak{}share} & Send feedback or a transcript & Sends conversation content — see Table 23 \\
\texttt{/\allowbreak{}install-\allowbreak{}github-\allowbreak{}app} · \texttt{/\allowbreak{}install-\allowbreak{}slack-\allowbreak{}app} & Integration installers &  \\
\texttt{/\allowbreak{}autofix-\allowbreak{}pr} & Watch a PR and push fixes & Can trigger comment-driven automation \\
\texttt{/\allowbreak{}powerup} & Interactive feature lessons &  \\
\texttt{/\allowbreak{}exit} & Exit & Alias \texttt{/\allowbreak{}quit} \\
}%
\def\tblbody{\begin{minipage}{\textwidth}
{\scriptsize\begin{tabular}{L{48.6mm}L{40.8mm}L{50.5mm}}
\toprule
\textbf{Command} & \textbf{Purpose} & \textbf{Note} \\
\midrule
\tblrows
\bottomrule\end{tabular}}\end{minipage}}%
\begingroup
\def\sloppy{\tolerance 9999\emergencystretch 3em\hfuzz 200pt\vfuzz 200pt}%
\hbadness=10000\vbadness=10000\hfuzz=200pt\vfuzz=200pt
\global\setbox\tblbox=\hbox{\tblbody}%
\endgroup
\par\addvspace{2.6mm}
\ifdim\dimexpr\ht\tblbox+\dp\tblbox\relax>0.55\textheight
  \typeout{HANDBOOK-TABLE broken \the\dimexpr\ht\tblbox+\dp\tblbox\relax}%
  
  {\scriptsize\begin{longtable}{L{48.6mm}L{40.8mm}L{50.5mm}}
  \toprule
\textbf{Command} & \textbf{Purpose} & \textbf{Note} \\
\midrule\endfirsthead
  \multicolumn{3}{@{}l@{}}{%
  \sffamily\footnotesize\itshape\color{inkgrey}Continued}\\[1.2mm]
  \toprule
\textbf{Command} & \textbf{Purpose} & \textbf{Note} \\
\midrule\endhead
  \bottomrule\endfoot
  \bottomrule\endlastfoot
  \tblrows
  \end{longtable}}%
\else
  \typeout{HANDBOOK-TABLE atomic \the\dimexpr\ht\tblbox+\dp\tblbox\relax}%
  \noindent\tblbody
\fi
\par\addvspace{2.6mm}
\endgroup

\FloatBarrier
\setcounter{chapter}{1}
\setcounter{section}{0}
\appendixchapter{B}{Shell commands and flags}
Run \texttt{claude -\allowbreak{}-\allowbreak{}help} and \texttt{claude mcp -\allowbreak{}-\allowbreak{}help} before scripting: interactive commands and shell flags are different interfaces, and a few flags are deliberately absent from \texttt{-\allowbreak{}-\allowbreak{}help} while their feature is in preview. Verified against [62] on 23 August 2026.


\texttt{claude -\allowbreak{}-\allowbreak{}help} reports 62 long options in option position and 13 subcommands, identical between 2.1.241 and 2.1.246. Three further flags — \texttt{-\allowbreak{}-\allowbreak{}allowed-\allowbreak{}tools}, \texttt{-\allowbreak{}-\allowbreak{}background} and \texttt{-\allowbreak{}-\allowbreak{}disallowed-\allowbreak{}tools} — are named only inside other flags' descriptions rather than in option position, and are not counted among the 62. The listing below is a practitioner's selection from that set, not a reproduction of it, and \texttt{-\allowbreak{}-\allowbreak{}effort} is one of the flags whose accepted values exceed what \texttt{-\allowbreak{}-\allowbreak{}help} enumerates (Table 3).


\begin{codeblock}{8.0}{9.4}
\cl{claude~~~~~~~~~~~~~~~~~~~~~~~~~~~~~~~~~~\#~start~interactive}
\cl{claude~--version~~~~~~~~~~~~~~~~~~~~~~~~\#~version}
\cl{claude~doctor~~~~~~~~~~~~~~~~~~~~~~~~~~~\#~diagnostics}
\cl{claude~-p~"<prompt>"~~~~~~~~~~~~~~~~~~~~\#~non-interactive~(alias~--print)}
\cl{claude~-p~--output-format~json~~~~~~~~~~\#~structured~result:~JSON~or~stream-json}
\cl{claude~--continue~~~~~~~~~~~~~~~~~~~~~~~\#~resume~most~recent~here~(alias~-c)}
\cl{claude~--resume~[name|id]~~~~~~~~~~~~~~~\#~picker,~or~resume~directly~(alias~-r)}
\cl{claude~--from-pr~<number>~~~~~~~~~~~~~~~\#~picker~filtered~to~a~pull~request}
\cl{claude~-n~<name>~~~~~~~~~~~~~~~~~~~~~~~~\#~name~the~session~(alias~--name)}
\cl{claude~--fork-session~~~~~~~~~~~~~~~~~~~\#~fork~on~resume,~in~a~new~process}
\cl{claude~--no-session-persistence~~~~~~~~~\#~suppress~transcript~writes~for~this~run}
\cl{claude~--model~<model>~~~~~~~~~~~~~~~~~~\#~model~for~this~session}
\cl{claude~--effort~ultracode~~~~~~~~~~~~~~~\#~xhigh~plus~automatic~workflows~(2.1.203+)}
\cl{claude~--permission-mode~plan~~~~~~~~~~~\#~start~in~a~specific~permission~mode}
\cl{claude~--dangerously-skip-permissions~~~\#~isolated~environments~only}
\cl{claude~--add-dir~<path>~~~~~~~~~~~~~~~~~\#~additional~working~directory}
\cl{claude~--settings~<file>~~~~~~~~~~~~~~~~\#~explicit~settings~payload}
\cl{claude~--setting-sources~<list>~~~~~~~~~\#~restrict~which~settings~sources~load}
\cl{claude~--append-system-prompt~"<text>"~~\#~one-invocation~system~prompt~addition}
\cl{claude~--agent~<name>~~~~~~~~~~~~~~~~~~~\#~run~the~session~as~a~subagent~definition}
\cl{claude~--agents~\textquotesingle{}<json>\textquotesingle{}~~~~~~~~~~~~~~~~\#~define~session-scoped~subagents}
\cl{claude~--bg~"<prompt>"~~~~~~~~~~~~~~~~~~\#~background~session~(alias~--background)}
\cl{claude~agents~~~~~~~~~~~~~~~~~~~~~~~~~~~\#~agent~view:~monitor~background~sessions}
\cl{claude~--worktree~<name>~~~~~~~~~~~~~~~~\#~isolated~git~worktree~and~branch}
\cl{claude~--teammate-mode~auto~~~~~~~~~~~~~\#~agent-team~display~mode~(experimental)}
\cl{claude~--chrome~/~--no-chrome~~~~~~~~~~~\#~browser~integration~for~this~session}
\cl{claude~--remote-control~[name]~~~~~~~~~~\#~expose~this~local~session~(alias~--rc)}
\cl{claude~remote-control~~~~~~~~~~~~~~~~~~~\#~server~mode}
\cl{claude~--cloud~"<task>"~~~~~~~~~~~~~~~~~\#~create~a~cloud~session}
\cl{claude~--cloud~<session-id>~~~~~~~~~~~~~\#~attach,~or~with~-p~queue~a~message}
\cl{claude~--teleport~[session-id]~~~~~~~~~~\#~pull~a~cloud~session~into~the~terminal}
\cl{claude~--channels~plugin:<name>@<mkt>~~~\#~enable~channel~servers~for~this~session}
\cl{claude~--plugin-dir~<path|zip>~~~~~~~~~~\#~load~a~local~plugin}
\cl{claude~--plugin-url~<url>~~~~~~~~~~~~~~~\#~load~a~plugin~archive~for~this~session}
\cl{claude~--mcp-config~<file>~~~~~~~~~~~~~~\#~session~MCP~configuration}
\cl{claude~mcp~add~--transport~http~<n>~<u>~\#~add~an~MCP~server}
\cl{claude~mcp~list~|~get~|~remove~|~login~|~logout}
\cl{claude~plugin~init~|~validate~|~details~<name>}
\cl{claude~plugin~marketplace~add~<owner/repo>}
\cl{claude~auth~login~~~~~~~~~~~~~~~~~~~~~~~\#~account~authentication}
\cl{claude~setup-token~~~~~~~~~~~~~~~~~~~~~~\#~long-lived~token~(blocks~Chrome,~artifacts)}
\end{codeblock}

\texttt{claude -\allowbreak{}-\allowbreak{}effort auto} is refused: the flag warns and the session runs at the default effort. \texttt{claude -\allowbreak{}-\allowbreak{}permission-\allowbreak{}mode} is the one flag whose value set this book states in full, in Table 5.


Flags not restored on \texttt{-\allowbreak{}-\allowbreak{}resume}: \texttt{-\allowbreak{}-\allowbreak{}mcp-\allowbreak{}config}, \texttt{-\allowbreak{}-\allowbreak{}settings}, \texttt{-\allowbreak{}-\allowbreak{}plugin-\allowbreak{}dir}, \texttt{-\allowbreak{}-\allowbreak{}fallback-\allowbreak{}model}, \texttt{-\allowbreak{}-\allowbreak{}add-\allowbreak{}dir} [30].


\FloatBarrier
\setcounter{chapter}{2}
\setcounter{section}{0}
\appendixchapter{C}{Files, directories and scope}
\texttt{\textasciitilde{}} is the current user's home directory. Windows paths resolve under the user profile; use the live settings documentation for exact platform expansion [63], [92].


\begingroup
\def\tblrows{%
\texttt{CLAUDE.\allowbreak{}md} & Project instructions & Yes, after review \\
\texttt{.\allowbreak{}claude/\allowbreak{}CLAUDE.\allowbreak{}md} & Alternative project instruction location & Yes, after review \\
\texttt{CLAUDE.\allowbreak{}local.\allowbreak{}md} & Private project instructions & No; keep ignored \\
\texttt{\textasciitilde{}/\allowbreak{}.\allowbreak{}claude/\allowbreak{}CLAUDE.\allowbreak{}md} & User-wide instructions & No \\
\texttt{/\allowbreak{}Library/\allowbreak{}Application\allowbreak{} Support/\allowbreak{}Claude\allowbreak{}Code/\allowbreak{}CLAUDE.\allowbreak{}md} (macOS) · \texttt{/\allowbreak{}etc/\allowbreak{}claude-\allowbreak{}code/\allowbreak{}CLAUDE.\allowbreak{}md} (Linux, WSL) · \texttt{C:\allowbreak{}\textbackslash{}Program Fil\allowbreak{}es\textbackslash{}Claude\allowbreak{}Code\textbackslash{}CLAUDE.\allowbreak{}md} (Windows) & Managed policy instructions; cannot be excluded & Deployed by MDM \\
\texttt{.\allowbreak{}claude/\allowbreak{}rules/\allowbreak{}*.\allowbreak{}md} & Modular project rules, optionally path-scoped & Usually yes \\
\texttt{\textasciitilde{}/\allowbreak{}.\allowbreak{}claude/\allowbreak{}rules/\allowbreak{}*.\allowbreak{}md} & User rules; load before project rules & No \\
\texttt{.\allowbreak{}claude/\allowbreak{}settings.\allowbreak{}json} & Shared project settings & Yes, if secret-free \\
\texttt{.\allowbreak{}claude/\allowbreak{}settings.\allowbreak{}local.\allowbreak{}json} & Personal project settings and saved approvals & No \\
\texttt{\textasciitilde{}/\allowbreak{}.\allowbreak{}claude/\allowbreak{}settings.\allowbreak{}json} & User settings & No \\
\texttt{managed-\allowbreak{}settings.\allowbreak{}json} and \texttt{managed-\allowbreak{}settings.\allowbreak{}d/\allowbreak{}*.\allowbreak{}json} in the system directory & Organisation policy & Deployed, not committed \\
\texttt{managed-\allowbreak{}mcp.\allowbreak{}json} & Managed MCP configuration & Deployed \\
\texttt{.\allowbreak{}mcp.\allowbreak{}json} & Shared project MCP configuration & Yes only if secret-free and reviewed \\
\texttt{\textasciitilde{}/\allowbreak{}.\allowbreak{}claude.\allowbreak{}json} & User state; local and user MCP configuration & No \\
\texttt{.\allowbreak{}claude/\allowbreak{}skills/\allowbreak{}<name>/\allowbreak{}SKILL.\allowbreak{}md} & Project skill & Yes, after security review \\
\texttt{\textasciitilde{}/\allowbreak{}.\allowbreak{}claude/\allowbreak{}skills/\allowbreak{}<name>/\allowbreak{}SKILL.\allowbreak{}md} & User skill & No \\
\texttt{.\allowbreak{}claude/\allowbreak{}commands/\allowbreak{}<name>.\allowbreak{}md} & Legacy project command & Migrate to a skill \\
\texttt{.\allowbreak{}claude/\allowbreak{}agents/\allowbreak{}<name>.\allowbreak{}md} & Project subagent & Yes, after tool review \\
\texttt{\textasciitilde{}/\allowbreak{}.\allowbreak{}claude/\allowbreak{}agents/\allowbreak{}<name>.\allowbreak{}md} & User subagent & No \\
\texttt{.\allowbreak{}claude/\allowbreak{}workflows/\allowbreak{}*.\allowbreak{}js} · \texttt{\textasciitilde{}/\allowbreak{}.\allowbreak{}claude/\allowbreak{}workflows/\allowbreak{}*.\allowbreak{}js} & Saved dynamic workflows & Project ones, after review \\
\texttt{.\allowbreak{}claude/\allowbreak{}output-\allowbreak{}styles/\allowbreak{}} · \texttt{\textasciitilde{}/\allowbreak{}.\allowbreak{}claude/\allowbreak{}output-\allowbreak{}styles/\allowbreak{}} & Output styles & Project ones, yes \\
\texttt{.\allowbreak{}claude/\allowbreak{}loop.\allowbreak{}md} · \texttt{\textasciitilde{}/\allowbreak{}.\allowbreak{}claude/\allowbreak{}loop.\allowbreak{}md} & Default prompt for a bare \texttt{/\allowbreak{}loop} & Project one, after review \\
\texttt{.\allowbreak{}claude/\allowbreak{}hooks/\allowbreak{}} & Hook scripts & Yes, after review \\
\texttt{\textasciitilde{}/\allowbreak{}.\allowbreak{}claude/\allowbreak{}projects/\allowbreak{}<project>/\allowbreak{}<session-\allowbreak{}id>.\allowbreak{}jsonl} & Session transcripts, plaintext JSONL & No \\
\texttt{\textasciitilde{}/\allowbreak{}.\allowbreak{}claude/\allowbreak{}projects/\allowbreak{}<project>/\allowbreak{}memory/\allowbreak{}} & Auto memory; \textbf{excluded from the retention sweep} & No \\
\texttt{\textasciitilde{}/\allowbreak{}.\allowbreak{}claude/\allowbreak{}agent-\allowbreak{}memory/\allowbreak{}<name>/\allowbreak{}} · \texttt{.\allowbreak{}claude/\allowbreak{}agent-\allowbreak{}memory[-\allowbreak{}local]/\allowbreak{}<name>/\allowbreak{}} & Subagent persistent memory & Project one, deliberately \\
\texttt{\textasciitilde{}/\allowbreak{}.\allowbreak{}claude/\allowbreak{}teams/\allowbreak{}<team>/\allowbreak{}} · \texttt{\textasciitilde{}/\allowbreak{}.\allowbreak{}claude/\allowbreak{}tasks/\allowbreak{}<team>/\allowbreak{}} & Agent-team config, mailboxes and task list & No \\
\texttt{\textasciitilde{}/\allowbreak{}.\allowbreak{}claude/\allowbreak{}channels/\allowbreak{}<name>/\allowbreak{}.\allowbreak{}env} & Channel plugin credentials & \textbf{Never} \\
\texttt{\textasciitilde{}/\allowbreak{}.\allowbreak{}claude/\allowbreak{}feedback-\allowbreak{}bundles/\allowbreak{}} & Local feedback archives on third-party providers & No \\
\texttt{.\allowbreak{}env} & Local environment values & \textbf{Never} \\
\texttt{.\allowbreak{}env.\allowbreak{}example} & Variable names only & Yes \\
\texttt{.\allowbreak{}gitignore} · \texttt{.\allowbreak{}worktreeinc\allowbreak{}lude} & Exclusion and worktree rules & Yes \\
}%
\def\tblbody{\begin{minipage}{\textwidth}
{\scriptsize\begin{tabular}{L{76.9mm}L{38.3mm}L{24.6mm}}
\toprule
\textbf{Path} & \textbf{Scope and purpose} & \textbf{Commit?} \\
\midrule
\tblrows
\bottomrule\end{tabular}}\end{minipage}}%
\begingroup
\def\sloppy{\tolerance 9999\emergencystretch 3em\hfuzz 200pt\vfuzz 200pt}%
\hbadness=10000\vbadness=10000\hfuzz=200pt\vfuzz=200pt
\global\setbox\tblbox=\hbox{\tblbody}%
\endgroup
\par\addvspace{2.6mm}
\ifdim\dimexpr\ht\tblbox+\dp\tblbox\relax>0.55\textheight
  \typeout{HANDBOOK-TABLE broken \the\dimexpr\ht\tblbox+\dp\tblbox\relax}%
  
  {\scriptsize\begin{longtable}{L{76.9mm}L{38.3mm}L{24.6mm}}
  \toprule
\textbf{Path} & \textbf{Scope and purpose} & \textbf{Commit?} \\
\midrule\endfirsthead
  \multicolumn{3}{@{}l@{}}{%
  \sffamily\footnotesize\itshape\color{inkgrey}Continued}\\[1.2mm]
  \toprule
\textbf{Path} & \textbf{Scope and purpose} & \textbf{Commit?} \\
\midrule\endhead
  \bottomrule\endfoot
  \bottomrule\endlastfoot
  \tblrows
  \end{longtable}}%
\else
  \typeout{HANDBOOK-TABLE atomic \the\dimexpr\ht\tblbox+\dp\tblbox\relax}%
  \noindent\tblbody
\fi
\par\addvspace{2.6mm}
\endgroup

\FloatBarrier
\section*{Settings precedence}
\markboth{Settings precedence}{Settings precedence}

Highest to lowest [63]:


\begin{codeblock}{9.0}{10.6}
\cl{Managed~settings~~~(managed-settings.json,~MDM,~or~the~claude.ai~console)}
\cl{~~~~~~~~↓}
\cl{Command~line~~~~~~~(claude~--settings)}
\cl{~~~~~~~~↓}
\cl{Project~local~~~~~~(.claude/settings.local.json)}
\cl{~~~~~~~~↓}
\cl{Shared~project~~~~~(.claude/settings.json)}
\cl{~~~~~~~~↓}
\cl{User~~~~~~~~~~~~~~~(\textasciitilde{}/.claude/settings.json)}
\end{codeblock}

Two cautions. \textbf{No level, including command-line arguments, overrides a managed permission rule} [28]. And permission-rule evaluation has its own precedence — deny, then ask, then allow — which is independent of settings precedence: do not assume a narrower allow defeats a broader deny (Section 7.2).


\FloatBarrier
\setcounter{chapter}{3}
\setcounter{section}{0}
\appendixchapter{D}{Settings and variables of security consequence}
A working set for a security review. Full definitions are in [100] and [101].


\begingroup
\def\tblrows{%
\texttt{permissions.\allowbreak{}deny} / \texttt{ask} / \texttt{allow} & Tool access. Deny → ask → allow \\
\texttt{permissions.\allowbreak{}disable\allowbreak{}Bypass\allowbreak{}Permissions\allowbreak{}Mode} & Prevents \texttt{bypass\allowbreak{}Permissions} \\
\texttt{allow\allowbreak{}Managed\allowbreak{}Permission\allowbreak{}Rules\allowbreak{}Only} & Only managed permission rules apply \\
\texttt{allow\allowbreak{}Managed\allowbreak{}Hooks\allowbreak{}Only} & Blocks user, project, local and plugin hooks — and \texttt{/\allowbreak{}goal} \\
\texttt{allow\allowbreak{}Managed\allowbreak{}Mcp\allowbreak{}Servers\allowbreak{}Only} · \texttt{allowed\allowbreak{}Mcp\allowbreak{}Servers} · \texttt{denied\allowbreak{}Mcp\allowbreak{}Servers} & MCP allow/deny at organisation level \\
\texttt{disable\allowbreak{}All\allowbreak{}Hooks} & Turns hooks off, subject to precedence; cannot disable managed hooks \\
\texttt{allowed\allowbreak{}Http\allowbreak{}Hook\allowbreak{}Urls} · \texttt{http\allowbreak{}Hook\allowbreak{}Allowed\allowbreak{}Env\allowbreak{}Vars} & Restrict HTTP hook destinations and header interpolation \\
\texttt{sandbox.\allowbreak{}enabled} · \texttt{filesystem.\allowbreak{}deny\allowbreak{}Read} / \texttt{deny\allowbreak{}Write} / \texttt{allow\allowbreak{}Read} & Filesystem isolation \\
\texttt{sandbox.\allowbreak{}credentials.\allowbreak{}files} / \texttt{.\allowbreak{}env\allowbreak{}Vars} & Credential deny or mask \\
\texttt{sandbox.\allowbreak{}network.\allowbreak{}allow\allowbreak{}Managed\allowbreak{}Domains\allowbreak{}Only} · \texttt{strict\allowbreak{}Allowlist} & Network allowlist locks \\
\texttt{auto\allowbreak{}Allow\allowbreak{}Bash\allowbreak{}If\allowbreak{}Sandboxed} & Defaults to \texttt{true}; set \texttt{false} to prompt anyway \\
\texttt{strict\allowbreak{}Known\allowbreak{}Marketplace\allowbreak{}s} · \texttt{blocked\allowbreak{}Marketplace\allowbreak{}s} · \texttt{plugin\allowbreak{}Trust\allowbreak{}Message} & Plugin supply chain \\
\texttt{strict\allowbreak{}Plugin\allowbreak{}Only\allowbreak{}Customizati\allowbreak{}on} & Blocks skills, agents, hooks and MCP from user and project sources \\
\texttt{disable\allowbreak{}Sideload\allowbreak{}Flags} & Rejects \texttt{-\allowbreak{}-\allowbreak{}plugin-\allowbreak{}dir}, \texttt{-\allowbreak{}-\allowbreak{}plugin-\allowbreak{}url}, \texttt{-\allowbreak{}-\allowbreak{}agents}, \texttt{-\allowbreak{}-\allowbreak{}mcp-\allowbreak{}config} \\
\texttt{disable\allowbreak{}Bundled\allowbreak{}Skills} · \texttt{skill\allowbreak{}Overrides} · \texttt{disable\allowbreak{}Skill\allowbreak{}Shell\allowbreak{}Execution} & Skill exposure and inline shell \\
\texttt{channels\allowbreak{}Enabled} · \texttt{allowed\allowbreak{}Channel\allowbreak{}Plugins} & Channels, which can relay permission decisions \\
\texttt{disable\allowbreak{}Artifact} · \texttt{CLAUDE\_\allowbreak{}CODE\_\allowbreak{}DISABLE\_\allowbreak{}ARTIFACT} & Artifact publishing \\
\texttt{disable\allowbreak{}Workflows} · \texttt{CLAUDE\_\allowbreak{}CODE\_\allowbreak{}DISABLE\_\allowbreak{}WORKFLOWS} · \texttt{workflow\allowbreak{}Size\allowbreak{}Guideline} & Dynamic workflows \\
\texttt{disable\allowbreak{}Remote\allowbreak{}Control} & Remote Control per device \\
\texttt{force\allowbreak{}Login\allowbreak{}Method} · \texttt{force\allowbreak{}Login\allowbreak{}Org\allowbreak{}UUID} & Route sessions into the intended organisation — required for ZDR \\
\texttt{available\allowbreak{}Models} · \texttt{enforce\allowbreak{}Available\allowbreak{}Models} & Model allowlist. A managed \texttt{model} is only a default \\
\texttt{cleanup\allowbreak{}Period\allowbreak{}Days} & Local transcript retention. Does \textbf{not} cover auto memory \\
\texttt{auto\allowbreak{}Memory\allowbreak{}Enabled} · \texttt{auto\allowbreak{}Memory\allowbreak{}Directory} · \texttt{CLAUDE\_\allowbreak{}CODE\_\allowbreak{}DISABLE\_\allowbreak{}AUTO\_\allowbreak{}MEMORY} & Auto memory \\
\texttt{claude\allowbreak{}Md} · \texttt{claude\allowbreak{}Md\allowbreak{}Excludes} & Managed instruction content and monorepo exclusions \\
\texttt{required\allowbreak{}Minimum\allowbreak{}Version} · \texttt{required\allowbreak{}Maximum\allowbreak{}Version} & Version floor and ceiling; fail open by design \\
\texttt{CLAUDE\_\allowbreak{}CODE\_\allowbreak{}EXPERIMENTA\allowbreak{}L\_\allowbreak{}AGENT\_\allowbreak{}TEAMS} & Agent teams; changes ordinary delegation when on \\
\texttt{CLAUDE\_\allowbreak{}CODE\_\allowbreak{}FORK\_\allowbreak{}SUBAGENT} & Fork mode; \textbf{on by default interactively} \\
\texttt{CLAUDE\_\allowbreak{}CODE\_\allowbreak{}MAX\_\allowbreak{}SUBAGENT\_\allowbreak{}SPAWN\_\allowbreak{}DEPTH} · \texttt{.\allowbreak{}.\allowbreak{}.\allowbreak{}\_\allowbreak{}MAX\_\allowbreak{}CONCURRENT\_\allowbreak{}SUBAGENTS} & Subagent nesting and concurrency \\
\texttt{CLAUDE\_\allowbreak{}CODE\_\allowbreak{}DISABLE\_\allowbreak{}CRON} & Disables the scheduler and \texttt{/\allowbreak{}loop} \\
\texttt{CLAUDE\_\allowbreak{}CODE\_\allowbreak{}SKIP\_\allowbreak{}PROMPT\_\allowbreak{}HISTORY} · \texttt{CLAUDE\_\allowbreak{}CONFIG\_\allowbreak{}DIR} · \texttt{CLAUDE\_\allowbreak{}CODE\_\allowbreak{}PROJECT\_\allowbreak{}DIR\_\allowbreak{}NAME} & Transcript writing and location \\
\texttt{DISABLE\_\allowbreak{}TELEMETRY} · \texttt{DISABLE\_\allowbreak{}ERROR\_\allowbreak{}REPORTING} · \texttt{DISABLE\_\allowbreak{}FEEDBACK\_\allowbreak{}COMMAND} · \texttt{CLAUDE\_\allowbreak{}CODE\_\allowbreak{}DISABLE\_\allowbreak{}FEEDBACK\_\allowbreak{}SURVEY} & Individual egress flows \\
\texttt{CLAUDE\_\allowbreak{}CODE\_\allowbreak{}DISABLE\_\allowbreak{}NONESSENTIA\allowbreak{}L\_\allowbreak{}TRAFFIC} & All of the above at once; also disables Remote Control feature flags and \texttt{/\allowbreak{}schedule} \\
\texttt{skip\allowbreak{}Web\allowbreak{}Fetch\allowbreak{}Preflight} & WebFetch hostname safety check \\
\texttt{OTEL\_\allowbreak{}LOG\_\allowbreak{}USER\_\allowbreak{}PROMPTS} · \texttt{OTEL\_\allowbreak{}LOG\_\allowbreak{}ASSISTANT\_\allowbreak{}RESPONSES} · \texttt{OTEL\_\allowbreak{}LOG\_\allowbreak{}TOOL\_\allowbreak{}CONTENT} · \texttt{OTEL\_\allowbreak{}LOG\_\allowbreak{}RAW\_\allowbreak{}API\_\allowbreak{}BODIES} & Content logging into telemetry. Off by default; leave them off unless you have decided otherwise with your privacy function \\
}%
\def\tblbody{\begin{minipage}{\textwidth}
{\footnotesize\begin{tabular}{H{78.7mm}L{64.3mm}}
\toprule
\textbf{Key or variable} & \textbf{Effect} \\
\midrule
\tblrows
\bottomrule\end{tabular}}\end{minipage}}%
\begingroup
\def\sloppy{\tolerance 9999\emergencystretch 3em\hfuzz 200pt\vfuzz 200pt}%
\hbadness=10000\vbadness=10000\hfuzz=200pt\vfuzz=200pt
\global\setbox\tblbox=\hbox{\tblbody}%
\endgroup
\par\addvspace{2.6mm}
\ifdim\dimexpr\ht\tblbox+\dp\tblbox\relax>0.55\textheight
  \typeout{HANDBOOK-TABLE broken \the\dimexpr\ht\tblbox+\dp\tblbox\relax}%
  
  {\footnotesize\begin{longtable}{H{78.7mm}L{64.3mm}}
  \toprule
\textbf{Key or variable} & \textbf{Effect} \\
\midrule\endfirsthead
  \multicolumn{2}{@{}l@{}}{%
  \sffamily\footnotesize\itshape\color{inkgrey}Continued}\\[1.2mm]
  \toprule
\textbf{Key or variable} & \textbf{Effect} \\
\midrule\endhead
  \bottomrule\endfoot
  \bottomrule\endlastfoot
  \tblrows
  \end{longtable}}%
\else
  \typeout{HANDBOOK-TABLE atomic \the\dimexpr\ht\tblbox+\dp\tblbox\relax}%
  \noindent\tblbody
\fi
\par\addvspace{2.6mm}
\endgroup

\FloatBarrier
\setcounter{chapter}{4}
\setcounter{section}{0}
\appendixchapter{E}{Security release checklist}
\FloatBarrier
\section*{Before the task}
\markboth{Before the task}{Before the task}

\begin{codeblock}{7.0}{8.3}
\cl{[~]~Correct~machine,~account,~organisation,~project,~repository~and~branch~confirmed}
\cl{[~]~Working~directory~is~narrow;~not~a~home~directory}
\cl{[~]~Clean~baseline~or~recoverable~backup~exists}
\cl{[~]~No~secret~in~prompts,~in~files~in~scope,~or~in~version~control}
\cl{[~]~Permission~rules~follow~least~privilege,~and~you~have~tested~that~a~deny~actually~blocks}
\cl{[~]~Sandboxing~and~credential~protection~enabled~where~supported}
\cl{[~]~External~services~use~test~or~least-privilege~accounts}
\cl{[~]~Third-party~skills,~plugins,~marketplaces,~MCP~servers,~channels,~hooks~and~scripts~reviewed}
\cl{[~]~Human~approval~gates~written~into~the~plan}
\cl{[~]~Token,~time,~turn~and~agent~bounds~explicit}
\cl{[~]~Fork~mode~and~permission-mode~inheritance~understood~for~any~delegated~reviewer}
\end{codeblock}

\FloatBarrier
\section*{During the task}
\markboth{During the task}{During the task}

\begin{codeblock}{9.0}{10.6}
\cl{[~]~Initial~tool~calls~match~the~intended~direction}
\cl{[~]~Every~permission~request~understood~before~approval}
\cl{[~]~Unexpected~network,~credential,~package~or~account~access~stops~the~task}
\cl{[~]~Browser~actions~stay~within~the~approved~profile~and~domain~set}
\cl{[~]~Parallel~sessions~have~isolated~ownership~or~worktrees}
\cl{[~]~Material~decisions~and~findings~are~being~written~to~files}
\cl{[~]~No~stale~goal~or~loop~is~competing~with~the~current~task}
\cl{[~]~Inter-agent~messages~are~treated~as~coordination,~never~as~approval}
\end{codeblock}

\FloatBarrier
\section*{Before release}
\markboth{Before release}{Before release}

\begin{codeblock}{9.0}{10.6}
\cl{[~]~Diff~reviewed~file~by~file}
\cl{[~]~Automated~checks~actually~ran,~and~their~output~was~inspected}
\cl{[~]~Behaviour~or~rendering~observed~directly}
\cl{[~]~Independent~review~completed,~and~confirmed~to~be~independent}
\cl{[~]~External~side~effects~enumerated,~including~anything~published~or~messaged}
\cl{[~]~Links~and~citations~checked}
\cl{[~]~No~secret~in~output,~logs,~artifacts,~or~transcript~excerpts~quoted~into~files}
\cl{[~]~Publication,~merge,~push,~payment~or~account~change~separately~approved}
\cl{[~]~Rollback~route~and~monitoring~owner~known}
\end{codeblock}

\FloatBarrier
\setcounter{chapter}{5}
\setcounter{section}{0}
\appendixchapter{F}{Threat model and control map}
\begingroup
\def\tblrows{%
Indirect prompt injection & Web page, tool output, PR comment, fire payload, channel message & Isolated web-fetch context; untrusted-input wrappers; classifier review of inter-agent messages & Explicit objective and domain allowlist; human gate on external writes & Chapter 20, Chapter 13, Chapter 24 \\
Credential exposure & \texttt{.\allowbreak{}env} read, command output, transcript, artifact & Deny rules; sandbox credential deny/mask & Least-privilege keys; rotation; output discipline & Chapter 14 \\
Malicious or careless third-party code & Skill, plugin, marketplace, MCP server, channel plugin & Review before install; pinned versions; \texttt{strict\allowbreak{}Known\allowbreak{}Marketplace\allowbreak{}s} & Isolated test profile; \texttt{disable\allowbreak{}Sideload\allowbreak{}Flags} & Chapter 15, Chapter 18, Chapter 13 \\
Unreviewed autonomous action & Routine, \texttt{/\allowbreak{}loop}, goal, auto-fix, agent team & Human gates in the prompt; least-privilege connectors; bounded turns & Notifications; inspection of the first several runs & Chapter 23, Chapter 24, Chapter 21 \\
Loss of oversight through delegation & Fork inheriting authoring context; teammate inheriting permissive parent & Explicit agent definitions; verify independence & \texttt{permission\allowbreak{}Mode}; parent mode discipline & Chapter 16, Chapter 30 \\
Data leaving the boundary & Telemetry, feedback, artifacts, cloud sessions, Remote Control transcript & Egress opt-outs; ZDR where eligible & \texttt{force\allowbreak{}Login\allowbreak{}Org\allowbreak{}UUID}; artifact and sharing controls & Chapter 32 \\
Irreversible external action & Push, publish, payment, deploy, account change & Human gate; branch protection; \texttt{ask} rules & Separate approval step; scoped credentials & Chapter 8, Appendix E \\
Concurrent-edit corruption & Two sessions, one checkout & Worktrees or file ownership & Sequential integration & Chapter 29 \\
Configuration drift and shadow policy & User settings, plugin settings, unselected managed source & \texttt{allow\allowbreak{}Managed*Onl\allowbreak{}y} locks; \texttt{/\allowbreak{}status} verification & \texttt{Config\allowbreak{}Change} hooks; \texttt{claude doct\allowbreak{}or} & Chapter 31 \\
Over-trust in agent self-report & "Done" without evidence & Evidence-based completion conditions & Independent review; release checklist & Chapter 21, Chapter 34 \\
}%
\def\tblbody{\begin{minipage}{\textwidth}
{\scriptsize\begin{tabular}{L{25.3mm}L{21.8mm}L{38.7mm}L{33.7mm}L{13.9mm}}
\toprule
\textbf{Threat} & \textbf{Vector} & \textbf{Primary control} & \textbf{Secondary control} & \textbf{Chapter} \\
\midrule
\tblrows
\bottomrule\end{tabular}}\end{minipage}}%
\begingroup
\def\sloppy{\tolerance 9999\emergencystretch 3em\hfuzz 200pt\vfuzz 200pt}%
\hbadness=10000\vbadness=10000\hfuzz=200pt\vfuzz=200pt
\global\setbox\tblbox=\hbox{\tblbody}%
\endgroup
\par\addvspace{2.6mm}
\ifdim\dimexpr\ht\tblbox+\dp\tblbox\relax>0.55\textheight
  \typeout{HANDBOOK-TABLE broken \the\dimexpr\ht\tblbox+\dp\tblbox\relax}%
  
  {\scriptsize\begin{longtable}{L{25.3mm}L{21.8mm}L{38.7mm}L{33.7mm}L{13.9mm}}
  \toprule
\textbf{Threat} & \textbf{Vector} & \textbf{Primary control} & \textbf{Secondary control} & \textbf{Chapter} \\
\midrule\endfirsthead
  \multicolumn{5}{@{}l@{}}{%
  \sffamily\footnotesize\itshape\color{inkgrey}Continued}\\[1.2mm]
  \toprule
\textbf{Threat} & \textbf{Vector} & \textbf{Primary control} & \textbf{Secondary control} & \textbf{Chapter} \\
\midrule\endhead
  \bottomrule\endfoot
  \bottomrule\endlastfoot
  \tblrows
  \end{longtable}}%
\else
  \typeout{HANDBOOK-TABLE atomic \the\dimexpr\ht\tblbox+\dp\tblbox\relax}%
  \noindent\tblbody
\fi
\par\addvspace{2.6mm}
\endgroup

\FloatBarrier
\setcounter{chapter}{6}
\setcounter{section}{0}
\appendixchapter{G}{Troubleshooting}
\begingroup
\def\tblrows{%
\texttt{claude:\allowbreak{} command no\allowbreak{}t found} & Installer path not on \texttt{PATH} & Restart the terminal; run \texttt{claude doct\allowbreak{}or}; see [102] \\
Login repeats, or the wrong account appears & Stored credentials, or an API key taking precedence over a claude.ai login & \texttt{/\allowbreak{}status}, then \texttt{/\allowbreak{}login}. Check \texttt{ANTHROPIC\_\allowbreak{}API\_\allowbreak{}KEY}, \texttt{ANTHROPIC\_\allowbreak{}AUTH\_\allowbreak{}TOKEN}, \texttt{api\allowbreak{}Key\allowbreak{}Helper} \\
A documented command is missing & Authentication path, provider, plan, or an organisation policy & \texttt{/\allowbreak{}status}; check whether the feature needs a claude.ai login \\
Claude cannot find files & Wrong project folder or working directory & \texttt{pwd}; use \texttt{/\allowbreak{}add-\allowbreak{}dir} only after confirming the target \\
Repeated permission prompts & Rule too narrow, not persistent, or saved in another scope & \texttt{/\allowbreak{}permissions}; inspect the source settings file \\
A \texttt{CLAUDE.\allowbreak{}md} edit had no effect & Instructions are read at session start and cached & Start a new session; see [60] \\
An instruction disappeared after \texttt{/\allowbreak{}compact} & It was given only in conversation & Add it to \texttt{CLAUDE.\allowbreak{}md}; project-root files survive compaction \\
A rule is not loading & Path-scoped rule has not matched a file yet & \texttt{/\allowbreak{}context}; add an \texttt{Instruction\allowbreak{}s\allowbreak{}Loaded} hook \\
Managed policy is not applying & A higher-priority managed source was selected, or the file is invalid & \texttt{/\allowbreak{}status} \texttt{Setting sou\allowbreak{}rces} line; \texttt{claude doct\allowbreak{}or} \\
An MCP server is disconnected & Session did not reload, OAuth expired, or transport changed & New session; \texttt{/\allowbreak{}mcp}; \texttt{claude mcp \allowbreak{}get <name>} \\
A project MCP server never connects & Awaiting approval, or the workspace is untrusted & \texttt{claude mcp \allowbreak{}list} for \texttt{⏸ Pending a\allowbreak{}pproval}; run \texttt{claude} interactively \\
An MCP tool is absent & Tool Search deferred it, or the scope differs & Ask for it explicitly; inspect \texttt{/\allowbreak{}mcp} \\
\texttt{/\allowbreak{}agents} does not open an editor & Behaviour changed in 2.1.198 & Edit \texttt{.\allowbreak{}claude/\allowbreak{}agents/\allowbreak{}} directly \\
A "reviewer" already knows the argument & It was a fork & Check \texttt{CLAUDE\_\allowbreak{}CODE\_\allowbreak{}FORK\_\allowbreak{}SUBAGENT}; use an explicit agent definition \\
A subagent ignores its \texttt{permission\allowbreak{}Mode} & Parent is in \texttt{accept\allowbreak{}Edits}, \texttt{bypass\allowbreak{}Permissions} or \texttt{auto} & Change the parent's mode \\
Claude spawned teammates unasked & Agent teams enabled; a named subagent launches as a teammate & Set \texttt{CLAUDE\_\allowbreak{}CODE\_\allowbreak{}EXPERIMENTA\allowbreak{}L\_\allowbreak{}AGENT\_\allowbreak{}TEAMS=\allowbreak{}0}; check settings precedence \\
A workflow looks expensive & Too many lanes, or duplicated work & \texttt{/\allowbreak{}workflows}, \texttt{/\allowbreak{}usage}; pause; set \texttt{workflow\allowbreak{}Size\allowbreak{}Guideline} \\
A loop did not run on time & Session busy or closed, jitter, expiry, or a missed interval & Avoid \texttt{:\allowbreak{}00} and \texttt{:\allowbreak{}30}; use a durable scheduler if timing matters \\
A loop is running something you did not ask for & A bare \texttt{/\allowbreak{}loop} started the maintenance prompt & \texttt{Esc} while waiting; check for \texttt{loop.\allowbreak{}md} \\
A session vanished from the picker & Its first prompt was \texttt{/\allowbreak{}loop} & \texttt{claude -\allowbreak{}-\allowbreak{}resume <id>}, or \texttt{Ctrl+A} in the picker \\
A routine "succeeded" but did nothing & Green means no infrastructure error & Open the run and read the transcript \\
A routine cannot reach local software & It runs in cloud infrastructure & Use a Desktop scheduled task or a secured local scheduler \\
Chrome will not connect & Extension, native messaging host, or an API-key session & \texttt{/\allowbreak{}chrome} → Reconnect; restart Chrome; confirm \texttt{/\allowbreak{}login} \\
An artifact created a second page & The session had no URL for the original & Attach with \texttt{/\allowbreak{}artifacts}, or pass the URL \\
A secret was printed & A tool read it, or a command echoed it & Stop, rotate, purge, then fix the control (Section 14.3) \\
Two sessions changed the same file & Shared checkout collision & Stop writers; inspect git; move to worktrees \\
}%
\def\tblbody{\begin{minipage}{\textwidth}
{\scriptsize\begin{tabular}{L{29.7mm}L{35.7mm}L{74.4mm}}
\toprule
\textbf{Symptom} & \textbf{Likely cause} & \textbf{Safe next check} \\
\midrule
\tblrows
\bottomrule\end{tabular}}\end{minipage}}%
\begingroup
\def\sloppy{\tolerance 9999\emergencystretch 3em\hfuzz 200pt\vfuzz 200pt}%
\hbadness=10000\vbadness=10000\hfuzz=200pt\vfuzz=200pt
\global\setbox\tblbox=\hbox{\tblbody}%
\endgroup
\par\addvspace{2.6mm}
\ifdim\dimexpr\ht\tblbox+\dp\tblbox\relax>0.55\textheight
  \typeout{HANDBOOK-TABLE broken \the\dimexpr\ht\tblbox+\dp\tblbox\relax}%
  
  {\scriptsize\begin{longtable}{L{29.7mm}L{35.7mm}L{74.4mm}}
  \toprule
\textbf{Symptom} & \textbf{Likely cause} & \textbf{Safe next check} \\
\midrule\endfirsthead
  \multicolumn{3}{@{}l@{}}{%
  \sffamily\footnotesize\itshape\color{inkgrey}Continued}\\[1.2mm]
  \toprule
\textbf{Symptom} & \textbf{Likely cause} & \textbf{Safe next check} \\
\midrule\endhead
  \bottomrule\endfoot
  \bottomrule\endlastfoot
  \tblrows
  \end{longtable}}%
\else
  \typeout{HANDBOOK-TABLE atomic \the\dimexpr\ht\tblbox+\dp\tblbox\relax}%
  \noindent\tblbody
\fi
\par\addvspace{2.6mm}
\endgroup

Use \texttt{claude doct\allowbreak{}or}, \texttt{/\allowbreak{}doctor}, \texttt{/\allowbreak{}hooks}, \texttt{/\allowbreak{}context} and [72] before applying a fix found in a forum, and [103] to look up an exact error string.


\FloatBarrier
\setcounter{chapter}{7}
\setcounter{section}{0}
\appendixchapter{H}{Standards and regulatory crosswalk}
This crosswalk maps the controls described in this book to seventeen external frameworks. It is offered as a starting point for an assurance conversation, not as a certification claim: a control being present is not evidence that it is effective, and mapping is a matter of interpretation. Use it to answer the question an auditor, a regulator or a customer security questionnaire will actually ask — \emph{which of your controls addresses this requirement, and what evidence do you hold?}


\begin{calloutbox}{palegrey}{quoterule}{1.2mm}
\textbf{UNVERIFIED — ISO/IEC 27001:2022 Annex A identifiers.} The NIST, IEEE, SSDF and ISO/IEC 42001 clause identifiers in this crosswalk were confirmed against primary or publisher sources. The Annex A control numbers were not: ISO standards are paywalled, and only secondary summaries were available at the verification date. Confirm them against your own copy of the standard before relying on this mapping in an audit response. Note also that ISO/IEC 27001:2022 carries Amendment 1:2024 [12].
\end{calloutbox}

\FloatBarrier
\section*{Control-to-framework map}
\markboth{Control-to-framework map}{Control-to-framework map}

\begingroup
\def\tblrows{%
Permission rules and modes & Chapter 7 & GOVERN 1.2, MANAGE 2.2 & 8.1 operational planning and control & A.8.2 privileged access; A.8.3 information access restriction & NIST SP 800-207 §3 policy enforcement point [8] \\
Sandboxing and isolation & Chapter 7 & MANAGE 2.3 & 8.1 & A.8.22 segregation; A.8.31 separation of environments & ASD Essential Eight: application control [18] \\
Managed policy tier & Chapter 31 & GOVERN 1.1, 2.1 & 5.2 policy; 9.2 internal audit & A.5.1 policies; A.8.9 configuration management & NIST CSF 2.0 GV.PO, PR.PS [10] \\
Credential handling and rotation & Chapter 14 & MANAGE 2.2 & 8.1 & A.5.17 authentication information; A.8.24 cryptography & Essential Eight: restrict admin privileges [18] \\
Supply-chain review of skills, plugins, MCP, channels & Chapter 15;Chapter 18;Chapter 13;Chapter 27 & MAP 4.1, MANAGE 3.1, 3.2 & 8.3 AI risk treatment & A.5.19–A.5.23 supplier and cloud services; A.8.30 outsourced development & NIST SP 800-218 PO.3, PS.3, PW.4 [9]; SLSA provenance [104] \\
Untrusted-input handling and injection defence & Chapter 20;Chapter 13;Chapter 24 & MAP 5.1, MEASURE 2.7, MANAGE 4.1 & 6.1.2 risk assessment & A.8.26 application security requirements & OWASP LLM01 [4]; MITRE ATLAS AML.T0051 [5]; [22] \\
Human approval gates & Chapter 34 & GOVERN 3.2, MANAGE 4.1 & 8.4 AI system impact assessment & A.5.4 management responsibilities & EU AI Act Art. 14 human oversight [17]; IEEE 7000-2021 [16] \\
Evidence-Gated Delivery and verification & Chapter 34;“Scope and method” & MEASURE 2.5, 2.9, 3.3 & 8.1; 9.1 monitoring & A.8.29 security testing; A.8.32 change management & IEEE Std 1012-2016 V\&V [15]; ISO/IEC 25010 [14] \\
Version control and rollback & Chapter 8 & MANAGE 4.3 & 8.1 & A.8.32 change management; A.8.13 backup & NIST SP 800-218 PS.1, PS.2 [9] \\
Telemetry, audit and traceability & Chapter 33 & MEASURE 1.3, 2.8, MANAGE 4.1 & 9.1; 10.2 nonconformity & A.8.15 logging; A.8.16 monitoring & NIST CSF 2.0 DE.CM, DE.AE [10]; EU AI Act Art. 12 record-keeping [17] \\
Data retention, residency and ZDR & Chapter 32 & GOVERN 6.1, MAP 4.2 & 7.5 documented information & A.5.33 protection of records; A.5.34 privacy and PII & GDPR Arts. 5, 25, 32 [91]; Australian Privacy Principles 8, 11 [90] \\
Threat modelling of the agentic environment & Appendix F & MAP 1.1, 2.3, 5.1 & 6.1.2 & A.8.25 secure development lifecycle & ISO/IEC 23894 §6 [13]; MITRE ATLAS [5] \\
Accessibility of the toolchain & Chapter 33 & GOVERN 5.1, MAP 1.6 & 4.2 interested parties & — & IEEE 7000-2021 value elicitation [16] \\
Re-verification against a moving baseline & Appendix K & MEASURE 2.4, MANAGE 4.2 & 9.3 management review; 10.1 improvement & A.8.8 technical vulnerability management & Secure AI system development, "secure operation" [105] \\
}%
\def\tblbody{\begin{minipage}{\textwidth}\boxcaption{Table 26 — Book controls mapped to external frameworks}
{\scriptsize\begin{tabular}{L{25.7mm}L{17.1mm}L{12.2mm}L{22.2mm}L{27.4mm}L{25.6mm}}
\toprule
\textbf{Control in this book} & \textbf{Chapter} & \textbf{NIST AI RMF 1.0 [6]} & \textbf{ISO/IEC 42001 [11]} & \textbf{ISO/IEC 27001:2022 Annex A [12]} & \textbf{Other} \\
\midrule
\tblrows
\bottomrule\end{tabular}}\end{minipage}}%
\begingroup
\def\sloppy{\tolerance 9999\emergencystretch 3em\hfuzz 200pt\vfuzz 200pt}%
\hbadness=10000\vbadness=10000\hfuzz=200pt\vfuzz=200pt
\global\setbox\tblbox=\hbox{\tblbody}%
\endgroup
\par\addvspace{2.6mm}
\ifdim\dimexpr\ht\tblbox+\dp\tblbox\relax>0.55\textheight
  \typeout{HANDBOOK-TABLE broken \the\dimexpr\ht\tblbox+\dp\tblbox\relax}%
  \tabcaption{Table 26 — Book controls mapped to external frameworks}
  {\scriptsize\begin{longtable}{L{25.7mm}L{17.1mm}L{12.2mm}L{22.2mm}L{27.4mm}L{25.6mm}}
  \toprule
\textbf{Control in this book} & \textbf{Chapter} & \textbf{NIST AI RMF 1.0 [6]} & \textbf{ISO/IEC 42001 [11]} & \textbf{ISO/IEC 27001:2022 Annex A [12]} & \textbf{Other} \\
\midrule\endfirsthead
  \multicolumn{6}{@{}l@{}}{%
  \sffamily\footnotesize\itshape\color{inkgrey}Table 26 — Book controls mapped to external frameworks \textemdash\ continued}\\[1.2mm]
  \toprule
\textbf{Control in this book} & \textbf{Chapter} & \textbf{NIST AI RMF 1.0 [6]} & \textbf{ISO/IEC 42001 [11]} & \textbf{ISO/IEC 27001:2022 Annex A [12]} & \textbf{Other} \\
\midrule\endhead
  \bottomrule\endfoot
  \bottomrule\endlastfoot
  \tblrows
  \end{longtable}}%
\else
  \typeout{HANDBOOK-TABLE atomic \the\dimexpr\ht\tblbox+\dp\tblbox\relax}%
  \noindent\tblbody
\fi
\par\addvspace{2.6mm}
\endgroup

\FloatBarrier
\section*{Reading the map for three common obligations}
\markboth{Reading the map for three common obligations}{Reading the map for three common obligations}

\textbf{EU AI Act human oversight (Art. 14).} Article 14(4) enumerates five capabilities the natural person must be enabled to have: to understand the system's capacities and limitations; to remain aware of automation bias — "the possible tendency of automatically relying or over-relying on the output"; to correctly interpret the output; \textbf{to decide, in any particular situation, not to use the system or to otherwise disregard, override or reverse its output}; and to intervene or interrupt the system through a stop control or equivalent [17]. In the terms of this book: Part 0 addresses capacities and limitations; the evidence-label convention and the UNVERIFIED class address automation bias directly; the completion-condition discipline of Chapter 21 addresses correct interpretation of output; the human release gate of Chapter 34 is precisely the fourth capability, the decision to disregard or override an agent's output; and \texttt{Esc}, \texttt{/\allowbreak{}goal clear}, \texttt{/\allowbreak{}rewind} and the managed \texttt{disable\allowbreak{}Bypass\allowbreak{}Permissions\allowbreak{}Mode} key address intervention and stop. The residual gap is that an agent's external actions — a push, a published artifact, a sent message — are not reversible by any local control, which is why they are gated rather than monitored.


\textbf{ISO/IEC 42001 clause 8.1, operational planning and control.} Clause 8 is \emph{Operation}, and its sub-clauses are 8.1 operational planning and control, 8.2 AI risk assessment, 8.3 AI risk treatment and 8.4 AI system impact assessment [11]. The obligation under 8.1 is to plan, implement and control the processes needed to meet the management system's requirements. The managed-policy tier of Chapter 31 is the only mechanism in the product that survives a user's contrary intent, so any 42001 control that must hold organisation-wide is implemented there or is not implemented. Appendix L sequences this into maturity levels.


\textbf{Secure development (NIST SP 800-218).} SSDF practice groups map cleanly: PO.3 supporting toolchains covers the plugin and MCP review of Chapter 18 and Chapter 13; PS.1–PS.3 protecting software covers version control and provenance in Chapter 8; PW.4 reusing existing software covers the supply-chain gate on third-party skills; and RV.1 identifying vulnerabilities covers \texttt{/\allowbreak{}security-\allowbreak{}review}, the security-guidance plugin and the re-verification procedure [9], [69].


\FloatBarrier
\setcounter{chapter}{8}
\setcounter{section}{0}
\appendixchapter{I}{Glossary}
\begin{deflist}
\item \textbf{Agent} \textemdash\ a model operating in a loop with tools, state and a task.  
\item \textbf{Agent team} \textemdash\ a lead session coordinating peer sessions with a shared task list and mailboxes.  
\item \textbf{Allow / ask / deny} \textemdash\ permission-rule outcomes, evaluated in that reverse order: deny first.  
\item \textbf{Artifact} \textemdash\ a live page published from a session to a URL on claude.ai.  
\item \textbf{Auto memory} \textemdash\ repository-scoped notes Claude maintains for future sessions.  
\item \textbf{Auto mode} \textemdash\ permission mode in which a classifier reviews actions instead of you.  
\item \textbf{Channel} \textemdash\ an MCP server that pushes events into a running session.  
\item \textbf{Checkpoint} \textemdash\ saved conversation and file state used by rewind.  
\item \textbf{Connector} \textemdash\ a claude.ai account integration, distinct from a locally configured MCP server.  
\item \textbf{Context engineering} \textemdash\ deliberate selection, organisation, persistence and removal of context.  
\item \textbf{Effort} \textemdash\ configured reasoning allocation for a model.  
\item \textbf{Fork} \textemdash\ a subagent that inherits the parent conversation; the interactive default.  
\item \textbf{Goal} \textemdash\ a persistent completion condition evaluated after each turn by a separate model.  
\item \textbf{Handoff} \textemdash\ compact state prepared for another session or agent.  
\item \textbf{Headless mode} \textemdash\ non-interactive execution with \texttt{claude -\allowbreak{}p}.  
\item \textbf{Hook} \textemdash\ an executable action bound to a lifecycle event.  
\item \textbf{Human gate} \textemdash\ a point that cannot proceed without explicit human confirmation.  
\item \textbf{Managed settings} \textemdash\ organisation policy applied above every other settings level.  
\item \textbf{MCP} \textemdash\ Model Context Protocol, an open standard for exposing tools and resources to AI clients.  
\item \textbf{Output style} \textemdash\ a modification of the system prompt setting role, tone and format.  
\item \textbf{Prompt injection} \textemdash\ untrusted content attempting to redirect a model or its tool use.  
\item \textbf{Routine} \textemdash\ a saved cloud configuration fired by schedule, API call or GitHub event.  
\item \textbf{Sandbox} \textemdash\ operating-system-enforced filesystem and network constraints on tool execution.  
\item \textbf{Subagent} \textemdash\ a specialised instance delegated a bounded task in its own context window.  
\item \textbf{Surface} \textemdash\ an interface through which Claude Code runs: terminal, IDE, Desktop, web, mobile.  
\item \textbf{Teleport} \textemdash\ pulling a cloud session and its branch into the local terminal.  
\item \textbf{Tool Search} \textemdash\ deferred loading of tool definitions until they are needed.  
\item \textbf{Transcript} \textemdash\ locally stored conversation and tool activity, plaintext JSONL.  
\item \textbf{Ultracode} \textemdash\ \texttt{xhigh} effort combined with automatic dynamic-workflow orchestration.  
\item \textbf{Workflow} \textemdash\ a script that orchestrates subagents, executed by a runtime outside the conversation.  
\item \textbf{Worktree} \textemdash\ a separate git working tree and branch isolating concurrent change.  
\item \textbf{Zero Data Retention} \textemdash\ an enterprise configuration under which inference data is not stored.  
\end{deflist}

\FloatBarrier
\setcounter{chapter}{9}
\setcounter{section}{0}
\appendixchapter{J}{Verification and evidence ledger}
\FloatBarrier
\section*{Verification baseline}
\markboth{Verification baseline}{Verification baseline}

\begingroup
\def\tblrows{%
Installed CLI & \texttt{2.\allowbreak{}1.\allowbreak{}241 (Claude\allowbreak{} Code)}, 23 Aug. 2026; re-verified against \texttt{2.\allowbreak{}1.\allowbreak{}246}, 26 Aug. 2026 & High \\
Official documentation & Pages under \texttt{code.\allowbreak{}claude.\allowbreak{}com/\allowbreak{}docs/\allowbreak{}en/\allowbreak{}}, enumerated from the published index [1] & High for behaviour as of the date \\
Traceable-origin material & Introductory video course, publication date and duration confirmed; used only to attribute the origin of superseded claims [42] & High \\
Practitioner technique & Labelled PRACTITIONER NOTE throughout; judgement, not product behaviour & Medium unless independently verified \\
Third-party components & None installed or endorsed; risk depends on implementation and version & Insufficient until individually reviewed \\
}%
\def\tblbody{\begin{minipage}{\textwidth}
{\footnotesize\begin{tabular}{L{32.8mm}L{80.4mm}L{26.6mm}}
\toprule
\textbf{Item} & \textbf{Observed result} & \textbf{Confidence} \\
\midrule
\tblrows
\bottomrule\end{tabular}}\end{minipage}}%
\begingroup
\def\sloppy{\tolerance 9999\emergencystretch 3em\hfuzz 200pt\vfuzz 200pt}%
\hbadness=10000\vbadness=10000\hfuzz=200pt\vfuzz=200pt
\global\setbox\tblbox=\hbox{\tblbody}%
\endgroup
\par\addvspace{2.6mm}
\ifdim\dimexpr\ht\tblbox+\dp\tblbox\relax>0.55\textheight
  \typeout{HANDBOOK-TABLE broken \the\dimexpr\ht\tblbox+\dp\tblbox\relax}%
  
  {\footnotesize\begin{longtable}{L{32.8mm}L{80.4mm}L{26.6mm}}
  \toprule
\textbf{Item} & \textbf{Observed result} & \textbf{Confidence} \\
\midrule\endfirsthead
  \multicolumn{3}{@{}l@{}}{%
  \sffamily\footnotesize\itshape\color{inkgrey}Continued}\\[1.2mm]
  \toprule
\textbf{Item} & \textbf{Observed result} & \textbf{Confidence} \\
\midrule\endhead
  \bottomrule\endfoot
  \bottomrule\endlastfoot
  \tblrows
  \end{longtable}}%
\else
  \typeout{HANDBOOK-TABLE atomic \the\dimexpr\ht\tblbox+\dp\tblbox\relax}%
  \noindent\tblbody
\fi
\par\addvspace{2.6mm}
\endgroup

\FloatBarrier
\section*{Method}
\markboth{Method}{Method}

Each atomic product claim was adjudicated against the raw text of the primary documentation page that governs it, not against a summary of that page. Where an automated summarisation pass and the raw text disagreed — which happened, notably over the existence of \texttt{/\allowbreak{}sandbox}, \texttt{/\allowbreak{}workflows} and \texttt{-\allowbreak{}-\allowbreak{}worktree} — the raw text was taken as authoritative and the summary discarded. Claims that could be settled only against a running CLI are not asserted \textbf{from documentation}. Where one was later settled by exercising the product, it is recorded below as configuration-specific; where it was not, it appears as UNVERIFIED in the open list at the end of this appendix, with the local check that closes it.


\FloatBarrier
\section*{Misconception ledger}
\markboth{Misconception ledger}{Misconception ledger}

Each row corresponds to an entry in Part 0. The verdict records \emph{how} the claim failed, which is the point of the table: a claim that is simply false and a claim that is true but imprecise call for different responses from a reader who holds it.


\begingroup
\def\tblrows{%
\textbf{REJECTED} & The claim is false. The mechanism does something else entirely, and acting on the claim produces an outcome the reader did not intend. \\
\textbf{CORRECTED} & The claim names a real mechanism but states it wrongly — a detail, a quantifier or a release. The correction is a matter of precision. \\
\textbf{QUALIFIED} & The claim is true only under a condition the reader is unlikely to have met, and false as a general statement. \\
\textbf{REFRAMED} & The claim is factually accurate but its framing understates a risk that governs how the feature should be used. \\
}%
\def\tblbody{\begin{minipage}{\textwidth}
{\footnotesize\begin{tabular}{L{18.4mm}L{124.5mm}}
\toprule
\textbf{Verdict} & \textbf{What it means} \\
\midrule
\tblrows
\bottomrule\end{tabular}}\end{minipage}}%
\begingroup
\def\sloppy{\tolerance 9999\emergencystretch 3em\hfuzz 200pt\vfuzz 200pt}%
\hbadness=10000\vbadness=10000\hfuzz=200pt\vfuzz=200pt
\global\setbox\tblbox=\hbox{\tblbody}%
\endgroup
\par\addvspace{2.6mm}
\ifdim\dimexpr\ht\tblbox+\dp\tblbox\relax>0.55\textheight
  \typeout{HANDBOOK-TABLE broken \the\dimexpr\ht\tblbox+\dp\tblbox\relax}%
  
  {\footnotesize\begin{longtable}{L{18.4mm}L{124.5mm}}
  \toprule
\textbf{Verdict} & \textbf{What it means} \\
\midrule\endfirsthead
  \multicolumn{2}{@{}l@{}}{%
  \sffamily\footnotesize\itshape\color{inkgrey}Continued}\\[1.2mm]
  \toprule
\textbf{Verdict} & \textbf{What it means} \\
\midrule\endhead
  \bottomrule\endfoot
  \bottomrule\endlastfoot
  \tblrows
  \end{longtable}}%
\else
  \typeout{HANDBOOK-TABLE atomic \the\dimexpr\ht\tblbox+\dp\tblbox\relax}%
  \noindent\tblbody
\fi
\par\addvspace{2.6mm}
\endgroup

\begingroup
\def\tblrows{%
M01 & Run \texttt{/\allowbreak{}loop} with no arguments to see your scheduled tasks & REJECTED & [24] & Chapter 23 \\
M02 & Closing the session stops a loop & REJECTED & [24] & Chapter 23 \\
M03 & A subagent gives you an independent second opinion & QUALIFIED & [25] & Chapter 16 \\
M04 & \texttt{.\allowbreak{}env} keeps credentials out of Claude's reach & REJECTED & [26]–[28] & Chapter 14 \\
M05 & Scheduling jitter is unpredictable & CORRECTED & [24] & Chapter 23 \\
M06 & \texttt{CLAUDE\_\allowbreak{}CODE\_\allowbreak{}DISABLE\_\allowbreak{}AUTO\_\allowbreak{}MEMORY} is how you turn auto memory off & CORRECTED & [26] & Chapter 10 \\
M07 & The 30-day cleanup period governs local Claude Code state & CORRECTED & [26], [29] & Chapter 10 \\
M08 & Session history is capped at one month & CORRECTED & [30] & Chapter 12 \\
M09 & \texttt{bypass\allowbreak{}Permissions} is a reasonable productivity setting & REFRAMED & [19], [29], [31] & Chapter 7 \\
M10 & Every MCP tool definition always consumes context & CORRECTED & [32], [33] & Chapter 13 \\
M11 & Copying \texttt{.\allowbreak{}mcp.\allowbreak{}json} transfers an authenticated server & REJECTED & [32] & Chapter 13 \\
M12 & A popular third-party skill is effectively safe & REJECTED & [34] & Chapter 15 \\
M13 & \texttt{/\allowbreak{}agents} opens an agent editor & CORRECTED & [25], [35] & Chapter 16 \\
M14 & Workflows are beta and cannot be monitored in an IDE & CORRECTED & [19] & Chapter 22 \\
M15 & Workflow agents inherit the session's permission mode & CORRECTED & [19] & Chapter 22 \\
M16 & Claude in Chrome is a beta feature & CORRECTED & [21], [36] & Chapter 20 \\
M17 & \texttt{/\allowbreak{}loop} is equivalent to operating-system cron & REJECTED & [24] & Chapter 23 \\
M18 & Plugins are convenient feature packages & REFRAMED & [37], [38] & Chapter 18 \\
M19 & Routines are simply sessions on a schedule & CORRECTED & [20] & Chapter 24 \\
M20 & A green status on a routine run means the task succeeded & CORRECTED & [20] & Chapter 24 \\
M21 & Remote Control keeps everything local & QUALIFIED & [39]–[41] & Chapter 25 \\
}%
\def\tblbody{\begin{minipage}{\textwidth}\boxcaption{Table 27 — Misconception adjudication, with verdict and primary source}
{\footnotesize\begin{tabular}{L{7.6mm}L{71.6mm}L{18.4mm}L{16.9mm}L{18.8mm}}
\toprule
\textbf{} & \textbf{Claim} & \textbf{Verdict} & \textbf{Primary source} & \textbf{Treated in} \\
\midrule
\tblrows
\bottomrule\end{tabular}}\end{minipage}}%
\begingroup
\def\sloppy{\tolerance 9999\emergencystretch 3em\hfuzz 200pt\vfuzz 200pt}%
\hbadness=10000\vbadness=10000\hfuzz=200pt\vfuzz=200pt
\global\setbox\tblbox=\hbox{\tblbody}%
\endgroup
\par\addvspace{2.6mm}
\ifdim\dimexpr\ht\tblbox+\dp\tblbox\relax>0.55\textheight
  \typeout{HANDBOOK-TABLE broken \the\dimexpr\ht\tblbox+\dp\tblbox\relax}%
  \tabcaption{Table 27 — Misconception adjudication, with verdict and primary source}
  {\footnotesize\begin{longtable}{L{7.6mm}L{71.6mm}L{18.4mm}L{16.9mm}L{18.8mm}}
  \toprule
\textbf{} & \textbf{Claim} & \textbf{Verdict} & \textbf{Primary source} & \textbf{Treated in} \\
\midrule\endfirsthead
  \multicolumn{5}{@{}l@{}}{%
  \sffamily\footnotesize\itshape\color{inkgrey}Table 27 — Misconception adjudication, with verdict and primary source \textemdash\ continued}\\[1.2mm]
  \toprule
\textbf{} & \textbf{Claim} & \textbf{Verdict} & \textbf{Primary source} & \textbf{Treated in} \\
\midrule\endhead
  \bottomrule\endfoot
  \bottomrule\endlastfoot
  \tblrows
  \end{longtable}}%
\else
  \typeout{HANDBOOK-TABLE atomic \the\dimexpr\ht\tblbox+\dp\tblbox\relax}%
  \noindent\tblbody
\fi
\par\addvspace{2.6mm}
\endgroup

This table and the register of Table 1 are generated from a single record, so a claim is stated once and cannot drift between them.


\FloatBarrier
\section*{Reference and standards verification}
\markboth{Reference and standards verification}{Reference and standards verification}

A dedicated verification pass on 24 August 2026 checked the two reference classes a link checker cannot see: entries cited by document number or DOI, and assertions about what a standard or regulation actually requires. Twenty-three primary-source checks were made against publisher and issuing-body sources.


\textbf{Confirmed without change.} NIST AI 100-1, AI 600-1, SP 800-207 (including author list), SP 800-218 (authors and date), and CSWP 29; ISO/IEC 25010:2023, 23894:2023 and 42001:2023 as cited; the title, edition and publication date of ISO/IEC 27001; IEEE Std 1012-2016 and 7000-2021, including both DOIs against their IEEE Xplore document numbers; Greshake \emph{et al.}, Perez and Ribeiro, and Amodei \emph{et al.}; \textbf{all twenty-six NIST AI RMF category identifiers} used in the crosswalk; and \textbf{all six SSDF practice identifiers} — PO.3, PS.1, PS.2, PS.3, PW.4 and RV.1.


\textbf{Corrected.}


\begingroup
\def\tblrows{%
V1 & ISO/IEC 42001 clause numbers were wrong in seven crosswalk rows and the worked example. Clause 8 is \emph{Operation}: 8.1 operational planning and control, 8.2 AI risk assessment, 8.3 AI risk treatment, 8.4 AI system impact assessment & Renumbered throughout Appendix H; the sub-clause structure is now stated explicitly \\
V2 & The Article 14 mapping gave four of the five capabilities in Art. 14(4), omitting the decision not to use the system or to disregard, override or reverse its output & Restored, and mapped to the human release gate of Chapter 34 \\
V3 & ISO/IEC 27001:2022 carries Amendment 1:2024 & Noted in the reference entry \\
V4 & The OWASP LLM Top 10 has moved to the OWASP GenAI Security Project, with a 2026 edition published before this book's verification date & Reference updated to the current project, edition and canonical URL \\
V5 & SLSA cited at v1.0; current is v1.2, maintained by a vendor-neutral steering committee within the OpenSSF ecosystem & Reference corrected \\
V6 & Two self-describing counts were stale: the automated check count and the number of frameworks in the crosswalk & Corrected to 96 and seventeen \\
V7 & One internal reference pointed at the pre-renumbering appendix letter & Corrected \\
}%
\def\tblbody{\begin{minipage}{\textwidth}
{\footnotesize\begin{tabular}{L{4.8mm}L{99.7mm}L{35.3mm}}
\toprule
\textbf{ID} & \textbf{Finding} & \textbf{Resolution} \\
\midrule
\tblrows
\bottomrule\end{tabular}}\end{minipage}}%
\begingroup
\def\sloppy{\tolerance 9999\emergencystretch 3em\hfuzz 200pt\vfuzz 200pt}%
\hbadness=10000\vbadness=10000\hfuzz=200pt\vfuzz=200pt
\global\setbox\tblbox=\hbox{\tblbody}%
\endgroup
\par\addvspace{2.6mm}
\ifdim\dimexpr\ht\tblbox+\dp\tblbox\relax>0.55\textheight
  \typeout{HANDBOOK-TABLE broken \the\dimexpr\ht\tblbox+\dp\tblbox\relax}%
  
  {\footnotesize\begin{longtable}{L{4.8mm}L{99.7mm}L{35.3mm}}
  \toprule
\textbf{ID} & \textbf{Finding} & \textbf{Resolution} \\
\midrule\endfirsthead
  \multicolumn{3}{@{}l@{}}{%
  \sffamily\footnotesize\itshape\color{inkgrey}Continued}\\[1.2mm]
  \toprule
\textbf{ID} & \textbf{Finding} & \textbf{Resolution} \\
\midrule\endhead
  \bottomrule\endfoot
  \bottomrule\endlastfoot
  \tblrows
  \end{longtable}}%
\else
  \typeout{HANDBOOK-TABLE atomic \the\dimexpr\ht\tblbox+\dp\tblbox\relax}%
  \noindent\tblbody
\fi
\par\addvspace{2.6mm}
\endgroup

\textbf{Recorded as unverified.} The ISO/IEC 27001:2022 Annex A control identifiers used in the crosswalk could not be checked: ISO standards are paywalled and only secondary summaries were reachable. A caveat now appears in Appendix H rather than an implication that they were verified. The page range for Greshake \emph{et al.} is unconfirmed — the ACM Digital Library refuses automated clients — though title, venue and DOI are confirmed. The official expansion of the MITRE ATLAS acronym varies between the project site and its repository.


\textbf{A methodological result.} During this pass an automated summarisation of a national cyber security agency's page asserted that an international secure-AI guideline had been issued by one agency alone. The issuing bodies' own announcements record a joint publication by two national agencies with seventeen further international partners; the original citation was correct and the summary was wrong. This is the second confirmed false negative from automated summarisation in the preparation of this book, and it is why “Scope and method” treats raw-source adjudication as a rule rather than a preference.


\FloatBarrier
\section*{Corrections applied at re-verification}
\markboth{Corrections applied at re-verification}{Corrections applied at re-verification}

The re-verification of 26 August 2026 against Claude Code 2.1.246 changed the statements below. Each was adjudicated against the raw text of the page named, and each is a case where the 23 August 2026 adjudication had reached a different answer or had stopped short of a condition the page attaches. They are listed here because the verification statement in the front matter tells the reader that anything \textbf{not} listed rests on the baseline adjudication; that sentence is only true if this table exists.


\begingroup
\def\tblrows{%
R1 & Chapter 7 stated that \texttt{bypass\allowbreak{}Permissions} propagates into a workflow's subagents. They always run in \texttt{accept\allowbreak{}Edits} and inherit your tool allowlist, whatever the session's mode & [19] \\
R2 & \texttt{/\allowbreak{}stats} is a documented alias of \texttt{/\allowbreak{}usage} that opens on the Stats tab, closing the last of the three UNVERIFIED items this book shipped with & [35] \\
R3 & The \texttt{Shift+Tab} cycle was recorded from a session capture and was wrong in order and in kind: the base cycle is three modes, and two more slot in only when each is enabled & [31] \\
R4 & The \texttt{Large workf\allowbreak{}low} warning fires above 25 agents rather than at 25, and a size guideline chosen by the reader replaces that number with its own & [19] \\
R5 & A resumed session has three permission-mode exceptions rather than one, conditions on restoring the model, and tasks and mid-session directories that are not restored at all & [30] \\
R6 & \texttt{auto\allowbreak{}Memory\allowbreak{}Directory} set in a project settings file is honoured only under the workspace-trust rule that governs hooks in settings files & [26] \\
R7 & A parent in \texttt{auto} also defeats a subagent's own \texttt{permission\allowbreak{}Mode}, and \texttt{auto} is the mode sessions start in on Pro, Max and Team plans & [25] \\
}%
\def\tblbody{\begin{minipage}{\textwidth}\boxcaption{Table 28 — Corrections applied at the 26 August 2026 re-verification}
{\footnotesize\begin{tabular}{L{4.1mm}L{118.8mm}L{16.9mm}}
\toprule
\textbf{} & \textbf{What the re-verification changed} & \textbf{Primary source} \\
\midrule
\tblrows
\bottomrule\end{tabular}}\end{minipage}}%
\begingroup
\def\sloppy{\tolerance 9999\emergencystretch 3em\hfuzz 200pt\vfuzz 200pt}%
\hbadness=10000\vbadness=10000\hfuzz=200pt\vfuzz=200pt
\global\setbox\tblbox=\hbox{\tblbody}%
\endgroup
\par\addvspace{2.6mm}
\ifdim\dimexpr\ht\tblbox+\dp\tblbox\relax>0.55\textheight
  \typeout{HANDBOOK-TABLE broken \the\dimexpr\ht\tblbox+\dp\tblbox\relax}%
  \tabcaption{Table 28 — Corrections applied at the 26 August 2026 re-verification}
  {\footnotesize\begin{longtable}{L{4.1mm}L{118.8mm}L{16.9mm}}
  \toprule
\textbf{} & \textbf{What the re-verification changed} & \textbf{Primary source} \\
\midrule\endfirsthead
  \multicolumn{3}{@{}l@{}}{%
  \sffamily\footnotesize\itshape\color{inkgrey}Table 28 — Corrections applied at the 26 August 2026 re-verification \textemdash\ continued}\\[1.2mm]
  \toprule
\textbf{} & \textbf{What the re-verification changed} & \textbf{Primary source} \\
\midrule\endhead
  \bottomrule\endfoot
  \bottomrule\endlastfoot
  \tblrows
  \end{longtable}}%
\else
  \typeout{HANDBOOK-TABLE atomic \the\dimexpr\ht\tblbox+\dp\tblbox\relax}%
  \noindent\tblbody
\fi
\par\addvspace{2.6mm}
\endgroup

\FloatBarrier
\section*{Closed by direct observation}
\markboth{Closed by direct observation}{Closed by direct observation}

Three items carried as UNVERIFIED in earlier drafts were reopened by exercising the product directly on Claude Code 2.1.246 on 26 August 2026, and settled against the documentation that governs each. Each verdict is \textbf{configuration-specific}: it records what one installation did on one date, on the surface named, and asserts nothing more general. What each observation did not reach stays in the open list below.


\begingroup
\def\tblrows{%
The \texttt{/\allowbreak{}effort} vocabulary & Seven levels, five accepted everywhere and two differing by surface (Table 3) & Terminal CLI, VS Code extension, and the \texttt{claude -\allowbreak{}-\allowbreak{}effort} flag & CONFIRMED (configuration-specific) \\
The \texttt{Shift+Tab} permission-mode cycle & The cycle never reaches \texttt{dont\allowbreak{}Ask}; order and conditions are given in Table 5 & Terminal CLI, then adjudicated against [31] & CONFIRMED \\
The \texttt{claude -\allowbreak{}-\allowbreak{}help} flag set & 62 long options in option position and 13 subcommands, identical between 2.1.241 and 2.1.246 & The pinned binaries of both releases & CONFIRMED \\
}%
\def\tblbody{\begin{minipage}{\textwidth}\boxcaption{Table 29 — Items closed by direct observation on Claude Code 2.1.246, 26 August 2026}
{\footnotesize\begin{tabular}{H{30.7mm}L{34.0mm}L{26.7mm}H{45.1mm}}
\toprule
\textbf{Item} & \textbf{Observation} & \textbf{Surface observed} & \textbf{Verdict} \\
\midrule
\tblrows
\bottomrule\end{tabular}}\end{minipage}}%
\begingroup
\def\sloppy{\tolerance 9999\emergencystretch 3em\hfuzz 200pt\vfuzz 200pt}%
\hbadness=10000\vbadness=10000\hfuzz=200pt\vfuzz=200pt
\global\setbox\tblbox=\hbox{\tblbody}%
\endgroup
\par\addvspace{2.6mm}
\ifdim\dimexpr\ht\tblbox+\dp\tblbox\relax>0.55\textheight
  \typeout{HANDBOOK-TABLE broken \the\dimexpr\ht\tblbox+\dp\tblbox\relax}%
  \tabcaption{Table 29 — Items closed by direct observation on Claude Code 2.1.246, 26 August 2026}
  {\footnotesize\begin{longtable}{H{30.7mm}L{34.0mm}L{26.7mm}H{45.1mm}}
  \toprule
\textbf{Item} & \textbf{Observation} & \textbf{Surface observed} & \textbf{Verdict} \\
\midrule\endfirsthead
  \multicolumn{4}{@{}l@{}}{%
  \sffamily\footnotesize\itshape\color{inkgrey}Table 29 — Items closed by direct observation on Claude Code 2.1.246, 26 August 2026 \textemdash\ continued}\\[1.2mm]
  \toprule
\textbf{Item} & \textbf{Observation} & \textbf{Surface observed} & \textbf{Verdict} \\
\midrule\endhead
  \bottomrule\endfoot
  \bottomrule\endlastfoot
  \tblrows
  \end{longtable}}%
\else
  \typeout{HANDBOOK-TABLE atomic \the\dimexpr\ht\tblbox+\dp\tblbox\relax}%
  \noindent\tblbody
\fi
\par\addvspace{2.6mm}
\endgroup

The \texttt{claude -\allowbreak{}-\allowbreak{}help} row closes an item that named a single release. Comparing the two releases directly was cheaper than re-capturing one, and the zero delta is the more useful result: it records that the flag surface did not move across the interval, which is what a reader re-verifying against a later release actually needs to know.


\FloatBarrier
\section*{A help string is not an enumeration}
\markboth{A help string is not an enumeration}{A help string is not an enumeration}

Exercising those two vocabularies produced a result worth recording as method rather than as content. \textbf{Twice, the product accepted a value its own help text does not list.} \texttt{claude -\allowbreak{}-\allowbreak{}effort} accepts \texttt{ultracode} while \texttt{-\allowbreak{}-\allowbreak{}help} enumerates five levels. \texttt{claude -\allowbreak{}-\allowbreak{}permission-\allowbreak{}mode} accepts \texttt{default} while the error message it prints when refusing a value enumerates six and omits it — so the under-reporting appeared in the most authoritative-looking enumeration the command line offers, the one shown at the moment of refusal.


The second case is the more instructive, because adjudication later explained it. \texttt{manual} is a documented alias for \texttt{default} (2.1.200 or later), and the validator lists the alias rather than the canonical value that settings files and the SDK use [31]. The enumeration was never wrong; it was incomplete in a way no amount of staring at it could reveal. Had this book treated the omission of \texttt{default} as evidence that \texttt{default} was not accepted, it would have printed a false statement drawn from an authoritative-looking source — which is precisely the failure the rule below prevents.


The consequence is a rule rather than a caveat. \textbf{A help string is evidence that a value is supported; it is never evidence that an unlisted value is refused.} Every statement in this book of the form \emph{the accepted values are X, Y and Z} is therefore stated for the surface it was observed on, and an absence from an enumeration is recorded as unestablished rather than as a refusal — unless the refusal was itself observed, as \texttt{claude -\allowbreak{}-\allowbreak{}effort auto} was. This strengthens the UNVERIFIED discipline of “Scope and method” rather than weakening it: the two vocabularies were larger than the documentation of them, and only exercising the product showed it.


\FloatBarrier
\section*{Residual uncertainty}
\markboth{Residual uncertainty}{Residual uncertainty}

Claude Code changes weekly. Model availability, interface labels, plan entitlements and preview features may all have moved since the verification date. Third-party MCP servers, plugins, skills and channels require separate per-version assessment; this book does not assert that the absence of a documented incident makes any component safe. Two classes of statement in this edition are explicitly weaker than the rest: anything labelled PRACTITIONER NOTE, which is judgement rather than product behaviour, and anything labelled UNVERIFIED, listed below.


\textbf{Open UNVERIFIED items at publication:}


\begingroup
\def\tblrows{%
The exact \texttt{/\allowbreak{}effort} menu on a given plan & Menu contents are plan- and model-dependent. Table 3 closes this for one configuration only & Observe \texttt{/\allowbreak{}effort} locally on your own plan and model \\
Whether the terminal \texttt{/\allowbreak{}effort} selector offers \texttt{auto} & The selector was observed to exist and to accept \texttt{ultracode}; its full contents were not captured & Run a bare \texttt{/\allowbreak{}effort} in the terminal and read the selector \\
Whether the VS Code \texttt{/\allowbreak{}effort} accepts \texttt{ultracode} & Absent from that surface's usage string, which — see above — is not evidence of refusal & Type \texttt{/\allowbreak{}effort ultr\allowbreak{}acode} in the extension and observe the result \\
}%
\def\tblbody{\begin{minipage}{\textwidth}
{\footnotesize\begin{tabular}{L{31.4mm}L{66.2mm}L{42.2mm}}
\toprule
\textbf{Item} & \textbf{Why it is open} & \textbf{How to close it} \\
\midrule
\tblrows
\bottomrule\end{tabular}}\end{minipage}}%
\begingroup
\def\sloppy{\tolerance 9999\emergencystretch 3em\hfuzz 200pt\vfuzz 200pt}%
\hbadness=10000\vbadness=10000\hfuzz=200pt\vfuzz=200pt
\global\setbox\tblbox=\hbox{\tblbody}%
\endgroup
\par\addvspace{2.6mm}
\ifdim\dimexpr\ht\tblbox+\dp\tblbox\relax>0.55\textheight
  \typeout{HANDBOOK-TABLE broken \the\dimexpr\ht\tblbox+\dp\tblbox\relax}%
  
  {\footnotesize\begin{longtable}{L{31.4mm}L{66.2mm}L{42.2mm}}
  \toprule
\textbf{Item} & \textbf{Why it is open} & \textbf{How to close it} \\
\midrule\endfirsthead
  \multicolumn{3}{@{}l@{}}{%
  \sffamily\footnotesize\itshape\color{inkgrey}Continued}\\[1.2mm]
  \toprule
\textbf{Item} & \textbf{Why it is open} & \textbf{How to close it} \\
\midrule\endhead
  \bottomrule\endfoot
  \bottomrule\endlastfoot
  \tblrows
  \end{longtable}}%
\else
  \typeout{HANDBOOK-TABLE atomic \the\dimexpr\ht\tblbox+\dp\tblbox\relax}%
  \noindent\tblbody
\fi
\par\addvspace{2.6mm}
\endgroup

\FloatBarrier
\setcounter{chapter}{10}
\setcounter{section}{0}
\appendixchapter{K}{Re-verification and version decay}
\FloatBarrier
\section*{Why this appendix exists}
\markboth{Why this appendix exists}{Why this appendix exists}

Between the release the course was recorded against and the release this book was verified against, documented behaviour changed at 2.1.154, 2.1.160, 2.1.178, 2.1.181, 2.1.186, 2.1.193, 2.1.196, 2.1.198, 2.1.199, 2.1.200, 2.1.202, 2.1.203, 2.1.205, 2.1.206, 2.1.207, 2.1.208, 2.1.209, 2.1.210, 2.1.211, 2.1.213, 2.1.214, 2.1.216, 2.1.217, 2.1.219, 2.1.221, 2.1.223, 2.1.225, 2.1.227, 2.1.228, 2.1.229, 2.1.232, 2.1.234, 2.1.236, 2.1.237 and 2.1.239 — thirty-five release-specific behaviour notes appear in the pages cited by this book, which is every release this book stamps inline and is checked against them. A reference work in this domain has a half-life measured in weeks, and the honest response is a procedure rather than a disclaimer.


\FloatBarrier
\section*{The procedure}
\markboth{The procedure}{The procedure}

Perform this quarterly, or before relying on the book for a consequential decision.


\begin{enumerate}
\item \textbf{Establish the delta.} Read the changelog and the weekly digest from the verification date forward [36], [106].
\item \textbf{Re-fetch the index.} Compare the current documentation index against the reference list of this edition [1]. New pages indicate new capability the book does not cover; removed pages indicate deprecation.
\item \textbf{Re-verify the version-stamped claims.} Every statement in this book carrying a release number is a candidate for change.
\item \textbf{Re-run the local evidence script} below, and reconcile it with Appendix A and Appendix B.
\item \textbf{Re-test the security controls that matter to you.} A deny rule that blocked last quarter should be tested again, not assumed.
\item \textbf{Update the ledger.} Record what changed, in Appendix J, with the source and the date.
\end{enumerate}

\FloatBarrier
\section*{Local evidence script}
\markboth{Local evidence script}{Local evidence script}

Run this on the machine whose behaviour you want to record. It writes a plain-text bundle you can attach to a review. It reads and reports; it changes nothing.


\begin{codeblock}{6.0}{7.1}
\cl{\#!/usr/bin/env~bash}
\cl{\#~collect-claude-evidence.sh~—~record~the~locally~observable~state~of~a~Claude~Code~install.}
\cl{set~-uo~pipefail}
\cl{OUT="claude-evidence-\$(date~-u~+\%Y\%m\%dT\%H\%M\%SZ).txt"}
\cl{\{}
\cl{~~echo~"===~collected~(UTC)~===";~date~-u}
\cl{~~echo;~echo~"===~claude~--version~===";~claude~--version~2>\&1}
\cl{~~echo;~echo~"===~claude~doctor~===";~claude~doctor~2>\&1}
\cl{~~echo;~echo~"===~claude~--help~===";~claude~--help~2>\&1}
\cl{~~echo;~echo~"===~claude~mcp~--help~===";~claude~mcp~--help~2>\&1}
\cl{~~echo;~echo~"===~claude~mcp~list~===";~claude~mcp~list~2>\&1}
\cl{~~echo;~echo~"===~claude~plugin~--help~===";~claude~plugin~--help~2>\&1}
\cl{~~echo;~echo~"===~settings~files~present~==="}
\cl{~~for~f~in~\textasciitilde{}/.claude/settings.json~.claude/settings.json~.claude/settings.local.json~\textbackslash{}}
\cl{~~~~~~~~~~~/etc/claude-code/managed-settings.json~\textbackslash{}}
\cl{~~~~~~~~~~~"/Library/Application~Support/ClaudeCode/managed-settings.json";~do}
\cl{~~~~[~-f~"\$f"~]~\&\&~\{~echo~"---~\$f";~sed~-e~\textquotesingle{}s/\textbackslash{}("[A-Za-z\_]*[Tt]oken"[[:space:]]*:[[:space:]]*\textbackslash{}).*/\textbackslash{}1"<redacted>"/\textquotesingle{}~\textbackslash{}}
\cl{~~~~~~~~~~~~~~~~~~~~~~~~~~~~~~~~~~~~~~-e~\textquotesingle{}s/\textbackslash{}("[A-Za-z\_]*[Kk]ey"[[:space:]]*:[[:space:]]*\textbackslash{}).*/\textbackslash{}1"<redacted>"/\textquotesingle{}~"\$f";~\}}
\cl{~~done}
\cl{~~echo;~echo~"===~instruction~and~memory~files~==="}
\cl{~~ls~-la~\textasciitilde{}/.claude/rules/~.claude/rules/~2>/dev/null}
\cl{~~ls~-la~\textasciitilde{}/.claude/projects/*/memory/~2>/dev/null~|~head~-40}
\cl{~~echo;~echo~"===~skills,~agents,~workflows~==="}
\cl{~~ls~-la~\textasciitilde{}/.claude/skills/~.claude/skills/~\textasciitilde{}/.claude/agents/~.claude/agents/~\textbackslash{}}
\cl{~~~~~~~~~\textasciitilde{}/.claude/workflows/~.claude/workflows/~2>/dev/null}
\cl{\}~>~"\$OUT"~2>\&1}
\cl{echo~"Wrote~\$OUT~—~review~it~for~secrets~before~sharing."}
\end{codeblock}

Then, inside a session, capture what only the running product can tell you: \texttt{/\allowbreak{}status}, \texttt{/\allowbreak{}permissions}, \texttt{/\allowbreak{}context}, \texttt{/\allowbreak{}memory}, \texttt{/\allowbreak{}hooks}, \texttt{/\allowbreak{}mcp}, \texttt{/\allowbreak{}usage}, and the \texttt{/\allowbreak{}effort} menu. Use \texttt{/\allowbreak{}export} to save the result [30], [35].


\begin{calloutbox}{palebrass}{brassdark}{2.0mm}
\textbf{CAUTION:} the redaction in the script above is a convenience, not a guarantee. Read the bundle before you send it anywhere.
\end{calloutbox}

\FloatBarrier
\setcounter{chapter}{11}
\setcounter{section}{0}
\appendixchapter{L}{Adoption maturity model}
A five-level model for organisational adoption of agentic development tooling. Each level states the controls that must be \emph{enforced} rather than requested, the evidence that demonstrates the level is real, and the role accountable for it. The model is the author's; it is offered as a sequencing aid, not as a standard.


The organising principle is stated in the executive summary and repeated here because it is the only rule that matters: \textbf{adopt the capability at the pace at which you can evidence its output.} A level is reached when the evidence exists, not when the intention is announced.


\begingroup
\def\tblrows{%
0 & \textbf{Unmanaged} & None. Individuals install and configure independently; defaults apply. & None available. The organisation cannot state which permission modes, connectors or plugins are in use. & Nobody \\
1 & \textbf{Visible} & Inventory only. No policy is deployed, but usage is known. & A list of users, surfaces and plans. \texttt{/\allowbreak{}status} output sampled across the fleet. Telemetry enabled with content logging off. & Engineering manager \\
2 & \textbf{Bounded} & Managed settings deployed and verified: deny rules for credential paths, \texttt{disable\allowbreak{}Bypass\allowbreak{}Permissions\allowbreak{}Mode}, \texttt{allow\allowbreak{}Managed\allowbreak{}Permission\allowbreak{}Rules\allowbreak{}Only}, an MCP allowlist, a marketplace policy. Sandboxing enabled where the platform supports it. & \texttt{/\allowbreak{}status} shows the intended managed source on sampled machines. \texttt{claude doct\allowbreak{}or} shows no dropped entries. A tested deny rule demonstrably blocks. & Security engineering \\
3 & \textbf{Evidenced} & Level 2, plus: telemetry exported to a collector with retention defined; audit review of \texttt{tool\_\allowbreak{}decision}, \texttt{permission\_\allowbreak{}mode\_\allowbreak{}changed}, \texttt{plugin\_\allowbreak{}installed} and \texttt{mcp\_\allowbreak{}server\_\allowbreak{}connection} events; a release gate for anything an agent publishes, merges or sends. & Traceability from a single prompt identifier to every action it caused. A completed release checklist per consequential deliverable. A named reviewer independent of the author. & Head of engineering, with security assurance \\
4 & \textbf{Governed} & Level 3, plus: data residency and retention decided per cohort with \texttt{force\allowbreak{}Login\allowbreak{}Org\allowbreak{}UUID} deployed; autonomous execution — routines, unattended goals, agent teams — permitted only under named approval; supply-chain review with pinned versions and a re-review trigger; accessibility position stated. & A control-to-framework map maintained (Appendix H). A quarterly re-verification record (Appendix K). An incident procedure exercised at least once. A register of approved autonomous workloads with owners. & Executive risk owner, reported to the board \\
}%
\def\tblbody{\begin{minipage}{\textwidth}\boxcaption{Table 30 — Adoption maturity levels for agentic development tooling}
{\scriptsize\begin{tabular}{L{9.7mm}L{14.8mm}L{52.8mm}L{34.6mm}L{21.4mm}}
\toprule
\textbf{Level} & \textbf{Name} & \textbf{Enforced controls} & \textbf{Evidence that the level is real} & \textbf{Accountable} \\
\midrule
\tblrows
\bottomrule\end{tabular}}\end{minipage}}%
\begingroup
\def\sloppy{\tolerance 9999\emergencystretch 3em\hfuzz 200pt\vfuzz 200pt}%
\hbadness=10000\vbadness=10000\hfuzz=200pt\vfuzz=200pt
\global\setbox\tblbox=\hbox{\tblbody}%
\endgroup
\par\addvspace{2.6mm}
\ifdim\dimexpr\ht\tblbox+\dp\tblbox\relax>0.55\textheight
  \typeout{HANDBOOK-TABLE broken \the\dimexpr\ht\tblbox+\dp\tblbox\relax}%
  \tabcaption{Table 30 — Adoption maturity levels for agentic development tooling}
  {\scriptsize\begin{longtable}{L{9.7mm}L{14.8mm}L{52.8mm}L{34.6mm}L{21.4mm}}
  \toprule
\textbf{Level} & \textbf{Name} & \textbf{Enforced controls} & \textbf{Evidence that the level is real} & \textbf{Accountable} \\
\midrule\endfirsthead
  \multicolumn{5}{@{}l@{}}{%
  \sffamily\footnotesize\itshape\color{inkgrey}Table 30 — Adoption maturity levels for agentic development tooling \textemdash\ continued}\\[1.2mm]
  \toprule
\textbf{Level} & \textbf{Name} & \textbf{Enforced controls} & \textbf{Evidence that the level is real} & \textbf{Accountable} \\
\midrule\endhead
  \bottomrule\endfoot
  \bottomrule\endlastfoot
  \tblrows
  \end{longtable}}%
\else
  \typeout{HANDBOOK-TABLE atomic \the\dimexpr\ht\tblbox+\dp\tblbox\relax}%
  \noindent\tblbody
\fi
\par\addvspace{2.6mm}
\endgroup

\FloatBarrier
\section*{Level transitions, and the failure at each one}
\markboth{Level transitions, and the failure at each one}{Level transitions, and the failure at each one}

\textbf{0 → 1.} The failure is believing an inventory exists because a procurement record exists. Ask for \texttt{/\allowbreak{}status} output, not a licence count.


\textbf{1 → 2.} The failure is writing standards into \texttt{CLAUDE.\allowbreak{}md}. That documents an intention; a managed settings file deploys a control (Figure 1). The second failure is deploying a policy that is silently ignored because a higher-priority managed source is present (Section 31.1).


\textbf{2 → 3.} The failure is enabling telemetry without deciding who reads it. An unread audit trail is an expense, not a control. The second failure is an "independent" reviewer that is a fork (Section 16.1).


\textbf{3 → 4.} The failure is permitting autonomous execution because it is available rather than because a named person accepted the residual risk. A routine runs with your connectors, under your identity, with no approval prompt (Chapter 24); somebody must own that.


\textbf{A note on speed.} Level 4 is not the goal for every organisation. A two-person research group operating at Level 2 with disciplined release gates is in better shape than a large enterprise at Level 3 that has never tested a deny rule. The model sequences capability against evidence; it does not argue that more governance is always better.


\FloatBarrier
\setcounter{chapter}{12}
\setcounter{section}{0}
\appendixchapter{M}{Worked solutions to the exercises}

Each of the thirteen exercises has one worked solution here. They are defensible answers, not the only answers: every solution shows a route that survives inspection and states the evidence that would demonstrate the exercise was completed correctly rather than merely attempted. The exercises were written so that the failure modes are the lesson, so read the solution after your own attempt, not before it.


Three conventions run through all of them. Where an instruction can be \emph{enforced} rather than \emph{requested}, the solution enforces it — a permission rule, a hook, a sandbox — because a sentence in a prompt is a preference and a deny rule is a control (Chapter 7). Where a solution produces an artefact, the artefact has a stated schema, so that "done" is decidable by someone who was not watching. And every solution ends with a check that can fail; a check that cannot fail is a description, not a test.


\FloatBarrier
\section*{Chapter 1 — a four-element brief you can run}
\markboth{Chapter 1 — a four-element brief you can run}{Chapter 1 — a four-element brief you can run}

\textbf{Asked:} write a four-line brief in Outcome / Context / Boundaries / Evidence form and keep it for Chapter 5.


The four elements each close a specific failure mode: an unbounded outcome, an assumed context, an unstated prohibition, and an unverifiable result. A worked version, for a practice workspace holding \texttt{notes/\allowbreak{}}, \texttt{sources/\allowbreak{}} and an empty \texttt{drafts/\allowbreak{}}:


\begin{codeblock}{7.0}{8.3}
\cl{Outcome:~~~~A~brief~of~at~most~600~words,~built~only~from~the~files~in~notes/,~answering~two}
\cl{~~~~~~~~~~~~questions:~what~does~this~material~claim,~and~what~does~it~not~establish?}
\cl{Context:~~~~notes/~holds~fourteen~markdown~files~written~between~March~and~August~2026.}
\cl{~~~~~~~~~~~~sources/~holds~the~PDFs~they~cite.~Treat~notes/~as~claims~and~sources/~as~evidence.}
\cl{~~~~~~~~~~~~Identify~evidence~gaps;~do~not~fill~them.}
\cl{Boundaries:~Read-only~outside~drafts/.~Do~not~edit,~move,~rename~or~delete~anything~in~notes/}
\cl{~~~~~~~~~~~~or~sources/.~No~network~access.~No~shell~command~that~writes~outside~drafts/.}
\cl{Evidence:~~~Write~drafts/brief-v1.md~containing~a~claim-to-source~table~with~one~row~per~claim}
\cl{~~~~~~~~~~~~and~the~columns:~Claim~|~notes/~file~and~line~|~sources/~file~and~page~|~Supported.}
\cl{~~~~~~~~~~~~Record~every~unsupported~claim~as~UNSUPPORTED~rather~than~omitting~it.}
\end{codeblock}

Three properties make this runnable rather than aspirational. The outcome carries a length bound \emph{and} a question, so completion is decidable. The boundary prohibits writes by name rather than asking for care, and it names the operations — edit, move, rename, delete — that a request to "not change anything" leaves ambiguous. The evidence clause fixes the artefact and its schema, so the session cannot end with a plausible summary and no audit trail. Note the deliberate instruction to record gaps rather than close them: an agent asked for a complete brief will produce one, and the missing evidence is the finding you actually want.


Enforce the boundary rather than trusting it. In \texttt{.\allowbreak{}claude/\allowbreak{}settings.\allowbreak{}json}:


\begin{codeblock}{9.0}{10.6}
\cl{\{}
\cl{~~"permissions":~\{}
\cl{~~~~"deny":~[}
\cl{~~~~~~"Edit(./notes/**)",}
\cl{~~~~~~"Edit(./sources/**)",}
\cl{~~~~~~"WebFetch",}
\cl{~~~~~~"WebSearch"}
\cl{~~~~]}
\cl{~~\}}
\cl{\}}
\end{codeblock}

\textbf{Check.} Ask one question of the brief: could two competent people disagree about whether it was completed? If they could, the element that permits the disagreement is the one to tighten. Then confirm the control fires — ask for an edit inside \texttt{notes/\allowbreak{}} and watch it refused. A deny rule you have never seen block anything is a rule you have not tested (Section 7.2).


\FloatBarrier
\section*{Chapter 3 — a read-only inspection in Plan mode}
\markboth{Chapter 3 — a read-only inspection in Plan mode}{Chapter 3 — a read-only inspection in Plan mode}

\textbf{Asked:} in Plan mode, have the workspace inspected and a smallest safe next step proposed, without modification and without leaving the folder.


Run the inspection from the practice folder and nowhere else. The prompt in the exercise is already well formed; what the exercise tests is whether you can tell the difference between a plan that respected the boundary and one that merely reported respecting it.


A good returned plan has four properties. It describes structure it could only have learned by reading — file counts, naming patterns, date ranges — rather than restating your prompt back. It separates what is missing into things you must decide and things it can determine. Its proposed next step is genuinely small: one file, one directory, one reversible operation. And it states what it did \emph{not} look at, which is the part most plans omit.


The mistake the exercise is built around is accepting a plan because it is well written. Plan mode constrains what Claude may do, not what it may claim (Section 6.1). Verify the boundary from outside the transcript:


\begin{codeblock}{9.0}{10.6}
\cl{\#~Nothing~under~the~workspace~changed~during~the~inspection.}
\cl{find~.~-newermt~\textquotesingle{}-15~minutes\textquotesingle{}~-type~f~-not~-path~\textquotesingle{}./.git/*\textquotesingle{}}
\cl{}
\cl{\#~And~with~Git~initialised,~the~stronger~statement:}
\cl{git~status~--porcelain}
\end{codeblock}

Both should be empty. If \texttt{find} reports files that Claude only read, check whether an editor or indexer touched them — atime is not mtime, and a non-empty result here is more often your own tooling than the agent.


\textbf{Check.} Three questions, answerable from the plan alone: which specific files does it cite as evidence for its description of the structure; what did it decline to inspect and why; and is the proposed next step reversible without a backup? A plan that cannot answer all three has not inspected the workspace, it has described it plausibly.


\FloatBarrier
\section*{Chapter 5 — running the brief, and steering without interrupting}
\markboth{Chapter 5 — running the brief, and steering without interrupting}{Chapter 5 — running the brief, and steering without interrupting}

\textbf{Asked:} run the Chapter 1 brief, queue one non-urgent steering message mid-task requiring findings to be ordered by decision relevance, and interrupt only if a boundary is crossed.


The exercise separates two things that novices conflate: \emph{steering}, which joins the queue and is read at the next turn boundary, and \emph{interruption}, which discards the current turn's remaining work. Steering is cheap; interruption throws away tokens already spent and forces re-derivation (Section 5.3).


Start the session with the brief exactly as written, then send this as a queued message while work is in progress — not after it stops:


\begin{codeblock}{7.0}{8.3}
\cl{Additional~requirement,~no~need~to~restart:~order~the~findings~by~decision~relevance~rather}
\cl{than~by~source~order.~A~finding~is~decision-relevant~if~acting~on~it~would~change~what~someone}
\cl{does~next.~Put~the~ones~that~change~a~decision~first,~and~say~for~each~what~decision~it~changes.}
\end{codeblock}

Two details make it a steering message rather than an interruption. It states explicitly that no restart is needed, which removes the ambiguity that causes an agent to discard completed work. And it defines the ordering criterion instead of naming it — "decision relevance" is otherwise interpreted as "importance", and you will get the same list in a different order.


Reserve interruption for boundary crossings: a write outside \texttt{drafts/\allowbreak{}}, a network call, an attempt to resolve an evidence gap by inference rather than recording it. When one occurs, interrupt and say what was crossed, rather than restating the boundary — the boundary was already stated, and repeating it teaches you nothing about why it failed.


\textbf{Check.} Two observations. First, did the steering message change the output without causing re-derivation of work already done? Compare the claim-to-source table before and after: the rows should be reordered and annotated, not regenerated. Second, count your interruptions. More than zero on a well-bounded brief means either the boundary was too loose or you interrupted for a disagreement rather than a violation, and the two have different remedies.


\FloatBarrier
\section*{Chapter 6 — a plan that satisfies the plan contract}
\markboth{Chapter 6 — a plan that satisfies the plan contract}{Chapter 6 — a plan that satisfies the plan contract}

\textbf{Asked:} plan a second brief that adds web research, requiring a query strategy, source-quality criteria, a claim ledger, a source register, a contradiction pass, and explicit approval before any external publication — and do not execute until the plan satisfies all eight points of Section 6.2.


The point of the exercise is that you review the plan against the contract, item by item, rather than against your impression of it. Ask for the plan in the contract's own shape so the review is mechanical:


\begin{codeblock}{8.0}{9.4}
\cl{Plan~only~—~do~not~execute.~This~brief~extends~drafts/brief-v1.md~with~web~research.}
\cl{}
\cl{Return~the~plan~under~these~headings,~and~under~each~state~what~you~will~do~and~how~I~would}
\cl{know~it~was~done:~Outcome;~Inputs~and~assumptions;~Steps~in~order;~Tools~and~permissions}
\cl{required;~What~could~go~wrong~and~the~detection~for~each;~What~you~will~not~do;~Evidence}
\cl{produced;~Approval~gates.}
\cl{}
\cl{Requirements~the~plan~must~satisfy:}
\cl{-~A~query~strategy:~the~search~terms,~why~those~terms,~and~what~would~make~you~change~them.}
\cl{-~Source-quality~criteria~stated~before~searching,~not~chosen~after~seeing~results.}
\cl{-~A~claim~ledger:~one~row~per~claim,~with~the~source~that~supports~it~and~its~status.}
\cl{-~A~source~register:~one~row~per~source,~with~its~class,~date~and~why~it~qualifies.}
\cl{-~A~contradiction~pass:~an~explicit~step~that~looks~for~sources~that~disagree,~and~records}
\cl{~~the~disagreement~rather~than~resolving~it~silently.}
\cl{-~No~external~publication,~message,~commit~or~push~without~separate~approval.}
\end{codeblock}

Reviewing it, the two contract items most often failed are \emph{what could go wrong and the detection for each} — plans list risks without detections, which is a list of worries, not a control — and \emph{what you will not do}, which is usually absent because it is the only item that requires the plan to constrain itself. Send the plan back for those two rather than accepting a plan that is strong on the first six.


The source-quality criteria must be fixed before searching. Criteria chosen after seeing results are a rationalisation of the results, and the failure is invisible in the finished brief. State them concretely: primary or issuing-body sources preferred over commentary; publication or revision date required; a named author or institution; and for any standard or regulation, the issuing body's own text rather than a summary of it (“Scope and method”).


\textbf{Check.} Go through the eight contract points and mark each satisfied or not, in writing. If you cannot mark a point without re-reading the plan twice, that point is not satisfied — the contract is a checklist precisely so that satisfaction is legible at a glance. Then confirm the approval gate is real by checking that publication appears in the plan as a step \emph{you} perform, not as a step Claude performs after asking.


\FloatBarrier
\section*{Chapter 8 — a conservative \texttt{.\allowbreak{}gitignore}, explained pattern by pattern}
\markboth{Chapter 8 — a conservative \texttt{.\allowbreak{}gitignore}, explained pattern by pattern}{Chapter 8 — a conservative \texttt{.\allowbreak{}gitignore}, explained pattern by pattern}

\textbf{Asked:} initialise Git in the practice workspace, have a conservative \texttt{.\allowbreak{}gitignore} drafted with every pattern explained, show the staged diff, and commit only after confirming that nothing sensitive or unrelated is included.


Order matters: \texttt{.\allowbreak{}gitignore} must exist before the first commit, because \texttt{.\allowbreak{}gitignore} does not untrack what is already tracked (Section 8.1). A file committed once remains in history after you ignore it, and removing it from history is a rewrite, not a fix.


\begin{codeblock}{9.0}{10.6}
\cl{git~init}
\cl{\#~draft~.gitignore,~review~it,~and~only~then:}
\cl{git~add~-A}
\cl{git~status~--porcelain~~~~~~~~~~\#~what~will~be~committed,~as~a~list}
\cl{git~diff~--staged~~~~~~~~~~~~~~~\#~what~will~be~committed,~as~content}
\cl{git~commit~-m~"Initial~commit:~practice~workspace"}
\end{codeblock}

A conservative starting point for this workspace:


\begin{codeblock}{8.0}{9.4}
\cl{.env~~~~~~~~~~~~~~~~~~~~~~~\#~local~secrets,~never~committed}
\cl{.env.*~~~~~~~~~~~~~~~~~~~~~\#~per-environment~secrets:~.env.local,~.env.production}
\cl{!.env.example~~~~~~~~~~~~~~\#~the~template~is~safe~and~belongs~in~the~repository}
\cl{.claude/settings.local.json~\#~machine-local~permission~overrides,~not~shared~policy}
\cl{CLAUDE.local.md~~~~~~~~~~~~\#~personal~project~notes,~not~team~instructions}
\cl{secrets/~~~~~~~~~~~~~~~~~~~\#~anything~filed~here~is~excluded~by~location,~not~by~name}
\cl{drafts/*.tmp~~~~~~~~~~~~~~~\#~intermediate~output;~keep~the~finished~drafts}
\cl{.DS\_Store~~~~~~~~~~~~~~~~~~\#~macOS~directory~metadata}
\cl{\_\_pycache\_\_/~~~~~~~~~~~~~~~\#~build~products,~regenerable}
\cl{*.log~~~~~~~~~~~~~~~~~~~~~~\#~logs~may~quote~transcripts,~tokens~or~paths}
\end{codeblock}

Two patterns deserve the explanation the exercise asks for. \texttt{!.\allowbreak{}env.\allowbreak{}example} is a negation, and its position matters: a negation cannot resurrect a file whose \emph{parent directory} is excluded, so \texttt{secrets/\allowbreak{}} plus \texttt{!secrets/\allowbreak{}example.\allowbreak{}json} does not work — the directory is never descended into. And \texttt{secrets/\allowbreak{}} excludes by location rather than by name, which is the more robust of the two strategies: a rule that depends on someone naming a file correctly fails the first time someone does not.


The instruction to have every pattern explained is not politeness. An unexplained \texttt{.\allowbreak{}gitignore} is usually copied from a template for a different stack, and it will either exclude something you need or fail to exclude something you have. Ask for the explanation in the file itself, as above, so it survives into the repository.


\textbf{Check.} Before committing, run \texttt{git diff -\allowbreak{}-\allowbreak{}staged -\allowbreak{}-\allowbreak{}stat} and read the file list end to end. Ask of each entry: do I know why this is here? Then run a negative test — create \texttt{.\allowbreak{}env} with a dummy value and confirm \texttt{git status -\allowbreak{}-\allowbreak{}porcelain} does not list it. Finally, confirm push is not part of the flow: push is a separate, human-gated action, never a step inside an agent's plan (Section 8.2).


\FloatBarrier
\section*{Chapter 9 — a \texttt{CLAUDE.\allowbreak{}md} that is an index, not a warehouse}
\markboth{Chapter 9 — a \texttt{CLAUDE.\allowbreak{}md} that is an index, not a warehouse}{Chapter 9 — a \texttt{CLAUDE.\allowbreak{}md} that is an index, not a warehouse}

\textbf{Asked:} create the project \texttt{CLAUDE.\allowbreak{}md}, start a new session, and have Claude state the directory rules without searching every file. Then determine which instructions are stable and which belong in a path-scoped rule or a skill.


The exercise tests one property: whether the file is an index into the workspace or a container for everything you know about it (Section 9.4). A warehouse file is loaded in full at every session start and consumes context that the task needs.


\begin{codeblock}{9.0}{10.6}
\cl{\#~Practice~research~workspace}
\cl{}
\cl{\#\#~Layout}
\cl{-~\textasciigrave{}notes/\textasciigrave{}~—~source~claims.~Read-only.~Never~edited~by~an~agent.}
\cl{-~\textasciigrave{}sources/\textasciigrave{}~—~primary~evidence,~PDFs.~Read-only.}
\cl{-~\textasciigrave{}drafts/\textasciigrave{}~—~the~only~writable~directory.~Output~goes~here.}
\cl{}
\cl{\#\#~Rules}
\cl{-~Cite~by~file~and~line~for~notes,~by~file~and~page~for~sources.}
\cl{-~Record~an~unsupported~claim~as~UNSUPPORTED;~do~not~resolve~it~by~inference.}
\cl{-~No~network~access~unless~the~task~brief~grants~it~explicitly.}
\cl{-~Publication,~commit~and~push~are~human~actions,~not~steps~in~a~plan.}
\cl{}
\cl{\#\#~Where~to~look}
\cl{-~Claim-to-source~schema:~\textasciigrave{}drafts/SCHEMA.md\textasciigrave{}}
\cl{-~Source-quality~criteria:~\textasciigrave{}drafts/CRITERIA.md\textasciigrave{}}
\cl{-~Prior~briefs~and~their~review~notes:~\textasciigrave{}drafts/archive/\textasciigrave{}}
\end{codeblock}

The three "where to look" pointers are what make it an index. The schema and criteria are stable documents that change rarely and are needed only by tasks that produce briefs; inlining them costs context in every session, including the ones that never write a brief.


Verify in a clean session, without leading the answer:


\begin{codeblock}{8.0}{9.4}
\cl{Without~reading~the~contents~of~notes/~or~sources/,~state~the~rules~that~apply~to~this}
\cl{workspace~and~where~each~one~comes~from.}
\end{codeblock}

If the answer arrives immediately and cites \texttt{CLAUDE.\allowbreak{}md}, the file is doing its job. If Claude searches first, the rules are either not in the file or are buried below the point where they read as instructions.


The second half of the exercise is the classification. A rule is stable if it would still be true after the current task changes; it belongs in \texttt{CLAUDE.\allowbreak{}md}. A rule that applies only inside a subtree belongs in a path-scoped rule so it loads only when that subtree is in play. A rule that is really a \emph{procedure} — a sequence you repeat with variation — is a skill, not an instruction, and leaving it in \texttt{CLAUDE.\allowbreak{}md} both wastes context and makes it probabilistic (Chapter 15).


\begingroup
\def\tblrows{%
Directory layout and write boundary & Yes & \texttt{CLAUDE.\allowbreak{}md} \\
Citation format for notes and sources & Yes & \texttt{CLAUDE.\allowbreak{}md} \\
"Never edit a PDF in place" & Yes, but scoped & Path-scoped rule on \texttt{sources/\allowbreak{}} \\
The eight-step brief procedure & It is a procedure & A skill (Section 15.1) \\
"Use British spelling in drafts" & Yes, but scoped & Path-scoped rule on \texttt{drafts/\allowbreak{}} \\
}%
\def\tblbody{\begin{minipage}{\textwidth}
{\footnotesize\begin{tabular}{L{72.0mm}L{24.7mm}L{43.1mm}}
\toprule
\textbf{Instruction} & \textbf{Stable} & \textbf{Belongs in} \\
\midrule
\tblrows
\bottomrule\end{tabular}}\end{minipage}}%
\begingroup
\def\sloppy{\tolerance 9999\emergencystretch 3em\hfuzz 200pt\vfuzz 200pt}%
\hbadness=10000\vbadness=10000\hfuzz=200pt\vfuzz=200pt
\global\setbox\tblbox=\hbox{\tblbody}%
\endgroup
\par\addvspace{2.6mm}
\ifdim\dimexpr\ht\tblbox+\dp\tblbox\relax>0.55\textheight
  \typeout{HANDBOOK-TABLE broken \the\dimexpr\ht\tblbox+\dp\tblbox\relax}%
  
  {\footnotesize\begin{longtable}{L{72.0mm}L{24.7mm}L{43.1mm}}
  \toprule
\textbf{Instruction} & \textbf{Stable} & \textbf{Belongs in} \\
\midrule\endfirsthead
  \multicolumn{3}{@{}l@{}}{%
  \sffamily\footnotesize\itshape\color{inkgrey}Continued}\\[1.2mm]
  \toprule
\textbf{Instruction} & \textbf{Stable} & \textbf{Belongs in} \\
\midrule\endhead
  \bottomrule\endfoot
  \bottomrule\endlastfoot
  \tblrows
  \end{longtable}}%
\else
  \typeout{HANDBOOK-TABLE atomic \the\dimexpr\ht\tblbox+\dp\tblbox\relax}%
  \noindent\tblbody
\fi
\par\addvspace{2.6mm}
\endgroup

\textbf{Check.} Two measurements. Count the lines: a project file that has grown past roughly two hundred lines has almost certainly absorbed something that belongs elsewhere (Section 9.2). Then run \texttt{/\allowbreak{}context} in a fresh session and look at what the instruction load costs before any work has been done. If the project file is a visible fraction of the window at turn zero, move the stable-but-narrow material into path-scoped rules and the procedures into skills.


\FloatBarrier
\section*{Chapter 11 — a handoff that survives a clean session}
\markboth{Chapter 11 — a handoff that survives a clean session}{Chapter 11 — a handoff that survives a clean session}

\textbf{Asked:} after drafting the brief, run \texttt{/\allowbreak{}context}, save a handoff, start a clean session, and have the new session continue — verifying that it can name completed work, open issues and the next check without receiving the old transcript.


This is the exercise that distinguishes context \emph{engineering} from context \emph{hoarding}. The test is not whether the new session can be brought up to speed; it is whether the handoff you wrote carries the state, or whether you unconsciously relied on the transcript.


Run \texttt{/\allowbreak{}context} first and record the numbers, because the comparison afterwards is the finding (Section 11.1). Then ask for the handoff with a schema, not as a summary:


\begin{codeblock}{8.0}{9.4}
\cl{Write~drafts/HANDOFF.md~for~a~session~that~will~have~none~of~this~conversation.~Include:}
\cl{}
\cl{1.~Goal~—~the~outcome~in~one~sentence,~in~the~terms~of~the~original~brief.}
\cl{2.~Done~—~completed~work,~each~item~with~the~artefact~that~evidences~it.}
\cl{3.~Not~done~—~remaining~work,~in~the~order~it~should~be~attempted.}
\cl{4.~Open~issues~—~decisions~I~must~make,~each~with~the~options~and~what~each~would~cost.}
\cl{5.~Rejected~—~approaches~already~tried~and~abandoned,~with~the~reason.}
\cl{6.~Next~check~—~the~single~next~verification,~and~what~its~passing~would~prove.}
\cl{7.~Boundaries~—~the~constraints~in~force,~restated~in~full.}
\cl{}
\cl{Write~no~history.~A~new~session~must~not~need~to~know~what~happened,~only~what~is~true.}
\end{codeblock}

Item 5 is the one that is routinely omitted and the one that pays for the file. Without it the new session re-proposes a rejected approach, and you spend the saved context re-rejecting it. Item 7 is restated in full rather than referenced because a boundary that lives only in the old transcript does not survive the handoff — which is precisely the failure the exercise induces.


Start the clean session and ask it to continue from \texttt{drafts/\allowbreak{}HANDOFF.\allowbreak{}md} alone. Do not paste the brief, do not summarise, do not correct it until it has committed to a next step.


\textbf{Check.} Three answers from the new session, before you supply anything: what is done and what evidences it; what decision is currently blocking; what is the next check and what would its passing prove. Any of the three that requires you to fill in a gap identifies a missing section in the handoff. Run \texttt{/\allowbreak{}context} again in the new session and compare with the number you recorded: a handoff that reproduces the working state at a fraction of the context cost is the demonstration the exercise is after.


\FloatBarrier
\section*{Chapter 13 — a threat model before an MCP connection}
\markboth{Chapter 13 — a threat model before an MCP connection}{Chapter 13 — a threat model before an MCP connection}

\textbf{Asked:} without installing anything, produce a threat model for a hypothetical Notion connection covering data reachable, write operations exposed, credential storage, scope, approval gates, logging and revocation — then replace each assumption with a statement from the provider's documentation before approving installation.


The exercise is deliberately ordered so that the model is written while you still have no stake in the answer. A threat model written after installation is a justification.


Ask for it as a table with an explicit assumption column, because the assumptions are the deliverable:


\begin{codeblock}{8.0}{9.4}
\cl{Do~not~install~anything.~Produce~a~threat~model~for~connecting~a~Notion~MCP~server~to~this}
\cl{workspace,~as~a~table~with~the~columns:~Question~|~Current~answer~|~Is~this~an~assumption~or}
\cl{a~documented~fact~|~Source.~Cover:~what~data~becomes~reachable;~which~write~operations~are}
\cl{exposed;~where~the~credential~is~stored~and~what~can~read~it;~what~scope~the~credential}
\cl{carries;~which~operations~require~approval~and~which~do~not;~what~is~logged~and~where;~and}
\cl{how~access~is~revoked~and~how~I~would~confirm~revocation~took~effect.}
\cl{}
\cl{Mark~every~row~you~cannot~support~from~the~provider\textquotesingle{}s~own~documentation~as~ASSUMPTION.}
\cl{Do~not~fill~a~gap~by~inference.}
\end{codeblock}

The rows that most often come back as assumptions, and which matter most:


\begin{itemize}
\item \textbf{Scope.} Integration tokens frequently carry workspace-wide reach rather than page-level reach, and "connected to a page" describes what you selected in a picker, not what the token permits.
\item \textbf{Write operations.} Read-only is a property of the token, not of the connector's name.
\item \textbf{Revocation.} Revoking in the provider's console and removing the server locally are two different operations, and either alone leaves the other in place.
\item \textbf{Logging.} Whether the provider records the agent's reads, and whether you can retrieve that record, are separate questions.
\end{itemize}

Then do the second half, which is the actual work: replace each ASSUMPTION with a quotation from the provider's documentation, with its URL and the date you read it. Where the documentation does not answer the question, the row stays ASSUMPTION — an unanswered question recorded as unanswered is a finding, and treating silence as permission is the failure mode this book returns to repeatedly (Section 13.1).


An MCP server is untrusted input as well as a capability: its tool descriptions and its returned content both reach the model, so a compromised or careless server is an injection surface, not only a data-access surface (Section 13.4).


\textbf{Check.} Count the rows still marked ASSUMPTION. If any of scope, write operations or revocation remains an assumption, the connection is not ready to approve regardless of how useful it would be. Then state, in one sentence and without looking, what the worst outcome of a compromised token would be for this workspace. If you cannot, the model has been produced but not read.


\FloatBarrier
\section*{Chapter 15 — a skill improved only by rules that generalise}
\markboth{Chapter 15 — a skill improved only by rules that generalise}{Chapter 15 — a skill improved only by rules that generalise}

\textbf{Asked:} install the example skill locally, invoke \texttt{/\allowbreak{}research-\allowbreak{}brief} on a second source set, record every place it made an assumption, and improve the skill only with a rule that generalises.


The constraint in the last clause is the whole exercise. The natural response to a bad output is to add a sentence forbidding exactly that output, and a skill maintained that way becomes a list of past failures that grows without ever getting better.


Run it on a second source set — one with a different shape from the first, ideally with a missing date, a source that contradicts another, and one document that is not what its filename says. Record assumptions as they happen, in a file, with the evidence:


\begin{codeblock}{7.0}{8.3}
\cl{|~\#~|~Where~it~assumed~|~What~it~assumed~|~What~it~should~have~done~|}
\cl{|---|---|---|---|}
\cl{|~1~|~sources/report.pdf~|~that~the~undated~file~was~current~|~recorded~the~date~as~unknown~|}
\cl{|~2~|~two~disagreeing~notes~|~picked~the~more~recent~one~silently~|~recorded~the~disagreement~|}
\cl{|~3~|~a~filename~ending~.md~|~that~the~content~matched~the~name~|~checked~before~classifying~|}
\end{codeblock}

Now apply the test. For each row, ask: is the fix a rule about \emph{this input}, or a rule about a \emph{class of input}? Row 1's specific fix — "treat report.pdf as undated" — generalises to a rule about missing metadata. Row 2's specific fix — "prefer the note from August" — is a rule about two files and generalises to a rule about contradiction. Row 3 generalises to a rule about trusting declared type over observed content.


The generalised rules, which is what actually goes into the skill:


\begin{codeblock}{7.0}{8.3}
\cl{\#\#~Evidence~rules}
\cl{-~A~source~with~no~determinable~date~is~recorded~as~date-unknown~and~may~not~be~used~to}
\cl{~~establish~precedence~over~a~dated~source.}
\cl{-~When~two~sources~disagree,~record~the~disagreement~in~the~claim~ledger~with~both~positions.}
\cl{~~Do~not~resolve~it~by~recency,~by~source~order,~or~by~which~is~easier~to~summarise.}
\cl{-~Classify~a~source~by~inspecting~its~content,~not~by~its~filename~or~extension.~Where~the~two}
\cl{~~conflict,~record~the~conflict.}
\end{codeblock}

Each of those would have prevented the observed failure \emph{and} prevents failures you have not seen yet. That is the test for whether a rule belongs in a skill at all: a rule that can only fire on the input that produced it belongs in the task brief, not in the reusable procedure.


Keep the skill's frontmatter description precise while you are at it. A skill is selected by its description, so a vague one is either never invoked or invoked for the wrong task (Section 15.3).


\textbf{Check.} Re-run \texttt{/\allowbreak{}research-\allowbreak{}brief} on the \emph{first} source set after the change. The output should be unchanged except where the new rules legitimately apply — a rule that alters the first result in some unrelated way has changed the procedure's behaviour rather than closed a gap. Then count: if the skill grew by more lines than the number of generalised rules you identified, something specific has been smuggled in.


\FloatBarrier
\section*{Chapter 17 — a project-only \texttt{Stop} hook that records completion}
\markboth{Chapter 17 — a project-only \texttt{Stop} hook that records completion}{Chapter 17 — a project-only \texttt{Stop} hook that records completion}

\textbf{Asked:} create a project-only \texttt{Stop} hook that appends a timestamped, non-sensitive completion record to a local file, making no network call, approving no permission, and including no conversation content. Test success, failure and removal.


Follow the creation sequence rather than writing the configuration first: define the event and matcher, write a standalone script with deterministic input and output, test it against benign fixtures by hand, and only then wire it in (Section 17.4).


The script, tested on its own before Claude Code ever calls it:


\begin{codeblock}{8.0}{9.4}
\cl{\#!/usr/bin/env~bash}
\cl{\#~.claude/hooks/record-stop.sh~—~append~a~completion~record.~No~network.~No~content.}
\cl{set~-euo~pipefail}
\cl{LOG="\$\{CLAUDE\_PROJECT\_DIR\}/drafts/session-log.tsv"}
\cl{mkdir~-p~"\$(dirname~"\$LOG")"}
\cl{printf~\textquotesingle{}\%s\textbackslash{}t\%s\textbackslash{}t\%s\textbackslash{}n\textquotesingle{}~\textbackslash{}}
\cl{~~"\$(date~-u~+\%Y-\%m-\%dT\%H:\%M:\%SZ)"~\textbackslash{}}
\cl{~~"stop"~\textbackslash{}}
\cl{~~"\$(basename~"\$\{CLAUDE\_PROJECT\_DIR\}")"~>>~"\$LOG"}
\cl{exit~0}
\end{codeblock}

Three deliberate omissions. It reads nothing from standard input, so no conversation content can reach the log by accident — the requirement is met structurally rather than by care. It writes a fixed three-field record, so the log cannot grow a field that later carries something sensitive. And it exits 0 unconditionally, because a \texttt{Stop} hook that exits 2 \emph{prevents the session from stopping}, which turns a logging failure into a session that will not end.


The configuration, in the project's \texttt{.\allowbreak{}claude/\allowbreak{}settings.\allowbreak{}json} so that it is project-only and visible in review:


\begin{codeblock}{9.0}{10.6}
\cl{\{}
\cl{~~"hooks":~\{}
\cl{~~~~"Stop":~[}
\cl{~~~~~~\{}
\cl{~~~~~~~~"hooks":~[}
\cl{~~~~~~~~~~\{}
\cl{~~~~~~~~~~~~"type":~"command",}
\cl{~~~~~~~~~~~~"command":~"\$\{CLAUDE\_PROJECT\_DIR\}/.claude/hooks/record-stop.sh",}
\cl{~~~~~~~~~~~~"timeout":~10,}
\cl{~~~~~~~~~~~~"statusMessage":~"Recording~completion…"}
\cl{~~~~~~~~~~\}}
\cl{~~~~~~~~]}
\cl{~~~~~~\}}
\cl{~~~~]}
\cl{~~\}}
\cl{\}}
\end{codeblock}

\texttt{\$\{CLAUDE\_\allowbreak{}PROJECT\_\allowbreak{}DIR\}} rather than a relative path, because it stays at the project root even inside a worktree, and a relative path in a hook resolves against a working directory you do not control (Section 17.2).


Now the three tests the exercise names.


\emph{Success.} Run any short task, let it finish, and inspect \texttt{drafts/\allowbreak{}session-\allowbreak{}log.\allowbreak{}tsv}. One new row, three fields, no prose.


\emph{Failure.} Make the script fail deterministically — \texttt{chmod -\allowbreak{}x} it, or point \texttt{LOG} at an unwritable path — and run another task. Confirm the session still ends and the failure is visible rather than silent. This is the test people skip, and it is the one that matters: a logging hook that fails closed will eventually stop your work, and one that fails invisibly will leave you believing you have a record you do not have.


\emph{Removal.} Delete the \texttt{Stop} block, start a new session, and confirm with \texttt{/\allowbreak{}hooks} that no Stop hook is configured — then run a task and confirm no new row appears. Verifying removal from the configuration file alone is insufficient, because hooks arrive from several sources and settings precedence decides which one is live (Section 17.5).


\textbf{Check.} Read the log file as an adversary would. Does any field carry a path, a prompt, a filename or a token? Then confirm the negative properties directly rather than by inspection of intent: the script contains no \texttt{curl}, \texttt{wget} or redirect to a URL, and returns no \texttt{permission\allowbreak{}Decision}. A hook that \emph{can} approve a permission is privileged code, and this one must not be able to.


\FloatBarrier
\section*{Chapter 20 — a bounded browser task, verified from outside the transcript}
\markboth{Chapter 20 — a bounded browser task, verified from outside the transcript}{Chapter 20 — a bounded browser task, verified from outside the transcript}

\textbf{Asked:} on a harmless public page or a local fixture, extract a small table with its location cited and take no action; then navigate to one linked page; then inspect browser history and the transcript to confirm the exact scope of what happened.


Use a local fixture if you have one — it removes the network from the first half of the exercise entirely, and it lets you know the ground truth of what the page contains. Whatever you use, the browser profile is the control surface that matters: a browsing agent operating in your everyday profile is operating with your logged-in sessions (Section 20.2).


The two-stage prompt, kept deliberately separate so that the second stage is a decision:


\begin{codeblock}{8.0}{9.4}
\cl{Stage~1.~Read~the~table~on~this~page.~Return~it~as~markdown,~and~for~each~row~state~where~on}
\cl{the~page~it~came~from.~Take~no~action:~no~click,~no~form~entry,~no~navigation,~no~download.}
\cl{}
\cl{Stage~2.~(Send~only~after~reviewing~stage~1.)~Navigate~to~the~single~link~labelled~"…"~and}
\cl{report~the~page~title~and~heading~structure.~Do~not~follow~any~further~link.}
\end{codeblock}

The verification is the point of the exercise, and it must come from outside the transcript, because the transcript is a record of what was reported, not of what occurred:


\begin{codeblock}{9.0}{10.6}
\cl{\#~Chrome~history,~read~from~a~copy~—~the~live~database~is~locked~while~Chrome~runs.}
\cl{cp~"\$HOME/.config/google-chrome/Default/History"~/tmp/h.db}
\cl{sqlite3~/tmp/h.db~\textbackslash{}}
\cl{~~"SELECT~datetime(last\_visit\_time/1000000-11644473600,\textquotesingle{}unixepoch\textquotesingle{}),~url}
\cl{~~~FROM~urls~ORDER~BY~last\_visit\_time~DESC~LIMIT~20;"}
\end{codeblock}

Adjust the profile path for your platform. Compare that list against the transcript: exactly one navigation should appear for stage 2, and none for stage 1. Extra entries are the finding — commonly a redirect, a prefetch, or a link that resolved somewhere other than its label.


Page content is untrusted input. Text on a page can address the agent directly, and an instruction that arrives inside fetched content is indistinguishable in form from one you wrote (Section 20.5). Stage 1 asks for extraction and nothing else precisely so that a page attempting to induce an action produces a visible refusal rather than a quiet compliance.


\textbf{Check.} Three comparisons. The extracted table against the page — is every row present and correctly located, or has a merged cell been silently flattened? The history against the transcript — same count, same URLs. And the transcript against the stage boundary — did anything in stage 1 attempt an action, and if so, what on the page induced it?


\FloatBarrier
\section*{Chapter 22 — a small workflow, and whether it earned its cost}
\markboth{Chapter 22 — a small workflow, and whether it earned its cost}{Chapter 22 — a small workflow, and whether it earned its cost}

\textbf{Asked:} run \texttt{/\allowbreak{}deep-\allowbreak{}research} on a non-sensitive topic with a \texttt{small} size guideline, requiring direct sources, a contradiction pass and a source register; compare the result with a single-agent attempt; and use \texttt{/\allowbreak{}usage} to decide whether the workflow added decision value or only cost.


Set the guideline before the run, not during it — a run cannot be resized once started, and \texttt{small} bounds it below five agents (Section 22.2). Set \texttt{workflow\allowbreak{}Size\allowbreak{}Guideline} to \texttt{small} in \texttt{/\allowbreak{}config} or in a settings file.


Then commit a clean baseline. This is not general hygiene; it is specific to workflows. The subagents a workflow spawns always run in \texttt{accept\allowbreak{}Edits} and inherit your tool allowlist, whatever the session's mode [19]. Your mode governs only the launch prompt. File edits are auto-approved for the duration of the run.


\begin{codeblock}{9.0}{10.6}
\cl{git~add~-A~\&\&~git~commit~-m~"Baseline~before~workflow~run"}
\end{codeblock}

Pick a topic with genuine disagreement in the literature, because a contradiction pass over a settled topic proves nothing. Then:


\begin{codeblock}{8.0}{9.4}
\cl{/deep-research~What~is~the~current~evidence~on~the~effectiveness~of~<topic>?~Requirements:}
\cl{cite~issuing~bodies~and~primary~studies~rather~than~commentary~about~them;~produce~a~source}
\cl{register~with~one~row~per~source~giving~its~class,~date~and~why~it~qualifies;~and~run~an}
\cl{explicit~contradiction~pass~that~reports~sources~which~disagree~and~states~the~disagreement}
\cl{rather~than~resolving~it.~Where~the~evidence~does~not~settle~a~question,~say~so.}
\end{codeblock}

Watch \texttt{/\allowbreak{}workflows} while it runs. Then run the same question as a single agent in a fresh session, with the same requirements, and put the two outputs side by side.


The comparison has a specific shape. Ask three questions of the pair, in this order. Did the workflow surface a source the single agent did not, or the same sources organised differently? Did the contradiction pass find a real disagreement, or manufacture one by quoting two sources that are not in conflict? And would a decision-maker choose differently given one output rather than the other? Only the third is decisive. The first two can both be yes while the answer to the third is no, and that is the outcome the exercise is designed to make visible.


Then read \texttt{/\allowbreak{}usage} for the session (Section 4.4). Divide the cost by the answer to the third question. A workflow that produced a better-organised version of the same conclusion at several times the cost did not add decision value; it added confidence, which is not the same thing and is sometimes worse than nothing.


If you stop the run mid fan-out, note what happens on resume: replay follows the order agents \emph{started}, cached results stop at the first agent that did not finish, and every agent that started after it runs again even if it had completed (Section 22.3). Stopping is therefore expensive, and it is worth stopping once deliberately so that the cost is a measurement rather than a surprise.


\textbf{Check.} Write two sentences before looking at the cost: what the workflow told you that the single agent did not, and what decision that changes. Then look at \texttt{/\allowbreak{}usage}. If the two sentences were hard to write, the run was an exercise in throughput rather than in research — which is a legitimate finding, and the one this exercise most often produces.


\FloatBarrier
\section*{Chapter 25 — a read-only status report from a second device}
\markboth{Chapter 25 — a read-only status report from a second device}{Chapter 25 — a read-only status report from a second device}

\textbf{Asked:} start Remote Control on the practice workspace, request a read-only status report from a second device, confirm the output appears locally, then disconnect — granting no new permissions during the exercise.


The clause that carries the lesson is the last one. Remoteness creates pressure to widen access — the review surface is worse, so the temptation is to approve rather than inspect. The exercise exists to establish, once, that you can complete useful remote work without loosening anything (Section 25.1).


Start from the practice directory and name the session so that a phone cannot confuse it with another project:


\begin{codeblock}{9.0}{10.6}
\cl{cd~\textasciitilde{}/practice-workspace}
\cl{claude~--remote-control~research-brief}
\end{codeblock}

Confirm before connecting the second device: correct machine, correct directory, correct branch, correct account. On a small screen these are the four facts you will not otherwise re-establish, and the phone shows a conversation, not a working directory.


From the second device, ask for something that cannot require a permission decision:


\begin{codeblock}{7.0}{8.3}
\cl{Read-only~status~report.~Do~not~edit,~create~or~delete~any~file,~and~do~not~run~any~command}
\cl{that~writes.~Report:~current~branch~and~whether~the~tree~is~clean;~which~files~in~drafts/}
\cl{changed~most~recently~and~when;~any~claim~in~drafts/brief-v1.md~still~marked~UNSUPPORTED;~and}
\cl{the~single~next~check~from~drafts/HANDOFF.md.~If~anything~you~would~need~requires~a~write~or~a}
\cl{new~permission,~stop~and~say~what~it~was.}
\end{codeblock}

The final sentence is the instrument. If the report comes back complete with no such statement, the task genuinely stayed read-only. If it comes back with a request, you have learned where the read-only boundary actually sits for this workspace — and the correct response during this exercise is to decline it and note it, not to approve it from a phone.


Confirm the output appears in the local terminal as well as on the device. Execution stays local: code, files and commands never leave the machine, and what travels is the conversation, relayed and stored while connected so it can sync across devices. Two consequences follow that are easy to miss — the local process must keep running, because closing it ends the execution path even though the mobile client stays open; and the relayed transcript is cloud-stored data for as long as you are connected, so anything you would not put in a cloud-stored conversation does not belong in a remote session (Chapter 32).


Then disconnect deliberately rather than by closing a lid. Close remote control when supervision is over, and confirm the local session's state is what you expect afterwards.


\textbf{Check.} Four confirmations. The report's facts match what you observe locally — branch, cleanliness, file times. \texttt{git status -\allowbreak{}-\allowbreak{}porcelain} is unchanged from before the exercise. No permission was granted: check your settings files for anything added during the session, since a permission approved "just for now" is written down somewhere. And the session is actually closed, verified on the machine rather than by the absence of the session on the phone.


\FloatBarrier
\section*{A note on the solutions}
\markboth{A note on the solutions}{A note on the solutions}

Eight of the thirteen solutions end in a verification performed \emph{outside} the transcript — a filesystem check, a browser history, a \texttt{git status}, a settings file or a produced artefact read directly. That is not a stylistic preference. The transcript records what was reported, and every failure mode this book treats as serious is one in which the report and the reality differ. Where an exercise can be marked complete from the conversation alone, it is testing comprehension; where it requires an external check, it is testing the practice. The second kind is the one worth repeating on real work.


\FloatBarrier
\setcounter{chapter}{13}
\setcounter{section}{0}
\appendixchapter{N}{Cross-reference matrix and index}
\FloatBarrier
\section*{Chapter dependency matrix}
\markboth{Chapter dependency matrix}{Chapter dependency matrix}

\begingroup
\def\tblrows{%
1 & What Claude Code is, and what it is not & 5, 21 \\
2 & Install, verify and authenticate & — \\
3 & Surfaces, projects and boundaries & — \\
4 & Modes, models, effort and usage & 7, 16, 22 \\
5 & Start and steer a session & 1, 7, 16 \\
6 & Plan, comment, approve, execute & 7 \\
7 & Permissions, sandboxing and checkpoints & 8, 17, 22, 31 \\
8 & Protect work with Git and GitHub & 17 \\
9 & Build an effective \texttt{CLAUDE.\allowbreak{}md} & 7, 17 \\
10 & Global instructions and auto memory & 15, 17 \\
11 & Engineer context deliberately & 13 \\
12 & Sessions: resume, name, retain & 10, 23, 29, 32 \\
13 & Connect external tools with MCP & — \\
14 & Environment variables and secrets & 7, 12, 33 \\
15 & Turn repeatable work into skills & — \\
16 & Delegate bounded work to subagents & 22, 30 \\
17 & Event-driven behaviour with hooks & 21 \\
18 & Evaluate and install plugins & — \\
19 & Output styles: Claude Code beyond software engineering & — \\
20 & Use Claude in Chrome safely & — \\
21 & Keep work moving with goals & — \\
22 & Dynamic workflows and ultracode & — \\
23 & Session-scoped scheduling with \texttt{/\allowbreak{}loop} & — \\
24 & Cloud routines and desktop scheduled tasks & — \\
25 & Continue local work with Remote Control & 32 \\
26 & Cloud sessions, \texttt{-\allowbreak{}-\allowbreak{}cloud} and \texttt{-\allowbreak{}-\allowbreak{}teleport} & — \\
27 & Push events into a session with channels & — \\
28 & Publish work as artifacts & — \\
29 & Parallel sessions without collisions & 12, 16, 22, 30 \\
30 & Agent teams & 22, 29, 34 \\
31 & Enterprise configuration and managed policy & 7, 18, 21, 32 \\
32 & Data governance, residency and retention & 10, 14, 30 \\
33 & Observability, cost control and accessibility & 11, 16, 19, 22, 30, 31 \\
34 & Build a verified research-brief agent & 1, 3, 4, 5, 6, 7, 8, 9, 10, 11, 12, 13, 14, 15, 16, 17, 18, 19, 20, 21, 22, 23, 24, 25, 27, 29, 30, 31, 32, 33 \\
}%
\def\tblbody{\begin{minipage}{\textwidth}
{\footnotesize\begin{tabular}{L{16.9mm}L{97.0mm}L{25.9mm}}
\toprule
\textbf{Chapter} & \textbf{Title} & \textbf{Depends on / refers to} \\
\midrule
\tblrows
\bottomrule\end{tabular}}\end{minipage}}%
\begingroup
\def\sloppy{\tolerance 9999\emergencystretch 3em\hfuzz 200pt\vfuzz 200pt}%
\hbadness=10000\vbadness=10000\hfuzz=200pt\vfuzz=200pt
\global\setbox\tblbox=\hbox{\tblbody}%
\endgroup
\par\addvspace{2.6mm}
\ifdim\dimexpr\ht\tblbox+\dp\tblbox\relax>0.55\textheight
  \typeout{HANDBOOK-TABLE broken \the\dimexpr\ht\tblbox+\dp\tblbox\relax}%
  
  {\footnotesize\begin{longtable}{L{16.9mm}L{97.0mm}L{25.9mm}}
  \toprule
\textbf{Chapter} & \textbf{Title} & \textbf{Depends on / refers to} \\
\midrule\endfirsthead
  \multicolumn{3}{@{}l@{}}{%
  \sffamily\footnotesize\itshape\color{inkgrey}Continued}\\[1.2mm]
  \toprule
\textbf{Chapter} & \textbf{Title} & \textbf{Depends on / refers to} \\
\midrule\endhead
  \bottomrule\endfoot
  \bottomrule\endlastfoot
  \tblrows
  \end{longtable}}%
\else
  \typeout{HANDBOOK-TABLE atomic \the\dimexpr\ht\tblbox+\dp\tblbox\relax}%
  \noindent\tblbody
\fi
\par\addvspace{2.6mm}
\endgroup

\FloatBarrier
\section*{Index}
\markboth{Index}{Index}

\begin{multicols}{2}\footnotesize
\begin{deflist}
\item \textbf{accessibility} \textemdash\ 12.2, 33.3, 34.10, H, L, N  
\item \textbf{agent teams} \textemdash\ 16.6, 22, 30, 30.1, 30.7, D, G, L, N  
\item \textbf{artifacts} \textemdash\ 2.1, 2.3, 2.4, 7.5, 8.5, 28, 28.4, 32.3, 32.5, 34.8, A, F  
\item \textbf{auto memory} \textemdash\ 10, 10.2, 10.3, 10.4, 16.1, 32.2, 32.5, 34.3, A, C, D, I  
\item \textbf{auto mode} \textemdash\ 3.2, 3.3, 6.4, 7.1, 16.3, 17.1, 19.1, 19.3, 21.3, 29.3, I  
\item \textbf{bypassPermissions} \textemdash\ 7.1, 7.5, 12.1, 12.4, 16.3, 16.7, 19.1, 34.6, D, G, H, J  
\item \textbf{channels} \textemdash\ 2.1, 12.4, 27, 27.1, 27.2, 27.3, 32.5, 34.10, C, D, H, J  
\item \textbf{checkpoint} \textemdash\ 5.3, 7.5, 11.3, A, I, N  
\item \textbf{CLAUDE.md} \textemdash\ 1.3, 2.3, 3.3, 4.4, 4.5, 7, 9, 9.1, 9.3, 9.5, 9.6, 10  
\item \textbf{cleanupPeriodDays} \textemdash\ 10.2, 10.4, 12.3, 30.2, 32.2, 32.5, D  
\item \textbf{cloud environment} \textemdash\ 10.2  
\item \textbf{compaction} \textemdash\ 9.6, 11.1, 11.2, 21.3, G  
\item \textbf{context window} \textemdash\ 3.3, 7.3, 16, 22, 30, 30.2, 34.10, I  
\item \textbf{cron} \textemdash\ 23, 23.1, 23.2, 23.3, 24.1, D, J  
\item \textbf{deep-research} \textemdash\ 22.5, 22.6, A, M  
\item \textbf{deny rule} \textemdash\ 7.2, 7.5, 14.2, 14.3, 15.4, 17.5, 20.4, 23.3, A, F, K, L  
\item \textbf{dev container} \textemdash\ 7.4  
\item \textbf{effort} \textemdash\ 4.1, 4.2, 4.3, 4.5, 16.4, 22.1, 30.3, 33.2, A, B, I, J  
\item \textbf{fork} \textemdash\ 11.3, 11.4, 15.3, 16.1, 16.6, 16.7, 17.1, 19.3, 26.2, 29.1, 33.2, 34.6  
\item \textbf{goal} \textemdash\ 11.2, 12.1, 12.4, 17.5, 21, 21.1, 21.3, 21.4, 31.2, 34.8, A, D  
\item \textbf{hooks} \textemdash\ 10.2, 10.4, 12.3, 15.3, 15.4, 16.4, 17, 17.2, 17.3, 17.4, 17.5, 18.1  
\item \textbf{human gate} \textemdash\ 7.5, 13.3, 21.2, F, I  
\item \textbf{managed settings} \textemdash\ 7.2, 7.4, 9.1, 16.2, 17.5, 19.2, 21.1, 22.6, 27.2, 27.3, 30.1, 31  
\item \textbf{MCP} \textemdash\ 1.3, 4.4, 11.4, 12.1, 12.4, 13.1, 13.2, 13.3, 13.5, 16.4, 17.2, 18.1  
\item \textbf{output styles} \textemdash\ 18.1, 19, 19.1, 19.3, C, N  
\item \textbf{permission mode} \textemdash\ 4.1, 4.5, 6.4, 6.5, 7.1, 12.1, 19.1, 19.3, 21.3, 21.4, 25.1, 25.2  
\item \textbf{plan mode} \textemdash\ 3.3, 6.1, 6.5, 20.3, 20.6, 30.3, A, M  
\item \textbf{plugin} \textemdash\ 1.3, 4.4, 4.5, 8.4, 11.4, 12.1, 12.4, 16.2, 17.2, 17.5, 18, 18.1  
\item \textbf{prompt caching} \textemdash\ 4.4  
\item \textbf{prompt injection} \textemdash\ 13.3, 20.5, F, I  
\item \textbf{Remote Control} \textemdash\ 1.3, 25, 25.1, 25.2, 26, 27.3, 31.4, 32.3, 32.4, 32.5, D, F  
\item \textbf{routines} \textemdash\ 2.1, 2.3, 2.4, 24.1, 24.4, 24.5, 26.6, 34.10, A, J, L, N  
\item \textbf{rules} \textemdash\ 3.3, 7, 7.2, 7.4, 7.5, 9.1, 9.5, 9.6, 10, 10.1, 10.3, 10.4  
\item \textbf{sandbox} \textemdash\ 1.3, 2.1, 2.4, 3.2, 3.3, 6.4, 7.1, 7.4, 7.5, 14.2, 14.3, 21.3  
\item \textbf{skills} \textemdash\ 3.3, 15, 15.4, 16.1, 16.4, 18.1, 19.3, 22, 23.3, 30.3, 31.2, A  
\item \textbf{subagent} \textemdash\ 4.4, 7.1, 7.5, 9.3, 11.3, 15.3, 16, 16.1, 16.2, 16.3, 16.4, 16.5  
\item \textbf{teleport} \textemdash\ 2.3, 2.4, 26.2, 26.6, A, I, N  
\item \textbf{Tool Search} \textemdash\ 13.4, G, I  
\item \textbf{transcript} \textemdash\ 1.3, 5.4, 5.5, 10.2, 11.4, 12.3, 12.4, 14.1, 14.2, 14.3, 20.4, 20.6  
\item \textbf{ultracode} \textemdash\ 4.1, 4.2, 22.1, 22.2, 33.2, A, I, J, N  
\item \textbf{workflow} \textemdash\ 4.3, 7.1, 8.5, 12.2, 16.6, 22, 22.1, 22.2, 22.3, 22.5, 22.6, 26.4  
\item \textbf{worktree} \textemdash\ 3.3, 10.2, 10.4, 12.1, 12.4, 16.4, 17.1, 17.2, 29.1, 29.2, 29.4, 30.7  
\item \textbf{zero data retention} \textemdash\ 26.3, 28.4, 32, I  
\end{deflist}
\end{multicols}

\breakrule

\FloatBarrier
\backmatter
\frontchapter{References}

The reference list is numbered in order of first citation in the text. All online sources were accessed on 23 August 2026 unless stated otherwise. Every cited documentation page was re-fetched and hashed, and all 96 unique URLs in the text re-checked for reachability, on 26 August 2026. Product documentation describes behaviour observed on Claude Code 2.1.241 and is subject to change; see “Why this appendix exists” for the re-verification procedure.


\begin{refslist}
\item[1.] Anthropic, “Claude Code documentation index,” \emph{llms.txt}. [Online]. Available: \url{https://code.claude.com/docs/llms.txt}. [accessed 23 Aug. 2026].
\end{refslist}

\begin{refslist}
\item[2.] Anthropic, “Overview,” \emph{Claude Code Documentation}. [Online]. Available: \url{https://code.claude.com/docs/en/overview}. [accessed 23 Aug. 2026].
\end{refslist}

\begin{refslist}
\item[3.] Anthropic, “How Claude Code works,” \emph{Claude Code Documentation}. [Online]. Available: \url{https://code.claude.com/docs/en/how-claude-code-works}. [accessed 23 Aug. 2026].
\end{refslist}

\begin{refslist}
\item[4.] OWASP GenAI Security Project, \emph{OWASP Top 10 for LLM Applications}, 2026 ed., published 4 Aug. 2026. [Online]. Available: \url{https://genai.owasp.org/}. [accessed 23 Aug. 2026]. The earlier project page at owasp.org is retained by OWASP as a legacy entry point.
\end{refslist}

\begin{refslist}
\item[5.] MITRE, \emph{ATLAS: Adversarial Threat Landscape for Artificial-Intelligence Systems}. [Online]. Available: \url{https://atlas.mitre.org/}. [accessed 23 Aug. 2026].
\end{refslist}

\begin{refslist}
\item[6.] National Institute of Standards and Technology, \emph{Artificial Intelligence Risk Management Framework (AI RMF 1.0)}, NIST AI 100-1. Gaithersburg, MD, USA: NIST, Jan. 2023. doi: 10.6028/NIST.AI.100-1.
\end{refslist}

\begin{refslist}
\item[7.] National Institute of Standards and Technology, \emph{Artificial Intelligence Risk Management Framework: Generative Artificial Intelligence Profile}, NIST AI 600-1. Gaithersburg, MD, USA: NIST, Jul. 2024. doi: 10.6028/NIST.AI.600-1.
\end{refslist}

\begin{refslist}
\item[8.] S. Rose, O. Borchert, S. Mitchell and S. Connelly, \emph{Zero Trust Architecture}, NIST Special Publication 800-207. Gaithersburg, MD, USA: NIST, Aug. 2020. doi: 10.6028/NIST.SP.800-207.
\end{refslist}

\begin{refslist}
\item[9.] M. Souppaya, K. Scarfone and D. Dodson, \emph{Secure Software Development Framework (SSDF) Version 1.1: Recommendations for Mitigating the Risk of Software Vulnerabilities}, NIST Special Publication 800-218. Gaithersburg, MD, USA: NIST, Feb. 2022. doi: 10.6028/NIST.SP.800-218.
\end{refslist}

\begin{refslist}
\item[10.] National Institute of Standards and Technology, \emph{The NIST Cybersecurity Framework (CSF) 2.0}, NIST CSWP 29. Gaithersburg, MD, USA: NIST, Feb. 2024. doi: 10.6028/NIST.CSWP.29.
\end{refslist}

\begin{refslist}
\item[11.] International Organization for Standardization, \emph{ISO/IEC 42001:2023 — Information technology — Artificial intelligence — Management system}. Geneva, Switzerland: ISO, 2023.
\end{refslist}

\begin{refslist}
\item[12.] International Organization for Standardization, \emph{ISO/IEC 27001:2022 — Information security, cybersecurity and privacy protection — Information security management systems — Requirements}, 3rd ed. Geneva, Switzerland: ISO, 25 Oct. 2022. Amended by ISO/IEC 27001:2022/Amd 1:2024.
\end{refslist}

\begin{refslist}
\item[13.] International Organization for Standardization, \emph{ISO/IEC 23894:2023 — Information technology — Artificial intelligence — Guidance on risk management}. Geneva, Switzerland: ISO, 2023.
\end{refslist}

\begin{refslist}
\item[14.] International Organization for Standardization, \emph{ISO/IEC 25010:2023 — Systems and software engineering — Systems and software Quality Requirements and Evaluation (SQuaRE) — Product quality model}. Geneva, Switzerland: ISO, 2023.
\end{refslist}

\begin{refslist}
\item[15.] IEEE, \emph{IEEE Standard for System, Software, and Hardware Verification and Validation}, IEEE Std 1012-2016. New York, NY, USA: IEEE, 2017. doi: 10.1109/IEEESTD.2017.8055462.
\end{refslist}

\begin{refslist}
\item[16.] IEEE, \emph{IEEE Standard Model Process for Addressing Ethical Concerns during System Design}, IEEE Std 7000-2021. New York, NY, USA: IEEE, 2021. doi: 10.1109/IEEESTD.2021.9536679.
\end{refslist}

\begin{refslist}
\item[17.] European Parliament and Council of the European Union, “Regulation (EU) 2024/1689 laying down harmonised rules on artificial intelligence (Artificial Intelligence Act),” \emph{Official Journal of the European Union}, L series, 12 Jul. 2024.
\end{refslist}

\begin{refslist}
\item[18.] Australian Signals Directorate, Australian Cyber Security Centre, \emph{Essential Eight Maturity Model}. [Online]. Available: \url{https://www.cyber.gov.au/resources-business-and-government/}. essential-cyber-security/essential-eight. [accessed 23 Aug. 2026].
\end{refslist}

\begin{refslist}
\item[19.] Anthropic, “Orchestrate subagents at scale with dynamic workflows,” \emph{Claude Code Documentation}. [Online]. Available: \url{https://code.claude.com/docs/en/workflows}. [accessed 23 Aug. 2026].
\end{refslist}

\begin{refslist}
\item[20.] Anthropic, “Automate work with routines,” \emph{Claude Code Documentation}. [Online]. Available: \url{https://code.claude.com/docs/en/routines}. [accessed 23 Aug. 2026].
\end{refslist}

\begin{refslist}
\item[21.] Anthropic, “Use Claude Code with Chrome,” \emph{Claude Code Documentation}. [Online]. Available: \url{https://code.claude.com/docs/en/chrome}. [accessed 23 Aug. 2026].
\end{refslist}

\begin{refslist}
\item[22.] K. Greshake, S. Abdelnabi, S. Mishra, C. Endres, T. Holz and M. Fritz, “Not what you’ve signed up for: Compromising real-world LLM-integrated applications with indirect prompt injection,” in \emph{Proc. 16th ACM Workshop on Artificial Intelligence and Security (AISec ’23)}, Copenhagen, Denmark, Nov. 2023, pp. 79–90. doi: 10.1145/3605764.3623985.
\end{refslist}

\begin{refslist}
\item[23.] F. Perez and I. Ribeiro, “Ignore previous prompt: Attack techniques for language models,” in \emph{NeurIPS ML Safety Workshop}, New Orleans, LA, USA, Dec. 2022. arXiv:2211.09527.
\end{refslist}

\begin{refslist}
\item[24.] Anthropic, “Run prompts on a schedule,” \emph{Claude Code Documentation}. [Online]. Available: \url{https://code.claude.com/docs/en/scheduled-tasks}. [accessed 23 Aug. 2026].
\end{refslist}

\begin{refslist}
\item[25.] Anthropic, “Create custom subagents,” \emph{Claude Code Documentation}. [Online]. Available: \url{https://code.claude.com/docs/en/sub-agents}. [accessed 23 Aug. 2026].
\end{refslist}

\begin{refslist}
\item[26.] Anthropic, “How Claude remembers your project,” \emph{Claude Code Documentation}. [Online]. Available: \url{https://code.claude.com/docs/en/memory}. [accessed 23 Aug. 2026].
\end{refslist}

\begin{refslist}
\item[27.] Anthropic, “Configure the sandboxed Bash tool,” \emph{Claude Code Documentation}. [Online]. Available: \url{https://code.claude.com/docs/en/sandboxing}. [accessed 23 Aug. 2026].
\end{refslist}

\begin{refslist}
\item[28.] Anthropic, “Configure permissions,” \emph{Claude Code Documentation}. [Online]. Available: \url{https://code.claude.com/docs/en/permissions}. [accessed 23 Aug. 2026].
\end{refslist}

\begin{refslist}
\item[29.] Anthropic, “Orchestrate teams of Claude Code sessions,” \emph{Claude Code Documentation}. [Online]. Available: \url{https://code.claude.com/docs/en/agent-teams}. [accessed 23 Aug. 2026].
\end{refslist}

\begin{refslist}
\item[30.] Anthropic, “Manage sessions,” \emph{Claude Code Documentation}. [Online]. Available: \url{https://code.claude.com/docs/en/sessions}. [accessed 23 Aug. 2026].
\end{refslist}

\begin{refslist}
\item[31.] Anthropic, “Choose a permission mode,” \emph{Claude Code Documentation}. [Online]. Available: \url{https://code.claude.com/docs/en/permission-modes}. [accessed 23 Aug. 2026].
\end{refslist}

\begin{refslist}
\item[32.] Anthropic, “Connect Claude Code to tools via MCP,” \emph{Claude Code Documentation}. [Online]. Available: \url{https://code.claude.com/docs/en/mcp}. [accessed 23 Aug. 2026].
\end{refslist}

\begin{refslist}
\item[33.] Anthropic, “Explore the context window,” \emph{Claude Code Documentation}. [Online]. Available: \url{https://code.claude.com/docs/en/context-window}. [accessed 23 Aug. 2026].
\end{refslist}

\begin{refslist}
\item[34.] Anthropic, “Extend Claude with skills,” \emph{Claude Code Documentation}. [Online]. Available: \url{https://code.claude.com/docs/en/skills}. [accessed 23 Aug. 2026].
\end{refslist}

\begin{refslist}
\item[35.] Anthropic, “Commands,” \emph{Claude Code Documentation}. [Online]. Available: \url{https://code.claude.com/docs/en/commands}. [accessed 23 Aug. 2026].
\end{refslist}

\begin{refslist}
\item[36.] Anthropic, “What’s new,” \emph{Claude Code Documentation}. [Online]. Available: \url{https://code.claude.com/docs/en/whats-new/index}. [accessed 23 Aug. 2026].
\end{refslist}

\begin{refslist}
\item[37.] Anthropic, “Create plugins,” \emph{Claude Code Documentation}. [Online]. Available: \url{https://code.claude.com/docs/en/plugins}. [accessed 23 Aug. 2026].
\end{refslist}

\begin{refslist}
\item[38.] Anthropic, “Discover and install prebuilt plugins through marketplaces,” \emph{Claude Code Documentation}. [Online]. Available: \url{https://code.claude.com/docs/en/discover-plugins}. [accessed 23 Aug. 2026].
\end{refslist}

\begin{refslist}
\item[39.] Anthropic, “Continue local sessions from any device with Remote Control,” \emph{Claude Code Documentation}. [Online]. Available: \url{https://code.claude.com/docs/en/remote-control}. [accessed 23 Aug. 2026].
\end{refslist}

\begin{refslist}
\item[40.] Anthropic, “Zero data retention,” \emph{Claude Code Documentation}. [Online]. Available: \url{https://code.claude.com/docs/en/zero-data-retention}. [accessed 23 Aug. 2026].
\end{refslist}

\begin{refslist}
\item[41.] Anthropic, “Data usage,” \emph{Claude Code Documentation}. [Online]. Available: \url{https://code.claude.com/docs/en/data-usage}. [accessed 23 Aug. 2026].
\end{refslist}

\begin{refslist}
\item[42.] SOS Automazioni, \emph{Claude Code: il corso completo per fare tutto, non solo programmare}. [Video]. Published 10 Jul. 2026. Available: https://www.youtube.com/watch?v=w6LDEPJglHw. [accessed 23 Aug. 2026].
\end{refslist}

\begin{refslist}
\item[43.] Anthropic, “Platforms and integrations,” \emph{Claude Code Documentation}. [Online]. Available: \url{https://code.claude.com/docs/en/platforms}. [accessed 23 Aug. 2026].
\end{refslist}

\begin{refslist}
\item[44.] Anthropic, “Use Claude Code in VS Code,” \emph{Claude Code Documentation}. [Online]. Available: \url{https://code.claude.com/docs/en/vs-code}. [accessed 23 Aug. 2026].
\end{refslist}

\begin{refslist}
\item[45.] Anthropic, “JetBrains IDEs,” \emph{Claude Code Documentation}. [Online]. Available: \url{https://code.claude.com/docs/en/jetbrains}. [accessed 23 Aug. 2026].
\end{refslist}

\begin{refslist}
\item[46.] Anthropic, “Desktop application,” \emph{Claude Code Documentation}. [Online]. Available: \url{https://code.claude.com/docs/en/desktop}. [accessed 23 Aug. 2026].
\end{refslist}

\begin{refslist}
\item[47.] Anthropic, “Use Claude Code on the web,” \emph{Claude Code Documentation}. [Online]. Available: \url{https://code.claude.com/docs/en/claude-code-on-the-web}. [accessed 23 Aug. 2026].
\end{refslist}

\begin{refslist}
\item[48.] Anthropic, “Claude Code on mobile,” \emph{Claude Code Documentation}. [Online]. Available: \url{https://code.claude.com/docs/en/mobile}. [accessed 23 Aug. 2026].
\end{refslist}

\begin{refslist}
\item[49.] Anthropic, “Advanced setup,” \emph{Claude Code Documentation}. [Online]. Available: \url{https://code.claude.com/docs/en/setup}. [accessed 23 Aug. 2026].
\end{refslist}

\begin{refslist}
\item[50.] Anthropic, “Share session output as artifacts,” \emph{Claude Code Documentation}. [Online]. Available: \url{https://code.claude.com/docs/en/artifacts}. [accessed 23 Aug. 2026].
\end{refslist}

\begin{refslist}
\item[51.] Anthropic, “Push events into a running session with channels,” \emph{Claude Code Documentation}. [Online]. Available: \url{https://code.claude.com/docs/en/channels}. [accessed 23 Aug. 2026].
\end{refslist}

\begin{refslist}
\item[52.] Anthropic, “Authentication,” \emph{Claude Code Documentation}. [Online]. Available: \url{https://code.claude.com/docs/en/authentication}. [accessed 23 Aug. 2026].
\end{refslist}

\begin{refslist}
\item[53.] Anthropic, “Security,” \emph{Claude Code Documentation}. [Online]. Available: \url{https://code.claude.com/docs/en/security}. [accessed 23 Aug. 2026].
\end{refslist}

\begin{refslist}
\item[54.] Anthropic, “Set up Claude Code in a monorepo or large codebase,” \emph{Claude Code Documentation}. [Online]. Available: \url{https://code.claude.com/docs/en/large-codebases}. [accessed 23 Aug. 2026].
\end{refslist}

\begin{refslist}
\item[55.] Anthropic, “Model configuration,” \emph{Claude Code Documentation}. [Online]. Available: \url{https://code.claude.com/docs/en/model-config}. [accessed 23 Aug. 2026].
\end{refslist}

\begin{refslist}
\item[56.] Anthropic, “Speed up responses with fast mode,” \emph{Claude Code Documentation}. [Online]. Available: \url{https://code.claude.com/docs/en/fast-mode}. [accessed 23 Aug. 2026].
\end{refslist}

\begin{refslist}
\item[57.] Anthropic, “Escalate hard decisions with the advisor tool,” \emph{Claude Code Documentation}. [Online]. Available: \url{https://code.claude.com/docs/en/advisor}. [accessed 23 Aug. 2026].
\end{refslist}

\begin{refslist}
\item[58.] Anthropic, “Interactive mode,” \emph{Claude Code Documentation}. [Online]. Available: \url{https://code.claude.com/docs/en/interactive-mode}. [accessed 23 Aug. 2026].
\end{refslist}

\begin{refslist}
\item[59.] Anthropic, “Manage costs effectively,” \emph{Claude Code Documentation}. [Online]. Available: \url{https://code.claude.com/docs/en/costs}. [accessed 23 Aug. 2026].
\end{refslist}

\begin{refslist}
\item[60.] Anthropic, “How Claude Code uses prompt caching,” \emph{Claude Code Documentation}. [Online]. Available: \url{https://code.claude.com/docs/en/prompt-caching}. [accessed 23 Aug. 2026].
\end{refslist}

\begin{refslist}
\item[61.] Anthropic, “Voice dictation,” \emph{Claude Code Documentation}. [Online]. Available: \url{https://code.claude.com/docs/en/voice-dictation}. [accessed 23 Aug. 2026].
\end{refslist}

\begin{refslist}
\item[62.] Anthropic, “CLI reference,” \emph{Claude Code Documentation}. [Online]. Available: \url{https://code.claude.com/docs/en/cli-reference}. [accessed 23 Aug. 2026].
\end{refslist}

\begin{refslist}
\item[63.] Anthropic, “Claude Code settings,” \emph{Claude Code Documentation}. [Online]. Available: \url{https://code.claude.com/docs/en/settings}. [accessed 23 Aug. 2026].
\end{refslist}

\begin{refslist}
\item[64.] Anthropic, “Deploy managed settings,” \emph{Claude Code Documentation}. [Online]. Available: \url{https://code.claude.com/docs/en/managed-settings}. [accessed 23 Aug. 2026].
\end{refslist}

\begin{refslist}
\item[65.] Anthropic, “Choose a sandbox environment,” \emph{Claude Code Documentation}. [Online]. Available: \url{https://code.claude.com/docs/en/sandbox-environments}. [accessed 23 Aug. 2026].
\end{refslist}

\begin{refslist}
\item[66.] Anthropic, “Development containers,” \emph{Claude Code Documentation}. [Online]. Available: \url{https://code.claude.com/docs/en/devcontainer}. [accessed 23 Aug. 2026].
\end{refslist}

\begin{refslist}
\item[67.] Anthropic, “Checkpointing,” \emph{Claude Code Documentation}. [Online]. Available: \url{https://code.claude.com/docs/en/checkpointing}. [accessed 23 Aug. 2026].
\end{refslist}

\begin{refslist}
\item[68.] Anthropic, “Code Review,” \emph{Claude Code Documentation}. [Online]. Available: \url{https://code.claude.com/docs/en/code-review}. [accessed 23 Aug. 2026].
\end{refslist}

\begin{refslist}
\item[69.] Anthropic, “Catch security issues as Claude writes code,” \emph{Claude Code Documentation}. [Online]. Available: \url{https://code.claude.com/docs/en/security-guidance}. [accessed 23 Aug. 2026].
\end{refslist}

\begin{refslist}
\item[70.] Anthropic, “Scan your codebase for vulnerabilities,” \emph{Claude Code Documentation}. [Online]. Available: \url{https://code.claude.com/docs/en/claude-security}. [accessed 23 Aug. 2026].
\end{refslist}

\begin{refslist}
\item[71.] Anthropic, “Find bugs with ultrareview,” \emph{Claude Code Documentation}. [Online]. Available: \url{https://code.claude.com/docs/en/ultrareview}. [accessed 23 Aug. 2026].
\end{refslist}

\begin{refslist}
\item[72.] Anthropic, “Debug your configuration,” \emph{Claude Code Documentation}. [Online]. Available: \url{https://code.claude.com/docs/en/debug-your-config}. [accessed 23 Aug. 2026].
\end{refslist}

\begin{refslist}
\item[73.] Anthropic, “Hooks reference,” \emph{Claude Code Documentation}. [Online]. Available: \url{https://code.claude.com/docs/en/hooks}. [accessed 23 Aug. 2026].
\end{refslist}

\begin{refslist}
\item[74.] Anthropic, “Keep Claude working toward a goal,” \emph{Claude Code Documentation}. [Online]. Available: \url{https://code.claude.com/docs/en/goal}. [accessed 23 Aug. 2026].
\end{refslist}

\begin{refslist}
\item[75.] Model Context Protocol, “Introduction — Model Context Protocol specification.” [Online]. Available: \url{https://modelcontextprotocol.io/introduction}. [accessed 23 Aug. 2026].
\end{refslist}

\begin{refslist}
\item[76.] Anthropic, “Control MCP server access for your organization,” \emph{Claude Code Documentation}. [Online]. Available: \url{https://code.claude.com/docs/en/managed-mcp}. [accessed 23 Aug. 2026].
\end{refslist}

\begin{refslist}
\item[77.] GitHub, “Removing sensitive data from a repository,” \emph{GitHub Docs}. [Online]. Available: \url{https://docs.github.com/en/authentication/}. keeping-your-account-and-data-secure/removing-sensitive-data-from-a-repository. [accessed 23 Aug. 2026].
\end{refslist}

\begin{refslist}
\item[78.] Anthropic, “Plugins reference,” \emph{Claude Code Documentation}. [Online]. Available: \url{https://code.claude.com/docs/en/plugins-reference}. [accessed 23 Aug. 2026].
\end{refslist}

\begin{refslist}
\item[79.] Anthropic, “Output styles,” \emph{Claude Code Documentation}. [Online]. Available: \url{https://code.claude.com/docs/en/output-styles}. [accessed 23 Aug. 2026].
\end{refslist}

\begin{refslist}
\item[80.] Anthropic, “Let Claude use your computer from the CLI,” \emph{Claude Code Documentation}. [Online]. Available: \url{https://code.claude.com/docs/en/computer-use}. [accessed 23 Aug. 2026].
\end{refslist}

\begin{refslist}
\item[81.] Anthropic, “Tools reference,” \emph{Claude Code Documentation}. [Online]. Available: \url{https://code.claude.com/docs/en/tools-reference}. [accessed 23 Aug. 2026].
\end{refslist}

\begin{refslist}
\item[82.] Anthropic, “Schedule recurring tasks in Claude Code Desktop,” \emph{Claude Code Documentation}. [Online]. Available: \url{https://code.claude.com/docs/en/desktop-scheduled-tasks}. [accessed 23 Aug. 2026].
\end{refslist}

\begin{refslist}
\item[83.] Anthropic, “Configure cloud environments,” \emph{Claude Code Documentation}. [Online]. Available: \url{https://code.claude.com/docs/en/cloud-environments}. [accessed 23 Aug. 2026].
\end{refslist}

\begin{refslist}
\item[84.] Anthropic, “Run Claude Code programmatically,” \emph{Claude Code Documentation}. [Online]. Available: \url{https://code.claude.com/docs/en/headless}. [accessed 23 Aug. 2026].
\end{refslist}

\begin{refslist}
\item[85.] Anthropic, “Enterprise network configuration,” \emph{Claude Code Documentation}. [Online]. Available: \url{https://code.claude.com/docs/en/network-config}. [accessed 23 Aug. 2026].
\end{refslist}

\begin{refslist}
\item[86.] Anthropic, “Run agents in parallel,” \emph{Claude Code Documentation}. [Online]. Available: \url{https://code.claude.com/docs/en/agents}. [accessed 23 Aug. 2026].
\end{refslist}

\begin{refslist}
\item[87.] Anthropic, “Manage multiple agents with agent view,” \emph{Claude Code Documentation}. [Online]. Available: \url{https://code.claude.com/docs/en/agent-view}. [accessed 23 Aug. 2026].
\end{refslist}

\begin{refslist}
\item[88.] Anthropic, “Run parallel sessions with worktrees,” \emph{Claude Code Documentation}. [Online]. Available: \url{https://code.claude.com/docs/en/worktrees}. [accessed 23 Aug. 2026].
\end{refslist}

\begin{refslist}
\item[89.] Anthropic, “Message your other Claude Code sessions,” \emph{Claude Code Documentation}. [Online]. Available: \url{https://code.claude.com/docs/en/cross-session-messaging}. [accessed 23 Aug. 2026].
\end{refslist}

\begin{refslist}
\item[90.] Office of the Australian Information Commissioner, \emph{Australian Privacy Principles guidelines}, Privacy Act 1988 (Cth). [Online]. Available: \url{https://www.oaic.gov.au/}. [accessed 23 Aug. 2026].
\end{refslist}

\begin{refslist}
\item[91.] European Parliament and Council of the European Union, “Regulation (EU) 2016/679 (General Data Protection Regulation),” \emph{Official Journal of the European Union}, L 119, pp. 1–88, 4 May 2016.
\end{refslist}

\begin{refslist}
\item[92.] Anthropic, “Explore the .claude directory,” \emph{Claude Code Documentation}. [Online]. Available: \url{https://code.claude.com/docs/en/claude-directory}. [accessed 23 Aug. 2026].
\end{refslist}

\begin{refslist}
\item[93.] Anthropic, “Monitoring,” \emph{Claude Code Documentation}. [Online]. Available: \url{https://code.claude.com/docs/en/monitoring-usage}. [accessed 23 Aug. 2026].
\end{refslist}

\begin{refslist}
\item[94.] Anthropic, “Track team usage with analytics,” \emph{Claude Code Documentation}. [Online]. Available: \url{https://code.claude.com/docs/en/analytics}. [accessed 23 Aug. 2026].
\end{refslist}

\begin{refslist}
\item[95.] Anthropic, “Use Claude Code with a screen reader,” \emph{Claude Code Documentation}. [Online]. Available: \url{https://code.claude.com/docs/en/accessibility}. [accessed 23 Aug. 2026].
\end{refslist}

\begin{refslist}
\item[96.] Anthropic, “Customize keyboard shortcuts,” \emph{Claude Code Documentation}. [Online]. Available: \url{https://code.claude.com/docs/en/keybindings}. [accessed 23 Aug. 2026].
\end{refslist}

\begin{refslist}
\item[97.] Anthropic, “Customize your status line,” \emph{Claude Code Documentation}. [Online]. Available: \url{https://code.claude.com/docs/en/statusline}. [accessed 23 Aug. 2026].
\end{refslist}

\begin{refslist}
\item[98.] Anthropic, “Configure your terminal for Claude Code,” \emph{Claude Code Documentation}. [Online]. Available: \url{https://code.claude.com/docs/en/terminal-config}. [accessed 23 Aug. 2026].
\end{refslist}

\begin{refslist}
\item[99.] Anthropic, “Configure auto mode,” \emph{Claude Code Documentation}. [Online]. Available: \url{https://code.claude.com/docs/en/auto-mode-config}. [accessed 23 Aug. 2026].
\end{refslist}

\begin{refslist}
\item[100.] Anthropic, “Claude Code settings reference,” \emph{Claude Code Documentation}. [Online]. Available: \url{https://code.claude.com/docs/en/settings-reference}. [accessed 23 Aug. 2026].
\end{refslist}

\begin{refslist}
\item[101.] Anthropic, “Environment variables,” \emph{Claude Code Documentation}. [Online]. Available: \url{https://code.claude.com/docs/en/env-vars}. [accessed 23 Aug. 2026].
\end{refslist}

\begin{refslist}
\item[102.] Anthropic, “Troubleshoot installation and login,” \emph{Claude Code Documentation}. [Online]. Available: \url{https://code.claude.com/docs/en/troubleshoot-install}. [accessed 23 Aug. 2026].
\end{refslist}

\begin{refslist}
\item[103.] Anthropic, “Error reference,” \emph{Claude Code Documentation}. [Online]. Available: \url{https://code.claude.com/docs/en/errors}. [accessed 23 Aug. 2026].
\end{refslist}

\begin{refslist}
\item[104.] SLSA steering committee, \emph{Supply-chain Levels for Software Artifacts (SLSA)}, v1.2, a project within the Open Source Security Foundation. [Online]. Available: \url{https://slsa.dev/}. [accessed 23 Aug. 2026].
\end{refslist}

\begin{refslist}
\item[105.] National Cyber Security Centre (UK), Cybersecurity and Infrastructure Security Agency (US) and international partners, \emph{Guidelines for Secure AI System Development}, Nov. 2023. [Online]. Available: \url{https://www.ncsc.gov.uk/collection/guidelines-secure-ai-system-development}. [accessed 23 Aug. 2026].
\end{refslist}

\begin{refslist}
\item[106.] Anthropic, “Claude Code changelog,” \emph{Claude Code Documentation}. [Online]. Available: \url{https://code.claude.com/docs/en/changelog}. [accessed 23 Aug. 2026].
\end{refslist}


\breakrule

\FloatBarrier
\frontchapter{About the author}

\textbf{David Soldani} (Senior Member, IEEE) received the M.Sc. degree (\emph{magna cum laude}) in engineering from the University of Florence, Italy, in 1994, and the D.Sc. degree in technology (Hons.) from the Helsinki University of Technology, Finland, in 2006.


Throughout his career he has held academic positions including Visiting Professor at the University of Surrey, U.K. (2014), Industry Professor at the University of Technology Sydney (UTS), Australia (2016), and Adjunct Professor at the University of New South Wales (UNSW) (2018).


In his professional roles he has served as Chief Information and Security Officer (CISO) and SVP Innovation and Advanced Research at Rakuten; Chief Technology and Cyber Security Officer with Huawei Asia Pacific; Head of 5G Technology at Nokia; and Head of the Central Research Institute and VP Strategic Research and Innovation in Europe at Huawei European Research Center. He is currently SVP Advanced Research and Innovation with Rakuten Mobile Inc., Tokyo, Japan.


Recognised internationally as one of the pioneers of 5G and 6G, he works across multi-disciplinary frontier research and innovation, with expertise spanning 6G, 5G-Advanced, cyber security, cloud computing, eBPF, artificial intelligence and big data. He has received numerous awards for leadership and professional contribution across the organisations he has worked in. He was granted a Distinguished Talent visa by the Australian Government in 2016.


He can be reached at \href{https://www.linkedin.com/in/dr-david-soldani/}{linkedin.com/in/dr-david-soldani}.


\emph{The views expressed in this book are the author's own and do not represent those of any current or former employer. No employer sponsored, reviewed or endorsed this work.}


\end{document}